% !TeX program = pdflatex
% Tom I: rachunek prawdopodobienstwa. Rozklady, funkcje charakterystyczne i twierdzenia graniczne.
% Plik samodzielny; kolejne czesci mozna dopisywac przed \end{document}.
\documentclass[11pt,a4paper,oneside]{report}
\usepackage[T1]{fontenc}
\usepackage[utf8]{inputenc}
\usepackage[polish]{babel}
\usepackage{lmodern}
\usepackage{amsmath,amssymb,amsthm,mathtools}
\usepackage{microtype}
\usepackage{geometry}
\usepackage{enumitem}
\usepackage{xcolor}
\usepackage{tikz}
\usetikzlibrary{arrows.meta,positioning,calc,shapes.geometric}
\usepackage[unicode]{hyperref}
\geometry{margin=2.65cm}
\hypersetup{colorlinks=true,linkcolor=blue!45!black,urlcolor=blue!45!black}
\setlength{\parindent}{0pt}
\setlength{\parskip}{4pt}
\setlength{\emergencystretch}{2em}
\setlist[enumerate]{topsep=3pt,itemsep=2pt}
\setlist[itemize]{topsep=3pt,itemsep=2pt}
\numberwithin{equation}{chapter}

\theoremstyle{plain}
\newtheorem{twierdzenie}{Twierdzenie}[chapter]
\newtheorem{lemat}[twierdzenie]{Lemat}
\newtheorem{wniosek}[twierdzenie]{Wniosek}
\newtheorem{stwierdzenie}[twierdzenie]{Stwierdzenie}
\theoremstyle{definition}
\newtheorem{definicja}{Definicja}[chapter]
\newtheorem{przyklad}{Przykład}[chapter]
\theoremstyle{remark}
\newtheorem{uwaga}{Uwaga}[chapter]
\renewcommand{\proofname}{Dowód}

\newcommand{\R}{\mathbb{R}}
\newcommand{\N}{\mathbb{N}}
\newcommand{\Z}{\mathbb{Z}}
\newcommand{\Prob}{\mathbb{P}}
\newcommand{\E}{\mathbb{E}}
\newcommand{\Var}{\operatorname{Var}}
\newcommand{\Cov}{\operatorname{Cov}}
\newcommand{\Law}{\operatorname{Law}}
\newcommand{\supp}{\operatorname{supp}}
\newcommand{\ind}[1]{\mathbf{1}_{#1}}
\newcommand{\dd}{\mathrm{d}}
\newcommand{\Bin}{\operatorname{Bin}}
\newcommand{\Poi}{\operatorname{Poi}}
\newcommand{\Geom}{\operatorname{Geom}}
\newcommand{\Unif}{\operatorname{U}}
\newcommand{\Gam}{\operatorname{Gamma}}
\newcommand{\Betad}{\operatorname{Beta}}
\newcommand{\Normal}{\mathcal{N}}
\newcommand{\iid}{\mathrel{\overset{\mathrm{iid}}{\sim}}}

\begin{document}
\begin{titlepage}
\centering
\vspace*{4cm}
{\Large Podstawy teorii prawdopodobieństwa\par}
\vspace{1.0cm}
{\Huge\bfseries Tom I. Rachunek prawdopodobieństwa\par}
\vspace{0.7cm}
{\Large Od kombinatoryki do łańcuchów Markowa i martyngałów\par}
\vfill
{\large Wersja robocza do dalszej rozbudowy\par}
\vspace{1.4cm}
\end{titlepage}
\tableofcontents

\chapter*{Wprowadzenie}
\addcontentsline{toc}{chapter}{Wprowadzenie}
Projekt składa się z czterech tomów: \emph{Rachunek prawdopodobieństwa}, \emph{Procesy stochastyczne}, \emph{Analiza stochastyczna} oraz \emph{Statystyka matematyczna}.
Pierwszy tom buduje podstawy teorii na przestrzeniach probabilistycznych:
od metod zliczania zdarzeń elementarnych do opisu rozkładów zmiennych losowych.
W obecnej części tomu wprowadzimy również funkcje tworzące i charakterystyczne,
rodzaje zbieżności zmiennych losowych, słabe i mocne prawa wielkich liczb
oraz centralne twierdzenie graniczne. Następnie zbudujemy
warunkową wartość oczekiwaną z twierdzenia Radona--Nikodyma,
wykażemy jej własności projekcyjne w $L^2$ oraz rozszerzymy
teorię granic o koncentrację, schematy trójkątne,
metodę delty i zbieżność dystrybuant empirycznych.
Dalsze rozdziały zawierają teorię dyskretnych łańcuchów Markowa,
filtracji, czasów zatrzymania, martyngałów, podmartyngałów i nadmartyngałów.
W dowodach szczegółowo rozróżniamy ograniczoność czasu,
całkowalność oraz jednostajną całkowalność.

Przez \emph{model klasyczny} rozumiemy w szczególności model
Laplace'a oparty na skończonej liczbie jednakowo prawdopodobnych
wyników. Nie jest on alternatywą dla miary probabilistycznej,
lecz jej bardzo szczególnym przypadkiem. Zanim przejdziemy do
modelowania danych i nieznanych parametrów, oddzielimy aksjomaty
rachunku prawdopodobieństwa od interpretacji konkretnych doświadczeń.

Przyjmujemy znajomość podstaw analizy matematycznej, całki Lebesgue'a,
pojęcia przestrzeni mierzalnej i twierdzenia Radona--Nikodyma.
Dowody skupiają się na konstrukcjach i istotnych przejściach logicznych;
w miejscach wymagających zaawansowanej teorii miary zaznaczamy, z jakiego
uprzednio znanego twierdzenia korzystamy.

W tej części pozostajemy przy aksjomatycznym i klasycznym rachunku
prawdopodobieństwa. Wzór Bayesa wystąpi jako tożsamość wynikająca
z definicji prawdopodobieństwa warunkowego, lecz interpretację
bayesowską parametrów, rozkłady a priori i a posteriori,
wiarygodność i wnioskowanie statystyczne odłożymy na później.

Podstawowe pojęcia będziemy wprowadzać według trzech pytań: jaki
problem prowadzi do danej definicji, co ta definicja rzeczywiście
mówi o losowym doświadczeniu i jak wynikają z niej twierdzenia.
Obok rachunków skończonych zestawimy modele o nieskończenie
wielu wynikach, w których intuicja oparta wyłącznie na liczeniu
może być myląca.

\chapter{Kombinatoryka i modele skończone}

\section{Zasady zliczania}
Zacznijmy od najprostszego problemu: chcemy obliczyć szansę
pewnego wyniku, ale najpierw musimy ustalić, ile wyników w ogóle
rozważamy. Przy rzutach kostką, losowaniu kart albo układaniu haseł
istotne są trzy pytania: czy obiekty są rozróżnialne, czy kolejność
ma znaczenie i czy dopuszczamy powtórzenia. Zmiana choćby jednej
odpowiedzi prowadzi do innego modelu.

Zasada iloczynu nie wymaga probabilistycznej niezależności.
Jest twierdzeniem o liczbie możliwych ciągów wyborów, podczas gdy
niezależność jest własnością \emph{miary} nadanej zdarzeniom.
Dla skończonych modeli te dwa pojęcia często występują obok siebie,
ale nie wolno ich utożsamiać.

\begin{definicja}
Dla liczby naturalnej $n\geq 0$ definiujemy $n!:=\prod_{j=1}^n j$, przyjmując $0!:=1$.
Dla liczb całkowitych $n\geq0$ i $0\leq k\leq n$ definiujemy
$\binom nk:=n!/[k!(n-k)!]$; dla pozostałych całkowitych $k$ przyjmujemy
$\binom nk:=0$.
\end{definicja}

\begin{stwierdzenie}[Zasady sumy i iloczynu]
Niech $A_1,\dots,A_m$ będą skończonymi zbiorami.
Jeżeli są parami rozłączne, to $|\bigcup_{i=1}^m A_i|=\sum_{i=1}^m|A_i|$.
Bez żadnego założenia o rozłączności zachodzi
$|A_1\times\cdots\times A_m|=\prod_{i=1}^m|A_i|$.
\end{stwierdzenie}
\begin{proof}
Dla rozłącznych zbiorów odwzorowanie
$\coprod_{i=1}^m A_i\to\bigcup_{i=1}^m A_i$
wysyłające element z $i$-tego składnika na ten sam element
jest bijekcją, ponieważ każdy element sumy należy do
dokładnie jednego z $A_i$. Skończony zbiór jest więc
rozbity na rozłączne części o licznościach $|A_i|$.

Dla dwóch zbiorów $A,B$ rozbijmy $A\times B$ na włókna
$\{a\}\times B$ indeksowane przez $a\in A$.
Są parami rozłączne i każde ma dokładnie $|B|$ elementów.
Pierwsza zasada daje
$|A\times B|=\sum_{a\in A}|B|=|A|\,|B|$.
Dla $m>2$ traktujemy
$A_1\times\cdots\times A_m$ jako
$(A_1\times\cdots\times A_{m-1})\times A_m$
i stosujemy powyższy argument indukcyjnie.
Ważne, że jest to dowód czysto zbiorowy i w ogóle
nie korzysta z pojęcia prawdopodobieństwa.
\end{proof}

\begin{twierdzenie}[Permutacje, wariacje i kombinacje]
Niech $S$ będzie zbiorem $n$ różnych elementów i niech $0\leq k\leq n$.
Wówczas liczba bijekcji $S\to S$ wynosi $n!$, liczba uporządkowanych
$k$-elementowych ciągów bez powtórzeń wynosi
$(n)_k:=n!/(n-k)!$, a liczba $k$-elementowych podzbiorów $S$ wynosi
$\binom nk$. Liczba ciągów długości $k$ o wyrazach w $S$, z dopuszczeniem
powtórzeń, wynosi $n^k$.
\end{twierdzenie}
\begin{proof}
Uporządkowany ciąg $(s_1,\ldots,s_k)$ bez powtórzeń
budujemy kolejno. Pierwszy wyraz ma $n$ możliwych wartości.
Po wyborze pierwszych $j-1$ wyrazów pozostało dokładnie
$n-j+1$ wartości dla $s_j$. Z zasady iloczynu liczba
ciągów jest równa
$n(n-1)\cdots(n-k+1)=n!/(n-k)!$.
Przy $k=0$ otrzymujemy jeden ciąg pusty,
a przy $k=n$ wszystkie bijekcje, czyli permutacje.

Niech $\mathcal V_k$ będzie zbiorem tych ciągów, a
$\mathcal C_k$ zbiorem $k$-elementowych podzbiorów $S$.
Odwzorowanie $\pi:\mathcal V_k\to\mathcal C_k$,
które zapomina kolejność, jest surjekcją. Dla
każdego $C\in\mathcal C_k$ zbiór $\pi^{-1}(C)$
zawiera dokładnie $k!$ uporządkowań elementów $C$.
Wszystkie włókna mają więc tę samą liczność, skąd
$|\mathcal V_k|=k!|\mathcal C_k|$ i
$|\mathcal C_k|=\binom nk$.
Dla ciągów z powtórzeniami każde miejsce ma $n$
możliwych wartości bez względu na wcześniejsze wybory,
więc ich liczba to $n^k$.
\end{proof}

\begin{przyklad}[Kod dostępu]
Kod złożony z sześciu cyfr, w którym dopuszcza się początkowe zero,
może mieć $10^6$ postaci. Jeżeli cyfry nie mogą się powtarzać,
postaci jest $(10)_6$. Jeżeli dodatkowo kod ma zaczynać się cyfrą niezerową,
a pozostałe cyfry nie mogą się powtarzać, otrzymujemy
$9\cdot(9)_5$ postaci. Warunki dotyczące pierwszej pozycji rozpatruje się
zatem przed zastosowaniem zasady iloczynu.
\end{przyklad}

\section{Powtórzenia i rozdzielanie obiektów}
Kombinacje opisują wybór podzbioru różnych elementów. Nie nadają
się jednak bezpośrednio do opisu rozdzielania identycznych przedmiotów:
zamiana miejscami dwóch identycznych żetonów niczego nie zmienia.
Odwrotnie, jeśli mamy kilka kopii jednego symbolu, permutacje
sztucznie oznakowanych kopii wielokrotnie liczą ten sam wynik.
Stąd wynikają dwie osobne techniki: dzielenie przez liczbę
niewidocznych permutacji oraz kodowanie przydziałów
za pomocą gwiazdek i kresek.

\begin{twierdzenie}[Permutacje multizbioru]
Niech $n_1,\ldots,n_r$ będą nieujemnymi liczbami całkowitymi, a
$n=n_1+\cdots+n_r$. Liczba ciągów długości $n$, w których symbol $i$
występuje dokładnie $n_i$ razy dla każdego $i=1,\ldots,r$, wynosi
$n!/(n_1!\cdots n_r!)$.
\end{twierdzenie}
\begin{proof}
Wyobraźmy sobie, że wszystkie wystąpienia symbolu $i$
otrzymują dodatkowe etykiety $(i,1),\ldots,(i,n_i)$.
Powstaje $n$ różnych obiektów, których ustawień jest $n!$.
Po usunięciu etykiet otrzymujemy ciąg z zadanymi
licznościami symboli. Odwzorowanie „usuń etykiety”
obejmuje wszystkie dopuszczalne ciągi.

Ustalmy jeden taki ciąg. Etykiety symbolu $i$
można przypisać jego $n_i$ wystąpieniom na $n_i!$
sposobów. Te operacje dla różnych symboli są niezależnymi
operacjami zliczania, więc łącznie dostajemy
$\prod_{i=1}^r n_i!$ oznakowanych ciągów.
Każdy ciąg nieoznakowany ma tyle samo reprezentacji.
Dzieląc $n!$ przez rozmiar jednego włókna, otrzymujemy
$n!/(n_1!\cdots n_r!)$.
Dla $n_i=0$ czynnik $n_i!=0!=1$ nie zmienia rachunku.
\end{proof}

\begin{twierdzenie}[Metoda gwiazdek i kresek]
Dla liczb całkowitych $n\geq0$ i $r\geq1$ liczba rozwiązań
$(x_1,\ldots,x_r)\in\Z_{\geq0}^r$ równania $x_1+\cdots+x_r=n$ wynosi
$\binom{n+r-1}{r-1}$. Jeżeli dodatkowo $x_i\geq1$ dla każdego $i$,
a $n\geq r$, liczba rozwiązań wynosi $\binom{n-1}{r-1}$.
\end{twierdzenie}
\begin{proof}
Dla rozwiązania $(x_1,\ldots,x_r)$ z $x_i\geq0$
budujemy słowo, w którym umieszczamy $x_1$ gwiazdek,
następnie kreskę, potem $x_2$ gwiazdek i tak dalej,
aż do $x_r$ gwiazdek. Powstaje dokładnie $n$ gwiazdek
i $r-1$ kresek, a więc słowo długości $n+r-1$.
Dopuszczamy kreski obok siebie i na brzegach,
co odpowiada współrzędnym równym zero.

Odwrotnie, z każdego słowa zawierającego $n$ gwiazdek
i $r-1$ kresek odzyskujemy jednoznacznie liczby $x_i$
przez policzenie gwiazdek w kolejnych blokach.
Mamy więc bijekcję między rozwiązaniami a wyborami
$r-1$ pozycji kresek spośród $n+r-1$ pozycji.
Stąd pierwszy wzór. Gdy wymagamy $x_i\geq1$, podstawiamy
$y_i=x_i-1\geq0$. Równanie staje się
$y_1+\cdots+y_r=n-r$, więc stosujemy pierwszy wzór
do $n-r$ i otrzymujemy $\binom{n-1}{r-1}$.
\end{proof}

\begin{przyklad}[Rozdzielanie jednakowych przedmiotów]
Rozdzielamy 12 identycznych żetonów pomiędzy cztery rozróżnialne osoby.
Liczba przydziałów bez dodatkowych warunków wynosi $\binom{15}{3}$.
Jeżeli każda osoba ma otrzymać przynajmniej dwa żetony, podstawienie
$y_i=x_i-2$ daje $\sum_i y_i=4$, zatem $\binom73$ przydziałów.
Wyniki nie odnoszą się do 12 \emph{rozróżnialnych} żetonów:
wtedy bez ograniczeń mamy $4^{12}$ przydziałów.
\end{przyklad}

\begin{stwierdzenie}[Tożsamość Pascala i dwumian Newtona]
Dla $n\geq1$ oraz $1\leq k\leq n-1$ zachodzi
$\binom nk=\binom{n-1}{k-1}+\binom{n-1}k$. Dla dowolnych liczb
$x,y\in\R$ i $n\geq0$ zachodzi
$(x+y)^n=\sum_{k=0}^n\binom nk x^k y^{n-k}$.
\end{stwierdzenie}
\begin{proof}
Ustalmy zbiór $S=\{1,\ldots,n\}$ i wyróżnijmy element
$n$. Każdy $k$-elementowy podzbiór $S$ albo zawiera $n$,
albo go nie zawiera. W pierwszym przypadku pozostałe
$k-1$ elementów można wybrać z $n-1$ elementów na
$\binom{n-1}{k-1}$ sposobów. W drugim wybieramy wszystkie
$k$ elementów spośród $n-1$, na $\binom{n-1}{k}$
sposobów. Obie rodziny są rozłączne i wyczerpują wszystkie
przypadki; daje to tożsamość Pascala.

Aby rozwinąć $(x+y)^n$, zapiszmy iloczyn $n$ czynników
$(x+y)$. Każdy wyraz po wymnożeniu uzyskujemy przez
wybranie z każdego czynnika albo $x$, albo $y$.
Jeśli wybraliśmy $x$ dokładnie $k$ razy, powstaje
$x^ky^{n-k}$. Można to zrobić na $\binom nk$ sposobów.
Zsumowanie wkładów dla $k=0,\ldots,n$ daje wzór Newtona.
Ten argument działa dla wszystkich $x,y\in\R$, również
zerowych i ujemnych: jest to tożsamość algebraiczna.
\end{proof}

\section{Zasada włączeń i wyłączeń}
Dodawanie liczebności zbiorów działa tylko wtedy, gdy zbiory
nie zachodzą na siebie. Jeśli element spełnia kilka warunków,
pojawia się wielokrotne liczenie. Zasada włączeń i wyłączeń
naprawia to przez kolejne odejmowanie i dodawanie przecięć.
Jej charakterystyczne znaki naprzemienne nie są sztuczną sztuczką:
wynikają z tego, że element należący do $r$ zbiorów ma być
ostatecznie policzony dokładnie raz.

\begin{twierdzenie}[Zasada włączeń i wyłączeń]
Dla skończonych zbiorów $A_1,\ldots,A_m$ zachodzi
\begin{equation}\label{eq:wlaczenia}
\left|\bigcup_{i=1}^m A_i\right|
=\sum_{\varnothing\ne I\subseteq\{1,\ldots,m\}}
 (-1)^{|I|+1}\left|\bigcap_{i\in I}A_i\right|.
\end{equation}
\end{twierdzenie}
\begin{proof}
Niech $U=\bigcup_{i=1}^m A_i$. Wystarczy pokazać, że
po prawej stronie każdy element $x\in U$ jest policzony
dokładnie raz, a każdy element spoza $U$ zero razy.
Ustalmy dowolny $x$. Jeśli nie należy do żadnego $A_i$,
nie jest zawarty w żadnym przecięciu występującym
w prawej stronie. Jego całkowity wkład jest zerowy.

Jeśli $x$ należy dokładnie do $r\geq1$ spośród zbiorów,
to pojawia się w przecięciach tych $j$-elementowych
rodzin indeksów, które wybieramy z owych $r$ zbiorów.
Takich rodzin jest $\binom rj$. Ponieważ ich znak to
$(-1)^{j+1}$, łączny wkład punktu $x$ jest
\[
\sum_{j=1}^r(-1)^{j+1}\binom rj
=1-\sum_{j=0}^r(-1)^j\binom rj
=1-(1-1)^r=1.
\]
Każdy element sumy jest więc policzony dokładnie raz.
Zsumowanie po wszystkich elementach skończonego $U$
kończy dowód. Zauważmy, że naprzemienne znaki wynikają
z potrzeby usunięcia wielokrotnego liczenia.
\end{proof}

\begin{przyklad}[Permutacje bez punktów stałych]
Niech $D_n$ oznacza liczbę permutacji zbioru $\{1,\ldots,n\}$ bez punktów
stałych. Dla każdego $i$ oznaczmy przez $A_i$ zbiór permutacji
ustalających $i$. Dla $k$ ustalonych pozycji istnieje $(n-k)!$ permutacji,
więc z~\eqref{eq:wlaczenia} wynika
\[
D_n=\sum_{k=0}^{n}(-1)^k\binom nk(n-k)!
=n!\sum_{k=0}^{n}\frac{(-1)^k}{k!}.
\]
Dla losowej jednostajnej permutacji prawdopodobieństwo braku punktów
stałych wynosi $D_n/n!$ i zmierza do $e^{-1}$.
\end{przyklad}

\begin{przyklad}[Dokładnie jedna własność]
Dla trzech zbiorów $A,B,C$ liczba elementów należących do dokładnie
jednego z nich jest równa
$|A|+|B|+|C|-2(|A\cap B|+|A\cap C|+|B\cap C|)+3|A\cap B\cap C|$.
Istotnie, element należący do dokładnie $r\in\{0,1,2,3\}$ zbiorów
wnosi odpowiednio $0,1,0,0$ do wskazanej kombinacji.
\end{przyklad}

\section{Przykłady modeli kombinatorycznych}
W przykładach będziemy za każdym razem jawnie opisywać
przestrzeń wyników. To zasadniczy etap rozwiązania: pięć kart
losowanych równocześnie tworzy podzbiór, pięć kolejnych wyciągnięć
tworzy uporządkowany ciąg, a obie przestrzenie mają inne liczebności.
Jeśli każdy ciąg jest jednakowo prawdopodobny, odpowiednio opisane
zdarzenie może mieć w obu modelach tę samą szansę; nie oznacza to
jednak, że wolno zmieniać mianownik bez odpowiedniej zmiany licznika.

\begin{przyklad}[Losowanie kart]
Z talii 52 różnych kart losujemy pięć kart bez zwracania i nie
uwzględniamy kolejności. Przestrzeń wyników ma $\binom{52}5$ elementów.
Liczba układów z dokładnie dwiema kartami w ustalonym kolorze
(w znaczeniu jednego z czterech kolorów po 13 kart) wynosi
$\binom{13}2\binom{39}3$. Liczba układów zawierających dokładnie
jedną parę jednakowych rang wynosi
$13\binom42\binom{12}3 4^3$: wybieramy rangę pary,
dwie karty tej rangi, trzy inne różne rangi oraz po jednej z czterech
kart każdej z tych rang. Dzięki wymogowi różnych rang nie liczymy
dwóch par ani trójki.
\end{przyklad}

\begin{przyklad}[Problem urodzin]
Zakładamy, że daty urodzin $n$ osób są niezależne i jednostajne na
zbiorze 365 dni, bez uwzględniania lat przestępnych. Dla $n\leq365$
prawdopodobieństwo, że wszystkie daty są różne, jest równe
$\prod_{j=0}^{n-1}(1-j/365)$. Zatem prawdopodobieństwo przynajmniej
jednej zgodności wynosi
$1-\prod_{j=0}^{n-1}(1-j/365)$.
Dla $n=23$ przekracza ono $1/2$. Jest to wniosek z modelu,
nie twierdzenie o rzeczywistym rozkładzie dat urodzin.
\end{przyklad}

\begin{przyklad}[Drogi na siatce]
Rozważmy drogę z $(0,0)$ do $(a,b)$, złożoną z kroków $(1,0)$ i $(0,1)$,
gdzie $a,b\in\N\cup\{0\}$. Każda droga składa się z $a+b$ kroków,
w tym $a$ poziomych. Dróg jest więc $\binom{a+b}{a}$.
Po narzuceniu przejścia przez $(c,d)$, gdzie $0\leq c\leq a$,
$0\leq d\leq b$, liczba dróg jest iloczynem
$\binom{c+d}c\binom{a+b-c-d}{a-c}$.
\end{przyklad}

\section{Zliczanie podziałów}
Dotychczas przypisywaliśmy obiekty do ponumerowanych miejsc.
Teraz interesuje nas samo rozbicie zbioru na grupy, bez etykiet
grup. Tak powstają liczby Stirlinga drugiego rodzaju. Ich
rekurencja obrazuje często stosowaną metodę: wyróżniamy ostatni
element i pytamy, czy tworzy samodzielną grupę, czy dołącza
do istniejącej.

\begin{definicja}
Dla $n,k\geq0$ liczba Stirlinga drugiego rodzaju $S(n,k)$ oznacza
liczbę podziałów zbioru $n$-elementowego na dokładnie $k$
niepustych, nierozróżnialnych bloków. Ustalamy $S(0,0)=1$,
$S(n,0)=0$ dla $n>0$ oraz $S(0,k)=0$ dla $k>0$.
\end{definicja}
\begin{stwierdzenie}
Dla $n,k\geq1$ zachodzi
$S(n,k)=kS(n-1,k)+S(n-1,k-1)$.
Liczba surjekcji ze zbioru $n$-elementowego na zbiór $k$-elementowy
wynosi $k!S(n,k)=\sum_{j=0}^k(-1)^j\binom kj(k-j)^n$.
\end{stwierdzenie}
\begin{proof}
Rozważmy podział zbioru $\{1,\ldots,n\}$ na $k$
niepustych, nierozróżnialnych bloków. Blok zawierający
$n$ jest albo singletonem $\{n\}$, albo zawiera inny
element. W pierwszym przypadku po usunięciu tego bloku
zostaje podział $n-1$ elementów na $k-1$ bloków,
co można zrobić na $S(n-1,k-1)$ sposobów.

W drugim przypadku usunięcie $n$ pozostawia podział
$n-1$ elementów na $k$ niepustych bloków. Aby odtworzyć
wyjściowy podział, musimy dodatkowo wskazać jeden z $k$
bloków, do którego dołączamy $n$; jest więc
$kS(n-1,k)$ możliwości. Przypadki są rozłączne,
co daje rekurencję.

Surjekcja $f:\{1,\ldots,n\}\to\{1,\ldots,k\}$
wyznacza podział na $k$ niepustych włókien
$f^{-1}(j)$; odwrotnie, każdy podział na $k$ bloków
daje $k!$ surjekcji po przypisaniu etykiet włóknom.
Stąd $k!S(n,k)$. Alternatywnie, wszystkich funkcji jest
$k^n$, a te, które nie są surjekcjami, pomijają co najmniej
jeden element przeciwdziedziny. Dla ustalonego
$j$-elementowego zbioru pomijanych wartości pozostaje
$(k-j)^n$ funkcji. Z zasady włączeń i wyłączeń wynika
$\sum_{j=0}^k(-1)^j\binom kj(k-j)^n$.
\end{proof}

\begin{przyklad}[Obsadzenie wszystkich pudełek]
Mamy $n\geq1$ rozróżnialnych kul i $m\geq1$ rozróżnialnych pudełek.
Każdą kulę umieszczamy w jednym pudełku. Wszystkich przydziałów
jest $m^n$, a przydziałów bez pustego pudełka jest $m!S(n,m)$.
Ta sama liczba, policzona zasadą włączeń i wyłączeń, wynosi
\[
\sum_{j=0}^{m}(-1)^j\binom mj(m-j)^n.
\]
Gdy wszystkie $m^n$ przydziałów są jednakowo prawdopodobne,
prawdopodobieństwo obsadzenia wszystkich pudełek wynosi
$m!S(n,m)/m^n$. Przykładowo dla $n=4$ i $m=3$
mamy $S(4,3)=6$, więc prawdopodobieństwo wynosi $36/81=4/9$.
\end{przyklad}

\begin{przyklad}[Dokładnie dwie pary w pokerze]
Losujemy pięć kart z talii 52 kart, w której występuje 13 rang
po cztery karty każdej rangi. Układ z dokładnie dwiema parami
otrzymujemy następująco: wybieramy dwie rangi par na $\binom{13}2$
sposobów, wybieramy po dwie z czterech kart każdej rangi
na $\binom42^2$ sposobów, a ostatnią kartę spośród
$11\cdot4$ kart o pozostałych rangach. Wynikowe
prawdopodobieństwo to
$\binom{13}2\binom42^2\cdot44/\binom{52}5$.
Wymóg, aby ostatnia karta miała inną rangę, zapobiega
wielokrotnemu zliczaniu i wyklucza układ z trzema
kartami jednej rangi.
\end{przyklad}

\chapter{Przestrzeń probabilistyczna i miara prawdopodobieństwa}

\section{Zdarzenia i \texorpdfstring{$\sigma$}{sigma}-ciała}
W doświadczeniu losowym najpierw występuje pewien wynik
$\omega\in\Omega$. Z reguły nie chcemy jednak przyznawać
prawdopodobieństwa pojedynczym wynikom, lecz formułować zdania,
na przykład „suma oczek jest parzysta” lub „ciąg wyników ma
nieskończenie wiele jedynek”. Takie zdania odpowiadają podzbiorom
przestrzeni wyników. Operacje logiczne „nie”, „lub”, „i” odpowiadają
dopełnieniu, sumie i przecięciu zbiorów.

Dlaczego nie zawsze bierzemy wszystkie podzbiory $\Omega$?
Dla przestrzeni nieprzeliczalnych próba przypisania długości
każdemu podzbiorowi odcinka z zachowaniem standardowych
własności miary prowadzi do niemożliwości. Dlatego z góry
ustalamy rodzinę $\mathcal F$ dopuszczalnych zdarzeń.
Zamknięcie na operacje \emph{przeliczalne} pozwoli mówić
między innymi o zdarzeniu „co najmniej jedno z $A_n$ zachodzi”
i o zdarzeniach granicznych.

\begin{definicja}
Niech $\Omega$ będzie niepustym zbiorem. Rodzinę $\mathcal F\subseteq2^\Omega$
nazywamy $\sigma$-ciałem, gdy (i) $\Omega\in\mathcal F$,
(ii) z $A\in\mathcal F$ wynika $\Omega\setminus A\in\mathcal F$,
(iii) dla każdego ciągu $(A_n)_{n\geq1}\subseteq\mathcal F$ zachodzi
$\bigcup_{n\geq1}A_n\in\mathcal F$.
Parę $(\Omega,\mathcal F)$ nazywamy przestrzenią mierzalną,
a elementy $\mathcal F$ zdarzeniami.
\end{definicja}

\begin{stwierdzenie}
Jeżeli $\mathcal F$ jest $\sigma$-ciałem na $\Omega$, to
$\varnothing\in\mathcal F$ oraz $\bigcap_{n\geq1}A_n\in\mathcal F$
dla każdego ciągu $A_n\in\mathcal F$. Rodzina $\mathcal F$ jest
również zamknięta na skończone sumy, przekroje i różnice zbiorów.
\end{stwierdzenie}
\begin{proof}
Ponieważ $\Omega\in\mathcal F$ i $\mathcal F$ jest zamknięte
na dopełnienia, to $\varnothing=\Omega^c\in\mathcal F$.
Jeśli $A_1,A_2,\ldots\in\mathcal F$, to każde
$A_n^c$ również należy do $\mathcal F$. Wobec tego
$\bigcup_{n\geq1}A_n^c\in\mathcal F$, a przez ponowne
wzięcie dopełnienia i prawo de Morgana dostajemy
\[
\bigcap_{n\geq1}A_n
=\left(\bigcup_{n\geq1}A_n^c\right)^c\in\mathcal F.
\]
Sumy i przekroje skończone są przypadkami przeliczalnymi:
do skończonej rodziny dopisujemy na końcu zbiory puste
lub całe $\Omega$, odpowiednio. Różnicę zapisujemy jako
$A\setminus B=A\cap B^c$, co dowodzi pozostałych własności.
\end{proof}

\begin{definicja}
Dla przestrzeni topologicznej $E$ przez $\mathcal B(E)$ oznaczamy
najmniejsze $\sigma$-ciało zawierające wszystkie zbiory otwarte w $E$;
nazywamy je $\sigma$-ciałem borelowskim.
\end{definicja}

\begin{przyklad}
Dla zbioru skończonego $\Omega$ przykładem jest $2^\Omega$.
Dla dowolnego $\Omega$ przykładem jest $\{\varnothing,\Omega\}$.
Jeżeli $A$ jest niepustym właściwym podzbiorem $\Omega$, to
$\{\varnothing,A,\Omega\setminus A,\Omega\}$ jest $\sigma$-ciałem:
pozwala rozstrzygać jedynie, w której z dwóch grup leży wynik.
Na $\R$ najczęściej używamy $\mathcal B(\R)$, a nie całego $2^\R$.
\end{przyklad}

\section{Aksjomaty Kołmogorowa i ich konsekwencje}
Dla skończonego zbioru wyników wystarczy podać wagi punktów,
które sumują się do jedności. W przestrzeniach nieskończonych
potrzebujemy zasady zgodności również dla przeliczalnych sum
rozłącznych zdarzeń. Stąd przeliczalna addytywność: jeśli zdarzenie
rozłożono na przeliczalnie wiele wykluczających się przypadków,
jego masa musi być sumą ich mas. Nie zakładamy wszystkich znanych
wzorów rachunku prawdopodobieństwa, lecz \emph{wyprowadzamy} je
z dwóch warunków: normalizacji i przeliczalnej addytywności.

\begin{definicja}
Niech $(\Omega,\mathcal F)$ będzie przestrzenią mierzalną.
Miarą probabilistyczną na $\mathcal F$ nazywamy funkcję
$\Prob:\mathcal F\to[0,1]$, dla której $\Prob(\Omega)=1$
oraz dla każdego ciągu parami rozłącznych zdarzeń $A_n\in\mathcal F$
zachodzi
\begin{equation}\label{eq:sigma_add}
\Prob\!\left(\bigcup_{n=1}^\infty A_n\right)
=\sum_{n=1}^\infty\Prob(A_n).
\end{equation}
Trójkę $(\Omega,\mathcal F,\Prob)$ nazywamy przestrzenią probabilistyczną.
\end{definicja}

\begin{twierdzenie}[Podstawowe własności miary probabilistycznej]
Niech $(\Omega,\mathcal F,\Prob)$ będzie przestrzenią probabilistyczną.
Dla wszystkich $A,B\in\mathcal F$ oraz dowolnego ciągu
$(A_n)_{n\geq1}\subseteq\mathcal F$ zachodzą:
\begin{enumerate}[label=(\roman*)]
\item $\Prob(\varnothing)=0$ i $\Prob(A^c)=1-\Prob(A)$;
\item jeśli $A\subseteq B$, to $\Prob(A)\leq\Prob(B)$ oraz
      $\Prob(B\setminus A)=\Prob(B)-\Prob(A)$;
\item $\Prob(A\cup B)=\Prob(A)+\Prob(B)-\Prob(A\cap B)$;
\item $\Prob(\bigcup_{n\geq1}A_n)\leq\sum_{n\geq1}\Prob(A_n)$.
\end{enumerate}
\end{twierdzenie}
\begin{proof}
\emph{(i)} Z przeliczalnej addytywności zastosowanej do
rozłącznego ciągu $\Omega,\varnothing,\varnothing,\ldots$
otrzymujemy
$1=\Prob(\Omega)=1+\sum_{n\geq2}\Prob(\varnothing)$.
Nieujemność każdego składnika wymusza
$\Prob(\varnothing)=0$. Następnie
$\Omega=A\mathbin{\dot\cup}A^c$ daje
$1=\Prob(A)+\Prob(A^c)$.

\emph{(ii)} Jeśli $A\subseteq B$, to
$B=A\mathbin{\dot\cup}(B\setminus A)$.
Z addytywności mamy
$\Prob(B)=\Prob(A)+\Prob(B\setminus A)$, a ostatni
składnik jest nieujemny. Jednocześnie otrzymujemy wzór
na masę różnicy.

\emph{(iii)} Zapiszmy
$A\cup B=A\mathbin{\dot\cup}(B\setminus A)$ oraz
$B=(B\setminus A)\mathbin{\dot\cup}(A\cap B)$.
Stosując skończoną addytywność do obu równości i
eliminując $\Prob(B\setminus A)$, otrzymujemy
$\Prob(A\cup B)=\Prob(A)+\Prob(B)-\Prob(A\cap B)$.

\emph{(iv)} Dla ogólnych $A_1,A_2,\ldots$ nie możemy
od razu użyć addytywności, bo zbiory mogą się przecinać.
Budujemy więc ich rozłącznienie: $B_1=A_1$ oraz
$B_n=A_n\setminus\bigcup_{j=1}^{n-1}A_j$ dla $n\geq2$.
Wtedy zbiory $B_n$ są parami rozłączne,
$\bigcup_n B_n=\bigcup_n A_n$ i $B_n\subseteq A_n$.
Zatem dzięki monotoniczności i przeliczalnej addytywności
\[
\Prob\left(\bigcup_n A_n\right)
=\sum_n\Prob(B_n)\leq\sum_n\Prob(A_n).
\]
\end{proof}

\begin{twierdzenie}[Ciągłość miary]
Niech $(\Omega,\mathcal F,\Prob)$ będzie przestrzenią probabilistyczną.
Jeśli $A_n\uparrow A$, to $\Prob(A_n)\uparrow\Prob(A)$.
Jeśli $A_n\downarrow A$, to $\Prob(A_n)\downarrow\Prob(A)$.
Tutaj $A_n\uparrow A$ oznacza $A_n\subseteq A_{n+1}$ i
$A=\bigcup_{n\geq1}A_n$, a $A_n\downarrow A$ oznacza
$A_{n+1}\subseteq A_n$ i $A=\bigcap_{n\geq1}A_n$.
\end{twierdzenie}
\begin{proof}
Załóżmy najpierw, że $A_n\uparrow A$.
Wprowadźmy rozłączne zdarzenia
$D_1=A_1$ i $D_n=A_n\setminus A_{n-1}$ dla $n\geq2$.
Ponieważ ciąg jest rosnący, $A_k$ jest dokładnie sumą
$D_1,\ldots,D_k$, a $A$ jest sumą wszystkich $D_n$.
Z addytywności
$\Prob(A_k)=\sum_{n=1}^k\Prob(D_n)$ oraz
$\Prob(A)=\sum_{n=1}^\infty\Prob(D_n)$.
Granice sum częściowych nieujemnego szeregu dają
$\lim_{k\to\infty}\Prob(A_k)=\Prob(A)$.

Jeżeli $A_n\downarrow A$, to ciąg $A_n^c$ jest rosnący,
a z praw de Morgana
$\bigcup_n A_n^c=(\bigcap_n A_n)^c=A^c$.
Pierwsza część dowodu daje
$\Prob(A_n^c)\uparrow\Prob(A^c)$. Korzystając ze wzoru
$\Prob(E^c)=1-\Prob(E)$ dla $E\in\mathcal F$,
otrzymujemy
$\Prob(A_n)=1-\Prob(A_n^c)\downarrow1-\Prob(A^c)
=\Prob(A)$.
\end{proof}

\begin{uwaga}
Warunek $\Prob(\Omega)=1$ jest istotny dla ciągłości z góry
bez dodatkowych ograniczeń: dla miary ogólnej potrzeba
$\mu(A_1)<\infty$. W przestrzeni probabilistycznej jest on
spełniony automatycznie.
\end{uwaga}

\begin{stwierdzenie}[Włączenia i wyłączenia dla prawdopodobieństw]
Dla każdej przestrzeni probabilistycznej i zdarzeń $A_1,\ldots,A_m$
zachodzi
\[
\Prob\!\left(\bigcup_{i=1}^m A_i\right)
=\sum_{\varnothing\ne I\subseteq\{1,\ldots,m\}}
(-1)^{|I|+1}\Prob\!\left(\bigcap_{i\in I}A_i\right).
\]
\end{stwierdzenie}
\begin{proof}
Dowodzimy twierdzenia indukcyjnie względem liczby zbiorów.
Dla $m=1$ jest oczywiste, a dla $m=2$ jest to
wyprowadzona wcześniej tożsamość
$\Prob(A\cup B)=\Prob(A)+\Prob(B)-\Prob(A\cap B)$.
Załóżmy, że wzór obowiązuje dla $m-1$ zbiorów.
Oznaczmy $U=\bigcup_{i=1}^{m-1}A_i$. Wtedy
\[
\Prob(U\cup A_m)
=\Prob(U)+\Prob(A_m)-\Prob(U\cap A_m).
\]
Do pierwszego składnika stosujemy założenie indukcyjne
dla $A_1,\ldots,A_{m-1}$. Do ostatniego stosujemy
je dla zbiorów $A_1\cap A_m,\ldots,A_{m-1}\cap A_m$,
ponieważ
$U\cap A_m=\bigcup_{i=1}^{m-1}(A_i\cap A_m)$.
Po odjęciu tej drugiej sumy każdy przekrój indeksowany
zbiorem $I$ bez $m$ ma znak $(-1)^{|I|+1}$,
a każdy przekrój zawierający $m$ otrzymuje znak
$-(-1)^{|I\setminus\{m\}|+1}=(-1)^{|I|+1}$.
Pozostaje jeszcze $\Prob(A_m)$ odpowiadające $I=\{m\}$.
Są to wszystkie niepuste zbiory indeksów, co kończy indukcję.
\end{proof}


\begin{przyklad}[Prawie pewne zdarzenie nie musi zachodzić zawsze]
Niech $U$ ma rozkład jednostajny na $(0,1)$.
Wtedy zdarzenie $A=\{U\neq1/2\}$ ma prawdopodobieństwo
jeden, mimo że na przestrzeni wyników $\Omega=(0,1)$
nie jest równe całemu $\Omega$. Mówimy, że $A$ zachodzi
\emph{prawie na pewno}. Z drugiej strony zdarzenie
$\{U=1/2\}$ jest logicznie możliwe w tym modelu,
lecz ma prawdopodobieństwo zero.
\end{przyklad}


\section{Dodatkowe własności przestrzeni probabilistycznej}
Często przy formułowaniu twierdzeń używa się wyrażeń
„prawie na pewno”, „poza zbiorem zerowym” i „zdarzenie
niemal niemożliwe”. Pierwsze dwa określenia mają ścisłe
znaczenie matematyczne; trzecie jest tylko nieformalnym
opisem zdarzenia o małym, lecz dodatnim prawdopodobieństwie.

\begin{definicja}
Niech $(\Omega,\mathcal F,\Prob)$ będzie przestrzenią
probabilistyczną i $A\in\mathcal F$. Mówimy, że $A$
zachodzi \emph{prawie na pewno}, gdy $\Prob(A)=1$.
Mówimy, że $A$ jest zdarzeniem \emph{zerowym}, gdy
$\Prob(A)=0$. Własność określona dla punktów $\omega\in\Omega$
zachodzi prawie na pewno, jeśli zbiór punktów, na których
nie zachodzi, należy do $\mathcal F$ i ma prawdopodobieństwo zero.
\end{definicja}

\begin{stwierdzenie}[Zbiory zerowe i prawie pewne]
Niech $(\Omega,\mathcal F,\Prob)$ będzie przestrzenią
probabilistyczną. Jeśli $N_n\in\mathcal F$ i
$\Prob(N_n)=0$ dla każdego $n\geq1$, to
$\Prob(\bigcup_{n\geq1}N_n)=0$. Jeśli $A_n\in\mathcal F$
i $\Prob(A_n)=1$ dla każdego $n\geq1$, to
$\Prob(\bigcap_{n\geq1}A_n)=1$.
\end{stwierdzenie}
\begin{proof}
Z przeliczalnej podaddytywności otrzymujemy
$0\leq\Prob(\bigcup_nN_n)\leq\sum_n\Prob(N_n)=0$,
a więc pierwsza równość zachodzi. Dla drugiej zauważmy,
że $\Prob(A_n^c)=1-\Prob(A_n)=0$. Pierwsza część
dowodu daje
$\Prob(\bigcup_n A_n^c)=0$. Z praw de Morgana
$\bigcap_n A_n=(\bigcup_n A_n^c)^c$, dlatego
$\Prob(\bigcap_n A_n)=1$.
\end{proof}

\begin{przyklad}[Nieskończony ciąg rzutów i zdarzenie zerowe]
Niech $\Omega=\{0,1\}^{\mathbb N}$ z miarą niezależnych
uczciwych rzutów monetą. Dla $m\geq1$ niech $C_m$
oznacza zdarzenie, że pierwszych $m$ wyników to same orły.
Wtedy $C_{m+1}\subseteq C_m$ i $\Prob(C_m)=2^{-m}$.
Zdarzenie $C=\bigcap_{m\geq1}C_m$ oznacza nieskończony
ciąg samych orłów. Na mocy ciągłości miary z góry
$\Prob(C)=\lim_{m\to\infty}2^{-m}=0$.
Jednak $C$ nie jest puste: zawiera jeden konkretny wynik,
mianowicie nieskończony ciąg orłów. To ważny przykład,
w którym logicznie dopuszczalna realizacja ma masę zero.
\end{przyklad}

\begin{stwierdzenie}[Nierówność Boole'a dla zdarzeń]
Niech $A_1,\ldots,A_m$ będą zdarzeniami w przestrzeni
probabilistycznej, $m\geq1$. Wtedy
$\Prob(\bigcap_{i=1}^m A_i)
\geq 1-\sum_{i=1}^m\Prob(A_i^c)$.
\end{stwierdzenie}
\begin{proof}
Ponieważ $(\bigcap_i A_i)^c=\bigcup_i A_i^c$,
z dopełnienia i przeliczalnej podaddytywności otrzymujemy
$1-\Prob(\bigcap_i A_i)
=\Prob(\bigcup_i A_i^c)\leq\sum_i\Prob(A_i^c)$.
Przenosząc składniki, otrzymujemy żądaną nierówność.
Jest ona przydatna, gdy chcemy wykazać, że wiele
pożądanych własności zachodzi jednocześnie z dużym
prawdopodobieństwem, choć nie znamy zależności między nimi.
\end{proof}

\section{Niezależność zdarzeń}
Dwa zdarzenia mogą zachodzić jednocześnie, ale informacja
o jednym z nich nie musi zmieniać prawdopodobieństwa drugiego.
Tę sytuację formalizuje niezależność. Ponieważ nie chcemy
opierać definicji wyłącznie na warunkowaniu (które wymaga dodatniego
prawdopodobieństwa warunku), jako definicję podstawową przyjmujemy
równość dotyczącą przecięcia. Dla trzech i więcej zdarzeń
niezależność par nie wystarcza: zależność może być widoczna
dopiero po jednoczesnym uwzględnieniu wszystkich zdarzeń.

\begin{definicja}
Zdarzenia $A,B\in\mathcal F$ nazywamy niezależnymi, jeśli
$\Prob(A\cap B)=\Prob(A)\Prob(B)$.
Rodzinę zdarzeń $(A_i)_{i\in I}$ nazywamy wzajemnie niezależną,
jeśli dla każdej skończonej rodziny \emph{różnych} indeksów
$i_1,\ldots,i_m\in I$, $m\geq1$, zachodzi
$\Prob(\bigcap_{j=1}^m A_{i_j})=\prod_{j=1}^m\Prob(A_{i_j})$.
Rodzinę nazywamy parami niezależną, jeśli warunek ten obowiązuje
tylko dla par różnych indeksów.
\end{definicja}

\section{Zdarzenia zachodzące nieskończenie często}
W nieskończonym ciągu prób powstaje nowe pytanie: czy sukces
pojawi się choć raz, czy pojawi się po dowolnie późnym momencie,
czy może od pewnego miejsca będzie występował zawsze? Są to
trzy różne zdarzenia. Symbol $\limsup A_n$ oznacza, że
\emph{dla każdego} momentu $m$ znajdzie się późniejszy sukces.
Pierwszy lemat Borela--Cantellego wyraża zasadę, że zbyt mała
łączna suma szans sukcesów wyklucza nieskończenie wiele ich
wystąpień prawie na pewno. Drugi pokazuje, że przy niezależności
rozbieżność tej sumy działa w przeciwnym kierunku.

\begin{definicja}
Dla ciągu zdarzeń $(A_n)_{n\geq1}$ definiujemy
\begin{align*}
\limsup_{n\to\infty}A_n
&:=\bigcap_{m=1}^\infty\bigcup_{n=m}^\infty A_n,\\
\liminf_{n\to\infty}A_n
&:=\bigcup_{m=1}^\infty\bigcap_{n=m}^\infty A_n.
\end{align*}
Pierwsze zdarzenie oznacza, że zachodzi nieskończenie wiele zdarzeń $A_n$,
a drugie, że od pewnego miejsca zachodzą wszystkie zdarzenia $A_n$.
\end{definicja}


\begin{przyklad}[Kwantyfikatory ukryte w granicach zdarzeń]
Dla ustalonego $\omega\in\Omega$ zachodzi
\begin{align*}
\omega\in\limsup_{n\to\infty}A_n
&\iff (\forall m\geq1)(\exists n\geq m):\omega\in A_n,\\
\omega\in\liminf_{n\to\infty}A_n
&\iff (\exists m\geq1)(\forall n\geq m):\omega\in A_n.
\end{align*}
Jeżeli $A_n=\Omega$ dla parzystych $n$ i $A_n=\varnothing$
dla nieparzystych $n$, to $\limsup A_n=\Omega$, a
$\liminf A_n=\varnothing$. Jeśli zaś $A_n=A$ dla każdego $n$,
to obie granice są równe $A$. Zawsze mamy
$\liminf A_n\subseteq\limsup A_n$, ponieważ zdarzenie
zachodzące od pewnego momentu zawsze zachodzi nieskończenie często.
\end{przyklad}

\begin{twierdzenie}[Pierwszy lemat Borela--Cantellego]
Niech $(A_n)_{n\geq1}$ będzie ciągiem zdarzeń w przestrzeni
probabilistycznej $(\Omega,\mathcal F,\Prob)$.
Jeżeli $\sum_{n=1}^\infty\Prob(A_n)<\infty$, to
$\Prob(\limsup_{n\to\infty}A_n)=0$.
\end{twierdzenie}
\begin{proof}
Niech $L=\limsup_{n\to\infty}A_n$.
Jeżeli $\omega\in L$, to dla każdego ustalonego $m$
istnieje $n\geq m$ takie, że $\omega\in A_n$.
Zatem $L\subseteq U_m:=\bigcup_{n=m}^\infty A_n$.
Monotoniczność miary i nierówność sumy dają
\[
0\leq\Prob(L)\leq\Prob(U_m)
\leq\sum_{n=m}^\infty\Prob(A_n)
\qquad(m\geq1).
\]
Kluczowy jest ostatni wyraz: jest on ogonem 
\emph{zbieżnego szeregu nieujemnego}, toteż zmierza
do zera, gdy $m\to\infty$. Lewa strona nie zależy od $m$,
więc musi być zerem. Niezależność zdarzeń nie została
użyta w żadnym punkcie rozumowania.
\end{proof}


\begin{lemat}[Dopełnienia zdarzeń niezależnych]
Niech $A_1,\ldots,A_m$ będą wzajemnie niezależnymi
zdarzeniami w przestrzeni probabilistycznej, gdzie $m\geq1$.
Wówczas dla dowolnego $\varepsilon_1,\ldots,\varepsilon_m
\in\{0,1\}$ zdarzenia $B_i$, określone przez
$B_i=A_i$ gdy $\varepsilon_i=1$ i $B_i=A_i^c$
gdy $\varepsilon_i=0$, spełniają
$\Prob(\bigcap_{i=1}^m B_i)=\prod_{i=1}^m\Prob(B_i)$.
\end{lemat}
\begin{proof}
Dla dwóch zdarzeń korzystamy z rozłącznego rozkładu
$A=A\cap B\mathbin{\dot\cup}A\cap B^c$:
$\Prob(A\cap B^c)=\Prob(A)-\Prob(A\cap B)
=\Prob(A)(1-\Prob(B))$.
Dla dowolnej skończonej rodziny analogicznie zapisujemy
$C\cap A_j^c=C\setminus(C\cap A_j)$, gdzie $C$ jest
przecięciem już wybranych zdarzeń i ich dopełnień.
Zaczynając od wzoru na przecięcia samych $A_i$ i
kolejno zamieniając $A_j$ na $A_j^c$, uzyskujemy
czynnik $1-\Prob(A_j)$ w miejsce $\Prob(A_j)$.
Indukcja po liczbie dokonanych zamian daje tezę.
\end{proof}

\begin{twierdzenie}[Drugi lemat Borela--Cantellego]
Niech $(A_n)_{n\geq1}$ będzie ciągiem \emph{wzajemnie niezależnych}
zdarzeń w przestrzeni probabilistycznej $(\Omega,\mathcal F,\Prob)$.
Jeżeli $\sum_{n=1}^\infty\Prob(A_n)=\infty$, to
$\Prob(\limsup_{n\to\infty}A_n)=1$.
\end{twierdzenie}
\begin{proof}
Zauważmy najpierw, że
\[
(\limsup_{n\to\infty}A_n)^c
=\bigcup_{m=1}^\infty\bigcap_{n=m}^\infty A_n^c.
\]
Dopełnienie zdarzenia „$A_n$ występuje nieskończenie wiele
razy” oznacza bowiem „od pewnego momentu żadne $A_n$
już nie wystąpi”. Wystarczy więc pokazać, że dla każdego
$m$ zdarzenie $C_m:=\bigcap_{n=m}^\infty A_n^c$
ma prawdopodobieństwo zero.

Ustalmy $m$ i rozważmy skończone przekroje
$C_{m,N}=\bigcap_{n=m}^N A_n^c$ dla $N\geq m$.
Z niezależności i poprzedniego lematu otrzymujemy
\[
\Prob(C_{m,N})=\prod_{n=m}^N\bigl(1-\Prob(A_n)\bigr)
\leq \exp\!\left(-\sum_{n=m}^N\Prob(A_n)\right).
\]
Nierówność po prawej wynika ze
$1-u\leq e^{-u}$ dla $u\in[0,1]$, którą można otrzymać
na przykład z wypukłości funkcji wykładniczej.
Ponieważ cały szereg $\sum_n\Prob(A_n)$ jest rozbieżny,
jego ogon od ustalonego $m$ również ma sumy częściowe
dążące do $+\infty$. Prawa strona dąży więc do zera.
Zbiory $C_{m,N}$ maleją do $C_m$; dzięki ciągłości
miary z góry mamy
$\Prob(C_m)=\lim_{N\to\infty}\Prob(C_{m,N})=0$.
Na koniec przeliczalna podaddytywność daje
$\Prob((\limsup A_n)^c)
\leq\sum_{m=1}^\infty\Prob(C_m)=0$.
Oznacza to $\Prob(\limsup A_n)=1$.
\end{proof}

\begin{przyklad}[Niezależne próby]
W kolejnych niezależnych próbach zdarzenie sukcesu ma w próbie $n$
prawdopodobieństwo $p_n$. Jeśli $p_n=1/n^2$ dla $n\geq1$,
sukces następuje tylko skończenie wiele razy z prawdopodobieństwem jeden.
Jeśli $p_n=1/n$, sukces następuje nieskończenie wiele razy
z prawdopodobieństwem jeden. Pierwszy wniosek nie wymaga niezależności;
drugi jej wymaga.
\end{przyklad}

\begin{przyklad}[Dlaczego niezależność jest potrzebna]
Niech $A\in\mathcal F$ spełnia $\Prob(A)=1/2$ i niech $A_n=A$
dla wszystkich $n$. Wtedy $\sum_n\Prob(A_n)=\infty$,
lecz $\limsup A_n=A$, więc $\Prob(\limsup A_n)=1/2$.
Dla zdarzeń tych brak niezależności.
\end{przyklad}

\section{Konstrukcje przestrzeni probabilistycznych}
Abstrakcyjna definicja ma sens dopiero wtedy, gdy można zbudować
modele dla rzeczywistych doświadczeń. Zaczniemy od modelu dwóch
możliwości, potem przejdziemy do przestrzeni skończonych i
przeliczalnych, losowania punktu z odcinka oraz nieskończonego
ciągu rzutów monetą. Każda konstrukcja odpowiada na trzy pytania:
co może być wynikiem, jakie zdarzenia dopuszczamy i jak przypisujemy
im prawdopodobieństwa.


\begin{przyklad}[Wypadnie albo nie wypadnie: przestrzeń o prawdopodobieństwach $0$, $1/2$, $1$]
Załóżmy, że przy rzucie uczciwą monetą interesuje nas tylko,
czy wypadnie orzeł. Oznaczmy przez $o$ wynik „orzeł”, a przez
$r$ wynik „reszka”. Rozważamy
$\Omega=\{o,r\}$ oraz całe $\sigma$-ciało
$\mathcal F=2^\Omega=\{\varnothing,\{o\},\{r\},\Omega\}$.
Przy założeniu uczciwości monety definiujemy
$\Prob(\{o\})=\Prob(\{r\})=1/2$.
W konsekwencji pełna miara jest opisana przez tabelę
\[
\begin{array}{c|cccc}
A&\varnothing&\{o\}&\{r\}&\Omega\\ \hline
\Prob(A)&0&1/2&1/2&1
\end{array}
\]
Wartości prawdopodobieństw wszystkich dopuszczalnych zdarzeń
należą więc do zbioru $\{0,1/2,1\}$. Mamy cztery zdarzenia,
ale jedynie trzy różne wartości miary: dwóm różnym zdarzeniom
może odpowiadać to samo prawdopodobieństwo.

Sprawdźmy aksjomaty bez odwoływania się do intuicji. Suma
mas punktowych wynosi jeden, stąd $\Prob(\Omega)=1$.
Dla dowolnych parami rozłącznych $A_n\subseteq\Omega$ i
każdego $\omega\in\Omega$ punkt $\omega$ należy do
co najwyżej jednego $A_n$. Zatem
\[
\Prob\!\left(\bigcup_n A_n\right)
=\sum_{\omega\in\Omega}\frac12\ind{\{\omega\in\bigcup_n A_n\}}
=\sum_n\sum_{\omega\in A_n}\frac12
=\sum_n\Prob(A_n).
\]
Zamiana sum jest dopuszczalna, gdyż w tym przypadku tylko
skończenie wiele składników może być niezerowych.

\emph{Ważne:} sam opis „wypadnie albo nie wypadnie” nie oznacza,
że obie możliwości mają po $1/2$. Dla niesymetrycznego
mechanizmu z parametrem $p\in[0,1]$ możemy ustalić
$\Prob(\{o\})=p$ i $\Prob(\{r\})=1-p$.
Wtedy zbiór możliwych wartości miary to
$\{0,p,1-p,1\}$, który przy $p=1/2$ redukuje się
do $\{0,1/2,1\}$. Dwudzielność wyniku jest własnością
przestrzeni $\Omega$, a równość szans jest założeniem o miarze.
\end{przyklad}

\begin{stwierdzenie}[Charakteryzacja miar na przestrzeni dwupunktowej]
Niech $\Omega=\{o,r\}$, gdzie $o\ne r$, i niech
$\mathcal F=2^\Omega$. Każda miara probabilistyczna
$\Prob$ na $\mathcal F$ jest jednoznacznie określona przez
liczbę $p=\Prob(\{o\})\in[0,1]$ oraz ma postać
$\Prob(A)=p\ind{\{o\in A\}}+(1-p)\ind{\{r\in A\}}$
dla $A\in\mathcal F$. Odwrotnie, każda liczba $p\in[0,1]$
zadaje tym wzorem miarę probabilistyczną na $\mathcal F$.
\end{stwierdzenie}
\begin{proof}
Punkty tworzą rozłączny rozkład
$\Omega=\{o\}\mathbin{\dot\cup}\{r\}$, więc
$1=\Prob(\Omega)=\Prob(\{o\})+\Prob(\{r\})$.
Jeżeli pierwsza masa to $p$, druga musi wynosić $1-p$.
Każde zdarzenie jest sumą zawartych w nim singletonów, zatem
jego masa ma podaną postać. Na odwrót, dla ustalonego
$p\in[0,1]$ wzór ma wartości nieujemne, daje
$\Prob(\Omega)=p+(1-p)=1$ i jest przeliczalnie
addytywny, ponieważ każdy z dwóch punktów może leżeć
w co najwyżej jednym z parami rozłącznych zdarzeń.
\end{proof}

\begin{uwaga}[Prawdopodobieństwo zero a niemożliwość]
W modelu uczciwej monety jedynym zdarzeniem o masie zero jest
$\varnothing$, a więc prawdopodobieństwo zero oznacza tutaj
niemożliwość. Nie jest to ogólne prawo. W modelu jednostajnego
losowania $U$ z $[0,1]$ zdarzenie $\{U=1/2\}$ jest niepuste,
ale ma prawdopodobieństwo zero. Zatem implikacja
$A=\varnothing\Rightarrow\Prob(A)=0$ jest zawsze prawdziwa,
natomiast implikacja odwrotna może zawieść. Podobnie zdarzenie
o prawdopodobieństwie jeden nie musi być całym $\Omega$.
\end{uwaga}

\begin{przyklad}[Model skończony i model Laplace'a]
Niech $\Omega=\{\omega_1,\ldots,\omega_N\}$,
$\mathcal F=2^\Omega$ i niech $p_i\geq0$ oraz $\sum_{i=1}^N p_i=1$.
Wzór $\Prob(A)=\sum_{\omega_i\in A}p_i$ definiuje miarę
probabilistyczną: jest unormowana i addytywna przeliczalnie,
ponieważ w skończonym $\Omega$ tylko skończenie wiele rozłącznych
niepustych zbiorów może wystąpić. Przy $p_i=1/N$ otrzymujemy
$\Prob(A)=|A|/N$.
\end{przyklad}


\begin{przyklad}[Przestrzeń z czterema wynikami: dwa rzuty monetą]
Niech $\Omega=\{oo,or,ro,rr\}$ i $\mathcal F=2^\Omega$.
Zakładamy, że cztery wyniki mają masę $1/4$.
Wówczas zdarzenie $A$ „w pierwszym rzucie orzeł” to
$\{oo,or\}$, a zdarzenie $B$ „w drugim rzucie orzeł”
to $\{oo,ro\}$. Mamy $\Prob(A)=\Prob(B)=1/2$,
ale $\Prob(A\cap B)=1/4$, co odpowiada niezależności.
Zdarzenie $C=\{oo,rr\}$ również ma masę $1/2$,
choć oznacza inny warunek: dwa jednakowe wyniki.

Jeśli teraz obserwujemy jedynie, czy wyniki są jednakowe,
to możemy przejść do dwupunktowej przestrzeni wyników
$\Omega'=\{\text{jednakowe},\text{różne}\}$.
Zdarzenia $\{\text{jednakowe}\}$ i
$\{\text{różne}\}$ mają po $1/2$. Zmieniła się przestrzeń
wyników, ale prawdopodobieństwa zdarzeń wyrażalnych
przez tę obserwację pozostały takie same.
\end{przyklad}

\begin{przyklad}[Przestrzeń z niepełną informacją]
Niech $\Omega=\{1,2,3,4\}$ i
$\mathcal F=\{\varnothing,\{1,2\},\{3,4\},\Omega\}$.
Możemy pytać, czy wynik należy do pierwszej, czy drugiej
pary, lecz zdarzenie $\{1\}$ nie jest mierzalne.
Jeśli $\Prob(\{1,2\})=1/2$ oraz
$\Prob(\{3,4\})=1/2$, to prawdopodobieństwa wszystkich
\emph{dostępnych} zdarzeń ponownie należą do
$\{0,1/2,1\}$, mimo że sama przestrzeń wyników ma cztery
punkty. To pokazuje, że liczba wartości miary jest wspólną
cechą $\mathcal F$ oraz $\Prob$, nie samego $\Omega$.
\end{przyklad}

\begin{przyklad}[Model przeliczalny]
Niech $\Omega=\{1,2,\ldots\}$, $\mathcal F=2^\Omega$ i
$\Prob(A)=\sum_{n\in A}2^{-n}$. Ponieważ $\sum_{n\geq1}2^{-n}=1$,
jest to miara probabilistyczna. W szczególności
$\Prob(\{n\geq k\})=2^{-(k-1)}$ dla $k\geq1$.
\end{przyklad}

\begin{przyklad}[Jednostajny punkt odcinka]
Niech $\Omega=[0,1]$, $\mathcal F=\mathcal B([0,1])$,
a $\Prob$ będzie ograniczeniem jednowymiarowej miary Lebesgue'a
$\lambda_1$ do $\mathcal F$. Wtedy $\Prob([a,b])=b-a$
dla $0\leq a\leq b\leq1$. Wynikami elementarnymi są wszystkie
punkty odcinka, mimo że prawdopodobieństwo każdego pojedynczego
punktu wynosi zero. Dlatego skończonego wzoru Laplace'a nie można
tu stosować.
\end{przyklad}

\begin{przyklad}[Rzut monetą nieskończenie wiele razy]
Dla ustalonego $p\in[0,1]$ rozważamy $\Omega=\{0,1\}^{\N}$.
Zdarzenia cylindryczne określają skończoną liczbę współrzędnych.
Wymagamy, aby dla różnych indeksów $i_1,\ldots,i_m$ i
$a_j\in\{0,1\}$ zachodziło
\[
\Prob\{\omega:\omega_{i_j}=a_j,\ 1\leq j\leq m\}
=\prod_{j=1}^m p^{a_j}(1-p)^{1-a_j}.
\]
Na $\sigma$-ciele generowanym przez cylindryczne zdarzenia
istnieje dokładnie jedna miara probabilistyczna o tych wartościach;
jest to przeliczalny produkt miar Bernoulliego, uzyskiwany z twierdzenia
o istnieniu miary produktowej. Jego własności skończeniewymiarowe
opisują niezależne rzuty monetą, lecz pojedynczym wynikiem jest cały
nieskończony ciąg.
\end{przyklad}

\begin{uwaga}
Aksjomaty teorii prawdopodobieństwa nie przypisują wag dowolnym
podzbiorom nieprzeliczalnej przestrzeni $\Omega$. Zdarzenia muszą
należeć do wybranego $\sigma$-ciała. Ma to znaczenie choćby dla
nieborelowskich lub niemierzalnych w sensie Lebesgue'a podzbiorów $\R$.
\end{uwaga}

\chapter{Prawdopodobieństwo klasyczne, warunkowanie i niezależność}

\section{Model Laplace'a}
Najprostszy pomysł: jeśli istnieje $N$ równie możliwych wyników,
a zdarzenie sprzyjające zawiera $k$ z nich, przypisujemy mu
szansę $k/N$. Jest to jednak \emph{warunkowe} zastosowanie
symetrii: trzeba najpierw uzasadnić, dlaczego wybrane wyniki
elementarne są równoprawdopodobne. Przy dwóch kostkach są nimi
uporządkowane pary oczek, nie sumy oczek. Wzór Laplace'a jest
konsekwencją wcześniejszej konstrukcji miary jednostajnej,
a nie dodatkowym aksjomatem.

\begin{definicja}
Modelem Laplace'a nazywamy przestrzeń probabilistyczną
$(\Omega,2^\Omega,\Prob)$, w której $\Omega$ jest niepustym zbiorem
skończonym, a $\Prob(\{\omega\})=1/|\Omega|$ dla każdego
$\omega\in\Omega$. Dla dowolnego $A\subseteq\Omega$ zachodzi wówczas
$\Prob(A)=|A|/|\Omega|$.
\end{definicja}

\begin{uwaga}
Warunek jednakowego prawdopodobieństwa wyników jest założeniem
modelu, a nie wnioskiem z samej symetrii językowego opisu.
Przykładowo podczas rzucania dwiema uczciwymi kostkami
36 uporządkowanych par ma równe prawdopodobieństwa; 11 możliwych sum
oczek nie ma rozkładu jednostajnego.
\end{uwaga}

\begin{przyklad}[Dwie kostki]
Niech $\Omega=\{1,\ldots,6\}^2$ z rozkładem jednostajnym.
Zdarzenie sumy równej siedem składa się z sześciu par,
więc ma prawdopodobieństwo $6/36=1/6$. Zdarzenie sumy równej dwa
składa się tylko z $(1,1)$, toteż jego prawdopodobieństwo wynosi $1/36$.
\end{przyklad}

\begin{przyklad}[Losowanie bez zwracania]
W urnie znajduje się $N$ różnych kul, z których $K$ jest czerwonych.
Losujemy jednocześnie $n$ kul, gdzie $0\leq n\leq N$, i zakładamy
jednostajność wszystkich $n$-elementowych podzbiorów.
Prawdopodobieństwo dokładnie $k$ czerwonych kul wynosi
\[
\frac{\binom Kk\binom{N-K}{n-k}}{\binom Nn}.
\]
W liczniku niezależnie \emph{zliczamy wybory} kul obu kolorów;
nie należy mylić tej kombinatorycznej zasady iloczynu
ze stochastyczną niezależnością kolejnych losowań.
\end{przyklad}

\section{Prawdopodobieństwo warunkowe}
Otrzymanie informacji zmienia przestrzeń rozważanych przypadków.
Jeśli wiemy, że zaszło $B$, to interesujące nas zdarzenie $A$
może zajść jedynie przez $A\cap B$. Aby pozostała pewna
możliwość miała prawdopodobieństwo jeden, normalizujemy
wagi przez $\Prob(B)$. To intuicja ilorazu
$\Prob(A\cap B)/\Prob(B)$.

Warunek $\Prob(B)>0$ nie jest formalnością: dla zdarzeń
zerowych dostalibyśmy iloraz $0/0$. Warunkowanie względem
wartości ciągłej zmiennej losowej wymaga później innej konstrukcji.

\begin{definicja}
Niech $(\Omega,\mathcal F,\Prob)$ będzie przestrzenią probabilistyczną,
a $A,B\in\mathcal F$ oraz $\Prob(B)>0$.
Prawdopodobieństwo warunkowe zdarzenia $A$ przy założeniu $B$
definiujemy wzorem
\begin{equation}\label{eq:cond}
\Prob(A\mid B):=\frac{\Prob(A\cap B)}{\Prob(B)}.
\end{equation}
\end{definicja}

\begin{stwierdzenie}
Dla każdego zdarzenia $B\in\mathcal F$ z $\Prob(B)>0$
funkcja $\Prob_B(A):=\Prob(A\mid B)$, $A\in\mathcal F$,
jest miarą probabilistyczną na $(\Omega,\mathcal F)$.
Ponadto $\Prob(B\mid B)=1$ i $\Prob(A^c\mid B)=1-\Prob(A\mid B)$.
\end{stwierdzenie}
\begin{proof}
Musimy sprawdzić wszystkie aksjomaty miary, nie tylko
nieujemność. Dla $A\in\mathcal F$ mamy
$0\leq\Prob(A\cap B)\leq\Prob(B)$, stąd
$\Prob_B(A)\in[0,1]$. Dalej
$\Prob_B(\Omega)=\Prob(\Omega\cap B)/\Prob(B)=1$.

Jeśli $A_1,A_2,\ldots$ są parami rozłączne, to zbiory
$A_n\cap B$ są również parami rozłączne, a
$B\cap\bigcup_n A_n=\bigcup_n(B\cap A_n)$.
Przeliczalna addytywność miary $\Prob$ daje
\[
\Prob_B\!\left(\bigcup_n A_n\right)
=\frac{\sum_n\Prob(A_n\cap B)}{\Prob(B)}
=\sum_n\Prob_B(A_n).
\]
Wreszcie $\Prob(B\mid B)=\Prob(B)/\Prob(B)=1$, zaś
$B=(A\cap B)\mathbin{\dot\cup}(A^c\cap B)$.
Dzieląc równość dla mas obu części przez $\Prob(B)$,
otrzymujemy $\Prob(A^c\mid B)=1-\Prob(A\mid B)$.
\end{proof}


\begin{przyklad}[Warunkowanie jako nowa miara na przestrzeni dwupunktowej]
Wróćmy do uczciwej monety
$\Omega=\{o,r\}$, $\Prob(\{o\})=\Prob(\{r\})=1/2$.
Warunkujemy względem zdarzenia $B=\{o\}$.
Wtedy $\Prob_B(\{o\})=1$ i $\Prob_B(\{r\})=0$.
Wynik doświadczenia nie zmienił się wstecznie;
zmieniła się miara, która opisuje naszą informację
po dowiedzeniu się, że zaszło $B$.
Ta prosta konstrukcja ilustruje, dlaczego rozkład warunkowy
jest nadal miarą probabilistyczną i dlaczego prawdopodobieństwa
$0$ oraz $1$ mogą pojawić się po warunkowaniu.
\end{przyklad}

\begin{przyklad}[Warunkowanie nie jest symetryczne]
W urnie znajdują się dwie kule czerwone i jedna niebieska.
Losujemy dwie kule kolejno bez zwracania. Niech $A$ oznacza,
że pierwsza jest czerwona, a $B$, że druga jest czerwona.
Z symetrii $\Prob(A)=\Prob(B)=2/3$.
Natomiast
$\Prob(A\cap B)=(2/3)(1/2)=1/3$, więc
$\Prob(B\mid A)=(1/3)/(2/3)=1/2$.
W tym szczególnym przykładzie również $\Prob(A\mid B)=1/2$,
co wynika z symetrii kolejności losowań; nie jest to jednak
ogólna reguła. Jeśli zdarzenia mają różne masy, na przykład
$\Prob(A)=1/2$, $\Prob(B)=1/4$ i
$\Prob(A\cap B)=1/8$, to
$\Prob(A\mid B)=1/2$, lecz $\Prob(B\mid A)=1/4$.
Dane te realizuje na przykład przestrzeń ośmiu
jednakowo prawdopodobnych wyników $\Omega=\{1,\ldots,8\}$,
w której $A=\{1,2,3,4\}$ oraz $B=\{1,5\}$.
\end{przyklad}

\begin{twierdzenie}[Reguła mnożenia]
Niech $A_1,\ldots,A_n\in\mathcal F$, gdzie $n\geq2$, i załóżmy,
że $\Prob(A_1\cap\cdots\cap A_{j})>0$ dla $1\leq j<n$.
Wtedy
\[
\Prob\!\left(\bigcap_{j=1}^n A_j\right)
=\Prob(A_1)\prod_{j=2}^n
\Prob\!\left(A_j\,\middle|\,\bigcap_{i=1}^{j-1}A_i\right).
\]
\end{twierdzenie}
\begin{proof}
Wprowadźmy $C_j=\bigcap_{i=1}^j A_i$ dla
$j=1,\ldots,n$. Założenia zapewniają,
że $\Prob(C_{j-1})>0$ dla $j=2,\ldots,n$;
wszystkie występujące prawdopodobieństwa warunkowe
są więc prawidłowo zdefiniowane. Ponieważ
$C_j=C_{j-1}\cap A_j$, z definicji mamy
\[
\Prob(A_j\mid C_{j-1})
=\frac{\Prob(C_j)}{\Prob(C_{j-1})}.
\]
Mnożąc te równości dla $j=2,\ldots,n$ i dopisując
$\Prob(A_1)=\Prob(C_1)$, otrzymujemy iloczyn
\[
\Prob(C_1)\frac{\Prob(C_2)}{\Prob(C_1)}
\frac{\Prob(C_3)}{\Prob(C_2)}\cdots
\frac{\Prob(C_n)}{\Prob(C_{n-1})}=\Prob(C_n).
\]
To właśnie prawdopodobieństwo przecięcia wszystkich $A_j$.
Na żadnym etapie nie założyliśmy niezależności:
każdy czynnik odpowiada sytuacji po uwzględnieniu
poprzednich zdarzeń.
\end{proof}

\begin{przyklad}[Bez zwracania i ze zwracaniem]
W urnie znajduje się pięć czerwonych i trzy niebieskie kule.
Przy dwóch losowaniach bez zwracania prawdopodobieństwo dwóch
czerwonych kul wynosi $(5/8)(4/7)=5/14$.
Przy losowaniu ze zwracaniem wynosi $(5/8)^2=25/64$.
Różnica bierze się stąd, że po pierwszym losowaniu bez zwracania
zmienia się skład urny, a więc rozkład wyniku drugiego losowania.
\end{przyklad}

\section{Niezależność zdarzeń}
W kontekście warunkowania można teraz powiedzieć precyzyjniej:
przy $\Prob(B)>0$ zdarzenie $B$ nie wnosi informacji o $A$
wtedy i tylko wtedy, gdy $\Prob(A\mid B)=\Prob(A)$.
Przykład trzech zdarzeń zależnych wyłącznie jako trójka
wyjaśni, dlaczego żądanie niezależności wszystkich skończonych
przekrojów jest silniejsze od niezależności parami.

Niezależność zdarzeń została zdefiniowana przed sformułowaniem
lematów Borela--Cantellego. Rozróżnienie niezależności wzajemnej
i parami niezależnej ma znaczenie również dla skończonych rodzin.


\begin{stwierdzenie}
Niech $A,B\in\mathcal F$ i $\Prob(B)>0$.
Wówczas $A$ i $B$ są niezależne wtedy i tylko wtedy, gdy
$\Prob(A\mid B)=\Prob(A)$. Jeżeli $A$ i $B$ są niezależne,
to niezależne są także pary $A^c,B$ oraz $A^c,B^c$.
\end{stwierdzenie}
\begin{proof}
Pierwsza równoważność wynika przez podstawienie do \eqref{eq:cond}.
Ponadto $\Prob(A^c\cap B)=\Prob(B)-\Prob(A\cap B)
=(1-\Prob(A))\Prob(B)$, a więc $A^c$ i $B$ są niezależne.
Powtarzamy rozumowanie po zamianie $B$ na $B^c$.
\end{proof}

\begin{przyklad}[Niezależność parami nie wystarcza]
Przeprowadzamy dwa niezależne rzuty uczciwą monetą.
W przestrzeni $\Omega=\{0,1\}^2$ z prawdopodobieństwem jednostajnym
rozważamy $A=\{(1,0),(1,1)\}$,
$B=\{(0,1),(1,1)\}$ i $C=\{(0,0),(1,1)\}$.
Każde zdarzenie ma prawdopodobieństwo $1/2$, a każde przecięcie
dwóch różnych z nich ma prawdopodobieństwo $1/4$.
Są zatem parami niezależne. Jednak $\Prob(A\cap B\cap C)=1/4$,
podczas gdy $\Prob(A)\Prob(B)\Prob(C)=1/8$.
\end{przyklad}

\begin{uwaga}
Zdarzenia rozłączne o dodatnich prawdopodobieństwach \emph{nie są}
niezależne: ich przecięcie ma prawdopodobieństwo zero,
a iloczyn ich prawdopodobieństw jest dodatni.
\end{uwaga}

\section{Prawdopodobieństwo całkowite i wzór Bayesa}
Jeżeli doświadczenie przebiega dwustopniowo, na przykład najpierw
wybieramy urnę, a potem kulę, możemy obliczyć szansę zdarzenia
końcowego osobno dla każdego sposobu rozpoczęcia doświadczenia.
Wzór na prawdopodobieństwo całkowite jest sumowaniem tych
rozłącznych dróg. Wzór Bayesa jest natomiast czysto algebraicznym
odwróceniem warunku: znając szansę obserwacji przy poszczególnych
scenariuszach, obliczamy szansę scenariusza pod warunkiem obserwacji.
Nie wprowadzamy tutaj jeszcze odrębnego aparatu wnioskowania
bayesowskiego.

\begin{twierdzenie}[Wzór na prawdopodobieństwo całkowite]
Niech $(B_i)_{i\in I}$ będzie skończoną lub przeliczalną rodziną
parami rozłącznych zdarzeń z $\Prob(B_i)>0$ dla każdego $i\in I$
oraz $\Prob(\bigcup_{i\in I}B_i)=1$.
Wtedy dla każdego $A\in\mathcal F$
\begin{equation}\label{eq:totalprob}
\Prob(A)=\sum_{i\in I}\Prob(A\mid B_i)\Prob(B_i).
\end{equation}
\end{twierdzenie}
\begin{proof}
Niech $B=\bigcup_{i\in I}B_i$.
Z założenia $\Prob(B)=1$, a więc $\Prob(B^c)=0$.
Rozkładamy $A$ na rozłączne części
$A=(A\cap B)\mathbin{\dot\cup}(A\cap B^c)$.
Ponieważ $A\cap B^c\subseteq B^c$, monotoniczność
daje $\Prob(A\cap B^c)=0$ i
$\Prob(A)=\Prob(A\cap B)$.

W dalszym kroku korzystamy z rozłączności $B_i$:
$A\cap B=\bigcup_{i\in I}(A\cap B_i)$,
a zbiory $A\cap B_i$ są parami rozłączne. Zatem
\[
\Prob(A)=\sum_{i\in I}\Prob(A\cap B_i)
=\sum_{i\in I}\Prob(A\mid B_i)\Prob(B_i).
\]
Ostatnia równość to definicja prawdopodobieństwa
warunkowego zastosowana osobno do każdego $B_i$.
Jeśli $I$ jest nieskończony, użyliśmy przeliczalnej,
a nie tylko skończonej addytywności.
\end{proof}

\begin{twierdzenie}[Wzór Bayesa]
Niech $(B_i)_{i\in I}$ spełnia warunki poprzedniego twierdzenia,
a $A\in\mathcal F$ ma dodatnie prawdopodobieństwo.
Wtedy dla każdego $j\in I$
\begin{equation}\label{eq:bayes}
\Prob(B_j\mid A)
=\frac{\Prob(A\mid B_j)\Prob(B_j)}
{\sum_{i\in I}\Prob(A\mid B_i)\Prob(B_i)}.
\end{equation}
\end{twierdzenie}
\begin{proof}
Z definicji warunkowania, ponieważ $\Prob(A)>0$,
\[
\Prob(B_j\mid A)=\frac{\Prob(B_j\cap A)}{\Prob(A)}.
\]
Licznik możemy zapisać jako
$\Prob(A\mid B_j)\Prob(B_j)$, gdyż $\Prob(B_j)>0$.
Natomiast mianownik jest sumą wkładów rozłącznych
scenariuszy $B_i$ z poprzedniego twierdzenia:
$\Prob(A)=\sum_{i\in I}\Prob(A\mid B_i)\Prob(B_i)$.
Po podstawieniu obu równości otrzymujemy tezę.
Nie ma tu nowego aksjomatu ani nowej interpretacji
prawdopodobieństwa: wzór wynika bezpośrednio z
warunkowania i rozbicia przestrzeni na rozłączne przypadki.
\end{proof}

\begin{przyklad}[Test diagnostyczny]
Niech $D$ oznacza występowanie pewnej choroby, a $T$ dodatni wynik
testu. Przyjmijmy $\Prob(D)=0{,}01$,
$\Prob(T\mid D)=0{,}99$ i $\Prob(T\mid D^c)=0{,}05$.
Wtedy
\[
\Prob(D\mid T)
=\frac{0{,}99\cdot0{,}01}
{0{,}99\cdot0{,}01+0{,}05\cdot0{,}99}
=\frac16.
\]
Dodatni wynik nie oznacza więc prawdopodobieństwa choroby równego
$99\%$. Wartość predykcyjna wyniku zależy zarówno od jakości testu,
jak i częstości choroby w rozpatrywanej populacji.
\end{przyklad}

\begin{przyklad}[Trzy drzwi --- problem Monty'ego Halla]
Załóżmy, że nagroda znajduje się jednostajnie za jednymi z trzech drzwi.
Gracz wybiera drzwi, nie znając położenia nagrody. Prowadzący zna jej
położenie, zawsze otwiera inne drzwi, za którymi nie ma nagrody,
a gdy ma wybór, rozstrzyga go losowo i symetrycznie. Gracz może
pozostać przy swoim wyborze albo zamienić go na jedne pozostałe
zamknięte drzwi. Początkowy wybór był prawidłowy z prawdopodobieństwem
$1/3$; wtedy zamiana przegrywa. Był nieprawidłowy z prawdopodobieństwem
$2/3$; wtedy prowadzący musi wskazać drugie puste drzwi, a zamiana
wygrywa. Zatem strategia zamiany daje prawdopodobieństwo wygranej $2/3$.
Warunki dotyczące zachowania prowadzącego są częścią modelu;
po ich zmianie wynik może być inny.
\end{przyklad}

\begin{przyklad}[Wybór urny po obserwacji koloru]
Najpierw losujemy urnę: pierwszą z prawdopodobieństwem $1/3$,
drugą z prawdopodobieństwem $2/3$. Pierwsza zawiera trzy kule
czerwone i jedną niebieską, druga jedną czerwoną i trzy niebieskie.
Następnie losujemy jednostajnie jedną kulę z wybranej urny.
Niech $U_1$ oznacza wybranie pierwszej urny, a $R$ wylosowanie
czerwonej kuli. Wtedy $\Prob(R\mid U_1)=3/4$ oraz
$\Prob(R\mid U_1^c)=1/4$. Po zobaczeniu czerwonej kuli
otrzymujemy ze wzoru Bayesa
\[
\Prob(U_1\mid R)
=\frac{(1/3)(3/4)}{(1/3)(3/4)+(2/3)(1/4)}
=\frac35.
\]
Przed obserwacją pierwsza urna była mniej prawdopodobna;
po obserwacji czerwonej kuli prawdopodobieństwo jej wyboru wzrasta.
Jest to wyłącznie rachunek prawdopodobieństwa warunkowego.
\end{przyklad}

\section{Warunkowanie a pozorna zależność}
Zależności widoczne w poszczególnych podgrupach mogą zniknąć
albo nawet zmienić znak po złączeniu obserwacji. Wynika to z tego,
że średnie warunkowe ważymy różnymi prawdopodobieństwami podgrup.
Tak zwany paradoks Simpsona nie jest paradoksem logicznym:
wskazuje, że przy porównywaniu danych trzeba jawnie określić,
na jakiej populacji warunkujemy.

\begin{przyklad}[Paradoks Simpsona]
W dwóch grupach porównujemy skuteczność metod $A$ i $B$.
W grupie I metoda $A$ daje $9$ sukcesów na $10$ prób, a metoda $B$
$80$ sukcesów na $100$ prób; w grupie II metoda $A$ daje
$20$ sukcesów na $100$ prób, a metoda $B$ $1$ sukces na $10$ prób.
W każdej grupie metoda $A$ ma większą częstość sukcesu:
$9/10>80/100$ oraz $20/100>1/10$.
Po połączeniu grup metoda $A$ ma jednak $29/110$ sukcesów,
a metoda $B$ $81/110$, więc kierunek porównania się odwraca.
Zmieniły się proporcje obu grup w populacjach stosujących dane metody.
Nie jest to sprzeczność, lecz ostrzeżenie przed utożsamianiem
prawdopodobieństw warunkowych i bezwarunkowych.
\end{przyklad}

\chapter{Zmienne losowe i rozkłady prawdopodobieństwa}

\section{Zmienna losowa i jej rozkład}
Przestrzeń wyników może być bardzo złożona: pojedynczym wynikiem
może być pięć kart, pełna trajektoria urządzenia albo nieskończony
ciąg rzutów monetą. W praktyce interesuje nas często tylko jedna
liczba, na przykład liczba czerwonych kart lub czas oczekiwania.
Zmienna losowa jest więc funkcją, która tłumaczy wyniki doświadczenia
na wartości liczbowe. Jej rozkład opisuje, z jakimi
prawdopodobieństwami pojawiają się zbiory tych wartości.

Mierzalność zapewnia, że zdanie „$X$ przyjmuje wartość należącą
do borelowskiego zbioru $B$” jest dopuszczalnym zdarzeniem.
Nie jest to warunek ciągłości ani losowości samej funkcji.

\begin{definicja}
Niech $(\Omega,\mathcal F,\Prob)$ będzie przestrzenią probabilistyczną.
Zmienną losową rzeczywistą nazywamy funkcję
$X:\Omega\to\R$, dla której $X^{-1}(B)\in\mathcal F$
dla każdego $B\in\mathcal B(\R)$. Jej rozkładem nazywamy miarę
probabilistyczną $\mu_X$ na $\mathcal B(\R)$, zadaną wzorem
$\mu_X(B):=\Prob(X\in B)=\Prob(X^{-1}(B))$.
Stosujemy oznaczenie $\Law(X)=\mu_X$ oraz $X\sim\mu_X$.
\end{definicja}

\begin{stwierdzenie}[Kryterium mierzalności]
Funkcja $X:\Omega\to\R$ jest zmienną losową wtedy i tylko wtedy,
gdy $\{X\leq x\}\in\mathcal F$ dla każdego $x\in\R$.
\end{stwierdzenie}
\begin{proof}
Konieczność wynika z borelowskości przedziału $(-\infty,x]$.
Dla dostateczności zauważamy, że przedziały $(-\infty,x]$ generują
$\mathcal B(\R)$. Rodzina $\{B\subseteq\R:X^{-1}(B)\in\mathcal F\}$
jest $\sigma$-ciałem, gdyż przeciwobrazy zachowują dopełnienia
i sumy przeliczalne, a więc zawiera całe $\mathcal B(\R)$.
\end{proof}

\begin{przyklad}
W modelu dwóch kostek funkcja $X(i,j)=i+j$ jest zmienną losową,
której rozkład skupia się na $\{2,\ldots,12\}$.
Na $([0,1],\mathcal B([0,1]),\lambda_1)$ zmienna $U(\omega)=\omega$
ma rozkład jednostajny na $[0,1]$. Samo $\Omega$ nie przesądza więc
o postaci rozkładu: ważna jest również zmienna losowa.
\end{przyklad}

\begin{definicja}
Niech $\mu$ będzie borelowską miarą probabilistyczną na $\R$.
Nośnikiem topologicznym $\mu$ nazywamy zbiór
$\supp\mu:=\{x\in\R:\mu((x-\varepsilon,x+\varepsilon))>0
\text{ dla każdego }\varepsilon>0\}$.
\end{definicja}

\begin{uwaga}
Rozkład $\mu_X$ zawiera pełną informację o prawdopodobieństwach
zdarzeń wyrażalnych przez wartości $X$, ale nie o zależności $X$
od innych zmiennych. Dwie zmienne mogą mieć identyczne rozkłady,
choć jedna jest równa drugiej, a inna od niej niezależna.
\end{uwaga}

\section{Dystrybuanta}
Rozkład można opisywać przez prawdopodobieństwo każdego zbioru
borelowskiego, lecz w wymiarze jednym wygodniej zebrać
informację w jednej funkcji $F(x)=\Prob(X\leq x)$.
Przechowuje ona prawdopodobieństwa wszystkich lewych półprostych,
a z ich różnic odzyskujemy masy przedziałów. Skok dystrybuanty
w punkcie oznacza dodatnie prawdopodobieństwo dokładnie tej wartości.
Prawostronna ciągłość nie jest dodatkową umową: wynika z
przeliczalnej addytywności miary.

\begin{definicja}
Dystrybuantą rzeczywistej zmiennej losowej $X$ nazywamy funkcję
$F_X:\R\to[0,1]$ określoną wzorem
$F_X(x):=\Prob(X\leq x)=\mu_X((\! -\infty,x])$.
\end{definicja}

\begin{twierdzenie}[Własności dystrybuanty]
Niech $X$ będzie rzeczywistą zmienną losową o dystrybuancie $F$.
Wówczas $F$ jest niemalejąca, prawostronnie ciągła oraz
$\lim_{x\to-\infty}F(x)=0$ i $\lim_{x\to+\infty}F(x)=1$.
Dla dowolnych $a<b$ zachodzą ponadto
\[
\Prob(a<X\leq b)=F(b)-F(a),\qquad
\Prob(X=a)=F(a)-F(a-),
\]
gdzie $F(a-):=\lim_{x\uparrow a}F(x)$.
\end{twierdzenie}
\begin{proof}
Jeśli $x\leq y$, to każdy wynik spełniający $X\leq x$
spełnia też $X\leq y$. Zatem
$\{X\leq x\}\subseteq\{X\leq y\}$ i monotoniczność
miary daje $F(x)\leq F(y)$.

Ustalmy $x\in\R$ i weźmy dowolny ciąg $x_n\downarrow x$.
Zdarzenia $A_n=\{X\leq x_n\}$ maleją. Ich przekrój
to $\{X\leq x\}$: jeśli $X(\omega)>x$, to dla
dostatecznie dużego $n$ mamy $x_n<X(\omega)$,
więc $\omega\notin A_n$. Z ciągłości miary z góry
wynika $F(x_n)\to F(x)$. Jest to prawostronna ciągłość.

Dla $x_n=-n$ zbiory $\{X\leq-n\}$ maleją
do zbioru pustego, ponieważ $X$ przyjmuje tylko
wartości rzeczywiste. Ich prawdopodobieństwa dążą do zera.
Podobnie zbiory $\{X\leq n\}$ rosną do $\Omega$,
więc ich prawdopodobieństwa dążą do jedności.
Monotoniczność $F$ pozwala otrzymane granice dla
ciągów $\pm n$ przenieść na całe półproste.

Dla $a<b$ mamy rozłączny rozkład
$\{X\leq b\}=\{X\leq a\}\mathbin{\dot\cup}
\{a<X\leq b\}$, więc
$\Prob(a<X\leq b)=F(b)-F(a)$.
Wreszcie dla dowolnego ciągu $x_n\uparrow a$
zachodzi $\{X\leq x_n\}\uparrow\{X<a\}$.
Z ciągłości miary z dołu dostajemy
$F(a-)=\Prob(X<a)$. Ponieważ
$\{X\leq a\}=\{X<a\}\mathbin{\dot\cup}\{X=a\}$,
otrzymujemy $\Prob(X=a)=F(a)-F(a-)$.
\end{proof}

\begin{wniosek}
Dystrybuanta jest ciągła w punkcie $a$ wtedy i tylko wtedy,
gdy $\Prob(X=a)=0$. Zmiana dystrybuanty w pojedynczym punkcie
może zniszczyć jej prawostronną ciągłość i dlatego nie można
modyfikować dystrybuanty dowolnie na zbiorze miary zero.
\end{wniosek}

\begin{twierdzenie}[Charakteryzacja dystrybuant]
Funkcja $F:\R\to\R$ jest dystrybuantą pewnej rzeczywistej
zmiennej losowej wtedy i tylko wtedy, gdy jest niemalejąca,
prawostronnie ciągła, ma granicę $0$ przy $-\infty$
i granicę $1$ przy $+\infty$.
\end{twierdzenie}
\begin{proof}
Konieczność wszystkich warunków wynika z poprzedniego
twierdzenia. Pokażemy, że są one także wystarczające,
konstruując zmienną o zadanej dystrybuancie.
Niech $\Omega=(0,1)$, $\mathcal F=\mathcal B((0,1))$,
a $\Prob$ będzie miarą Lebesgue'a.
Dla $u\in(0,1)$ zdefiniujmy
$Q(u)=\inf\{x\in\R:F(x)\geq u\}$.
Z granic $F$ wynika, że dla każdego $u\in(0,1)$
istnieją $a,b\in\R$ z $F(a)<u\leq F(b)$.
Zbiór pod infimum jest więc niepusty i zawiera się
w $(a,\infty)$; liczba $Q(u)$ jest skończona.

Sprawdźmy kluczową równoważność. Jeżeli $u\leq F(x)$,
to $x$ należy do zbioru w definicji i $Q(u)\leq x$.
Odwrotnie, jeśli $Q(u)\leq x$, to z definicji infimum
istnieje ciąg $x_n\downarrow Q(u)$ taki, że
$F(x_n)\geq u$. Prawostronna ciągłość daje
$F(Q(u))\geq u$, a monotoniczność implikuje
$F(x)\geq F(Q(u))\geq u$.
Otrzymaliśmy zatem dla wszystkich $x\in\R$, $u\in(0,1)$
\[
Q(u)\leq x\quad\Longleftrightarrow\quad u\leq F(x).
\]
Funkcja $Q$ jest niemalejąca (większe $u$ wymaga
większego kwantyla), a więc mierzalna.
Dla $U(\omega)=\omega$ mamy
$\Prob(Q(U)\leq x)=\Prob(U\leq F(x))=F(x)$,
ponieważ $U$ jest jednostajna na $(0,1)$, a $0\leq F(x)\leq1$.
Zmienna $X=Q(U)$ realizuje więc żądaną dystrybuantę.
\end{proof}

\begin{definicja}
Dla zmiennej losowej $X$ o dystrybuancie $F_X$ i $u\in(0,1)$
kwantylem rzędu $u$ nazywamy
$Q_X(u):=\inf\{x\in\R:F_X(x)\geq u\}$.
Medianą nazywamy dowolną liczbę $m$ spełniającą
$\Prob(X\leq m)\geq1/2$ i $\Prob(X\geq m)\geq1/2$.
\end{definicja}

\begin{uwaga}
Kwantyl $Q_X(1/2)$ jest medianą, ale mediana nie musi być jedyna.
Przykładowo zmienna przyjmująca $0$ i $2$ z prawdopodobieństwami
po $1/2$ ma za medianę każdą liczbę z przedziału $[0,2]$.
\end{uwaga}

\section{Rozkłady dyskretne, ciągłe i singularne}
W przypadku dyskretnym przypisujemy masy poszczególnym punktom.
W przypadku rozkładu mającego gęstość zamiast mas punktowych
opisujemy prawdopodobieństwa \emph{przedziałów} przez całkę.
Wartość $f(x)$ nie jest prawdopodobieństwem, że $X=x$:
może przekraczać jeden, podczas gdy $\Prob(X=x)=0$.
Gęstość ma znaczenie lokalne: $f(x)h$ przy małym $h$
przybliża masę krótkiego odcinka tam, gdzie $f$ jest regularna.
Nie każda ciągła dystrybuanta posiada gęstość. Rozkład Cantora
wykaże, dlaczego to rozróżnienie jest ważne.

\begin{definicja}
Rozkład $\mu$ na $\R$ nazywamy dyskretnym, jeśli istnieje zbiór
co najwyżej przeliczalny $D\subseteq\R$ taki, że $\mu(D)=1$.
Dla $x\in\R$ piszemy wtedy $p(x):=\mu(\{x\})$, a funkcję $p$
nazywamy funkcją masy prawdopodobieństwa.
\end{definicja}

Jeżeli $\mu$ jest dyskretna i $D$ jest zbiorem jej atomów o dodatnich
masach, to $\sum_{x\in D}p(x)=1$ i
$\mu(B)=\sum_{x\in B\cap D}p(x)$ dla każdego $B\in\mathcal B(\R)$.
Z kolei dowolne nieujemne masy o sumie jeden definiują w ten sposób
rozkład dyskretny.

\begin{definicja}
Borelowską miarę probabilistyczną $\mu$ na $\R$ nazywamy
absolutnie ciągłą względem miary Lebesgue'a $\lambda_1$,
gdy $\lambda_1(B)=0$ implikuje $\mu(B)=0$ dla każdego
$B\in\mathcal B(\R)$. Wówczas na mocy twierdzenia Radona--Nikodyma
istnieje mierzalna funkcja $f:\R\to[0,\infty]$ taka, że
\begin{equation}\label{eq:density}
\mu(B)=\int_B f(x)\,\dd\lambda_1(x)
\quad(B\in\mathcal B(\R)),\qquad
\int_\R f(x)\,\dd\lambda_1(x)=1.
\end{equation}
Funkcję $f$ nazywamy gęstością prawdopodobieństwa; jest ona
wyznaczona jednoznacznie z dokładnością do równości
$\lambda_1$-prawie wszędzie.
\end{definicja}

\begin{stwierdzenie}
Jeżeli $X$ ma gęstość $f$, to
$F_X(x)=\int_{(-\infty,x]} f(t)\,\dd\lambda_1(t)$.
Dystrybuanta $F_X$ jest ciągła, a $F_X'(x)=f(x)$
dla $\lambda_1$-prawie każdego $x$.
\end{stwierdzenie}
\begin{proof}
Z definicji gęstości otrzymujemy bezpośrednio
\[
F_X(x)=\mu_X((\!-\infty,x])
=\int_{(-\infty,x]}f(t)\,\dd\lambda_1(t).
\]
Pojedynczy punkt $\{x\}$ ma miarę Lebesgue'a zero,
więc z absolutnej ciągłości rozkładu
$\Prob(X=x)=\mu_X(\{x\})=0$.
Z twierdzenia o skokach dystrybuanty wiemy,
że $F_X(x)-F_X(x-)=\Prob(X=x)=0$.
Ponieważ dystrybuanta jest zawsze prawostronnie ciągła,
nie ma również nieciągłości z lewej strony; jest więc ciągła.

Funkcja $f$ jest całkowalna na $\R$, a więc lokalnie
całkowalna. Dla każdego $a\in\R$ możemy zapisać
$F_X(x)=F_X(a)+\int_a^x f(t)\,\dd\lambda_1(t)$,
przyjmując odpowiedni znak całki dla $x<a$.
Twierdzenie o różniczkowaniu całek Lebesgue'a
(fakt z analizy, którego tutaj nie dowodzimy) zapewnia,
że dla $\lambda_1$-prawie każdego $x$ pochodna ostatniej
całki jest równa $f(x)$. Stąd $F_X'(x)=f(x)$ prawie wszędzie.
Nie wolno jednak odwracać tego wniosku bez dodatkowej
absolutnej ciągłości: sama pochodna istniejąca prawie
wszędzie nie odzyskuje rozkładu singularnego.
\end{proof}

\begin{uwaga}
Sama ciągłość dystrybuanty nie oznacza istnienia gęstości.
Istnieją \emph{rozkłady singularne ciągłe}, na przykład rozkład
Cantora: jest skupiony na zbiorze Cantora o mierze Lebesgue'a zero,
a jednocześnie nie ma atomów. Jego dystrybuanta jest ciągła,
ale nie jest absolutnie ciągła. Ogólna miara probabilistyczna na $\R$
może zawierać jednocześnie część dyskretną, absolutnie ciągłą
i singularną ciągłą.
\end{uwaga}


\begin{twierdzenie}[Rozkład miary na część dyskretną, absolutnie ciągłą i singularną ciągłą]
Niech $\mu$ będzie borelowską miarą probabilistyczną na $\R$.
Wówczas istnieją jednoznaczne miary skończone
$\mu_{\mathrm d},\mu_{\mathrm{ac}},\mu_{\mathrm{sc}}$
takie, że $\mu=\mu_{\mathrm d}+\mu_{\mathrm{ac}}+
\mu_{\mathrm{sc}}$, gdzie $\mu_{\mathrm d}$ jest skupiona
na zbiorze co najwyżej przeliczalnym,
$\mu_{\mathrm{ac}}\ll\lambda_1$,
a $\mu_{\mathrm{sc}}\perp\lambda_1$ i nie ma atomów.
Każda z trzech miar może mieć masę mniejszą od jedności,
a nawet być miarą zerową.
\end{twierdzenie}
\begin{proof}
Najpierw zbiór atomów punktowych
$D=\{x\in\R:\mu(\{x\})>0\}$ jest co najwyżej
przeliczalny: dla każdego $m\geq1$ istnieje najwyżej
$m$ punktów o masie większej niż $1/m$, a
$D$ jest sumą tych skończonych zbiorów.
Definiujemy $\mu_{\mathrm d}(B)=\mu(B\cap D)$
oraz $\rho(B)=\mu(B\setminus D)$. Miara $\rho$
nie ma atomów punktowych.

Wyprowadzimy teraz potrzebny rozkład na część absolutnie
ciągłą i singularną z przyjętego twierdzenia
Radona--Nikodyma. Połóżmy $\nu=\rho+\lambda_1$.
Miara $\nu$ jest $\sigma$-skończona, ponieważ
$\rho$ jest skończona, a $\lambda_1$ jest
$\sigma$-skończona. Istnieją więc funkcje
$h=\dd\rho/\dd\nu$ i $g=\dd\lambda_1/\dd\nu$.
Dla $S=\{x:g(x)=0\}$ mamy $\lambda_1(S)=0$.
Definiujemy $\mu_{\mathrm{sc}}=\ind{S}\rho$ oraz
$\mu_{\mathrm{ac}}=\ind{\R\setminus S}\rho$.
Pierwsza miara jest singularna, bo skupia się na $S$.
Druga jest absolutnie ciągła względem $\lambda_1$:
jeśli $\lambda_1(B)=0$, to
$\int_B g\,\dd\nu=0$, a więc $g=0$
$\nu$-prawie wszędzie na $B$. Stąd
$\nu(B\cap\{g>0\})=0$ i
$\mu_{\mathrm{ac}}(B)=0$.
Obie części dziedziczą brak atomów od $\rho$.

Dla jednoznaczności zauważmy najpierw, że cała
masa punktowa musi należeć do $\mu_{\mathrm d}$,
gdyż dwie pozostałe części nie mają atomów.
Przypuśćmy następnie, że
$\rho=\alpha_1+\beta_1=\alpha_2+\beta_2$,
gdzie $\alpha_i\ll\lambda_1$, a $\beta_i$
są skupione na zbiorach $S_i$ o miarze Lebesgue'a zero.
Dla $S=S_1\cup S_2$ obie miary $\beta_i$
zanikają poza $S$, więc $\alpha_1$ i $\alpha_2$
są tam równe $\rho$. Na samym $S$ obie zanikają
z absolutnej ciągłości względem $\lambda_1$.
Zatem $\alpha_1=\alpha_2$ oraz
$\beta_1=\beta_2$. To kończy dowód jednoznaczności.
\end{proof}

\begin{przyklad}[Gęstość większa niż jeden i zerowa masa punktu]
Niech $X$ będzie jednostajna na $(0,1/4)$. Jej gęstość
to $f_X(x)=4\ind{\{0<x<1/4\}}$. Mamy $f_X(1/8)=4$,
co nie przeczy aksjomatom, ponieważ gęstość nie jest
prawdopodobieństwem zdarzenia $\{X=1/8\}$.
W rzeczywistości $\Prob(X=1/8)=0$, zaś
$\Prob(0<X<1/8)=\int_0^{1/8}4\,\dd\lambda_1(x)=1/2$.
Parametr gęstości ma jednostkę odwrotną do jednostki $x$;
liczbą z przedziału $[0,1]$ jest całka po zdarzeniu.
\end{przyklad}

\begin{przyklad}[Rozkład Cantora i brak gęstości]
Niech $(B_n)_{n\geq1}$ będą niezależne i
$\Prob(B_n=0)=\Prob(B_n=1)=1/2$. Szereg
$X=\sum_{n=1}^\infty 2B_n/3^n$ zbiega bezwzględnie,
a jego wartość należy do zbioru Cantora $C\subset[0,1]$.
Ponieważ $\lambda_1(C)=0$ oraz $\Prob(X\in C)=1$,
rozkład $X$ nie może mieć gęstości względem $\lambda_1$.

Pokażemy również, że nie ma atomów. Każde $m$ pierwszych
cyfr trójkowych $2B_1,\ldots,2B_m$ wyznacza przedział
konstrukcyjny długości $3^{-m}$ o masie $2^{-m}$.
Wartość $x\in C$ należy na każdym etapie do co najwyżej
dwóch takich domkniętych przedziałów (ewentualnie styka się
z ich końcami). Zatem $\Prob(X=x)\leq2\cdot2^{-m}$
dla każdego $m$, co po przejściu do granicy daje zero.
Dystrybuanta $X$ jest więc ciągła, choć jej rozkład
jest osobliwy względem miary Lebesgue'a.
\end{przyklad}

\begin{przyklad}[Rozkład mieszany]
Niech $B$ przyjmuje wartości $0$ i $1$ z prawdopodobieństwami po $1/2$,
a $U$ jest od niego niezależna i jednostajna na $(0,1)$.
Definiujemy $X=BU$. Wtedy $\Prob(X=0)=1/2$, a dla $0<x<1$
zachodzi $F_X(x)=1/2+x/2$. Łącznie
\[
F_X(x)=\begin{cases}
0,&x<0,\\
\frac12,&x=0,\\
\frac{1+x}{2},&0<x<1,\\
1,&x\geq1.
\end{cases}
\]
Rozkład nie ma gęstości względem $\lambda_1$, gdyż ma atom w zerze,
ale jego część na $(0,1)$ ma gęstość $1/2$.
\end{przyklad}

\section{Wektory losowe i rozkłady łączne}
Znajomość rozkładu wzrostu i rozkładu masy ciała osobno nie
mówi, czy wzrost i masa są ze sobą związane. Potrzebny jest
\emph{rozkład łączny} wektora dwóch cech. Rozkłady brzegowe
powstają przez „zapomnienie” o jednej współrzędnej, czyli
całkowanie gęstości łącznej po tej współrzędnej. Odwrotna
rekonstrukcja rozkładu łącznego z samych rozkładów brzegowych
na ogół nie jest możliwa; niezależność stanowi szczególny
przypadek, gdy taka rekonstrukcja jest możliwa przez iloczyn.

\begin{definicja}
Niech $d\geq1$. Wektorem losowym w $\R^d$ nazywamy funkcję
$X=(X_1,\ldots,X_d):\Omega\to\R^d$ mierzalną względem
$\mathcal F$ i $\mathcal B(\R^d)$. Jej rozkładem łącznym jest
$\mu_X(B)=\Prob(X\in B)$ dla $B\in\mathcal B(\R^d)$.
Rozkład $X_i$ nazywamy $i$-tym rozkładem brzegowym.
\end{definicja}

\begin{stwierdzenie}
Wektor $X=(X_1,\ldots,X_d)$ jest mierzalny wtedy i tylko wtedy,
gdy każda współrzędna $X_i$ jest rzeczywistą zmienną losową.
\end{stwierdzenie}
\begin{proof}
Rzuty współrzędne są funkcjami ciągłymi, co daje jeden kierunek.
W przeciwnym kierunku przeciwobraz prostokąta
$B_1\times\cdots\times B_d$, gdzie $B_i\in\mathcal B(\R)$,
jest przecięciem zdarzeń $\{X_i\in B_i\}$. Ponieważ
borelowskie prostokąty generują $\mathcal B(\R^d)$,
otrzymujemy mierzalność wektora.
\end{proof}

\begin{definicja}
Rzeczywiste zmienne losowe $X_1,\ldots,X_d$ nazywamy niezależnymi,
jeżeli dla wszystkich $B_1,\ldots,B_d\in\mathcal B(\R)$ zachodzi
$\Prob(X_1\in B_1,\ldots,X_d\in B_d)=\prod_{i=1}^d\Prob(X_i\in B_i)$.
Dla dowolnej rodziny zmiennych losowych niezależność oznacza
spełnienie tej własności przez każdy skończony podzbiór rodziny.
\end{definicja}

\begin{definicja}
Jeśli rozkład wektora $X\in\R^d$ jest absolutnie ciągły względem
$d$-wymiarowej miary Lebesgue'a $\lambda_d$, funkcję
$f_X:\R^d\to[0,\infty]$ spełniającą
$\Prob(X\in B)=\int_B f_X(x)\,\dd\lambda_d(x)$ nazywamy
gęstością łączną.
\end{definicja}

\begin{twierdzenie}[Gęstości brzegowe i niezależność]
Niech $(X,Y)$ będzie wektorem losowym o gęstości łącznej
$f_{X,Y}:\R^2\to[0,\infty]$. Wówczas funkcje
\[
f_X(x)=\int_\R f_{X,Y}(x,y)\,\dd\lambda_1(y),\qquad
f_Y(y)=\int_\R f_{X,Y}(x,y)\,\dd\lambda_1(x)
\]
są gęstościami rozkładów brzegowych (z dokładnością do wersji
prawie wszędzie). Zmienne $X,Y$ są niezależne wtedy i tylko wtedy,
gdy $f_{X,Y}(x,y)=f_X(x)f_Y(y)$ dla $\lambda_2$-prawie każdego
$(x,y)\in\R^2$.
\end{twierdzenie}
\begin{proof}
Najpierw sprawdźmy, że podane całki są rzeczywiście
gęstościami brzegowymi. Dla borelowskiego $A\subseteq\R$
z definicji rozkładu łącznego
\[
\Prob(X\in A)=\Prob((X,Y)\in A\times\R)
=\int_{A\times\R}f_{X,Y}(x,y)\,\dd\lambda_2(x,y).
\]
Ponieważ funkcja podcałkowa jest nieujemna, możemy
zastosować twierdzenie Tonellego i scałkować najpierw po $y$:
\[
\Prob(X\in A)=\int_A
\left(\int_\R f_{X,Y}(x,y)\,\dd\lambda_1(y)\right)
\dd\lambda_1(x).
\]
Tak otrzymujemy gęstość $f_X$; argument dla $f_Y$
jest symetryczny.

Jeśli $f_{X,Y}=f_X f_Y$ prawie wszędzie, to dla
borelowskich $A,B$ twierdzenie Tonellego daje
$\Prob(X\in A,Y\in B)
=(\int_A f_X\,\dd\lambda_1)
 (\int_B f_Y\,\dd\lambda_1)
=\Prob(X\in A)\Prob(Y\in B)$.
Jest to definicja niezależności.

Odwrotnie, przy niezależności dla wszystkich prostokątów
$A\times B$ miara łączna i miara produktowa
$\mu_X\otimes\mu_Y$ mają te same wartości.
Przez jednoznaczność miar na $\sigma$-ciele generowanym
przez prostokąty są identyczne. Miara produktowa ma
gęstość $f_X(x)f_Y(y)$ względem $\lambda_2$
(ponownie Tonelli). Jednoznaczność pochodnej
Radona--Nikodyma daje
$f_{X,Y}=f_Xf_Y$ dla $\lambda_2$-prawie każdego punktu.
\end{proof}

\begin{przyklad}[Jednostajność na trójkącie]
Niech $(X,Y)$ ma gęstość $f_{X,Y}(x,y)=2$ na zbiorze
$T=\{(x,y):0<x<y<1\}$ i zero poza nim.
Ponieważ pole $T$ wynosi $1/2$, funkcja jest unormowana.
Dla $0<x<1$ mamy $f_X(x)=\int_x^1 2\,\dd\lambda_1(y)=2(1-x)$,
a dla $0<y<1$ mamy $f_Y(y)=\int_0^y 2\,\dd\lambda_1(x)=2y$.
Zatem obie gęstości brzegowe nie są jednostajne, a $X,Y$ nie są
niezależne, co widać także ze zdarzenia $\{X<Y\}$
o prawdopodobieństwie jeden.
\end{przyklad}

\section{Rozkłady warunkowe przy istnieniu gęstości}
Przy ciągłej zmiennej $Y$ zdarzenie $\{Y=y\}$ ma zwykle
prawdopodobieństwo zero. Wyrażenia $\Prob(X\in A\mid Y=y)$
nie można zatem bezpośrednio rozumieć jako ilorazu zdarzeń.
Gdy istnieje gęstość łączna, jej rozkład na czynniki
„gęstość brzegowa razy gęstość warunkowa” stanowi
naturalne rozszerzenie wzoru na warunkowanie. Na tym etapie
konstruujemy je bez wchodzenia w ogólną teorię regularnych
prawdopodobieństw warunkowych.

\begin{stwierdzenie}[Warunkowa gęstość]
Niech $(X,Y)$ będzie wektorem losowym o gęstości łącznej $f_{X,Y}$,
a $f_Y$ niech oznacza gęstość brzegową $Y$.
Wybierzmy borelowskie wersje $f_{X,Y}$ i $f_Y$; po ewentualnej
zmianie na zbiorze wartości $y$ o zerowej mierze możemy przyjąć
$f_Y(y)=\int_\R f_{X,Y}(x,y)\,\dd\lambda_1(x)$,
gdy całka jest skończona. Dla $y$ z $0<f_Y(y)<\infty$
połóżmy $f_{X\mid Y=y}(x)=f_{X,Y}(x,y)/f_Y(y)$.
Dla pozostałych $y$ wybierzmy dowolną ustaloną gęstość,
np. standardową normalną. Tak wybrana funkcja
$x\mapsto f_{X\mid Y=y}(x)$ jest gęstością
probabilistyczną dla każdego $y$, po zmianie wersji na zbiorze
wyjątkowym o $\Law(Y)$-mierze zero.
Ponadto dla borelowskich $A,B\subseteq\R$
\begin{equation}\label{eq:joint_cond}
\Prob(X\in A,Y\in B)
=\int_B\left(\int_A f_{X\mid Y=y}(x)\,\dd\lambda_1(x)\right)
 f_Y(y)\,\dd\lambda_1(y),
\end{equation}
przy dowolnym wyborze gęstości probabilistycznej dla tych wartości $y$,
dla których iloraz nie jest określony.
\end{stwierdzenie}
\begin{proof}
Z twierdzenia Tonellego $f_Y(y)=\int_\R f_{X,Y}(x,y)\,
\dd\lambda_1(x)$ dla prawie wszystkich $y$, więc dla
$0<f_Y(y)<\infty$ całka z ilorazu po $x$ wynosi jeden.
Wzór \eqref{eq:joint_cond} wynika z przemnożenia przez $f_Y(y)$
i zastosowania twierdzenia Fubiniego--Tonellego.
\end{proof}

\begin{uwaga}
Symbol $\Prob(X\in A\mid Y=y)$, rozumiany przez powyższy wzór,
nie jest zwykłym ilorazem $\Prob(X\in A,Y=y)/\Prob(Y=y)$,
gdyż przy rozkładzie ciągłym obie wielkości często wynoszą zero.
Wersje gęstości warunkowej mogą się różnić na zbiorze wartości $y$
o zerowym prawdopodobieństwie.
\end{uwaga}

\begin{przyklad}[Warunkowanie na ustaloną wartość zmiennej ciągłej]
Wróćmy do rozkładu jednostajnego na trójkącie
$0<x<y<1$, o gęstości łącznej $f_{X,Y}(x,y)=2$.
Wykazaliśmy, że $f_Y(y)=2y$ dla $0<y<1$.
Zatem dla każdej takiej wartości $y$ możemy wybrać
wersję gęstości warunkowej
$f_{X\mid Y=y}(x)=y^{-1}\ind{\{0<x<y\}}$.
Jest ona unormowana, gdyż $\int_0^y y^{-1}\,\dd\lambda_1(x)=1$.
Na przykład przy $0<a<y<1$ otrzymujemy
$\Prob(X\leq a\mid Y=y)=a/y$ w sensie tej wersji
rozkładu warunkowego. Wniosek ma prostą interpretację:
jeśli obserwujemy wartość $Y=y$, to $X$ jest jednostajna
na przedziale $(0,y)$. Nie obliczaliśmy tu ilorazu
prawdopodobieństw zdarzeń o zerowej mierze;
posłużyliśmy się rozkładem łącznym i jego faktoryzacją.
\end{przyklad}



\begin{twierdzenie}[Warunkowa wartość oczekiwana obliczana z gęstości]
Niech $(X,Y)$ będzie rzeczywistym wektorem losowym
mającym borelowską gęstość łączną $f_{X,Y}$.
Niech $k(y,x)$ będzie gęstością warunkową
$f_{X\mid Y=y}(x)$ z poprzedniego stwierdzenia,
przy ustalonym borelowskim wyborze wersji.
Jeśli $h:\R\to\R$ jest funkcją borelowską i
$\E|h(X)|<\infty$, to funkcja
$g(y)=\int_\R h(x)k(y,x)\,\dd\lambda_1(x)$,
przy dowolnej wartości na zbiorze, gdzie całka nie jest
bezwzględnie zbieżna, jest borelowska i spełnia
\[
\E[h(X)\mid\sigma(Y)]=g(Y)\quad\text{p.n.}
\]
\end{twierdzenie}
\begin{proof}
Najpierw rozważmy $h\geq0$. Z mierzalności wspólnej
funkcji $(y,x)\mapsto h(x)k(y,x)$ oraz twierdzenia
Tonellego wynika borelowska mierzalność $g$.
Dla dowolnego zbioru borelowskiego $B\subseteq\R$
z poprzedniego wzoru na gęstość warunkową
otrzymujemy
\[
\E[h(X)\ind{\{Y\in B\}}]
=\int_B\!\left(\int_\R h(x)k(y,x)\,\dd\lambda_1(x)\right)
\dd\Law(Y)(y).
\]
Tożsamość dla funkcji nieujemnej można uzyskać przez
aproksymację rosnącą funkcjami prostymi, a następnie
zastosowanie zbieżności monotonicznej.
Dla ogólnej całkowalnej $h$ stosujemy ten argument
osobno do $h^+$ i $h^-$. Po scałkowaniu przy $B=\R$
widzimy, że obie odpowiednie funkcje warunkowe są
skończone dla $\Law(Y)$-prawie każdego $y$.
Różnica jest zatem dobrze określona prawie wszędzie,
a jej złożenie z $Y$ jest $\sigma(Y)$-mierzalne.
Zdarzenia $\{Y\in B\}$, $B\in\mathcal B(\R)$,
stanowią dokładnie $\sigma(Y)$, więc uzyskana
równość całek jest definicją warunkowej wartości
oczekiwanej. Jednoznaczność wersji kończy dowód.
\end{proof}
\begin{przyklad}[Obliczanie warunkowej średniej i wariancji]
W modelu jednostajnym na trójkącie $0<x<y<1$
mamy $k(y,x)=y^{-1}\ind{\{0<x<y\}}$ dla $0<y<1$.
Zatem
$\E[X\mid Y=y]=y^{-1}\int_0^y x\,\dd\lambda_1(x)=y/2$
oraz
$\E[X^2\mid Y=y]=y^{-1}\int_0^y x^2\,\dd\lambda_1(x)=y^2/3$.
Wariancja warunkowa jest równa $y^2/3-(y/2)^2=y^2/12$.
Stąd, w sensie wersji warunkowej,
$\E[X\mid\sigma(Y)]=Y/2$ i
$\Var(X\mid\sigma(Y))=Y^2/12$.
To przykład, w którym pojedyncza obserwacja $Y$
zmienia nie tylko prognozę, lecz również pozostałą
niepewność co do $X$.
\end{przyklad}

\section{Przekształcanie zmiennych i suma niezależnych zmiennych}
Jeśli znamy rozkład $X$, to często chcemy poznać rozkład
$g(X)$, na przykład kwadratu, logarytmu lub odwrotności.
Przy zmianie zmiennych gęstość musi być skorygowana o to,
jak $g$ rozciąga albo ściska małe przedziały: stąd pochodna
funkcji odwrotnej. Przy sumie niezależnych zmiennych na
ustaloną wartość $s$ składają się wszystkie pary $(x,s-x)$;
„sumowanie” ich wkładów prowadzi do splotu.

\begin{stwierdzenie}[Transformacja monotoniczna]
Niech $X$ ma gęstość $f_X$, a $g:I\to J$ będzie dyfeomorfizmem
przedziałów otwartych $I,J\subseteq\R$. Załóżmy, że
$\Prob(X\in I)=1$. Wtedy $Y=g(X)$ ma gęstość
$f_Y(y)=f_X(g^{-1}(y))\,|(g^{-1})'(y)|$ dla $y\in J$
i gęstość zero poza $J$.
\end{stwierdzenie}
\begin{proof}
Dla dowolnego borelowskiego zbioru $B\subseteq\R$
rozpoczynamy od definicji rozkładu zmiennej $Y=g(X)$:
\[
\Prob(Y\in B)=\Prob(X\in g^{-1}(B\cap J))
=\int_{g^{-1}(B\cap J)}f_X(x)\,\dd\lambda_1(x).
\]
Ostatnia równość wynika z założenia $\Prob(X\in I)=1$.
Ponieważ $g:I\to J$ jest dyfeomorfizmem, stosujemy
jednowymiarowe twierdzenie o zmianie zmiennych
$y=g(x)$ i otrzymujemy
\[
\Prob(Y\in B)=\int_{B\cap J}
 f_X(g^{-1}(y))\,|(g^{-1})'(y)|\,\dd\lambda_1(y).
\]
Zatem wskazana funkcja jest gęstością rozkładu $Y$.
Wartość bezwzględna przy pochodnej jest konieczna,
ponieważ odwracanie orientacji przedziału nie może
prowadzić do ujemnej gęstości.
\end{proof}

\begin{stwierdzenie}[Skończenie wiele gałęzi]
Niech $X$ ma gęstość $f_X$ i niech rozłączne przedziały otwarte
$I_1,\ldots,I_r$ spełniają $\Prob(X\in\bigcup_jI_j)=1$.
Załóżmy, że dla każdego $j$ funkcja $g|_{I_j}$ jest
$C^1$-dyfeomorfizmem na przedział $J_j$.
Wtedy $Y=g(X)$ ma gęstość
\[
f_Y(y)=\sum_{j:y\in J_j}
\frac{f_X((g|_{I_j})^{-1}(y))}
{|g'((g|_{I_j})^{-1}(y))|}
\quad\text{dla prawie wszystkich }y\in\R.
\]
\end{stwierdzenie}
\begin{proof}
Niech $B\subseteq\R$ będzie borelowski. Ponieważ
$\Prob(X\in\bigcup_j I_j)=1$ i przedziały $I_j$
są rozłączne, możemy rozłożyć zdarzenie $\{g(X)\in B\}$
według gałęzi:
\[
\Prob(g(X)\in B)=\sum_{j=1}^r
\int_{I_j\cap g^{-1}(B)} f_X(x)\,\dd\lambda_1(x).
\]
Na każdym $I_j$ funkcja $g$ jest dyfeomorfizmem
na $J_j$, więc poprzednie twierdzenie pozwala
przejść od $x$ do $y=g(x)$. Otrzymujemy
\[
\Prob(g(X)\in B)=\sum_{j=1}^r
\int_{B\cap J_j}
\frac{f_X((g|_{I_j})^{-1}(y))}
{|g'((g|_{I_j})^{-1}(y))|}\,\dd\lambda_1(y).
\]
Przenosimy sumę skończoną pod całkę, uzyskując
podany wzór. Wszystkie gałęzie prowadzące do tej samej
wartości $y$ muszą zostać dodane: w przeciwieństwie
do przypadku monotonicznego, pojedynczy wynik $Y=y$
może mieć kilka wartości przeciwobrazu.
\end{proof}

\begin{twierdzenie}[Splot]
Niech $X,Y$ będą niezależnymi rzeczywistymi zmiennymi losowymi
o gęstościach $f_X,f_Y$. Wtedy $S=X+Y$ ma gęstość
\begin{equation}\label{eq:convolution}
 f_S(s)=(f_X*f_Y)(s)
 :=\int_\R f_X(x)f_Y(s-x)\,\dd\lambda_1(x)
\end{equation}
dla $\lambda_1$-prawie każdego $s$.
\end{twierdzenie}
\begin{proof}
Z niezależności rozkład pary $(X,Y)$ jest iloczynem
rozkładów brzegowych, a więc ma gęstość
$f_{X,Y}(x,y)=f_X(x)f_Y(y)$.
Niech $B\subseteq\R$ będzie dowolnym zbiorem borelowskim.
Z definicji $S=X+Y$ i twierdzenia Tonellego
\[
\Prob(S\in B)=\int_{\R^2}\ind{\{x+y\in B\}}
 f_X(x)f_Y(y)\,\dd\lambda_2(x,y).
\]
Rozważmy liniową zmianę współrzędnych
$(x,y)\mapsto (x,s=x+y)$, której macierz ma wyznacznik
równy jeden. Jej odwrotność to $(x,s)\mapsto(x,s-x)$.
Po zmianie zmiennych otrzymujemy
\[
\Prob(S\in B)=\int_{\R^2}\ind{\{s\in B\}}
 f_X(x)f_Y(s-x)\,\dd\lambda_2(x,s).
\]
Ponieważ funkcja podcałkowa jest nieujemna,
możemy zmienić kolejność całkowania:
\[
\Prob(S\in B)=\int_B\left(
\int_\R f_X(x)f_Y(s-x)\,\dd\lambda_1(x)
\right)\dd\lambda_1(s).
\]
Z definicji gęstości wynika wzór na splot. W szczególności
całka z prawej strony po całej $\R$ wynosi jeden,
a splot nieujemnych gęstości jest znów gęstością.
\end{proof}

\begin{przyklad}[Suma dwóch jednostajnych zmiennych]
Jeśli $U,V$ są niezależne i jednostajne na $(0,1)$,
to dla $S=U+V$ ze wzoru na splot otrzymujemy
\[
f_S(s)=\begin{cases}
s,&0<s\leq1,\\
2-s,&1<s<2,\\
0,&\text{w pozostałych przypadkach}.
\end{cases}
\]
Rozkład sumy nie jest jednostajny: dla $s$ bliskiego $1$
istnieje więcej par $(u,v)$ o sumie $s$ niż dla $s$ bliskiego $0$.
\end{przyklad}

\begin{przyklad}[Bez niezależności wzór na splot może zawieść]
Niech $U\sim\Unif(0,1)$ oraz $V=1-U$.
Wtedy również $V\sim\Unif(0,1)$, lecz $U+V=1$
prawie na pewno. Rozkład sumy jest więc skupiony
w jednym punkcie i nie ma gęstości względem miary
Lebesgue'a, mimo że obie zmienne mają gęstości.
Nie jest to sprzeczne z twierdzeniem o splocie,
ponieważ $U$ i $V$ nie są niezależne.
\end{przyklad}


\chapter{Wartość oczekiwana, momenty i podstawowe nierówności}

\section{Całka względem rozkładu}
Prawdopodobieństwo odpowiada na pytanie, jak często wydarza się
zdarzenie. Wartość oczekiwana odpowiada na inne pytanie:
jaka jest przeciętna wielkość liczbowego wyniku, jeśli ważymy
możliwości ich prawdopodobieństwami? W modelu skończonym jest
to ważona suma. Całka względem miary probabilistycznej jest
jej konsekwentnym uogólnieniem na dowolne zmienne losowe.
Nazwa „oczekiwana” nie oznacza wartości najbardziej prawdopodobnej;
może to być wartość, której zmienna nigdy nie przyjmuje.

\begin{definicja}
Niech $X$ będzie rzeczywistą zmienną losową na
$(\Omega,\mathcal F,\Prob)$. Jeśli $X\geq0$, definiujemy
$\E X:=\int_\Omega X(\omega)\,\dd\Prob(\omega)\in[0,\infty]$.
Jeśli $\E|X|<\infty$, nazywamy $X$ całkowalną i definiujemy
$\E X:=\int_\Omega X(\omega)\,\dd\Prob(\omega)\in\R$.
W przypadku ogólnym, gdy przynajmniej jedna z liczb
$\E X^+,\E X^-$ jest skończona, dopuszczamy rozszerzoną
wartość oczekiwaną $\E X=\E X^+-\E X^-$.
\end{definicja}

\begin{twierdzenie}[Przenoszenie całkowania na rozkład]
Niech $X$ będzie rzeczywistą zmienną losową o rozkładzie $\mu_X$,
a $h:\R\to\R$ mierzalną funkcją borelowską.
Jeżeli $h\geq0$ albo $\int_\R|h(x)|\,\dd\mu_X(x)<\infty$, to
\begin{equation}\label{eq:lotus}
\E[h(X)]=\int_\R h(x)\,\dd\mu_X(x).
\end{equation}
W szczególności, dla dyskretnego $X$ o masach $p(x)$ prawa strona
jest sumą $\sum_xh(x)p(x)$, a dla $X$ o gęstości $f_X$
jest całką $\int_\R h(x)f_X(x)\,\dd\lambda_1(x)$.
\end{twierdzenie}
\begin{proof}
Dowód polega na przejściu od funkcji najprostszych
do ogólnych; wyjaśnia on zarazem, dlaczego rozkład
$\mu_X$ pozwala obliczać wartości oczekiwane funkcji $X$.

Dla funkcji wskaźnikowej $h=\ind B$, gdzie
$B\in\mathcal B(\R)$, mamy
$\E[h(X)]=\Prob(X\in B)=\mu_X(B)
=\int_\R\ind B(x)\,\dd\mu_X(x)$.
Dla nieujemnej funkcji prostej
$h=\sum_{j=1}^r a_j\ind{B_j}$, gdzie $a_j\geq0$,
równość wynika z liniowości całki i poprzedniego przypadku.

Dla dowolnej nieujemnej funkcji mierzalnej $h$
wybieramy ciąg nieujemnych funkcji prostych $h_n\uparrow h$.
Wówczas $h_n(X)\uparrow h(X)$ na $\Omega$.
Stosując twierdzenie o zbieżności monotonicznej
po obu stronach już udowodnionych równości, dostajemy
\[
\E[h(X)]=\lim_n\E[h_n(X)]
=\lim_n\int_\R h_n\,\dd\mu_X
=\int_\R h\,\dd\mu_X.
\]
Jeśli $h$ jest całkowalna względem $\mu_X$,
rozkładamy $h=h^+-h^-$ i stosujemy pierwszy przypadek
do obu nieujemnych części. Obie odpowiadające całki
są skończone, więc ich różnica jest dobrze określona.
\end{proof}


\begin{przyklad}[Średnia nie musi być możliwym wynikiem]
Niech $X$ przyjmuje wartości $0$ i $1$ z prawdopodobieństwami
po $1/2$. Wówczas $\E X=1/2$, ale $\Prob(X=1/2)=0$.
Średnia jest środkiem ciężkości rozkładu, a nie typową
zrealizowaną wartością. Dla bardzo nierównomiernych rozkładów
może ona znajdować się daleko od większości realizacji.
\end{przyklad}

\begin{definicja}
Dla całkowalnej zmiennej losowej $X$ definiujemy średnią
$m_X:=\E X$. Jeżeli $\E X^2<\infty$, definiujemy
wariancję $\Var(X):=\E[(X-\E X)^2]$.
Odchyleniem standardowym nazywamy $\sqrt{\Var(X)}$.
Dla zmiennych $X,Y$ o skończonych drugich momentach definiujemy
kowariancję $\Cov(X,Y):=\E[(X-\E X)(Y-\E Y)]$.
\end{definicja}

\begin{stwierdzenie}[Algebra momentów]
Dla rzeczywistych zmiennych losowych $X,Y$ o skończonych
drugich momentach oraz $a,b\in\R$ zachodzą
\begin{align*}
\Var(X)&=\E X^2-(\E X)^2, &
\Var(aX+b)&=a^2\Var(X),\\
\Cov(X,Y)&=\E[XY]-\E X\,\E Y, &
\Var(X+Y)&=\Var(X)+\Var(Y)+2\Cov(X,Y).
\end{align*}
Jeżeli $X$ i $Y$ są niezależne, to $\E[XY]=\E X\,\E Y$,
a zatem $\Cov(X,Y)=0$ i $\Var(X+Y)=\Var(X)+\Var(Y)$.
\end{stwierdzenie}
\begin{proof}
Wszystkie podane wyrażenia są skończone,
ponieważ nierówność Cauchy'ego--Schwarza zapewnia
$\E|XY|\leq(\E X^2)^{1/2}(\E Y^2)^{1/2}<\infty$.
Dzięki temu możemy swobodnie korzystać z liniowości całki.
Oznaczając $m=\E X$, rozwijamy
$\E(X-m)^2=\E X^2-2m\E X+m^2=\E X^2-m^2$.
Ponadto $\E(aX+b)=a\E X+b$, więc
$aX+b-\E(aX+b)=a(X-\E X)$,
a stąd $\Var(aX+b)=a^2\Var(X)$.

Dla $m_X=\E X$, $m_Y=\E Y$ mamy
$\Cov(X,Y)=\E[(X-m_X)(Y-m_Y)]
=\E[XY]-m_Xm_Y$.
Ponieważ
$X+Y-\E(X+Y)=(X-m_X)+(Y-m_Y)$,
rozwinięcie kwadratu i całkowanie daje
$\Var(X+Y)=\Var(X)+\Var(Y)+2\Cov(X,Y)$.

Jeśli $X,Y$ są niezależne, ich rozkład łączny to
$\mu_X\otimes\mu_Y$. Przez Fubiniego, którego
stosowanie usprawiedliwia całkowalność iloczynu,
\[
\E[XY]=\int_{\R^2}xy\,\dd(\mu_X\otimes\mu_Y)(x,y)
=\left(\int_\R x\,\dd\mu_X(x)\right)
 \left(\int_\R y\,\dd\mu_Y(y)\right)
=\E X\,\E Y.
\]
To daje zerową kowariancję i wzór na wariancję sumy.
Odwrotna implikacja nie jest prawdziwa i nie wolno
jej wykorzystywać bez dodatkowych założeń.
\end{proof}

\begin{uwaga}
Zerowa kowariancja nie oznacza na ogół niezależności.
Jeśli $X$ jest jednostajna na $[-1,1]$, a $Y=X^2$, to
$\Cov(X,Y)=\E X^3-\E X\,\E X^2=0$,
ale $Y$ jest funkcją $X$, a więc zmienne nie są niezależne.
\end{uwaga}

\begin{stwierdzenie}[Momenty zmiennej wskaźnikowej]
Jeżeli $A\in\mathcal F$ i $I=\ind{A}$, to
$\E I=\Prob(A)$ oraz
$\Var(I)=\Prob(A)(1-\Prob(A))$.
Dla dowolnych zdarzeń $A,B$ zachodzi
$\Cov(\ind{A},\ind{B})=\Prob(A\cap B)-\Prob(A)\Prob(B)$.
\end{stwierdzenie}
\begin{proof}
Zmienna $I$ przyjmuje wartości zero i jeden, a $I^2=I$.
Ponadto $\ind{A}\ind{B}=\ind{A\cap B}$.
Wstawiamy te równości do definicji momentów.
\end{proof}

\section{Nierówności Markowa i Czebyszewa}
Nie zawsze umiemy wyznaczyć cały rozkład, ale czasem znamy
średnią albo wariancję. Czy można na tej podstawie oszacować
prawdopodobieństwo dużego wyniku lub dużego odchylenia od średniej?
Nierówność Markowa opiera się tylko na dodatniości, a nierówność
Czebyszewa stosuje tę samą ideę do kwadratu odchylenia.
Ich siła tkwi w uniwersalności: nie wymagają normalności,
jednostajności ani nawet znajomości gęstości.

\begin{twierdzenie}[Nierówność Markowa]
Jeżeli $X\geq0$ prawie na pewno i $a>0$, to
$\Prob(X\geq a)\leq\E X/a$, gdzie dopuszczamy $\E X=\infty$.
\end{twierdzenie}
\begin{proof}
Na zdarzeniu $A=\{X\geq a\}$ zachodzi $X\geq a$,
a poza nim $a\ind A=0\leq X$.
Możemy więc zapisać punktową nierówność
$a\ind A\leq X$ dla każdego wyniku poza ewentualnym
zbiorem zerowym z założenia $X\geq0$ prawie na pewno.
Z monotoniczności całki mamy
$a\Prob(A)=\E[a\ind A]\leq\E X$.
Dzielenie przez $a>0$ daje tezę. Nierówność pozostaje
prawdziwa, choć może być bezużyteczna, gdy $\E X=\infty$.
Dowód pokazuje również, kiedy stała jest optymalna:
gdy $X$ przyjmuje wyłącznie wartości $0$ i $a$,
zachodzi równość.
\end{proof}

\begin{twierdzenie}[Nierówność Czebyszewa]
Niech $X$ będzie rzeczywistą zmienną losową o skończonym
drugim momencie. Wtedy dla każdego $\varepsilon>0$ zachodzi
\[
\Prob(|X-\E X|\geq\varepsilon)
\leq\frac{\Var(X)}{\varepsilon^2}.
\]
\end{twierdzenie}
\begin{proof}
Wprowadźmy nieujemną zmienną
$Y=(X-\E X)^2$. Jej wartość oczekiwana jest równa
$\E Y=\Var(X)<\infty$. Ponadto dla $\varepsilon>0$
zdarzenia
$\{Y\geq\varepsilon^2\}$ oraz
$\{|X-\E X|\geq\varepsilon\}$ są identyczne.
Stosując właśnie udowodnioną nierówność Markowa
do zmiennej $Y$ i progu $\varepsilon^2$, otrzymujemy
\[
\Prob(|X-\E X|\geq\varepsilon)
=\Prob(Y\geq\varepsilon^2)
\leq\frac{\E Y}{\varepsilon^2}
=\frac{\Var(X)}{\varepsilon^2}.
\]
Nie używaliśmy żadnych informacji o kształcie rozkładu $X$.
\end{proof}

\begin{przyklad}[Kontrola odchylenia sumy]
Niech $X_1,\ldots,X_n$ będą niezależnymi zmiennymi losowymi
o wspólnej średniej $m$ i wariancji $\sigma^2<\infty$,
gdzie $n\geq1$. Dla średniej $\overline X_n:=n^{-1}\sum_{i=1}^nX_i$
mamy $\E\overline X_n=m$ i $\Var(\overline X_n)=\sigma^2/n$.
Z nierówności Czebyszewa wynika
$\Prob(|\overline X_n-m|\geq\varepsilon)
\leq\sigma^2/(n\varepsilon^2)$.
Jest to zalążek słabego prawa wielkich liczb, którego
pełne sformułowania i uogólnienia pojawią się w dalszej części tomu.
\end{przyklad}

\section{Momenty wyższych rzędów i sens funkcji tworzących}

Dwa pierwsze momenty często nie wystarczają, aby odróżnić
rozkłady. Zmienna może mieć tę samą średnią i wariancję co inna,
a jednocześnie zupełnie inny kształt ogonów. Chcielibyśmy zatem
przechowywać informacje o całej rodzinie momentów, a najlepiej
w jednej funkcji, na której można wykonywać zwykłe operacje
analityczne. Dwa użyteczne pomysły polegają na obliczaniu
$\E z^X$ dla zmiennych całkowitoliczbowych i $\E e^{tX}$
dla zmiennych rzeczywistych. Trzeci --- wykorzystujący
$e^{itX}$ --- pozwala uniknąć problemu rozbieżności.

\begin{definicja}
Niech $X$ będzie rzeczywistą zmienną losową i $r\in\N$.
Jeśli $\E|X|^r<\infty$, to $m_r(X):=\E X^r$
nazywamy momentem zwykłym rzędu $r$. Jeśli $\E X$
istnieje i $\E|X|^r<\infty$, to
$\mu_r(X):=\E(X-\E X)^r$ nazywamy momentem centralnym.
Momentami absolutnymi nazywamy $\E|X|^r$.
\end{definicja}

\begin{uwaga}
Dla nieparzystego $r$ całka $\E X^r$ może nie istnieć,
nawet gdy symetria rozkładu formalnie sugeruje zero.
Istnienie momentu jako skończonej liczby wymaga
$\E|X|^r<\infty$. Gdy $r>s>0$, z nierówności
$|x|^s\leq 1+|x|^r$ wynika, że istnienie momentu
absolutnego rzędu $r$ zapewnia istnienie wszystkich momentów
absolutnych mniejszych rzędów.
\end{uwaga}

\section{Funkcja tworząca prawdopodobieństwa}

Dla zmiennej $N$ przyjmującej wartości w $\N_0:=\{0,1,2,\ldots\}$
rozkład jest ciągiem $p_k=\Prob(N=k)$. Możemy potraktować
$p_k$ jako współczynniki szeregu potęgowego. Wówczas
mnożenie funkcji tworzących odpowiada dodawaniu niezależnych
liczb wystąpień --- dokładnie tak, jak splot ciągów mas
opisuje rozkład ich sumy.

\begin{definicja}
Niech $N:\Omega\to\N_0$ będzie zmienną losową.
Jej funkcją tworzącą prawdopodobieństwa nazywamy
$G_N(z):=\E z^N=\sum_{k=0}^\infty \Prob(N=k)z^k$
dla $|z|\leq1$, z konwencją $0^0=1$.
Dla argumentów zespolonych stosujemy tę samą definicję.
\end{definicja}

\begin{twierdzenie}[Odzyskiwanie rozkładu i momentów silniowych]
Niech $N$ przyjmuje wartości w $\N_0$, a $G_N$ będzie jej
funkcją tworzącą. Wtedy $G_N(1)=1$, a dla $k\in\N_0$
zachodzi $\Prob(N=k)=G_N^{(k)}(0)/k!$.
Dla $r\in\N$ i $0\leq z<1$ mamy
\[
G_N^{(r)}(z)=\E\bigl[(N)_r z^{N-r}\ind{\{N\geq r\}}\bigr],
\qquad (n)_r:=n(n-1)\cdots(n-r+1).
\]
W powyższym wzorze całe wyrażenie pod wartością oczekiwaną
ma wartość $0$ na $\{N<r\}$, bez obliczania $z^{N-r}$;
dla $z=0$ i $N=r$ przyjmujemy $0^0=1$.
W szczególności $G_N^{(r)}(1-)=\E(N)_r\in[0,\infty]$.
Jeżeli $\E N^2<\infty$, to
$\E N=G_N'(1-)$ i
$\Var(N)=G_N''(1-)+G_N'(1-)-[G_N'(1-)]^2$.
\end{twierdzenie}
\begin{proof}
Dla $|z|<1$ szereg $\sum_k p_kz^k$ jest bezwzględnie
zbieżny i można go różniczkować dowolną liczbę razy
wewnątrz koła jednostkowego: na każdym $|z|\leq\rho<1$
pochodne kolejnych wyrazów są ograniczone przez
$(k)_r\rho^{k-r}$, a szereg tych ograniczeń jest zbieżny,
ponieważ $p_k\leq1$. Współczynnik przy $z^k$ jest więc
równy $G_N^{(k)}(0)/k!$. Z drugiej strony, podstawiając
$z=1$, otrzymujemy $G_N(1)=\sum_kp_k=1$.
Dla $z\uparrow1$ każdy składnik $(k)_rz^{k-r}$ jest
niemalejący i nieujemny. Twierdzenie o zbieżności
monotonicznej dla sum daje
$G_N^{(r)}(1-)=\sum_{k\geq r}(k)_rp_k=\E(N)_r$,
bez uprzedniego zakładania skończoności tej liczby.
Wreszcie $N^2=N(N-1)+N$; podstawiając tę tożsamość
do wzoru $\Var(N)=\E N^2-(\E N)^2$, otrzymujemy tezę.
\end{proof}

\begin{twierdzenie}[Suma niezależnych liczników]
Jeśli $N_1,N_2$ są niezależnymi zmiennymi o wartościach
w $\N_0$, to $G_{N_1+N_2}(z)=G_{N_1}(z)G_{N_2}(z)$
dla $|z|\leq1$.
\end{twierdzenie}
\begin{proof}
Dla $|z|\leq1$ wszystkie zmienne $z^{N_j}$ są ograniczone
co do modułu przez jeden. Niezależność $N_1,N_2$ implikuje
niezależność ich mierzalnych funkcji, więc przez Fubiniego
$\E[z^{N_1}z^{N_2}]=\E z^{N_1}\,\E z^{N_2}$.
Lewa strona jest równa $\E z^{N_1+N_2}$.
\end{proof}

\begin{przyklad}[Cztery klasyczne funkcje tworzące]
Dla $B\sim\operatorname{Bern}(p)$ zachodzi
$G_B(z)=1-p+pz$. Dla $N\sim\Bin(n,p)$ mamy
$G_N(z)=(1-p+pz)^n$. Dla $N\sim\Poi(\lambda)$:
$G_N(z)=\exp(\lambda(z-1))$.
Jeśli $T\sim\Geom(p)$ oznacza numer pierwszej udanej próby,
to $N=T-1$ liczy poprzedzające ją porażki i
$G_N(z)=p/[1-(1-p)z]$ dla $|z|\leq1$.
Dla liczby porażek przed $r$-tym sukcesem otrzymujemy
$G(z)=[p/(1-(1-p)z)]^r$.
Porównanie tych wzorów od razu wyjaśnia stabilność
rozkładu Poissona przy sumowaniu i konstrukcję
ujemnego dwumianowego jako sumy czasów oczekiwania.
\end{przyklad}

\section{Funkcja tworząca momenty i kumulanty}

\begin{definicja}
Dla rzeczywistej zmiennej losowej $X$ definiujemy
$M_X(t):=\E e^{tX}\in(0,\infty]$ dla $t\in\R$.
Dziedziną skończoności jest $D_X:=\{t\in\R:M_X(t)<\infty\}$;
zawsze $0\in D_X$, lecz $D_X$ nie musi zawierać
żadnego otoczenia zera. Jeśli $M_X$ jest skończona
w otoczeniu zera, określamy również funkcję tworzącą
kumulanty $K_X(t):=\log M_X(t)$.
\end{definicja}

\begin{twierdzenie}[Istnienie pochodnych i odzyskiwanie momentów]
Niech $X$ będzie rzeczywistą zmienną losową oraz
$\delta>0$ takie, że $M_X(-\delta)<\infty$ i
$M_X(\delta)<\infty$. Wtedy $M_X$ jest nieskończenie
wiele razy różniczkowalna na $(-\delta,\delta)$ oraz
dla każdego $r\in\N_0$ i $|t|<\delta$ zachodzi
\[
M_X^{(r)}(t)=\E[X^r e^{tX}].
\]
W szczególności $\E|X|^r<\infty$ i
$M_X^{(r)}(0)=\E X^r$ dla każdego $r\in\N_0$.
\end{twierdzenie}
\begin{proof}
Najpierw zwróćmy uwagę, że
$e^{\delta|X|}\leq e^{\delta X}+e^{-\delta X}$.
Zatem $\E e^{\delta|X|}<\infty$.
Ustalmy $0<\eta<\delta$ i $r\geq0$. Funkcja
$u\mapsto u^re^{-(\delta-\eta)u}$ jest ograniczona na
$[0,\infty)$, więc istnieje $C_{r,\eta,\delta}<\infty$
taka, że dla $|t|\leq\eta$:
$|X|^re^{tX}\leq C_{r,\eta,\delta}e^{\delta|X|}$.
Prawa strona jest całkowalna i nie zależy od $t$.
Możemy więc stosować twierdzenie o zbieżności zdominowanej
do ilorazów różnicowych, a następnie do wyższych pochodnych:
każde różniczkowanie $e^{tX}$ mnoży wyrażenie przez $X$.
Dostajemy wzór na wszystkie pochodne. Dla $t=0$
otrzymujemy momenty, które są skończone na mocy
powyższego oszacowania dla $r$ i $\eta=0$.
\end{proof}

\begin{twierdzenie}[Dodawanie kumulantów]
Niech $X,Y$ będą niezależnymi rzeczywistymi zmiennymi losowymi,
a ich funkcje $M_X,M_Y$ będą skończone w pewnym wspólnym
otoczeniu zera. Wtedy w tym otoczeniu
$M_{X+Y}(t)=M_X(t)M_Y(t)$ i
$K_{X+Y}(t)=K_X(t)+K_Y(t)$.
Ponadto $K_X'(0)=\E X$, $K_X''(0)=\Var(X)$
oraz $K_X'''(0)=\E(X-\E X)^3$.
\end{twierdzenie}
\begin{proof}
Nieujemność $e^{tX}$ i $e^{tY}$ oraz niezależność pozwalają
zastosować Tonellego do miary produktowej:
$\E e^{t(X+Y)}=\E e^{tX}\E e^{tY}$.
Obie wartości są skończone w podanym otoczeniu, więc
możemy wziąć logarytm. Ponieważ $M_X(0)=1$, reguły
różniczkowania logarytmu dają
$K_X'=M_X'/M_X$ i
$K_X''=M_X''/M_X-(M_X'/M_X)^2$.
Dla trzeciej pochodnej otrzymujemy przy $t=0$
$K_X'''(0)=\E X^3-3\E X\,\E X^2+2(\E X)^3$,
co jest rozwinięciem trzeciego momentu centralnego.
\end{proof}

\begin{przyklad}[Kumulanty trzech rodzin]
Dla $X\sim\Normal(\mu,\sigma^2)$ mamy
$K_X(t)=\mu t+\sigma^2t^2/2$: wszystkie pochodne
rzędu co najmniej trzeciego znikają. Dla
$N\sim\Poi(\lambda)$ zachodzi $K_N(t)=\lambda(e^t-1)$,
a więc każda kumulanta dodatniego rzędu jest równa $\lambda$.
Dla zmiennej Bernoulliego $K_B(t)=\log(1-p+pe^t)$.
\end{przyklad}

\begin{stwierdzenie}[Jednoznaczność funkcji tworzącej momenty]
Niech $X,Y$ będą rzeczywistymi zmiennymi losowymi.
Jeśli ich funkcje $M_X,M_Y$ są skończone na
$(-\delta,\delta)$ dla pewnego $\delta>0$
i są tam równe, to $\Law(X)=\Law(Y)$.
\end{stwierdzenie}
\begin{proof}
Rozważmy transformaty zespolone
$L_X(z)=\E e^{zX}$ i $L_Y(z)=\E e^{zY}$
na pasie $\{z\in\mathbb C:|\operatorname{Re}z|<\delta\}$.
Są tam holomorficzne: dla dowolnych $0<\eta<\delta'
<\delta$, na pasie $|\operatorname{Re}z|\leq\eta$
moduły wyrazów $X^re^{zX}$ można oszacować przez
$C_r e^{\delta'|X|}$, którego wartość oczekiwana
jest skończona, ponieważ $M_X(\delta')$ i
$M_X(-\delta')$ są skończone. To zapewnia zbieżność
zdominowaną dla kolejnych pochodnych zespolonych.
Z zasady identyczności funkcji holomorficznych równość
$L_X(t)=L_Y(t)$ na $(-\delta,\delta)$ implikuje
równość na całym pasie, w szczególności dla $z=it$.
Zatem $\varphi_X(t)=\varphi_Y(t)$ dla każdego $t\in\R$.
Twierdzenie o jednoznaczności funkcji charakterystycznej,
udowodnione w następnej sekcji, daje $\Law(X)=\Law(Y)$.
\end{proof}

\begin{przyklad}[Wszystkie momenty nie muszą wyznaczać rozkładu]
Niech $Z\sim\Normal(0,1)$ oraz $X=e^Z$.
Zmienna $X$ ma wszystkie dodatnie całkowite momenty
$\E X^n=e^{n^2/2}$, ale dla każdego $t>0$ zachodzi
$M_X(t)=\infty$ (w całce po $z\to\infty$ wykładnik
$t e^z-z^2/2$ zmierza do $+\infty$).
Co więcej, przyjmijmy gęstość $f_0(x)=
\ind{\{x>0\}}e^{-(\log x)^2/2}/(x\sqrt{2\pi})$
i dla $|\varepsilon|\leq1$ połóżmy
$f_\varepsilon(x)=f_0(x)[1+\varepsilon\sin(2\pi\log x)]$.
Są to nieujemne gęstości: dla każdego $n\in\N_0$
\[
\int_0^\infty x^n f_0(x)\sin(2\pi\log x)
\,\dd\lambda_1(x)
=\operatorname{Im}\exp((n+2\pi i)^2/2)=0,
\]
gdyż $\sin(2\pi n)=0$. Dla $n=0$ otrzymujemy
normalizację, a dla wszystkich $n$ --- identyczne momenty.
Dla różnych $\varepsilon$ gęstości są jednak różne prawie
wszędzie. Ten przykład pokazuje, dlaczego sama lista momentów
bez warunków zapewniających jednoznaczność może być niewystarczająca.
\end{przyklad}

\section{Funkcja charakterystyczna i transformata Fouriera}

Kluczowa różnica między $e^{tX}$ a $e^{itX}$ polega na tym,
że $|e^{itX}|=1$ nawet dla bardzo ciężkich ogonów rozkładu.
Całkowanie zespolonej wykładniczej nie wymaga więc istnienia
żadnego momentu. Zmienna faza $e^{itx}$ rejestruje zarazem
informację o całym rozkładzie, nie tylko o jego średniej.

\begin{definicja}
Funkcją charakterystyczną rzeczywistej zmiennej losowej $X$
nazywamy $\varphi_X:\R\to\mathbb C$ zadaną wzorem
$\varphi_X(t):=\E e^{itX}=\int_\R e^{itx}\,\dd\mu_X(x)$.
\end{definicja}

\begin{twierdzenie}[Własności funkcji charakterystycznej]
Dla każdej rzeczywistej zmiennej losowej $X$ i
$t,s,a,b\in\R$ zachodzą:
\begin{enumerate}
\item $\varphi_X(0)=1$, $|\varphi_X(t)|\leq1$
 i $\varphi_X(-t)=\overline{\varphi_X(t)}$;
\item $\varphi_X$ jest jednostajnie ciągła na $\R$;
\item $\varphi_{aX+b}(t)=e^{itb}\varphi_X(at)$;
\item dla niezależnych $X,Y$ mamy
 $\varphi_{X+Y}(t)=\varphi_X(t)\varphi_Y(t)$;
\item dla dowolnych $n\in\N$, $t_1,\ldots,t_n\in\R$,
 $c_1,\ldots,c_n\in\mathbb C$ zachodzi
 $\sum_{j,k=1}^n c_j\overline{c_k}\varphi_X(t_j-t_k)\geq0$.
\end{enumerate}
Jeżeli $\E|X|^r<\infty$ dla pewnego $r\in\N$,
to $\varphi_X$ jest $r$ razy różniczkowalna i
$\varphi_X^{(k)}(t)=\E[(iX)^k e^{itX}]$
dla $0\leq k\leq r$.
\end{twierdzenie}
\begin{proof}
Ponieważ $|e^{itX}|=1$, całka zawsze istnieje,
a nierówność trójkąta daje $|\varphi_X(t)|\leq1$.
Pozostałe tożsamości z punktu pierwszego wynikają
z podstawienia $t=0$ i sprzężenia całki.
Dla ciągu $h_n\to0$ wyrażenie
$|\varphi_X(t+h_n)-\varphi_X(t)|$
jest nie większe niż $\E|e^{ih_nX}-1|$, niezależnie od $t$.
Funkcja pod wartością oczekiwaną dąży punktowo do zera
i nie przekracza $2$, więc zbieżność zdominowana daje
jednostajną ciągłość. Punkt trzeci wynika z równości
$e^{it(aX+b)}=e^{itb}e^{i(at)X}$.
Dla czwartego korzystamy z niezależności i Fubiniego:
$\E[e^{itX}e^{itY}]=\E e^{itX}\E e^{itY}$.
Dla ostatniego punktu obliczamy
\[
\sum_{j,k}c_j\overline{c_k}\varphi_X(t_j-t_k)
=\E\left|\sum_{j=1}^n c_je^{it_jX}\right|^2\geq0.
\]
Wreszcie $|(iX)^k e^{itX}|=|X|^k$ jest całkowalna
przy $k\leq r$. Zbieżność zdominowana stosowana
kolejno do ilorazów różnicowych uzasadnia różniczkowanie
pod całką i podany wzór na pochodne.
\end{proof}

\begin{lemat}[Całka Fouriera dla rozkładu Gaussa]
Jeśli $Z\sim\Normal(0,1)$, to
$\varphi_Z(t)=e^{-t^2/2}$ dla $t\in\R$.
Jeśli $Y\sim\Normal(m,\sigma^2)$, to
$\varphi_Y(t)=\exp(imt-\sigma^2t^2/2)$.
\end{lemat}
\begin{proof}
Niech $g(x)=(2\pi)^{-1/2}e^{-x^2/2}$.
Zdominowana zbieżność pozwala różniczkować
$\varphi_Z(t)=\int_\R e^{itx}g(x)\,\dd\lambda_1(x)$.
Ponieważ $g'(x)=-xg(x)$, całkowanie przez części daje
$\int_\R xe^{itx}g(x)\,\dd\lambda_1(x)
=it\int_\R e^{itx}g(x)\,\dd\lambda_1(x)$;
wyrazy brzegowe znikają dzięki wykładniczemu zanikowi $g$.
Stąd $\varphi_Z'(t)=i(it\varphi_Z(t))=-t\varphi_Z(t)$.
Przy warunku początkowym $\varphi_Z(0)=1$ rozwiązanie tego
równania to $e^{-t^2/2}$. Drugą równość otrzymujemy
przez przekształcenie afiniczne $Y\stackrel d=m+\sigma Z$.
\end{proof}

\begin{twierdzenie}[Jednoznaczność funkcji charakterystycznej]
Jeżeli $X,Y$ są rzeczywistymi zmiennymi losowymi i
$\varphi_X(t)=\varphi_Y(t)$ dla każdego $t\in\R$,
to $\Law(X)=\Law(Y)$.
\end{twierdzenie}
\begin{proof}
Oznaczmy rozkłady przez $\mu$ i $\nu$, a przez
$g_\sigma(x)=(\sigma\sqrt{2\pi})^{-1}
\exp(-x^2/(2\sigma^2))$ gęstość $\Normal(0,\sigma^2)$
dla $\sigma>0$. Dzięki poprzedniemu lematowi i
całkowalności funkcji $t\mapsto e^{-\sigma^2t^2/2}$
możemy zastosować Fubiniego do wzoru Fouriera na gęstość
Gaussa i otrzymać dla każdego $x\in\R$
\[
\int_\R g_\sigma(x-y)\,\dd\mu(y)
=\frac1{2\pi}\int_\R
 e^{-itx}\varphi_\mu(t)e^{-\sigma^2t^2/2}
\,\dd\lambda_1(t).
\]
Tutaj użyliśmy elementarnej odwrotności transformaty
Fouriera Gaussa, wynikającej z udowodnionego wzoru na
$\varphi_Z$ (równoważnie z całki Gaussa po zamianie
ról zmiennych). To samo zachodzi dla $\nu$.
Równość funkcji charakterystycznych oznacza zatem,
że sploty $\mu*g_\sigma$ i $\nu*g_\sigma$ mają
identyczne gęstości, a więc są równymi miarami.

Dla dowolnej funkcji $h:\R\to\R$ ograniczonej i ciągłej
oraz niezależnej zmiennej $Z\sim\Normal(0,1)$ zachodzi
\[
\int_\R h(x)\,\dd(\mu*g_\sigma)(x)
=\int_\R\E[h(x+\sigma Z)]\,\dd\mu(x).
\]
Gdy $\sigma\downarrow0$, dla każdego ustalonego $x$
zmienna $h(x+\sigma Z)$ dąży prawie na pewno do $h(x)$,
a jej moduł jest ograniczony przez $\|h\|_\infty$.
Dwukrotne zastosowanie zbieżności zdominowanej daje
$\int h\,\dd(\mu*g_\sigma)\to\int h\,\dd\mu$;
podobnie dla $\nu$. Ponieważ sploty były równe dla
każdego $\sigma$, otrzymujemy
$\int h\,\dd\mu=\int h\,\dd\nu$ dla każdej
ograniczonej funkcji ciągłej $h$. Takie całki
jednoznacznie wyznaczają borelowskie miary probabilistyczne
na $\R$ (przez aproksymację wskaźników przedziałów
funkcjami ciągłymi). Wobec tego $\mu=\nu$.
\end{proof}

\begin{stwierdzenie}[Odwracanie transformaty przy całkowalności]
Niech $X$ będzie rzeczywistą zmienną losową taką,
że $\int_\R|\varphi_X(t)|\,\dd\lambda_1(t)<\infty$.
Wtedy rozkład $X$ ma ograniczoną, ciągłą gęstość
\[
f_X(x)=\frac1{2\pi}\int_\R e^{-itx}\varphi_X(t)
\,\dd\lambda_1(t),\qquad x\in\R.
\]
\end{stwierdzenie}
\begin{proof}
Niech $g_\sigma$ będzie gaussowską gęstością z
poprzedniego dowodu. Splot $\mu_X*g_\sigma$ ma gęstość
$f_\sigma(x)=(2\pi)^{-1}\int e^{-itx}\varphi_X(t)
 e^{-\sigma^2t^2/2}\,\dd\lambda_1(t)$.
Ponieważ $|e^{-\sigma^2t^2/2}|\leq1$, zbieżność zdominowana
względem $|\varphi_X(t)|\,\dd\lambda_1(t)$ daje
$f_\sigma\to f$ jednostajnie na całej $\R$.
Funkcja $f$ jest ciągła i ograniczona z dominacji
oraz ciągłości fazy $e^{-itx}$.
Dla nieujemnej funkcji ciągłej $h$ o zwartym nośniku
mamy $\int h f_\sigma\,\dd\lambda_1\to
\int h f\,\dd\lambda_1$ z jednostajnej zbieżności,
a jednocześnie $\int h f_\sigma\,\dd\lambda_1\to\E h(X)$
z argumentu z dowodu jednoznaczności.
Zatem $\int h f\,\dd\lambda_1=\E h(X)$
dla wszystkich takich $h$, co wyznacza miarę
$\mu_X=f\lambda_1$. W szczególności $f\geq0$
i $\int f\,\dd\lambda_1=1$.
\end{proof}

\begin{przyklad}[Ciężkie ogony a funkcja charakterystyczna]
Dla standardowego rozkładu Cauchy'ego zachodzi
$\varphi_X(t)=e^{-|t|}$. Udowodnimy tę tożsamość
bez odwoływania się do rachunku residuów.
Niech $Z_1,Z_2$ będą niezależnymi zmiennymi standardowo
normalnymi. Z przedstawienia rozkładu Studenta dla
jednego stopnia swobody wiemy, że $C=Z_1/|Z_2|$
ma rozkład Cauchy'ego. Korzystając z niezależności
i funkcji charakterystycznej rozkładu normalnego,
możemy całkować najpierw po $Z_1$:
\[
\varphi_C(t)=\E\exp\!\left(-\frac{t^2}{2Z_2^2}\right)
=\sqrt{\frac2\pi}\int_0^\infty
\exp\!\left(-\frac{z^2}{2}-\frac{t^2}{2z^2}\right)
\,\dd\lambda_1(z).
\]
Aby obliczyć tę całkę, dla $a>0$ i $b\geq0$ połóżmy
$I(a,b)=\int_0^\infty e^{-az^2-b/z^2}\,\dd\lambda_1(z)$.
Dla $b>0$ wolno różniczkować względem $b$ pod całką
(zdominowanie na dowolnym zwartym podprzedziale
$(0,\infty)$). Podstawienie $u=\sqrt{b/a}/z$ daje
\[
\frac{\partial I}{\partial b}(a,b)
=-\int_0^\infty z^{-2}e^{-az^2-b/z^2}\,\dd\lambda_1(z)
=-\sqrt{\frac ab}\,I(a,b).
\]
Zatem $\partial_b\log I(a,b)=-\sqrt{a/b}$,
czyli $I(a,b)=C(a)e^{-2\sqrt{ab}}$.
Zdominowana zbieżność przy $b\downarrow0$ i całka
Gaussa wyznaczają $C(a)=I(a,0)=\sqrt\pi/(2\sqrt a)$.
Podstawiając $a=1/2$ i $b=t^2/2$ do powyższego
wzoru na $\varphi_C$, otrzymujemy
$\varphi_C(t)=e^{-|t|}$ również dla $t=0$.
Funkcja charakterystyczna istnieje więc pomimo braku
wartości oczekiwanej rozkładu Cauchy'ego.
\end{przyklad}

\chapter{Podstawowe rodziny rozkładów prawdopodobieństwa}

\section{Rozkład Bernoulliego i dwumianowy}
Pierwszym modelem losowym jest doświadczenie z dwiema
możliwościami: sukces albo porażka. Nie ma potrzeby zakładać,
że są równie prawdopodobne. Parametr $p$ określa szansę sukcesu,
a wskaźnik sukcesu ma rozkład Bernoulliego. Licząc sukcesy
w $n$ niezależnych próbach o tym samym parametrze, otrzymujemy
rozkład dwumianowy. Współczynnik $\binom nk$ odpowiada
liczbie różnych rozmieszczeń sukcesów, nie nowej losowości.

\begin{definicja}
Niech $p\in[0,1]$. Mówimy, że zmienna losowa $X$ ma
\emph{rozkład Bernoulliego} z parametrem $p$, jeżeli
$\Prob(X=1)=p$ i $\Prob(X=0)=1-p$.
Piszemy $X\sim\operatorname{Bern}(p)$.
\end{definicja}

Zmienna o rozkładzie Bernoulliego jest wskaźnikiem zdarzenia
o prawdopodobieństwie $p$, toteż
$\E X=p$, $\Var(X)=p(1-p)$, a
$M_X(t)=1-p+pe^t$ dla $t\in\R$.

\begin{definicja}
Niech $n\in\N\cup\{0\}$ i $p\in[0,1]$.
Zmienna losowa $X$ ma \emph{rozkład dwumianowy} $\Bin(n,p)$,
gdy dla $k=0,\ldots,n$ zachodzi
\begin{equation}\label{eq:binomial}
\Prob(X=k)=\binom nk p^k(1-p)^{n-k},
\end{equation}
a poza tym zakresem $\Prob(X=k)=0$. Przy $p=0$ lub $p=1$
przyjmujemy standardową interpretację iloczynów z potęgą zerową,
w tym $0^0=1$ w tym wzorze.
\end{definicja}

\begin{twierdzenie}[Suma prób Bernoulliego]
Niech $X_1,\ldots,X_n$ będą niezależnymi zmiennymi o rozkładzie
$\operatorname{Bern}(p)$, gdzie $n\geq1$. Wtedy
$S_n:=X_1+\cdots+X_n\sim\Bin(n,p)$ oraz
\[
\E S_n=np,\qquad \Var(S_n)=np(1-p),\qquad
M_{S_n}(t)=(1-p+pe^t)^n.
\]
\end{twierdzenie}
\begin{proof}
Dla ustalonego $k\in\{0,\ldots,n\}$ określmy,
które spośród $n$ prób zakończyły się sukcesem.
Jeśli dokładnie zbiór indeksów $I\subseteq\{1,\ldots,n\}$,
$|I|=k$, odpowiada sukcesom, to zdarzenie to ma postać
$\bigcap_{i\in I}\{X_i=1\}\cap
\bigcap_{j\notin I}\{X_j=0\}$.
Z niezależności jego prawdopodobieństwo jest równe
$p^k(1-p)^{n-k}$.

Różne zbiory $I$ opisują rozłączne zdarzenia,
a zdarzenie $\{S_n=k\}$ jest ich sumą. Ponieważ
liczba $k$-elementowych podzbiorów indeksów wynosi
$\binom nk$, z addytywności dostajemy
$\Prob(S_n=k)=\binom nk p^k(1-p)^{n-k}$.
Po zsumowaniu tego wyrażenia dla wszystkich $k$
otrzymujemy przez dwumian Newtona
$(p+(1-p))^n=1$, co stanowi dodatkową kontrolę normalizacji.

Dla wskaźnika pojedynczej próby
$\E X_i=p$, $\Var(X_i)=p(1-p)$.
Liniowość wartości oczekiwanej i addytywność
wariancji niezależnych składników dają
$\E S_n=np$ i $\Var(S_n)=np(1-p)$.
Wreszcie $M_{X_i}(t)=1-p+pe^t$, więc niezależność daje
$M_{S_n}(t)=\prod_{i=1}^nM_{X_i}(t)
=(1-p+pe^t)^n$.
\end{proof}

\begin{przyklad}[Liczba wadliwych elementów]
Jeżeli każdy z 20 elementów jest wadliwy niezależnie z
prawdopodobieństwem $0{,}02$, to liczba wadliwych elementów
ma rozkład $\Bin(20,0{,}02)$. Prawdopodobieństwo
braku wad wynosi $0{,}98^{20}$, a dokładnie jednej wady
$20\cdot0{,}02\cdot0{,}98^{19}$.
Jeżeli elementy pobierano z małej partii \emph{bez zwracania},
rozkład dwumianowy nie jest na ogół dokładny.
\end{przyklad}

\section{Rozkład geometryczny i ujemny dwumianowy}
Rozkład dwumianowy ustala liczbę prób i pyta o liczbę sukcesów.
Można odwrócić punkt widzenia: ustalić liczbę sukcesów i
pytać, ile prób potrzeba, by je osiągnąć. Pierwszy sukces
prowadzi do rozkładu geometrycznego, a $r$-ty do ujemnego
dwumianowego. Ich postać odzwierciedla zdarzenie
„najpierw odpowiednia liczba porażek, a następnie sukces”.

\begin{definicja}
Niech $p\in(0,1]$. Rozkładem geometrycznym $\Geom(p)$
na $\{1,2,\ldots\}$ nazywamy rozkład o funkcji masy
$\Prob(T=k)=(1-p)^{k-1}p$, $k\geq1$.
Zmienna $T$ opisuje numer pierwszej udanej próby
w ciągu niezależnych prób Bernoulliego z prawdopodobieństwem sukcesu $p$.
\end{definicja}

\begin{stwierdzenie}
Dla $T\sim\Geom(p)$ zachodzi $\Prob(T>k)=(1-p)^k$ dla
$k\in\N\cup\{0\}$, a także
$\E T=1/p$ i $\Var(T)=(1-p)/p^2$.
Jeśli $p\in(0,1)$, to dla $m,n\geq0$
$\Prob(T>m+n\mid T>m)=\Prob(T>n)$.
\end{stwierdzenie}
\begin{proof}
Zdarzenie $\{T>k\}$ oznacza, że w pierwszych $k$
próbach nie wystąpił sukces. Każda próba kończy się
niepowodzeniem z prawdopodobieństwem $q=1-p$;
przez niezależność otrzymujemy
$\Prob(T>k)=q^k$.

Dla $q<1$ szereg geometryczny
$\sum_{k=0}^\infty q^k=1/(1-q)$ jest zbieżny
jednostajnie na każdym przedziale $[0,q_0]$ z $q_0<1$.
Również szeregi pochodnych, począwszy od
$\sum_{k=1}^\infty kq^{k-1}$, są jednostajnie zbieżne
na takich przedziałach (ich wyrazy można ograniczyć przez
$kq_0^{k-1}$). Pozwala to różniczkować wyraz po wyrazie.
Otrzymujemy
$\sum_{k=1}^\infty kq^{k-1}=1/(1-q)^2$.
Różniczkując jeszcze raz lub wykorzystując tożsamość
$k^2=k(k-1)+k$, dostajemy
$\sum_{k=1}^\infty k^2q^{k-1}=(1+q)/(1-q)^3$.
Stąd
$\E T=p/(1-q)^2=1/p$,
$\E T^2=p(1+q)/(1-q)^3=(1+q)/p^2$
i $\Var(T)=\E T^2-(\E T)^2=q/p^2$.

Dla $p\in(0,1)$ mamy $q>0$ i
$\Prob(T>m)>0$, więc warunkowanie jest zdefiniowane.
Ponieważ $\{T>m+n\}\subseteq\{T>m\}$,
\[
\Prob(T>m+n\mid T>m)
=\frac{q^{m+n}}{q^m}=q^n=\Prob(T>n).
\]
Przy $p=1$ zdarzenie $\{T>m\}$ dla $m\geq1$
ma masę zero i powyższe warunkowanie nie ma sensu,
co wyjaśnia ograniczenie w tezie.
\end{proof}

\begin{uwaga}
Spotyka się także rozkład geometryczny na $\{0,1,\ldots\}$,
gdzie zmienna oznacza liczbę niepowodzeń \emph{przed} pierwszym sukcesem.
Jest to rozkład $T-1$, a jego średnia wynosi $(1-p)/p$.
Dlatego zawsze należy podać konwencję numerowania.
\end{uwaga}

\begin{definicja}
Dla $r\in\N$ i $p\in(0,1]$ rozkładem ujemnym dwumianowym
(w konwencji liczby prób do $r$-tego sukcesu) nazywamy rozkład
zmiennej $T_r$ określonej przez
$\Prob(T_r=n)=\binom{n-1}{r-1}p^r(1-p)^{n-r}$, $n\geq r$.
\end{definicja}

\begin{stwierdzenie}
Jeżeli $T_r$ jest numerem próby, w której następuje $r$-ty sukces
w ciągu niezależnych prób Bernoulliego z parametrem $p>0$,
to ma powyższy rozkład, a ponadto
$\E T_r=r/p$ i $\Var(T_r)=r(1-p)/p^2$.
\end{stwierdzenie}
\begin{proof}
Aby $r$-ty sukces nastąpił dokładnie w próbie numer $n$,
próba $n$ musi zakończyć się sukcesem, a wśród pierwszych
$n-1$ prób musi znaleźć się dokładnie $r-1$ sukcesów.
Dla każdej konkretnej konfiguracji tych prób, przez
niezależność, prawdopodobieństwo wynosi
$p^r(1-p)^{n-r}$. Konfiguracji jest
$\binom{n-1}{r-1}$; są rozłączne, co daje wzór na masę.

Wykażmy starannie niezależność czasów oczekiwania,
której nie wolno wywnioskować wyłącznie z intuicji.
Niech $W_1$ będzie liczbą prób do pierwszego sukcesu,
a dla $j\geq2$ niech $W_j$ oznacza liczbę kolejnych prób
pomiędzy sukcesami numer $j-1$ i $j$, licząc próbę
zakończoną sukcesem. Dla dowolnych $k_1,\ldots,k_r\geq1$
zdarzenie $\{W_1=k_1,\ldots,W_r=k_r\}$ oznacza
jednoznacznie określony ciąg
$k_1+\cdots+k_r$ wyników: w każdym bloku $k_j-1$
niepowodzeń, po których następuje sukces.
Przez niezależność pojedynczych prób jego
prawdopodobieństwo jest równe
$\prod_{j=1}^r p(1-p)^{k_j-1}$.
Jest to iloczyn mas rozkładów $\Geom(p)$;
zatem $W_1,\ldots,W_r$ są niezależne i jednakowo
rozłożone. Ponieważ $T_r=W_1+\cdots+W_r$,
z liniowości wartości oczekiwanej i addytywności
wariancji otrzymujemy
$\E T_r=r/p$ i $\Var(T_r)=r(1-p)/p^2$.
\end{proof}

\section{Rozkład hipergeometryczny}
Liczba sukcesów w próbie bez zwracania różni się od liczby
sukcesów w niezależnych próbach Bernoulliego. Każdy kolejny
wynik zmienia skład dostępnej populacji. Zliczając możliwe
podzbiory zamiast ciągów niezależnych doświadczeń, otrzymujemy
rozkład hipergeometryczny. Ujemna kowariancja wyników kolejnych
losowań tłumaczy, dlaczego wariancja jest mniejsza niż w modelu
dwumianowym z taką samą średnią.

\begin{definicja}
Niech $N\in\N$, $K\in\{0,\ldots,N\}$
i $n\in\{0,\ldots,N\}$.
Zmienna losowa $X$ ma rozkład hipergeometryczny
$\operatorname{HGeom}(N,K,n)$, jeżeli
\begin{equation}\label{eq:hypergeom}
\Prob(X=k)=\frac{\binom Kk\binom{N-K}{n-k}}{\binom Nn}
\quad(k\in\Z),
\end{equation}
przy konwencji, że symbol Newtona jest zerem poza właściwym zakresem.
Jest to liczba wyróżnionych obiektów w próbie $n$ elementów
losowanej jednostajnie bez zwracania z populacji $N$ elementów,
w której $K$ obiektów jest wyróżnionych.
\end{definicja}

\begin{twierdzenie}
Jeżeli $X\sim\operatorname{HGeom}(N,K,n)$, to
$\E X=nK/N$. Dla $N>1$ zachodzi
\[
\Var(X)=n\frac KN\left(1-\frac KN\right)\frac{N-n}{N-1}.
\]
Dla $N=1$ wariancja wynosi zero.
\end{twierdzenie}
\begin{proof}
Realizujmy losowanie jako pierwszych $n$ pozycji jednostajnej
losowej permutacji $N$ różnych elementów. Dla
$j=1,\ldots,n$ niech $I_j$ oznacza, że na pozycji $j$
znajduje się obiekt wyróżniony. Wówczas
$X=\sum_{j=1}^n I_j$.
Każda pozycja ma ten sam rozkład i dokładnie $K$
z $N$ obiektów jest wyróżnionych, toteż
$\Prob(I_j=1)=K/N=:p$ i $\E I_j=p$.
Liniowość wartości oczekiwanej daje $\E X=np$.

Dla $N>1$ i różnych $i,j$ obliczamy prawdopodobieństwo
dwóch wyróżnionych obiektów przez warunkowanie:
$\Prob(I_i=I_j=1)=(K/N)((K-1)/(N-1))$,
przy czym dla $K=0$ lub $K=1$ wynik ten jest zerem.
Zatem
\[
\Cov(I_i,I_j)=\frac{K(K-1)}{N(N-1)}-
\left(\frac KN\right)^2
=-\frac{p(1-p)}{N-1}.
\]
Każdy wskaźnik ma wariancję $p(1-p)$, a w sumie
występuje $n$ wariancji i $n(n-1)$ kowariancji
uporządkowanych par różnych indeksów. Otrzymujemy
\[
\Var(X)=np(1-p)-\frac{n(n-1)}{N-1}p(1-p)
=np(1-p)\frac{N-n}{N-1}.
\]
Dla $N=1$ dopuszczalne $n$ to $0$ lub $1$,
a liczba obiektów wyróżnionych w próbie jest stała,
więc wariancja wynosi zero.
\end{proof}

\begin{uwaga}
Względem wariancji rozkładu $\Bin(n,K/N)$ pojawia się
czynnik korekcyjny $(N-n)/(N-1)$, odzwierciedlający ujemną
zależność losowań bez zwracania. Jeśli $N$ jest duże w porównaniu
z $n$, oba modele mogą dawać zbliżone wyniki.
\end{uwaga}


\begin{przyklad}[Taka sama średnia, ale różna wariancja]
W populacji znajduje się $N=10$ kul, z których $K=5$
jest czerwonych. Losujemy $n=4$ kule bez zwracania.
Liczba czerwonych ma rozkład hipergeometryczny,
a więc średnią $4\cdot(5/10)=2$ oraz wariancję
$4\cdot(1/2)(1/2)\cdot(10-4)/(10-1)=2/3$.
Gdybyśmy losowali ze zwracaniem, liczba czerwonych
miałaby rozkład $\Bin(4,1/2)$ o tej samej średniej $2$,
ale wariancji $4\cdot(1/2)(1/2)=1$.
Przy losowaniu bez zwracania pojawienie się czerwonej kuli
zmniejsza szansę czerwonej kuli w kolejnej próbie.
Ta ujemna zależność ogranicza rozrzut liczby sukcesów.
\end{przyklad}

\section{Rozkład Poissona}
Jeśli mamy bardzo wiele prób, lecz sukces w każdej z nich jest
bardzo rzadki, rozkład dwumianowy bywa niewygodny w obliczeniach.
Gdy liczba prób rośnie, prawdopodobieństwo sukcesu maleje, a
oczekiwana liczba sukcesów pozostaje stała, otrzymujemy
rozkład Poissona. Parametr $\lambda$ opisuje jednocześnie
średnią i wariancję. Samo występowanie zdarzeń „w czasie” nie
wystarcza jednak do przyjęcia rozkładu Poissona: odpowiednia
struktura zależności musi wynikać z modelu.

\begin{definicja}
Niech $\lambda>0$. Zmienna losowa $X$ ma rozkład Poissona
$\Poi(\lambda)$, jeżeli
$\Prob(X=k)=e^{-\lambda}\lambda^k/k!$ dla
$k\in\{0,1,2,\ldots\}$.
\end{definicja}

\begin{stwierdzenie}
Dla $X\sim\Poi(\lambda)$ zachodzi
$\E X=\lambda$, $\Var(X)=\lambda$ oraz
$M_X(t)=\exp(\lambda(e^t-1))$ dla $t\in\R$.
\end{stwierdzenie}
\begin{proof}
Najpierw sprawdzamy, że wzór definiuje rozkład:
\[
\sum_{k=0}^\infty e^{-\lambda}\frac{\lambda^k}{k!}
=e^{-\lambda}e^\lambda=1.
\]
Dla każdego $t\in\R$ funkcja tworząca momenty jest
skończona, ponieważ
\[
M_X(t)=\sum_{k=0}^\infty e^{tk}
 e^{-\lambda}\frac{\lambda^k}{k!}
=e^{-\lambda}\sum_{k=0}^\infty\frac{(\lambda e^t)^k}{k!}
=\exp\bigl(\lambda(e^t-1)\bigr).
\]
Ponieważ funkcja jest skończona w otoczeniu zera,
możemy odzyskać momenty przez różniczkowanie:
$M_X'(t)=\lambda e^tM_X(t)$,
a $M_X''(t)=(\lambda e^t+\lambda^2e^{2t})M_X(t)$.
W szczególności $\E X=M_X'(0)=\lambda$,
$\E X^2=M_X''(0)=\lambda+\lambda^2$,
stąd $\Var(X)=\lambda$.
\end{proof}

\begin{twierdzenie}[Prawo rzadkich zdarzeń]
Niech $\lambda>0$ i niech $X_n\sim\Bin(n,\lambda/n)$
dla całkowitych $n\geq\lceil\lambda\rceil$ z $n\geq1$.
Wtedy dla każdego ustalonego $k\in\N\cup\{0\}$
\[
\lim_{n\to\infty}\Prob(X_n=k)
=e^{-\lambda}\frac{\lambda^k}{k!}.
\]
\end{twierdzenie}
\begin{proof}
Dla ustalonego $k\geq0$ i dostatecznie dużego $n$
mamy $n\geq k$ oraz $\lambda/n\leq1$. Z definicji
rozkładu dwumianowego
\[
\Prob(X_n=k)=\binom nk\left(\frac\lambda n\right)^k
\left(1-\frac\lambda n\right)^{n-k}.
\]
Rozdzielamy to wyrażenie na czynniki o prostych granicach:
\[
\Prob(X_n=k)=\frac{\lambda^k}{k!}
\underbrace{\prod_{j=0}^{k-1}\left(1-\frac jn\right)}_{\longrightarrow1}
\underbrace{\left(1-\frac\lambda n\right)^n}_{\longrightarrow e^{-\lambda}}
\underbrace{\left(1-\frac\lambda n\right)^{-k}}_{\longrightarrow1}.
\]
Pierwszy iloczyn ma stałą liczbę czynników (bo $k$ nie
zależy od $n$). Granica środkowego czynnika jest
klasyczną granicą wykładniczą, możliwą też do
wyprowadzenia przez logarytmowanie:
$n\log(1-\lambda/n)\to-\lambda$.
Ostatni czynnik zmierza do jedności.
Iloczyn granic daje $e^{-\lambda}\lambda^k/k!$.
Twierdzenie stwierdza na tym etapie zbieżność
każdej ustalonej masy punktowej; przejście do
języka zbieżności według rozkładu omówimy osobno.
\end{proof}

\begin{przyklad}[Zliczanie defektów]
Jeśli w dużej liczbie niezależnych prób każdy defekt jest rzadki,
a oczekiwana liczba defektów wynosi około $\lambda$, to rozkład
Poissona może przybliżać rozkład liczby defektów.
Nie jest jednak automatycznym modelem dowolnych zliczeń:
przy klastrowaniu zdarzeń lub nadmiernej zmienności
wariancja może różnić się istotnie od średniej.
\end{przyklad}

\begin{stwierdzenie}[Dodawanie rozkładów Poissona]
Jeśli $X\sim\Poi(\lambda)$ i $Y\sim\Poi(\eta)$ są niezależne,
gdzie $\lambda,\eta>0$, to $X+Y\sim\Poi(\lambda+\eta)$.
\end{stwierdzenie}
\begin{proof}
Dla $k\geq0$ całkowitego z niezależności wynika
\[
\Prob(X+Y=k)
=\sum_{j=0}^k e^{-\lambda}\frac{\lambda^j}{j!}
 e^{-\eta}\frac{\eta^{k-j}}{(k-j)!}
=e^{-(\lambda+\eta)}\frac{(\lambda+\eta)^k}{k!},
\]
gdzie ostatnią równość daje dwumian Newtona.
\end{proof}

\begin{przyklad}[Oczekiwanie na zebranie wszystkich kuponów]
Niech $N\geq1$ i załóżmy, że w kolejnych niezależnych próbach
losujemy jednostajnie jedną z $N$ etykiet. Niech $T$ oznacza
liczbę prób potrzebną, aby zobaczyć wszystkie etykiety.
Po zebraniu $j$ różnych etykiet, gdzie $0\leq j<N$,
każda kolejna próba daje nową etykietę z prawdopodobieństwem
$p_j=(N-j)/N$. Liczba dalszych prób do kolejnej nowej etykiety
ma rozkład geometryczny $\Geom(p_j)$, a zatem wartość
oczekiwaną $N/(N-j)$. Sumując oczekiwane czasy przejść
między kolejnymi stanami liczby poznanych etykiet,
otrzymujemy
$\E T=\sum_{j=0}^{N-1}N/(N-j)=N\sum_{k=1}^N1/k$.
Nie trzeba przy tym zakładać, że wszystkie czasy oczekiwania
są z góry wzajemnie niezależne: do obliczenia średniej
wystarcza liniowość wartości oczekiwanej.
\end{przyklad}

\section{Rozkład jednostajny ciągły i wykładniczy}
W przestrzeni ciągłej nie da się przypisać jednakowej dodatniej
masy wszystkim punktom odcinka. Jednostajność oznacza więc,
że prawdopodobieństwo przedziału jest proporcjonalne do jego
długości. Rozkład wykładniczy odpowiada innej intuicji:
jeśli wydarzenie nie zaszło przez czas $s$, dodatkowy czas
oczekiwania ma mieć tę samą regułę co na początku.
Ta własność braku pamięci wyznacza wykładniczy kształt ogona.

\begin{definicja}
Dla $a,b\in\R$ z $a<b$ rozkładem jednostajnym
$\Unif(a,b)$ nazywamy rozkład o gęstości
$f(x)=\ind{\{a<x<b\}}/(b-a)$.
\end{definicja}

\begin{stwierdzenie}
Jeżeli $X\sim\Unif(a,b)$, to
\[
F_X(x)=\begin{cases}
0,&x\leq a,\\
\dfrac{x-a}{b-a},&a<x<b,\\
1,&x\geq b,
\end{cases}
\qquad
\E X=\frac{a+b}{2},\qquad
\Var(X)=\frac{(b-a)^2}{12}.
\]
\end{stwierdzenie}
\begin{proof}
Gęstość $f(x)=(b-a)^{-1}\ind{\{a<x<b\}}$
jest nieujemna i ma całkę równą jeden.
Dla $x\leq a$ przedział $(-\infty,x]$ nie przecina
nośnika gęstości z dodatnią miarą, więc $F_X(x)=0$.
Dla $a<x<b$ całkujemy stałą gęstość od $a$ do $x$,
co daje $F_X(x)=(x-a)/(b-a)$.
Dla $x\geq b$ całka obejmuje cały nośnik i wynosi jeden.

Momenty obliczamy z definicji:
\[
\E X=\frac1{b-a}\int_a^b x\,\dd\lambda_1(x)
=\frac{b^2-a^2}{2(b-a)}=\frac{a+b}2,
\]
\[
\E X^2=\frac1{b-a}\int_a^b x^2\,\dd\lambda_1(x)
=\frac{b^3-a^3}{3(b-a)}=\frac{a^2+ab+b^2}3.
\]
Po odjęciu $((a+b)/2)^2$ od drugiego momentu
i uproszczeniu otrzymujemy
$\Var(X)=(b-a)^2/12$.
\end{proof}

\begin{definicja}
Dla $\lambda>0$ rozkładem wykładniczym $\operatorname{Exp}(\lambda)$
z parametrem intensywności $\lambda$ nazywamy rozkład
o gęstości $f(x)=\lambda e^{-\lambda x}\ind{\{x\geq0\}}$.
\end{definicja}

\begin{stwierdzenie}
Jeżeli $X\sim\operatorname{Exp}(\lambda)$, to dla $x\geq0$
$F_X(x)=1-e^{-\lambda x}$, a
$\E X=1/\lambda$ i $\Var(X)=1/\lambda^2$.
Dla $s,t\geq0$ zachodzi
$\Prob(X>s+t\mid X>s)=\Prob(X>t)$.
\end{stwierdzenie}
\begin{proof}
Dla $x<0$ mamy $F_X(x)=0$. Dla $x\geq0$ całkowanie
gęstości daje
$F_X(x)=\int_0^x\lambda e^{-\lambda t}\,\dd\lambda_1(t)
=1-e^{-\lambda x}$.
Zatem ogon ma postać
$\Prob(X>x)=e^{-\lambda x}$.

Aby obliczyć momenty, wykonujemy zmianę zmiennych
$u=\lambda x$ i używamy rekurencji całki gamma
(albo dwukrotnego całkowania przez części):
\[
\E X^k=\int_0^\infty x^k\lambda e^{-\lambda x}
\,\dd\lambda_1(x)=\frac{k!}{\lambda^k}
\qquad(k=1,2).
\]
Otrzymujemy $\E X=1/\lambda$ i
$\Var(X)=2/\lambda^2-1/\lambda^2=1/\lambda^2$.

Wreszcie $\{X>s+t\}\subseteq\{X>s\}$ oraz
$\Prob(X>s)=e^{-\lambda s}>0$, więc
\[
\Prob(X>s+t\mid X>s)
=\frac{\Prob(X>s+t)}{\Prob(X>s)}
=\frac{e^{-\lambda(s+t)}}{e^{-\lambda s}}
=e^{-\lambda t}=\Prob(X>t).
\]
Interpretacja: wiedza o dotychczasowym czasie oczekiwania
nie zmienia rozkładu dodatkowego czasu.
\end{proof}

\begin{uwaga}
Rozkład wykładniczy jest ciągłym odpowiednikiem rozkładu
geometrycznego ze względu na własność braku pamięci.
W procesach stochastycznych pojawi się jako rozkład czasów
oczekiwania w procesie Poissona.
\end{uwaga}

\begin{przyklad}[Generowanie zmiennej wykładniczej z jednostajnej]
Niech $U\sim\Unif(0,1)$ i $\lambda>0$. Zmienna
$X=-\log(1-U)/\lambda$ ma rozkład $\operatorname{Exp}(\lambda)$,
ponieważ dla $x\geq0$ zachodzi
$\Prob(X\leq x)=\Prob(U\leq1-e^{-\lambda x})=1-e^{-\lambda x}$.
Jest to szczególny przypadek metody odwrotnej dystrybuanty,
opartej na kwantylu uogólnionym. Pozwala ona przejść od
jednostajnej liczby pseudolosowej do zadanej rodziny rozkładów.
\end{przyklad}

\section{Rozkład normalny}
Rozkład normalny wyróżnia jego postać analityczna oraz stabilność
względem dodawania niezależnych zmiennych normalnych. Służy też
jako rozkład graniczny dla odpowiednio unormowanych sum
wielu drobnych składników, co uzasadnimy później przy CTG.
Na razie nie będziemy motywować jego gęstości wyłącznie
rysunkiem dzwonu: wyprowadzimy normalizację przez całkę Gaussa
oraz sprawdzimy wzory na momenty i własności sum.

\begin{definicja}
Niech $\mu\in\R$ i $\sigma>0$. Rozkładem normalnym
$\Normal(\mu,\sigma^2)$ nazywamy rozkład o gęstości
\begin{equation}\label{eq:normal}
f(x)=\frac{1}{\sigma\sqrt{2\pi}}
\exp\!\left(-\frac{(x-\mu)^2}{2\sigma^2}\right),
\qquad x\in\R.
\end{equation}
Rozkład $\Normal(0,1)$ nazywamy standardowym normalnym.
Jego dystrybuantę oznaczamy przez $\Phi$.
\end{definicja}

\begin{lemat}[Całka Gaussa]
Zachodzi $\int_\R e^{-x^2/2}\,\dd\lambda_1(x)=\sqrt{2\pi}$.
\end{lemat}
\begin{proof}
Niech $I=\int_\R e^{-x^2/2}\,\dd\lambda_1(x)$.
Z twierdzenia Tonellego i przejścia do współrzędnych biegunowych
otrzymujemy
\[
I^2=\int_{\R^2}e^{-(x^2+y^2)/2}\,\dd\lambda_2(x,y)
=\int_0^{2\pi}\!\int_0^\infty e^{-r^2/2}r\,
\dd\lambda_1(r)\,\dd\lambda_1(\varphi)=2\pi.
\]
Ponieważ $I>0$, mamy $I=\sqrt{2\pi}$.
\end{proof}

\begin{stwierdzenie}[Momenty i funkcja tworząca]
Jeżeli $X\sim\Normal(\mu,\sigma^2)$, to
$\E X=\mu$, $\Var(X)=\sigma^2$ oraz
$M_X(t)=\exp(\mu t+\sigma^2t^2/2)$ dla wszystkich $t\in\R$.
Ponadto $aX+b\sim\Normal(a\mu+b,a^2\sigma^2)$
dla $a\ne0$ i $b\in\R$; dla $a=0$ rozkład jest skupiony
w punkcie $b$.
\end{stwierdzenie}
\begin{proof}
Niech $Z=(X-\mu)/\sigma$. Z twierdzenia o
przekształceniach zmiennych wynika, że $Z$ ma gęstość
$g(z)=(2\pi)^{-1/2}e^{-z^2/2}$.
Dla $t\in\R$ obliczamy
\[
M_X(t)=\frac{e^{\mu t}}{\sqrt{2\pi}}\int_\R
\exp\!\left(\sigma tz-\frac{z^2}2\right)
\,\dd\lambda_1(z).
\]
Tożsamość
$\sigma tz-z^2/2=-(z-\sigma t)^2/2+\sigma^2t^2/2$
pozwala wyjąć z całki czynnik
$e^{\sigma^2t^2/2}$. Pozostała całka z gęstości
Gaussa przesuniętej o $\sigma t$ wynosi jeden.
Otrzymujemy więc
$M_X(t)=\exp(\mu t+\sigma^2t^2/2)$
dla wszystkich $t\in\R$.
Ponieważ funkcja ta jest skończona na całej prostej,
możemy różniczkować pod całką. Z równości
$M_X'(0)=\mu$ i $M_X''(0)=\mu^2+\sigma^2$
wynikają wzory na średnią i wariancję.
Dla $a\ne0$ transformacja gęstości przez funkcję
$x\mapsto ax+b$ daje rozkład
$\Normal(a\mu+b,a^2\sigma^2)$;
gdy $a=0$, zmienna $aX+b$ jest stale równa $b$.
\end{proof}

\begin{stwierdzenie}[Suma niezależnych zmiennych normalnych]
Jeżeli $X\sim\Normal(\mu_1,\sigma_1^2)$ i
$Y\sim\Normal(\mu_2,\sigma_2^2)$ są niezależne,
gdzie $\sigma_1,\sigma_2>0$, to
$X+Y\sim\Normal(\mu_1+\mu_2,\sigma_1^2+\sigma_2^2)$.
\end{stwierdzenie}
\begin{proof}
Z niezależności i wzoru na splot wynika dla $s\in\R$
\[
f_{X+Y}(s)=\frac1{2\pi\sigma_1\sigma_2}
\int_\R\exp\!\left(-\frac{(x-\mu_1)^2}{2\sigma_1^2}
-\frac{(s-x-\mu_2)^2}{2\sigma_2^2}\right)\,\dd\lambda_1(x).
\]
Oznaczmy $v=\sigma_1^2+\sigma_2^2$ oraz
$c(s)=(\mu_1\sigma_2^2+(s-\mu_2)\sigma_1^2)/v$.
Uzupełnienie kwadratu daje
\[
\frac{(x-\mu_1)^2}{\sigma_1^2}
+\frac{(s-x-\mu_2)^2}{\sigma_2^2}
=\frac{v}{\sigma_1^2\sigma_2^2}(x-c(s))^2
+\frac{(s-\mu_1-\mu_2)^2}{v}.
\]
Całkując pierwszy czynnik za pomocą całki Gaussa, otrzymujemy
\[
f_{X+Y}(s)=\frac{1}{\sqrt{2\pi v}}
\exp\!\left(-\frac{(s-\mu_1-\mu_2)^2}{2v}\right),
\]
czyli żądaną gęstość.
\end{proof}

\begin{przyklad}[Standaryzacja]
Jeżeli $X\sim\Normal(\mu,\sigma^2)$, to
$Z:=(X-\mu)/\sigma\sim\Normal(0,1)$.
W konsekwencji dla $a<b$
$\Prob(a<X\leq b)=\Phi((b-\mu)/\sigma)-\Phi((a-\mu)/\sigma)$.
Dystrybuanta $\Phi$ nie ma pierwotnej wyrażalnej przez skończone
kombinacje standardowych funkcji elementarnych; jej wartości
wyznacza się numerycznie lub za pomocą funkcji błędu.
\end{przyklad}

\section{Funkcja gamma i funkcja beta}
Próby normalizowania funkcji takich jak
$x^{a-1}e^{-x}$ albo $x^{a-1}(1-x)^{b-1}$ prowadzą
do całek, które nie zawsze są elementarne. Funkcje gamma i beta
zapisują te całki w sposób systematyczny. Nie są jedynie
technicznym dodatkiem: ich tożsamości pozwalają wyprowadzać
wzory na momenty i zależności między rozkładami gamma,
chi-kwadrat, beta i t-Studenta.

\begin{definicja}
Dla $a,b>0$ definiujemy funkcję gamma i funkcję beta Eulera wzorami
\[
\Gamma(a):=\int_0^\infty t^{a-1}e^{-t}\,\dd\lambda_1(t),
\qquad
\mathrm B(a,b):=\int_0^1 t^{a-1}(1-t)^{b-1}\,\dd\lambda_1(t).
\]
Obie całki są skończone i dodatnie dla podanych parametrów.
\end{definicja}

\begin{lemat}[Podstawowe tożsamości]
Dla $a,b>0$ zachodzą
$\Gamma(a+1)=a\Gamma(a)$ oraz
$\mathrm B(a,b)=\Gamma(a)\Gamma(b)/\Gamma(a+b)$.
Ponadto $\Gamma(1)=1$, $\Gamma(1/2)=\sqrt\pi$
i $\Gamma(n+1)=n!$ dla $n\in\N\cup\{0\}$.
\end{lemat}
\begin{proof}
Najpierw uzasadnijmy zbieżność całki definiującej
$\Gamma(a)$. Przy zerze funkcja $t^{a-1}e^{-t}$
zachowuje się jak $t^{a-1}$, a ta jest całkowalna
na $(0,1)$ wtedy i tylko wtedy, gdy $a>0$.
Na $(1,\infty)$ czynnik wykładniczy zapewnia
całkowalność każdego wielomianowego czynnika.
Podobnie całka beta jest zbieżna, ponieważ
przy końcach przedziału czynniki zachowują się
jak $t^{a-1}$ i $(1-t)^{b-1}$.

Całkując przez części i wykorzystując znikanie
wyrazów brzegowych $t^a e^{-t}$ dla $t\to0+$
i $t\to\infty$, otrzymujemy
\[
\Gamma(a+1)=\int_0^\infty t^ae^{-t}\,\dd\lambda_1(t)
=a\int_0^\infty t^{a-1}e^{-t}\,\dd\lambda_1(t)
=a\Gamma(a).
\]
Ponieważ $\Gamma(1)=\int_0^\infty e^{-t}\,\dd\lambda_1(t)=1$,
indukcja daje $\Gamma(n+1)=n!$.
Podstawienie $t=u^2$ i symetria całki Gaussa dają
$\Gamma(1/2)=2\int_0^\infty e^{-u^2}\,\dd\lambda_1(u)
=\sqrt\pi$.

Aby wyprowadzić relację z funkcją beta, zapiszmy
przez Tonellego iloczyn dwóch dodatnich całek:
\[
\Gamma(a)\Gamma(b)
=\int_{(0,\infty)^2}u^{a-1}v^{b-1}e^{-u-v}
\,\dd\lambda_2(u,v).
\]
Wprowadzamy $r=u+v>0$ i $w=u/(u+v)\in(0,1)$.
Odwrotnie, $u=rw$, $v=r(1-w)$; wyznacznik macierzy
pochodnych odwzorowania $(r,w)\mapsto(u,v)$
ma moduł $r$. Funkcja podcałkowa wraz z jakobianem
przyjmuje postać
$r^{a+b-1}e^{-r}w^{a-1}(1-w)^{b-1}$.
Możemy więc ponownie rozdzielić całkę:
\[
\Gamma(a)\Gamma(b)
=\left(\int_0^\infty r^{a+b-1}e^{-r}\,\dd\lambda_1(r)\right)
\left(\int_0^1w^{a-1}(1-w)^{b-1}\,\dd\lambda_1(w)\right)
=\Gamma(a+b)\mathrm B(a,b).
\]
Dzielenie przez dodatnie $\Gamma(a+b)$ kończy dowód.
\end{proof}

\section{Rozkład gamma}
Jeśli czas do jednego zdarzenia ma rozkład wykładniczy, to
suma kilku niezależnych czasów oczekiwania prowadzi do rodziny
gamma. Parametr kształtu opisuje m.in. liczbę składników
w przypadkach całkowitych, a parametr intensywności reguluje
skalę czasową. Funkcję gamma można traktować jako naturalne
przedłużenie silni do argumentów dodatnich.

\begin{definicja}
Niech $\alpha,\lambda>0$. Rozkładem gamma $\Gam(\alpha,\lambda)$
o parametrze kształtu $\alpha$ i \emph{intensywności} $\lambda$
nazywamy rozkład o gęstości
\begin{equation}\label{eq:gamma}
f(x)=\frac{\lambda^\alpha}{\Gamma(\alpha)}
 x^{\alpha-1}e^{-\lambda x}\ind{\{x>0\}}.
\end{equation}
\end{definicja}

\begin{stwierdzenie}
Dla $X\sim\Gam(\alpha,\lambda)$ zachodzi
\[
\E X=\frac\alpha\lambda,\qquad
\Var(X)=\frac\alpha{\lambda^2},\qquad
M_X(t)=\left(\frac\lambda{\lambda-t}\right)^\alpha\quad(t<\lambda).
\]
Jeżeli $X\sim\Gam(\alpha,\lambda)$ i
$Y\sim\Gam(\beta,\lambda)$ są niezależne, to
$X+Y\sim\Gam(\alpha+\beta,\lambda)$.
\end{stwierdzenie}
\begin{proof}
Najpierw sprawdźmy normalizację gęstości. Dla
$\alpha,\lambda>0$, po podstawieniu $u=\lambda x$,
\[
\int_0^\infty\frac{\lambda^\alpha}{\Gamma(\alpha)}
 x^{\alpha-1}e^{-\lambda x}\,\dd\lambda_1(x)
=\frac1{\Gamma(\alpha)}\int_0^\infty
u^{\alpha-1}e^{-u}\,\dd\lambda_1(u)=1.
\]
Dla $r\geq0$ ten sam rachunek daje
$\E X^r=\Gamma(\alpha+r)/(\lambda^r\Gamma(\alpha))$.
Korzystając z rekurencji funkcji gamma dla $r=1,2$,
mamy $\E X=\alpha/\lambda$ i
$\E X^2=\alpha(\alpha+1)/\lambda^2$.
Różnica drugiego momentu i kwadratu średniej daje
$\Var(X)=\alpha/\lambda^2$.

Dla $t<\lambda$ obliczamy
\[
M_X(t)=\frac{\lambda^\alpha}{\Gamma(\alpha)}
\int_0^\infty x^{\alpha-1}e^{-(\lambda-t)x}
\,\dd\lambda_1(x)
=\left(\frac{\lambda}{\lambda-t}\right)^\alpha,
\]
gdzie użyliśmy kolejno zmiany $u=(\lambda-t)x$
i definicji funkcji gamma.

Dla sumy niezależnych zmiennych gamma z tą samą
intensywnością korzystamy ze splotu:
\[
f_{X+Y}(s)=\frac{\lambda^{\alpha+\beta}e^{-\lambda s}}
{\Gamma(\alpha)\Gamma(\beta)}
\int_0^s x^{\alpha-1}(s-x)^{\beta-1}\,\dd\lambda_1(x),
\quad s>0.
\]
Podstawiamy $x=sw$, co daje czynnik
$s^{\alpha+\beta-1}$ i całkę $\mathrm B(\alpha,\beta)$.
Tożsamość beta--gamma redukuje współczynnik do
$1/\Gamma(\alpha+\beta)$. Otrzymujemy dokładnie
gęstość $\Gam(\alpha+\beta,\lambda)$.
Wspólna intensywność jest istotna: dla różnych
intensywności ten prosty wzór na sumę nie obowiązuje.
\end{proof}

\begin{uwaga}
W literaturze używa się również parametru \emph{skali}
$\beta=1/\lambda$, zapisując gęstość jako
$x^{\alpha-1}e^{-x/\beta}/(\Gamma(\alpha)\beta^\alpha)$.
W tej książce drugi parametr rozkładu gamma jest zawsze
\emph{intensywnością}, a nie skalą.
Rozkład $\operatorname{Exp}(\lambda)$ jest równy $\Gam(1,\lambda)$.
\end{uwaga}

\section{Rozkład chi-kwadrat}
Suma kwadratów niezależnych standardowych zmiennych normalnych
jest naturalną miarą kwadratowego odchylenia wektora losowego
od początku układu. Jej rozkład należy do rodziny gamma i
nazywa się chi-kwadrat. Dla całkowitej liczby stopni swobody
ma ona bezpośredni sens geometryczny; definicję przez gamma
można natomiast rozszerzyć na dowolne dodatnie $\nu$.

\begin{definicja}
Niech $\nu>0$. Rozkładem chi-kwadrat z $\nu$ stopniami swobody,
oznaczanym $\chi^2_\nu$, nazywamy rozkład
$\Gam(\nu/2,1/2)$, a więc rozkład o gęstości
\begin{equation}\label{eq:chisq}
f_\nu(x)=\frac{1}{2^{\nu/2}\Gamma(\nu/2)}
 x^{\nu/2-1}e^{-x/2}\ind{\{x>0\}}.
\end{equation}
Dla naturalnego $\nu$ liczba $\nu$ jest liczbą stopni swobody
w poniższej reprezentacji za pomocą zmiennych normalnych.
\end{definicja}

\begin{twierdzenie}[Suma kwadratów zmiennych normalnych]
Niech $\nu\in\N$ i niech $Z_1,\ldots,Z_\nu$ będą niezależne
i mają rozkład $\Normal(0,1)$. Wówczas
$U:=\sum_{j=1}^\nu Z_j^2\sim\chi^2_\nu$.
Dla dowolnego rzeczywistego $\nu>0$, jeżeli $U\sim\chi^2_\nu$,
to $\E U=\nu$ oraz $\Var(U)=2\nu$.
\end{twierdzenie}
\begin{proof}
Wystarczy najpierw obliczyć rozkład $Z^2$ dla jednej
standardowej zmiennej normalnej $Z$.
Funkcja $g(z)=z^2$ jest różnowartościowa na
$(-\infty,0)$ i $(0,\infty)$, a dla $u>0$ ma
przeciwobrazy $z_\pm=\pm\sqrt u$. Pochodna ma
moduł $|g'(z_\pm)|=2\sqrt u$. Ze wzoru
na transformację o dwóch gałęziach
\[
f_{Z^2}(u)=\frac{\varphi(\sqrt u)}{2\sqrt u}
+\frac{\varphi(-\sqrt u)}{2\sqrt u}
=\frac{e^{-u/2}}{\sqrt{2\pi u}}
\qquad(u>0).
\]
Ponieważ $\Gamma(1/2)=\sqrt\pi$, wzór ten jest równy
$u^{-1/2}e^{-u/2}/(2^{1/2}\Gamma(1/2))$.
Jest to gęstość $\Gam(1/2,1/2)$.

Każda zmienna $Z_j^2$ ma więc taki sam rozkład gamma.
Z niezależności $Z_j$ wynika niezależność mierzalnych
funkcji $Z_j^2$. Własność dodawania niezależnych
rozkładów gamma o wspólnej intensywności daje
$\sum_{j=1}^\nu Z_j^2\sim\Gam(\nu/2,1/2)=\chi^2_\nu$.
Niezależnie od tego, czy $\nu$ jest całkowite,
przez definicję $\chi^2_\nu$ jako rozkładu gamma
obliczamy $\E U=(\nu/2)/(1/2)=\nu$ oraz
$\Var(U)=(\nu/2)/(1/2)^2=2\nu$.
\end{proof}

\begin{przyklad}[Trzy współrzędne normalne]
Jeżeli $Z_1,Z_2,Z_3$ są niezależne i standardowo normalne,
to $R^2:=Z_1^2+Z_2^2+Z_3^2\sim\chi^2_3$.
Średnia $R^2$ wynosi $3$, a wariancja $6$.
Nie należy mylić rozkładu $R^2$ z rozkładem $R$:
po spierwiastkowaniu powstaje inny rozkład, określany
jako rozkład chi.
\end{przyklad}

\section{Rozkład beta}
W przeciwieństwie do rozkładu gamma, który żyje na dodatniej
półprostej, rozkład beta jest skupiony na odcinku $(0,1)$.
Dzięki dwóm parametrom może być symetryczny, skupiony bliżej
jednego końca albo silnie wygięty w środku przedziału.
Jest więc użyteczny do modelowania proporcji i udziałów.
Jego normalizacja pochodzi wprost z funkcji beta Eulera.

\begin{definicja}
Niech $a,b>0$. Rozkładem beta $\Betad(a,b)$ na $(0,1)$
nazywamy rozkład o gęstości
\begin{equation}\label{eq:beta}
f(x)=\frac{x^{a-1}(1-x)^{b-1}}{\mathrm B(a,b)}
\ind{\{0<x<1\}}.
\end{equation}
\end{definicja}

\begin{stwierdzenie}
Jeżeli $X\sim\Betad(a,b)$, to
\[
\E X=\frac a{a+b},\qquad
\Var(X)=\frac{ab}{(a+b)^2(a+b+1)}.
\]
Rozkład $\Betad(1,1)$ jest jednostajny na $(0,1)$.
\end{stwierdzenie}
\begin{proof}
Najpierw sprawdzamy normalizację: gęstość jest
nieujemna, a jej całka po $(0,1)$ równa
$\mathrm B(a,b)/\mathrm B(a,b)=1$.
Dla $r\geq0$ mnożenie gęstości przez $x^r$
i wykorzystanie definicji funkcji beta daje
\[
\E X^r=\frac{\mathrm B(a+r,b)}{\mathrm B(a,b)}
=\frac{\Gamma(a+r)\Gamma(a+b)}
{\Gamma(a)\Gamma(a+b+r)}.
\]
Dla $r=1$ rekurencja $\Gamma(s+1)=s\Gamma(s)$
prowadzi do $\E X=a/(a+b)$. Dla $r=2$
otrzymujemy
$\E X^2=a(a+1)/[(a+b)(a+b+1)]$.
Stąd
\[
\Var(X)=\frac{a(a+1)}{(a+b)(a+b+1)}
-\frac{a^2}{(a+b)^2}
=\frac{ab}{(a+b)^2(a+b+1)}.
\]
Wreszcie $\mathrm B(1,1)=\int_0^1 1\,\dd\lambda_1(x)=1$,
więc dla $a=b=1$ gęstość jest stała i równa jeden
na $(0,1)$, tak jak w rozkładzie jednostajnym.
\end{proof}

\section{Rozkład t-Studenta}
Jeżeli zmienną normalną dzielimy przez niezależne losowe
odchylenie związane z rozkładem chi-kwadrat, mianownik może
przyjmować małe wartości. Powstają wówczas obserwacje znacznie
bardziej odległe od zera niż w rozkładzie normalnym. To źródło
cięższych ogonów rozkładu t-Studenta. Liczba stopni swobody
steruje częstością takich odchyleń i decyduje, które momenty
w ogóle istnieją.

\begin{definicja}
Niech $\nu>0$, $Z\sim\Normal(0,1)$ oraz
$U\sim\chi^2_\nu$ będą niezależne. Zmienną
$T:=Z/\sqrt{U/\nu}$ nazywamy zmienną o rozkładzie
t-Studenta z $\nu$ stopniami swobody i zapisujemy $T\sim t_\nu$.
\end{definicja}

\begin{twierdzenie}[Gęstość rozkładu Studenta]
Dla każdego $\nu>0$ zmienna $T\sim t_\nu$ ma gęstość
\begin{equation}\label{eq:student}
f_\nu(t)=
\frac{\Gamma((\nu+1)/2)}{\sqrt{\nu\pi}\,\Gamma(\nu/2)}
\left(1+\frac{t^2}{\nu}\right)^{-(\nu+1)/2},
\qquad t\in\R.
\end{equation}
Gęstość jest parzysta. Dla $\nu>1$ wartość oczekiwana
jest równa zero, a dla $\nu>2$ wariancja wynosi $\nu/(\nu-2)$.
Dla $0<\nu\leq1$ wartość oczekiwana nie istnieje,
a dla $0<\nu\leq2$ drugi moment jest nieskończony.
\end{twierdzenie}
\begin{proof}
Niech $\varphi(z)=(2\pi)^{-1/2}e^{-z^2/2}$ oznacza
gęstość standardowego rozkładu normalnego.
Z niezależności $(Z,U)$ ma na $\R\times(0,\infty)$
gęstość $\varphi(z)f_U(u)$, gdzie
$f_U(u)=u^{\nu/2-1}e^{-u/2}/(2^{\nu/2}\Gamma(\nu/2))$.
Dla ustalonego $u>0$ przechodzimy od
$(z,u)$ do $(t,u)$ przez $z=t\sqrt{u/\nu}$.
Moduł jakobianu odwrotnej zmiany jest
$\sqrt{u/\nu}$. Stąd gęstość pary $(T,U)$
wynosi $\varphi(t\sqrt{u/\nu})f_U(u)\sqrt{u/\nu}$.
Całkując po $u$ i grupując potęgi, otrzymujemy
\[
f_T(t)=\frac1{\sqrt{2\pi\nu}\,2^{\nu/2}\Gamma(\nu/2)}
\int_0^\infty u^{(\nu+1)/2-1}
\exp\!\left(-\frac{1+t^2/\nu}{2}u\right)
\,\dd\lambda_1(u).
\]
Oznaczmy $a=(\nu+1)/2$ i $c=(1+t^2/\nu)/2>0$.
Podstawienie $v=cu$ daje
$\int_0^\infty u^{a-1}e^{-cu}\,\dd\lambda_1(u)
=\Gamma(a)c^{-a}$.
Po wstawieniu i uproszczeniu potęg dwójki
otrzymujemy
\[
f_T(t)=\frac{\Gamma((\nu+1)/2)}
{\sqrt{\nu\pi}\,\Gamma(\nu/2)}
\left(1+\frac{t^2}{\nu}\right)^{-(\nu+1)/2}.
\]
Gęstość jest parzysta, ale symetria sama w sobie
nie gwarantuje istnienia wartości oczekiwanej.
Dla $|t|\geq1$ istnieją stałe dodatnie $c_1,c_2$
(zależne od $\nu$) takie, że
$c_1|t|^{-\nu-1}\leq f_T(t)
\leq c_2|t|^{-\nu-1}$. Zatem dla $r\geq0$
całka $\int_{|t|\geq1}|t|^rf_T(t)\,\dd\lambda_1(t)$
zachowuje się pod względem zbieżności jak
$\int_1^\infty t^{r-\nu-1}\,\dd\lambda_1(t)$,
a ta jest skończona dokładnie dla $r<\nu$.
W konsekwencji $\E|T|<\infty$ wtedy i tylko wtedy,
gdy $\nu>1$; w tym przypadku $\E T=0$ przez parzystość.
Dla $\nu>2$ drugi moment istnieje i
\[
\E T^2=\nu\E Z^2\,\E(U^{-1})
=\nu\frac{1}{\nu-2},
\]
ponieważ $\E Z^2=1$ oraz, z całki gamma,
$\E(U^{-1})=1/(\nu-2)$. Zatem
$\Var(T)=\nu/(\nu-2)$. Przy $\nu\leq2$ drugi
moment jest nieskończony, więc wariancja nie jest
skończoną liczbą; dla $\nu\leq1$ średnia nie istnieje.
\end{proof}

\begin{uwaga}
Parametr $\nu$ w definicji rozkładu Studenta może być dowolną
dodatnią liczbą rzeczywistą. W klasycznych zadaniach o estymacji
wariancji na podstawie próby normalnej występują zwykle całkowite
liczby stopni swobody, często $n-1$.
Dla $\nu\to\infty$ rozkład $t_\nu$ zbiega według rozkładu
do $\Normal(0,1)$; dowód znajdzie się w rozdziale
o centralnym twierdzeniu granicznym.
\end{uwaga}

\section{Rozkład Cauchy'ego}
Symetryczna, gładka gęstość nie musi opisywać rozkładu z dobrze
określoną średnią. Rozkład Cauchy'ego stanowi ważny kontrprzykład:
mimo symetrii jego dodatnia i ujemna część pierwszego momentu
są rozbieżne. Na nim zobaczymy, dlaczego formalne „znoszenie się”
nieskończoności nie jest dopuszczalnym sposobem obliczania
wartości oczekiwanej.

\begin{definicja}
Niech $x_0\in\R$ i $\gamma>0$. Rozkładem Cauchy'ego
o parametrze położenia $x_0$ i skali $\gamma$
nazywamy rozkład o gęstości
\begin{equation}\label{eq:cauchy}
f(x)=\frac1{\pi\gamma}
\frac{1}{1+((x-x_0)/\gamma)^2},\qquad x\in\R.
\end{equation}
\end{definicja}

\begin{stwierdzenie}
Jeżeli $X$ ma powyższy rozkład, to
\[
F_X(x)=\frac12+\frac1\pi
\arctan\!\left(\frac{x-x_0}{\gamma}\right).
\]
Dla każdego $r\geq1$ zachodzi $\E|X-x_0|^r=\infty$.
W szczególności wartość oczekiwana $\E X$ nie istnieje,
a wariancja nie jest zdefiniowana.
Rozkład standardowy Cauchy'ego ($x_0=0$, $\gamma=1$)
jest równy rozkładowi $t_1$.
\end{stwierdzenie}
\begin{proof}
Całka z gęstości \eqref{eq:cauchy} po $(-\infty,x]$
daje wskazaną dystrybuantę.
Dla dużych $|x|$ funkcja $|x-x_0|^r f(x)$
zachowuje się jak dodatnia stała razy $|x|^{r-2}$,
której całka po półprostej jest rozbieżna dla $r\geq1$.
W szczególności dodatnia i ujemna część
$X-x_0$ mają nieskończone wartości oczekiwane,
więc nie można definiować średniej jako ich różnicy.
Podstawienie $\nu=1$ do \eqref{eq:student} i tożsamości
$\Gamma(1)=1$, $\Gamma(1/2)=\sqrt\pi$ daje gęstość
$1/[\pi(1+x^2)]$ rozkładu standardowego Cauchy'ego.
\end{proof}

\begin{przyklad}[Wartość główna a wartość oczekiwana]
Dla standardowego rozkładu Cauchy'ego zachodzi
$\int_{-R}^{R}x/[\pi(1+x^2)]\,\dd\lambda_1(x)=0$
dla każdego $R>0$ z powodu nieparzystości funkcji.
Wartość główna Cauchy'ego całki wynosi zatem zero,
ale nie jest to wartość oczekiwana: całki z części dodatniej
i ujemnej są obie nieskończone. Nie można stosować
symetrycznego znoszenia się rozbieżnych całek w definicji
wartości oczekiwanej.
\end{przyklad}

\begin{przyklad}[Tangens kąta jednostajnego]
Jeśli $U\sim\Unif(-\pi/2,\pi/2)$ i $X=\tan U$,
to $X$ ma standardowy rozkład Cauchy'ego. Rzeczywiście,
$\tan$ jest dyfeomorfizmem tego przedziału na $\R$,
a pochodna funkcji odwrotnej wynosi
$(\arctan x)'=1/(1+x^2)$.
Wzór na transformację gęstości daje $f_X(x)=1/[\pi(1+x^2)]$.
\end{przyklad}


\begin{przyklad}[Porównanie ogonów rozkładu normalnego i Cauchy'ego]
Niech $Z\sim\Normal(0,1)$, $C$ ma standardowy rozkład
Cauchy'ego i niech $t>0$. Z nierówności Markowa
dla $Z^2$ wynika $\Prob(|Z|\geq t)\leq1/t^2$.
W rzeczywistości ogon normalny maleje znacznie szybciej:
dla każdego $u>0$, korzystając z
$M_Z(u)=e^{u^2/2}$, mamy
$\Prob(Z\geq t)=\Prob(e^{uZ}\geq e^{ut})
\leq e^{-ut+u^2/2}$.
Wybór $u=t$ daje $\Prob(Z\geq t)\leq e^{-t^2/2}$,
a przez symetrię $\Prob(|Z|\geq t)\leq2e^{-t^2/2}$.

Tymczasem dla $C$ zachodzi
$\Prob(|C|\geq t)=1-(2/\pi)\arctan t$.
Z asymptotyki $\pi/2-\arctan t\sim1/t$
otrzymujemy $\Prob(|C|\geq t)\sim 2/(\pi t)$.
Różnica między zanikiem wykładniczym a potęgowym
wyjaśnia, dlaczego w rozkładzie Cauchy'ego brak nawet
pierwszego bezwzględnego momentu.
\end{przyklad}

\section{Dalsze użyteczne rozkłady}
Nie wszystkie użyteczne rodziny wymagają nowej konstrukcji od zera.
Rozkład lognormalny powstaje przez wykładnicze przekształcenie
rozkładu normalnego. Rozkład Fishera--Snedecora powstaje z ilorazu
dwóch niezależnych odpowiednio unormowanych zmiennych chi-kwadrat.
To przykłady ważnej zasady: z prostych rozkładów tworzymy
bardziej złożone za pomocą działań na zmiennych losowych.

\begin{definicja}
Zmienna $X$ ma rozkład lognormalny z parametrami
$\mu\in\R$, $\sigma>0$, jeżeli $\log X\sim\Normal(\mu,\sigma^2)$
(prawie na pewno $X>0$). Ma wtedy gęstość
$\ind{\{x>0\}}\exp(- (\log x-\mu)^2/(2\sigma^2))
/[x\sigma\sqrt{2\pi}]$.
\end{definicja}

\begin{stwierdzenie}
Jeżeli $X$ ma rozkład lognormalny z parametrami $\mu,\sigma$,
to $\E X=e^{\mu+\sigma^2/2}$ oraz
$\Var(X)=(e^{\sigma^2}-1)e^{2\mu+\sigma^2}$.
\end{stwierdzenie}
\begin{uwaga}
Nie należy stąd wnioskować, że funkcja tworząca momenty
$M_X(t)=\E e^{tX}$ jest skończona w otoczeniu zera.
Dla lognormalnej zmiennej losowej $M_X(t)=\infty$
dla każdego $t>0$, choć wszystkie momenty
$\E X^r$ są skończone. Rozbieżność wynika z tego,
że dla $X=e^Y$, gdzie $Y$ jest normalna,
wykładnik $te^y-(y-\mu)^2/(2\sigma^2)$
rośnie do $+\infty$ przy $y\to+\infty$.
\end{uwaga}
\begin{proof}
Zapisujemy $X=e^Y$, gdzie $Y\sim\Normal(\mu,\sigma^2)$.
Wtedy $\E X=M_Y(1)$ i $\E X^2=M_Y(2)$,
co daje odpowiednio $e^{\mu+\sigma^2/2}$ i
$e^{2\mu+2\sigma^2}$. Odejmujemy kwadrat średniej.
\end{proof}

\begin{definicja}
Niech $d_1,d_2\in\N$ i niech $U\sim\chi^2_{d_1}$ oraz
$V\sim\chi^2_{d_2}$ będą niezależne. Rozkład zmiennej
$F:=(U/d_1)/(V/d_2)$ nazywamy rozkładem Fishera--Snedecora
i oznaczamy $F_{d_1,d_2}$.
\end{definicja}

\begin{stwierdzenie}
Rozkład $F_{d_1,d_2}$ ma gęstość
\[
f(x)=\frac{\Gamma((d_1+d_2)/2)}
{\Gamma(d_1/2)\Gamma(d_2/2)}
\left(\frac{d_1}{d_2}\right)^{d_1/2}
\frac{x^{d_1/2-1}}{(1+d_1x/d_2)^{(d_1+d_2)/2}}
\ind{\{x>0\}}.
\]
Jeżeli $T\sim t_\nu$ i $\nu\in\N$, to $T^2\sim F_{1,\nu}$.
\end{stwierdzenie}
\begin{proof}
Oznaczmy $a=d_1/2$ oraz $b=d_2/2$.
Z niezależności $U,V$ ich gęstość łączna na dodatnim
kwadrancie jest równa
\[
f_{U,V}(u,v)=\frac{u^{a-1}v^{b-1}e^{-(u+v)/2}}
{2^{a+b}\Gamma(a)\Gamma(b)}.
\]
Definicja $X=(U/d_1)/(V/d_2)$ prowadzi do zmiany
$u=(d_1/d_2)xv$, przy której moduł jakobianu
odwrotnej transformacji $(x,v)\mapsto(u,v)$ wynosi
$d_1v/d_2$. Podstawiając do gęstości łącznej,
a następnie całkując po $v>0$, dostajemy
\[
f_X(x)=\frac{(d_1/d_2)^a x^{a-1}}
{2^{a+b}\Gamma(a)\Gamma(b)}
\int_0^\infty v^{a+b-1}
\exp\!\left(-\frac{1+(d_1/d_2)x}{2}v\right)
\,\dd\lambda_1(v),\quad x>0.
\]
Po podstawieniu $w=(1+(d_1/d_2)x)v/2$
całka wynosi
$\Gamma(a+b)\,[2/(1+(d_1/d_2)x)]^{a+b}$.
Po uproszczeniu otrzymujemy podaną gęstość.

Dla $T\sim t_\nu$ mamy z definicji
$T^2=(Z^2/1)/(V/\nu)$, gdzie
$Z^2\sim\chi_1^2$, $V\sim\chi_\nu^2$ i są niezależne.
Zatem $T^2\sim F_{1,\nu}$. Ten argument działa
również dla dowolnego rzeczywistego $\nu>0$, jeśli
rozszerzymy definicję $F_{d_1,d_2}$ na rzeczywiste
dodatnie parametry; w tekście zasadnicza definicja
podana została dla całkowitych stopni swobody.
\end{proof}

\section{Porównanie rozkładów i wybór modelu}
Właściwy wybór modelu nie polega na dopasowaniu znanego wzoru
do przypadkowej tabeli. Najpierw trzeba rozpoznać mechanizm:
losowanie ze zwracaniem czy bez, ustalona liczba prób czy czekanie
na sukces, pojedyncza wielkość czy suma/iloraz wielkości losowych.
Dopiero potem parametryzuje się odpowiadający mu rozkład.

\begin{center}
\small
\begin{tabular}{@{}p{3.0cm}p{4.1cm}p{5.9cm}@{}}
\hline
Rozkład & Parametry & Typowe zastosowanie \\
\hline
Bernoulliego & $p\in[0,1]$ & Wynik pojedynczej próby binarnej \\
Dwumianowy & $n\in\N_0$, $p\in[0,1]$ & Liczba sukcesów w niezależnych próbach \\
Hipergeometryczny & $0\leq K,n\leq N$ & Losowanie bez zwracania \\
Poissona & $\lambda>0$ & Liczba rzadkich niezależnych zdarzeń \\
Geometryczny & $p\in(0,1]$ & Liczba prób do pierwszego sukcesu \\
Jednostajny & $a<b$ & Jednorodne losowanie punktu odcinka \\
Wykładniczy & $\lambda>0$ & Niepamiętający czas oczekiwania \\
Normalny & $\mu\in\R$, $\sigma>0$ & Błędy pomiaru, sumy licznych składników \\
Gamma & $\alpha,\lambda>0$ & Dodatnie czasy i sumy czasów oczekiwania \\
Chi-kwadrat & $\nu>0$ & Sumy kwadratów normalnych \\
Beta & $a,b>0$ & Udziały, proporcje i zmienne o wartościach w $(0,1)$ \\
t-Studenta & $\nu>0$ & Standaryzowane ilorazy normalnej i chi-kwadrat \\
Cauchy'ego & $x_0\in\R$, $\gamma>0$ & Model z ciężkimi ogonami, kontrprzykłady dla momentów \\
\hline
\end{tabular}
\end{center}

\begin{uwaga}
Zapis $\N_0$ w tabeli oznacza $\{0,1,2,\ldots\}$.
Przykłady zastosowań są wskazówkami modelowymi, a nie
warunkami definiującymi rozkłady. To, czy model jest odpowiedni,
wynika z przyjętych założeń, sposobu zbierania danych
i weryfikacji ich zgodności z modelem.
\end{uwaga}

\chapter{Zbieżność zmiennych losowych i rozkładów}
\label{ch:convergence}

Pojęcie granicy występuje w probabilistyce w kilku znaczeniach.
Przy wielokrotnych pomiarach chcielibyśmy wiedzieć, czy
ich średnia zbliży się do pewnej liczby; przy przybliżaniu
rozkładu dwumianowego rozkładem Poissona pytamy natomiast,
czy zbliżają się do siebie \emph{rozkłady}, niekoniecznie
realizacje zmiennych na jednej przestrzeni. Te zadania
wymagają rozróżnienia kilku topologii zbieżności.

\section{Zbieżność prawie na pewno, według prawdopodobieństwa i w~\texorpdfstring{$L^p$}{Lp}}

\begin{definicja}
Niech $(\Omega,\mathcal F,\Prob)$ będzie przestrzenią
probabilistyczną, a $X,X_1,X_2,\ldots:\Omega\to\R$
rzeczywistymi zmiennymi losowymi.
\begin{enumerate}
\item Mówimy, że $X_n$ jest zbieżny do $X$ \emph{prawie na pewno}
 i piszemy $X_n\to X$ p.n., gdy
 $\Prob(\{\omega:\lim_{n\to\infty}X_n(\omega)=X(\omega)\})=1$.
\item Mówimy, że $X_n$ jest zbieżny do $X$
 \emph{według prawdopodobieństwa}, $X_n\xrightarrow{\Prob}X$,
 gdy dla każdego $\varepsilon>0$ zachodzi
 $\lim_{n\to\infty}\Prob(|X_n-X|>\varepsilon)=0$.
\item Dla $p\geq1$ i $X,X_n\in L^p(\Omega,\mathcal F,\Prob)$
 mówimy o \emph{zbieżności w $L^p$}, gdy
 $\E|X_n-X|^p\to0$.
\end{enumerate}
\end{definicja}

Warunek „prawie na pewno” dotyczy niemal każdej
\emph{trajektorii} $n\mapsto X_n(\omega)$.
Zbieżność według prawdopodobieństwa pozwala na sporadyczne
duże odchylenia, o ile ich prawdopodobieństwo zanika.
Zbieżność w $L^p$ kontroluje natomiast średni koszt
odchylenia podniesionego do potęgi $p$.

\begin{twierdzenie}[Podstawowe implikacje]
Niech $X_n,X$ będą rzeczywistymi zmiennymi losowymi.
\begin{enumerate}
\item Jeśli $X_n\to X$ prawie na pewno, to
 $X_n\xrightarrow{\Prob}X$.
\item Jeśli $p\geq1$, $X_n,X\in L^p$ oraz
 $\E|X_n-X|^p\to0$, to $X_n\xrightarrow{\Prob}X$.
\item Jeśli $p\geq q\geq1$ i $X_n,X\in L^p$ oraz
 $X_n\to X$ w $L^p$, to $X_n\to X$ w $L^q$.
\end{enumerate}
\end{twierdzenie}
\begin{proof}
Dla pierwszej implikacji ustalmy $\varepsilon>0$.
Wskaźniki $\ind{\{|X_n-X|>\varepsilon\}}$ zbiegają
prawie na pewno do zera i są ograniczone przez jeden.
Zdominowana zbieżność daje zatem
$\Prob(|X_n-X|>\varepsilon)\to0$.
Dla drugiej implikacji do nieujemnej zmiennej
$|X_n-X|^p$ stosujemy nierówność Markowa:
\[
\Prob(|X_n-X|>\varepsilon)
\leq\frac{\E|X_n-X|^p}{\varepsilon^p}\longrightarrow0.
\]
Dla trzeciej wykorzystujemy nierówność Jensena
dla funkcji wklęsłej $u\mapsto u^{q/p}$ na
$[0,\infty)$ (albo nierówność H\"oldera przy
$\Prob(\Omega)=1$):
$\E|X_n-X|^q\leq(\E|X_n-X|^p)^{q/p}\to0$.
\end{proof}

\begin{przyklad}[Zbieżność w prawdopodobieństwie bez zbieżności p.n.]
Na $\Omega=[0,1)$ z miarą Lebesgue'a ustawmy kolejno
wszystkie przedziały diadyczne
$I_{m,j}=[j2^{-m},(j+1)2^{-m})$ dla
$m=0,1,2,\ldots$ oraz $0\leq j<2^m$.
Zdefiniujmy $X_{2^m+j}=\ind{I_{m,j}}$.
Wtedy $\Prob(X_{2^m+j}=1)=2^{-m}\to0$, a więc
$X_n\xrightarrow{\Prob}0$. Jednak dla każdego
$\omega\in[0,1)$ i każdego poziomu $m$ dokładnie
jeden ze wskaźników przedziałów tego poziomu wynosi
jeden. Dla $m\geq1$ co najmniej jeden wynosi także zero.
Zatem każda trajektoria przyjmuje wartości jeden i zero
nieskończenie wiele razy; $X_n$ nie jest zbieżny do zera
prawie na pewno. Jest to tzw. ciąg piszący na maszynie
(\emph{typewriter sequence}).
\end{przyklad}

\begin{przyklad}[Zbieżność w prawdopodobieństwie bez zbieżności w $L^1$]
Niech $U\sim\Unif(0,1)$ i
$X_n=n\ind{\{U\leq1/n\}}$ dla $n\geq1$.
Dla każdego $\varepsilon>0$ i $n>\varepsilon$ mamy
$\Prob(|X_n|>\varepsilon)=1/n\to0$.
Zatem $X_n\xrightarrow{\Prob}0$, lecz
$\E|X_n|=n\Prob(U\leq1/n)=1$.
Odchylenia występują coraz rzadziej, ale jednocześnie
stają się coraz większe: średni koszt nie znika.
\end{przyklad}

\begin{twierdzenie}[Kryterium podciągowe zbieżności według prawdopodobieństwa]
Dla rzeczywistych zmiennych losowych $X_n,X$ zachodzi
$X_n\xrightarrow{\Prob}X$ wtedy i tylko wtedy, gdy
z każdego podciągu $(X_{n_j})_{j\geq1}$ można wybrać
podciąg $(X_{n_{j_k}})_{k\geq1}$ zbieżny do $X$
prawie na pewno.
\end{twierdzenie}
\begin{proof}
Załóżmy zbieżność według prawdopodobieństwa.
Dla danego podciągu wybieramy rekurencyjnie indeksy
$j_1<j_2<\cdots$ tak, aby
$\Prob(|X_{n_{j_k}}-X|>2^{-k})\leq2^{-k}$.
Jest to możliwe, gdyż dla każdego ustalonego progu
prawdopodobieństwo odchylenia dąży do zera.
Suma prawdopodobieństw otrzymanych zdarzeń jest skończona.
Z pierwszego lematu Borela--Cantellego wynika, że
prawie na pewno występuje tylko skończenie wiele
odchyleń przekraczających $2^{-k}$; stąd zbieżność
podciągu do $X$ prawie na pewno.

Odwrotnie, gdyby zbieżność według prawdopodobieństwa
nie zachodziła, istniałyby $\varepsilon,\eta>0$ i
podciąg $(X_{n_j})$ taki, że
$\Prob(|X_{n_j}-X|>\varepsilon)\geq\eta$
dla każdego $j$. Żaden jego podciąg nie mógłby być
zbieżny do $X$ prawie na pewno, bo pierwsza implikacja
poprzedniego twierdzenia wymuszałaby zbieżność
prawdopodobieństw do zera. Otrzymujemy sprzeczność.
\end{proof}

\section{Zbieżność według rozkładu}

Zbieżność według rozkładu nie zakłada nawet istnienia
wspólnej przestrzeni probabilistycznej dla wszystkich
zmiennych. Interesują nas wyłącznie ich miary obrazowe.
Z tego względu jest właściwym językiem dla
przybliżeń Poissona i dla centralnego twierdzenia granicznego.

\begin{definicja}
Niech $X_n$ będą rzeczywistymi zmiennymi losowymi
(niekoniecznie na wspólnej przestrzeni), o dystrybuantach
$F_n$, i niech $X$ będzie zmienną losową o dystrybuancie
$F$. Mówimy, że $X_n$ jest zbieżny do $X$
\emph{według rozkładu}, i piszemy $X_n\xrightarrow{d}X$,
jeśli $F_n(x)\to F(x)$ dla każdego $x\in\R$ będącego
punktem ciągłości funkcji $F$.
\end{definicja}

Nie wymagamy zbieżności w punktach skoku $F$.
Dla przykładu, zmienna deterministyczna $X_n=1/n$
jest zbieżna do $X=0$, ale $F_{X_n}(0)=0$ dla
wszystkich $n$, podczas gdy $F_X(0)=1$.
Takie zachowanie nie przeczy definicji.

\begin{twierdzenie}[Zbieżność prawdopodobieństw a zbieżność rozkładów]
Niech $X_n,X$ będą rzeczywistymi zmiennymi losowymi
na wspólnej przestrzeni probabilistycznej.
Jeśli $X_n\xrightarrow{\Prob}X$, to
$X_n\xrightarrow{d}X$.
Jeśli $c\in\R$ jest stałą, to zbieżność
$X_n\xrightarrow{d}c$ implikuje
$X_n\xrightarrow{\Prob}c$.
\end{twierdzenie}
\begin{proof}
Niech $x$ będzie punktem ciągłości dystrybuanty $F_X$
i niech $\varepsilon>0$. Z inkluzji zdarzeń
\[
\{X\leq x-\varepsilon\}\cap\{|X_n-X|\leq\varepsilon\}
\subseteq\{X_n\leq x\}
\subseteq\{X\leq x+\varepsilon\}\cup
\{|X_n-X|>\varepsilon\}
\]
wynika, po przejściu z $n\to\infty$,
$F_X(x-\varepsilon)\leq\liminf_nF_n(x)
\leq\limsup_nF_n(x)\leq F_X(x+\varepsilon)$.
Po przejściu z $\varepsilon\downarrow0$ i użyciu
ciągłości $F_X$ w $x$ dostajemy $F_n(x)\to F_X(x)$.

Dla drugiego wniosku, jeśli $\varepsilon>0$,
to punkty $c-\varepsilon$ i $c+\varepsilon/2$
są punktami ciągłości dystrybuanty rozkładu
skupionego w $c$. Dlatego
\[
\Prob(|X_n-c|>\varepsilon)
\leq F_n(c-\varepsilon)
 +1-F_n(c+\varepsilon/2)\longrightarrow0.
\]
W rzeczywistości pierwsza nierówność można też
zapisać z ostrzejszymi granicami przedziałów,
ale podana wersja wystarcza do dowodu.
\end{proof}

\begin{przyklad}[Ten sam rozkład nie oznacza bliskości realizacji]
Niech $X$ przyjmuje wartości $-1$ i $1$ z
prawdopodobieństwami $1/2$, a $X_n=-X$ dla każdego $n$.
Wtedy $\Law(X_n)=\Law(X)$, więc $X_n\xrightarrow{d}X$.
Jednak $|X_n-X|=2$ prawie na pewno, toteż zbieżność
według prawdopodobieństwa do $X$ nie zachodzi.
\end{przyklad}

\begin{lemat}[Całki względem rozkładów zbieżnych]
Jeżeli $X_n\xrightarrow{d}X$, to dla każdej ograniczonej
funkcji ciągłej $h:\R\to\R$ zachodzi
$\E h(X_n)\to\E h(X)$.
W szczególności $\varphi_{X_n}(t)\to\varphi_X(t)$
dla każdego $t\in\R$.
\end{lemat}
\begin{proof}
Przytoczymy argument wyprowadzający własność całek
z definicji dystrybuant. Najpierw wybierzmy $R>0$ tak,
aby $-R$ i $R$ były punktami ciągłości $F_X$ oraz
$\Prob(|X|\geq R)$ było małe. Zbieżność dystrybuant
zapewnia, że dla wszystkich dostatecznie dużych $n$
prawdopodobieństwa $\Prob(|X_n|>R)$ również są małe
(z ewentualnym nieistotnym przesunięciem końców).
Na $[-R,R]$ funkcja $h$ jest jednostajnie ciągła.
Dzielimy odcinek na skończenie wiele przedziałów,
których końce są punktami ciągłości $F_X$, tak aby
oscylacja $h$ na każdym przedziale była mniejsza
od dowolnie zadanego $\varepsilon>0$.
Na takim przedziale $I=(a,b]$ zachodzi
$\Prob(X_n\in I)=F_n(b)-F_n(a)\to F_X(b)-F_X(a)$.
Zastępując $h$ przez funkcję schodkową stałą na
każdym z tych przedziałów, uzyskujemy zbieżność
całek funkcji schodkowych. Błąd zastąpienia nie
przekracza $\varepsilon$ na $[-R,R]$ i
$2\|h\|_\infty$ razy prawdopodobieństwo ogonów
poza nim. Możemy najpierw zwiększać $n$, następnie
zmniejszać $\varepsilon$, a na końcu zwiększać $R$;
to dowodzi tezy. Dla części rzeczywistej i urojonej
$h(x)=e^{itx}$ stosujemy wynik osobno, ponieważ
obie są ograniczone i ciągłe.
\end{proof}


\begin{przyklad}[Zbieżność według rozkładu nie oznacza zbieżności według prawdopodobieństwa]
Niech $Z$ przyjmuje wartości $-1$ i $1$ z jednakowymi
prawdopodobieństwami, a $X_n=(-1)^nZ$ i $X=Z$.
Dla każdego $n$ zmienna $X_n$ ma ten sam rozkład co $X$,
a więc $X_n\xrightarrow{d}X$. Jednak dla każdego
nieparzystego $n$ zachodzi $|X_n-X|=2$ prawie na pewno,
więc $\Prob(|X_n-X|>1)=1$.
Wobec tego nie zachodzi $X_n\xrightarrow{\Prob}X$.
Zbieżność według rozkładu nie porównuje wartości
$X_n(\omega)$ i $X(\omega)$ na tych samych trajektoriach;
porównuje wyłącznie miary obrazowe.
\end{przyklad}

\section{Twierdzenie Lévy'ego i przydatne operacje}

Odwrotne twierdzenie do ostatniego lematu jest jednym
z najważniejszych zastosowań funkcji charakterystycznych:
granicę rozkładów można rozpoznać po granicy ich
transformat Fouriera. Wyjaśnia to, dlaczego funkcje
charakterystyczne będą głównym narzędziem w dowodzie CTG.

\begin{lemat}[Zwartość rodzin ciasnych na prostej]
Niech $(\mu_n)_{n\geq1}$ będzie rodziną borelowskich miar
probabilistycznych na $\R$ taką, że dla każdego
$\varepsilon>0$ istnieje $R>0$ spełniające
$\sup_n\mu_n(\R\setminus[-R,R])<\varepsilon$.
Wówczas z każdego ciągu $(\mu_n)$ można wybrać
podciąg zbieżny według rozkładu do pewnej borelowskiej
miary probabilistycznej na $\R$.
\end{lemat}
\begin{proof}
\emph{Wybór podciągu.} Niech $F_n$ będą dystrybuantami
$\mu_n$. Uporządkujmy liczby wymierne jako
$\mathbb Q=\{q_1,q_2,\ldots\}$. Ciąg
$(F_n(q_1))_n$ jest ograniczony, zatem ma podciąg
zbieżny. Z tego podciągu wybierzmy następny, na którym
zbiega także $F_n(q_2)$, i kontynuujmy indukcyjnie.
Elementy przekątni tych podciągów wyznaczają jeden
podciąg $n_k$ taki, że dla każdej liczby wymiernej
$q$ istnieje granica $G(q)=\lim_kF_{n_k}(q)$.
Funkcja $G:\mathbb Q\to[0,1]$ jest niemalejąca.

\emph{Konstrukcja dystrybuanty granicznej.}
Dla każdego $x\in\R$ definiujemy
$F(x)=\inf\{G(q):q\in\mathbb Q,\ q>x\}$.
Wprost z definicji wynika, że $F$ jest niemalejąca.
Dla ciągu $x_j\downarrow x$ mamy $F(x_j)\geq F(x)$.
Z drugiej strony, jeśli ustalimy dowolne wymierne
$q>x$, to dla dostatecznie dużych $j$ zachodzi
$x_j<q$, skąd $F(x_j)\leq G(q)$. Biorąc najpierw
granicę po $j$, a potem infimum po $q>x$,
dostajemy $F(x_j)\to F(x)$. To dowodzi
prawostronnej ciągłości $F$.

\emph{Dlaczego masa nie ucieka do nieskończoności?}
Z ciasności dla ustalonego $\varepsilon>0$
wybieramy $R>0$ tak, aby dla wszystkich $n$
$\mu_n(\R\setminus[-R,R])<\varepsilon$.
Dla każdego wymiernego $q<-R$ mamy
$G(q)\leq\varepsilon$, a dla $q>R$ mamy
$G(q)\geq1-\varepsilon$. Stąd dla $x<-R$
$F(x)\leq\varepsilon$ (wybierając $q\in(x,-R)$),
a dla $x>R$ mamy $F(x)\geq1-\varepsilon$.
Ponieważ $\varepsilon$ jest dowolne,
$F(-\infty)=0$ i $F(+\infty)=1$.
Jest to zatem dystrybuanta pewnej miary
probabilistycznej $\mu$.

\emph{Zbieżność w punktach ciągłości.}
Niech $x$ będzie punktem ciągłości $F$.
Dla każdego wymiernego $q>x$ monotoniczność
$F_{n_k}$ daje
$\limsup_kF_{n_k}(x)\leq G(q)$, a więc
$\limsup_kF_{n_k}(x)\leq F(x)$.
Z drugiej strony, jeśli $y<x$, wybierzmy
wymierne $q\in(y,x)$. Wtedy
$\liminf_kF_{n_k}(x)\geq G(q)\geq F(y)$,
bo $F(y)$ jest infimum wartości $G$ w punktach
wymiernych większych od $y$. Zatem
$\liminf_kF_{n_k}(x)\geq F(x-)$, gdzie
$F(x-)=\lim_{y\uparrow x}F(y)$.
Ciągłość $F$ w $x$ daje $F(x-)=F(x)$, skąd
$F_{n_k}(x)\to F(x)$. Jest to zbieżność
podciągu według rozkładu.
\end{proof}

\begin{twierdzenie}[Twierdzenie Lévy'ego o ciągłości]
Niech $(\mu_n)_{n\geq1}$ będzie ciągiem miar
probabilistycznych na $\R$, a $\varphi_n$ ich
funkcjami charakterystycznymi. Jeśli dla każdego
$t\in\R$ zachodzi $\varphi_n(t)\to\varphi(t)$,
a funkcja $\varphi$ jest ciągła w zerze, to istnieje
jedyna miara probabilistyczna $\mu$ o funkcji
charakterystycznej $\varphi$ i $\mu_n$ zbiega
według rozkładu do $\mu$.
Odwrotnie, zbieżność według rozkładu implikuje
punktową zbieżność funkcji charakterystycznych
funkcji granicznego rozkładu.
\end{twierdzenie}
\begin{proof}
Odwrotna implikacja jest treścią poprzedniego lematu.
Dla właściwej części dowodu najpierw wykażemy ciasność.
Niech $X_n$ ma rozkład $\mu_n$ i ustalmy $T>0$.
Przez Fubiniego dla ograniczonych funkcji dostajemy
\[
\frac1{2T}\int_{-T}^T[1-\operatorname{Re}\varphi_n(t)]
\,\dd\lambda_1(t)
=\E\left[1-\frac{\sin(TX_n)}{TX_n}\right],
\]
przyjmując wartość ilorazu równą jeden dla $X_n=0$.
Jeżeli $|X_n|>2/T$, to
$|\sin(TX_n)/(TX_n)|\leq1/2$.
Ponieważ $1-\sin u/u\geq0$ dla wszystkich $u$, mamy
\[
\Prob(|X_n|>2/T)\leq\frac1T
\int_{-T}^T[1-\operatorname{Re}\varphi_n(t)]
\,\dd\lambda_1(t).
\]
Dla ustalonego $T$ prawa strona zbiega, przez zdominowaną
zbieżność, do $T^{-1}\int_{-T}^T
[1-\operatorname{Re}\varphi(t)]\,\dd\lambda_1(t)$.
Ponieważ $\varphi(0)=\lim_n\varphi_n(0)=1$
i $\varphi$ jest ciągła w zerze, ostatnia liczba jest
mała dla dostatecznie małego $T$. Wybierając następnie
$n$ dostatecznie duże, otrzymujemy jednolite
oszacowanie ogona dla wszystkich takich $n$;
dla skończenie wielu pozostałych miar zwiększamy $R$.
W ten sposób otrzymujemy ciasność całej rodziny.

Z lematu o zwartości wynika, że każdy podciąg $\mu_n$
ma dalszy podciąg zbieżny według rozkładu do pewnej
miary $\nu$. Z lematu o całkach ograniczonych ciągłych
funkcja charakterystyczna $\nu$ jest granicą
odpowiednich $\varphi_{n_k}$, a więc równa $\varphi$.
Jednoznaczność funkcji charakterystycznych zapewnia,
że wszystkie możliwe granice podciągów są równe tej
samej mierze $\mu$. Gdyby cały ciąg nie zbiegał do
$\mu$, istniałby punkt ciągłości jej dystrybuanty
i podciąg, na którym odchylenie dystrybuanty od granicy
byłoby większe od pewnej dodatniej stałej. Taki podciąg
ma jednak dalszy podciąg zbieżny do $\mu$, co daje
sprzeczność. Zatem $\mu_n\to\mu$ według rozkładu.
\end{proof}

\begin{twierdzenie}[Ciągłe przekształcenia i twierdzenie Slutsky'ego]
Niech $X_n,X,Y_n$ będą rzeczywistymi zmiennymi losowymi.
Jeżeli $X_n\xrightarrow{d}X$ i $g:\R\to\R$ jest ciągła,
to $g(X_n)\xrightarrow{d}g(X)$.
Jeżeli dodatkowo $X_n,Y_n$ są określone na wspólnej
przestrzeni, $Y_n\xrightarrow{\Prob}c$ dla pewnej stałej
$c\in\R$, to
$X_n+Y_n\xrightarrow{d}X+c$ oraz
$X_nY_n\xrightarrow{d}cX$; jeśli $c\ne0$, także
$X_n/Y_n\xrightarrow{d}X/c$, przy dowolnym określeniu
ilorazu na zdarzeniu $\{Y_n=0\}$.
\end{twierdzenie}
\begin{proof}
Przypomnijmy znaną charakteryzację zbieżności
według rozkładu przez całki ograniczonych funkcji
ciągłych, której jedna strona została dowiedziona
powyżej, a druga wynika przez aproksymację
wskaźników przedziałów funkcjami ciągłymi.
Dla ograniczonej ciągłej funkcji $h$ również
$h\circ g$ jest ograniczona i ciągła, więc
$\E h(g(X_n))\to\E h(g(X))$.

Dla twierdzenia Slutsky'ego zauważmy najpierw,
że zbieżność $X_n\xrightarrow d X$ zapewnia ciasność
rozkładów $X_n$. Na zdarzeniu $\{|X_n|\leq R,
|Y_n-c|\leq\delta\}$ każda z funkcji
$(x,y)\mapsto x+y$, $xy$ oraz $x/y$ (przy $c\ne0$
i $\delta<|c|/2$) różni się od odpowiadającej
funkcji z $y=c$ o wielkość dążącą jednostajnie do zera
gdy $\delta\downarrow0$. Z ciasności i zbieżności
$Y_n$ według prawdopodobieństwa wynika więc, że
różnice tych zmiennych od odpowiednio
$X_n+c$, $cX_n$ i $X_n/c$ dążą do zera według
prawdopodobieństwa. Pozostaje uzasadnić, że takie małe zaburzenia
nie zmieniają granicy rozkładów. Niech $A_n,B_n$
będą na jednej przestrzeni, $A_n-B_n\xrightarrow{\Prob}0$
oraz $B_n\xrightarrow d B$. Dla $\varepsilon>0$
inkluzje zdarzeń dają
\[
F_{B_n}(x-\varepsilon)-\Prob(|A_n-B_n|>\varepsilon)
\leq F_{A_n}(x)
\leq F_{B_n}(x+\varepsilon)+\Prob(|A_n-B_n|>\varepsilon).
\]
Jeśli $x$ jest punktem ciągłości dystrybuanty $B$,
wybieramy $\varepsilon$ tak, aby także
$x-\varepsilon$ i $x+\varepsilon$ były punktami
ciągłości (dystrybuanta ma co najwyżej przeliczalnie
wiele skoków). Przechodząc z $n\to\infty$, a potem
z $\varepsilon\downarrow0$, otrzymujemy
$F_{A_n}(x)\to F_B(x)$. Stosując ten fakt do
sumy, iloczynu i ilorazu, kończymy dowód.
\end{proof}

\section{Jednostajna całkowalność a zbieżność wartości oczekiwanych}

Zbieżność rozkładów i zbieżność według prawdopodobieństwa
nie muszą powodować zbieżności wartości oczekiwanych.
Przykład $X_n=n\ind{\{U\leq1/n\}}$ pokazuje mechanizm:
masa ucieka w ogony coraz rzadszych, ale coraz
większych wartości. Odpowiednim warunkiem wykluczającym
taką sytuację jest jednostajna całkowalność.

\begin{definicja}
Rodzinę całkowalnych rzeczywistych zmiennych losowych
$\mathcal H$ na wspólnej przestrzeni probabilistycznej
nazywamy \emph{jednostajnie całkowalną}, gdy
\[
\lim_{K\to\infty}\sup_{Y\in\mathcal H}
\E[|Y|\ind{\{|Y|>K\}}]=0.
\]
\end{definicja}

\begin{stwierdzenie}[Kryterium momentowe]
Jeżeli dla pewnego $p>1$ rodzina $\mathcal H$ spełnia
$\sup_{Y\in\mathcal H}\E|Y|^p<\infty$,
to jest jednostajnie całkowalna.
\end{stwierdzenie}
\begin{proof}
Na zdarzeniu $\{|Y|>K\}$ zachodzi
$|Y|\leq |Y|^p/K^{p-1}$. Zatem
$\E[|Y|\ind{\{|Y|>K\}}]
\leq K^{1-p}\E|Y|^p$, a prawa strona dąży
jednostajnie do zera, gdy $K\to\infty$.
\end{proof}

\begin{twierdzenie}[Zbieżność Vitaliego w $L^1$]
Niech $X_n,X$ będą rzeczywistymi zmiennymi losowymi
na jednej przestrzeni probabilistycznej.
Jeśli $X_n\xrightarrow{\Prob}X$ i rodzina
$\{X_n:n\geq1\}$ jest jednostajnie całkowalna,
to $X\in L^1$ oraz $\E|X_n-X|\to0$.
W szczególności $\E X_n\to\E X$.
\end{twierdzenie}
\begin{proof}
Wprowadzamy funkcję obcięcia
$\tau_K(x)=\max(-K,\min(x,K))$.
Jednostajna całkowalność zapewnia
$\sup_n\E|X_n|<\infty$: po rozdzieleniu wartości
oczekiwanej na $\{|X_n|\leq K\}$ i dopełnienie
pierwsza część jest nie większa niż $K$, a druga
jest jednostajnie mała dla dużych $K$.
Z kryterium podciągowego wybieramy podciąg
$X_{n_j}\to X$ prawie na pewno. Lemat Fatou
zapewnia $\E|X|\leq\liminf_j\E|X_{n_j}|<\infty$.
Dla $K>0$
mamy $\E|X_n-\tau_K(X_n)|\leq
\E[|X_n|\ind{\{|X_n|>K\}}]$, więc
po prawej stronie otrzymujemy wielkość zbieżną do
zera jednostajnie w $n$. Ponadto
$\E|\tau_K(X_n)-\tau_K(X)|\to0$:
różnica jest ograniczona przez $2K$ i zbiega do zera
według prawdopodobieństwa, zatem jej średnia dąży
do zera z oszacowania
$\E W\leq\varepsilon+2K\Prob(W>\varepsilon)$
dla $0\leq W\leq2K$.
Z całkowalności $X$ wynika również
$\E|X-\tau_K(X)|\to0$ przez zbieżność zdominowaną.
Rozkładamy różnicę $X_n-X$ na trzy składniki
$X_n-\tau_K(X_n)$,
$\tau_K(X_n)-\tau_K(X)$ i
$\tau_K(X)-X$. Dla dowolnego $\varepsilon>0$
wybieramy najpierw duże $K$ zmniejszające pierwszy
i trzeci składnik, a następnie dostatecznie duże $n$
zmniejszające drugi. Otrzymujemy $\E|X_n-X|\to0$.
\end{proof}

\begin{uwaga}
Implikacje między rodzajami zbieżności mają postać
$L^p\Rightarrow\Prob$ i p.n.$\Rightarrow\Prob$
$\Rightarrow d$. Brak strzałki między zbieżnością
prawie na pewno a zbieżnością w $L^p$ bez dodatkowych
założeń. Zbieżność w rozkładzie do \emph{stałej}
implikuje jednak zbieżność według prawdopodobieństwa
do tej stałej.
\end{uwaga}

% Rozszerzenia audytowe rozdzialow 7--9.

\begin{przyklad}[Ta sama granica rozkładów, różne granice średnich]
Niech $U\sim\Unif(0,1)$ i
$X_n=n\ind{\{U<1/n\}}$. Wtedy $X_n\to0$
prawie na pewno (dla każdego $U>0$ od pewnego
indeksu $1/n<U$), a zatem również według rozkładu.
Mimo to $\E X_n=1$ dla wszystkich $n$.
Różnica nie przeczy zbieżności słabej: funkcja
$x\mapsto x$ jest nieograniczona, więc nie jest
jedną z funkcji testowych w charakteryzacji przez
całki z funkcji ograniczonych i ciągłych.
Jednostajna całkowalność zapobiega dokładnie temu
rodzajowi ucieczki masy pierwszych momentów.
\end{przyklad}

\section{Kiedy zbieżność prawdopodobieństw wystarcza do zbieżności trajektorii?}

\begin{twierdzenie}[Kryterium sumowalnych odchyleń]
Niech $X,X_1,X_2,\ldots$ będą zmiennymi losowymi na jednej
przestrzeni probabilistycznej. Załóżmy, że dla każdego
$\varepsilon>0$ zachodzi
$\sum_{n=1}^\infty\Prob(|X_n-X|>\varepsilon)<\infty$.
Wówczas $X_n\to X$ prawie na pewno.
\end{twierdzenie}
\begin{proof}
Dla każdego $r\in\N$ połóżmy
$A_{n,r}=\{|X_n-X|>1/r\}$. Z pierwszego lematu
Borela--Cantellego otrzymujemy $\Prob(\limsup_n A_{n,r})=0$.
Przecięcie zbiorów
$\Omega\setminus\limsup_n A_{n,r}$ po $r\in\N$
ma prawdopodobieństwo jeden: korzystamy tu z przeliczalności
zbioru progów $1/r$. Dla ustalonego $\omega$ z tego
przecięcia i każdego $r$ nierówność
$|X_n(\omega)-X(\omega)|\leq1/r$ zachodzi od pewnego
indeksu, zależnego od $\omega$ i $r$. To jest dokładnie
definicja zbieżności liczbowej $X_n(\omega)\to X(\omega)$.
\end{proof}

\begin{uwaga}
Z kryterium wynika również, że
$\sum_n\E|X_n-X|^p<\infty$ dla pewnego $p>0$
implikuje zbieżność prawie na pewno: wystarczy
nierówność Markowa z progiem $\varepsilon^p$.
Odwrotność nie zachodzi. Zbieżność p.n. mówi,
że dla każdej prawie każdej trajektorii dużych odchyleń
jest skończenie wiele; nie mówi, że \emph{suma}
ich prawdopodobieństw musi być skończona.
\end{uwaga}

\begin{przyklad}[Zbieżność p.n. bez zbieżności w $L^1$]
Niech $U\sim\Unif(0,1)$ i $X_n=n\ind{\{U\leq1/n\}}$.
Dla każdego $U(\omega)>0$ od pewnego $n$ wskaźnik
jest zerowy. Zatem $X_n\to0$ p.n., chociaż
$\E|X_n|=1$ dla wszystkich $n$. Wcześniej ten sam
ciąg pokazał brak zbieżności w $L^1$ mimo zbieżności
według prawdopodobieństwa: tutaj widzimy mocniejszy fakt.
\end{przyklad}

\begin{przyklad}[Zbieżność w każdym $L^p$, lecz nie p.n.]
Dla ciągu wskaźników przedziałów diadycznych
z przykładu z \emph{ciągiem piszącym na maszynie}
mamy $\E|X_{2^m+j}|^p=2^{-m}$ dla każdego $p\geq1$.
Wszystkie te normy dążą do zera, mimo że trajektorie
oscylują między zerem i jedynką nieskończenie często.
Zatem nawet zbieżność w każdym ustalonym $L^p$
nie implikuje zbieżności p.n.
\end{przyklad}

\begin{przyklad}[Zbieżność w rozkładzie do zmiennej, ale nie do tej samej realizacji]
Niech $U\sim\Unif(0,1)$ i $X_n=1-U$, $X=U$.
Każda zmienna $X_n$ ma ten sam rozkład co $X$,
więc $X_n\xrightarrow{d}X$. Jednak
$\Prob(|X_n-X|>1/4)=\Prob(|1-2U|>1/4)=3/4$
dla każdego $n$. To ilustruje, że zbieżność
rozkładów porównuje rozkłady, a nie pary wyników
na wspólnej przestrzeni.
\end{przyklad}

\section{Dokładna charakteryzacja słabej zbieżności}

\begin{twierdzenie}[Kryterium funkcji ograniczonych ciągłych]
Dla borelowskich miar probabilistycznych $\mu_n,\mu$
na $\R$ następujące warunki są równoważne:
\begin{enumerate}
\item dystrybuanty $F_n$ zbiegają do $F$ w każdym
punkcie ciągłości $F$;
\item dla każdej ograniczonej ciągłej funkcji
$h:\R\to\R$ mamy $\int h\,\dd\mu_n\to\int h\,\dd\mu$.
\end{enumerate}
\end{twierdzenie}
\begin{proof}
Implikacja (1)$\Rightarrow$(2) została wykazana w lemacie
o całkach względem rozkładów zbieżnych. Dla kierunku
odwrotnego ustalmy punkt ciągłości $x$ dystrybuanty $F$
oraz $\delta>0$. Konstruujemy funkcje ciągłe
$h^-_\delta,h^+_\delta$ o wartościach w $[0,1]$:
$h^-_\delta=1$ na $(-\infty,x-\delta]$,
$h^-_\delta=0$ na $[x,\infty)$, liniowa pomiędzy;
$h^+_\delta=1$ na $(-\infty,x]$,
$h^+_\delta=0$ na $[x+\delta,\infty)$, również
liniowa pomiędzy. Zachodzi
$h^-_\delta\leq\ind{(-\infty,x]}\leq h^+_\delta$.
Całkując i korzystając z (2), dostajemy
\[
F(x-\delta)\leq\liminf_nF_n(x)
\leq\limsup_nF_n(x)\leq F(x+\delta).
\]
Gdy $\delta\downarrow0$, lewa i prawa strona
dążą do $F(x)$ na mocy ciągłości w $x$.
Jest to pełne uzasadnienie przejścia od całek
testowych do dystrybuant.
\end{proof}

\begin{przyklad}[Ograniczenie funkcji testowej ma znaczenie]
Niech $X_n=n\ind{\{U\leq1/n\}}$ i $X=0$.
Mamy $X_n\xrightarrow{d}0$, zatem
$\E h(X_n)\to h(0)$ dla każdej \emph{ograniczonej}
ciągłej $h$. Dla nieograniczonej funkcji ciągłej
$h(x)=x$ zachodzi jednak $\E h(X_n)=1$.
Nie można więc usunąć ograniczoności testów
bez dodatkowej kontroli ogonów.
\end{przyklad}

\section{Zbieżność wartości oczekiwanych przy zbieżności rozkładów}

\begin{twierdzenie}[Jednostajna całkowalność i zbieżność słaba]
Niech $X_n\xrightarrow{d}X$, gdzie zmienne mogą być
określone na różnych przestrzeniach. Jeśli rodzina
$\{X_n:n\geq1\}$ jest jednostajnie całkowalna
(w sensie jednostajnego oszacowania ich rozkładów),
to $X\in L^1$ i $\E X_n\to\E X$.
\end{twierdzenie}
\begin{proof}
Z jednostajnej całkowalności wynika
$M:=\sup_n\E|X_n|<\infty$. Dla każdego $K>0$
funkcja $x\mapsto|x|\wedge K$ jest ograniczona
i ciągła, więc kryterium funkcji testowych daje
$\E(|X|\wedge K)=\lim_n\E(|X_n|\wedge K)\leq M$.
Monotoniczna zbieżność po $K\uparrow\infty$ daje
$\E|X|\leq M$.

Niech $\tau_K(x)=(-K)\vee(x\wedge K)$.
Funkcja $\tau_K$ jest ograniczona i ciągła, a zatem
$\E\tau_K(X_n)\to\E\tau_K(X)$. Dodatkowo
$|x-\tau_K(x)|\leq |x|\ind{\{|x|>K\}}$,
co pozwala oszacować
\[
|\E X_n-\E X|\leq
\E|X_n-\tau_K(X_n)|
+|\E\tau_K(X_n)-\E\tau_K(X)|
+\E|X-\tau_K(X)|.
\]
Najpierw wybieramy $K$ tak duże, by pierwszy składnik
był mały jednostajnie w $n$, a trzeci mały dzięki
całkowalności $X$. Następnie wybieramy $n$ duże,
aby zmniejszyć środkowy składnik. To daje tezę.
\end{proof}

\begin{uwaga}
Twierdzenie dotyczy tylko wartości oczekiwanych.
Nie można z niego wywnioskować $\E|X_n-X|\to0$,
ponieważ przy zbieżności według rozkładu zmienne nie
muszą być nawet określone na wspólnej przestrzeni.
Zbieżność w $L^1$ jest wnioskiem z twierdzenia Vitaliego
dopiero wtedy, gdy mamy \emph{zbieżność według
prawdopodobieństwa} na tej samej przestrzeni.
\end{uwaga}

% Koniec dodatkow do rozdzialu 7.

\chapter{Prawa wielkich liczb}
\label{ch:lln}

Prawa wielkich liczb formalizują stabilizację częstości
obserwowanych w długiej serii doświadczeń. W rzucie monetą
względna liczba orłów nie musi być bliska $1/2$ po
kilkunastu próbach, lecz przy dużej liczbie niezależnych
rzutów odchylenia stają się rzadkie. Zależnie od tego,
czy badamy prawdopodobieństwo odchylenia dla dużych $n$,
czy całe nieskończone trajektorie średnich, otrzymujemy
słabe albo mocne prawo wielkich liczb.

\begin{lemat}[Tożsamość warstwowa dla wartości oczekiwanej]
Niech $Y$ będzie nieujemną rzeczywistą zmienną losową.
Wówczas, dopuszczając wartość nieskończoną,
\[
\E Y=\int_0^\infty \Prob(Y>t)\,\dd\lambda_1(t).
\]
W szczególności, gdy $\E Y<\infty$, to
$\sum_{k=1}^\infty\Prob(Y>k)\leq\E Y$.
\end{lemat}
\begin{proof}
Dla każdej liczby $y\geq0$ mamy tożsamość
$y=\int_0^\infty\ind{\{y>t\}}\,\dd\lambda_1(t)$.
Podstawiając $y=Y(\omega)$ i stosując twierdzenie
Tonellego do nieujemnej funkcji
$(\omega,t)\mapsto\ind{\{Y(\omega)>t\}}$,
otrzymujemy
$\E Y=\int_0^\infty \E\ind{\{Y>t\}}\,\dd\lambda_1(t)$,
a więc żądaną tożsamość. Ponieważ dla każdego $k\geq1$
i $t\in[k-1,k)$ zachodzi
$\Prob(Y>t)\geq\Prob(Y>k)$, suma całek po
kolejnych jednostkowych przedziałach daje nierówność.
\end{proof}

\section{Sformułowania i wersja wariancyjna}

\begin{definicja}
Niech $(X_n)_{n\geq1}$ będzie ciągiem całkowalnych
rzeczywistych zmiennych losowych na jednej przestrzeni
probabilistycznej i niech
$\overline X_n=n^{-1}\sum_{k=1}^nX_k$.
Dla $m\in\R$ mówimy, że zachodzi \emph{słabe prawo
wielkich liczb} (SPWL) z granicą $m$, jeśli
$\overline X_n\xrightarrow{\Prob}m$.
Mówimy, że zachodzi \emph{mocne prawo wielkich liczb}
(MPWL) z granicą $m$, jeśli
$\overline X_n\to m$ prawie na pewno.
\end{definicja}

Mocne prawo pociąga słabe, ponieważ zbieżność prawie na
pewno implikuje zbieżność według prawdopodobieństwa.
W drugą stronę implikacja nie zachodzi dla dowolnych ciągów
zmiennych losowych. Nawet przy zbieżności prawie na pewno
nie twierdzimy, że średnia od pewnego momentu \emph{zawsze}
jest równa $m$: tylko że jej odchylenie dąży do zera.

\begin{twierdzenie}[SPWL z kontroli wariancji]
Niech $(X_n)_{n\geq1}$ będą rzeczywistymi zmiennymi losowymi
o skończonych drugich momentach. Jeśli
\[
\frac1n\sum_{k=1}^n\E X_k\longrightarrow m,
\qquad
\frac1{n^2}\Var\!\left(\sum_{k=1}^nX_k\right)\longrightarrow0,
\]
to $n^{-1}\sum_{k=1}^nX_k\xrightarrow{\Prob}m$.
W szczególności wystarcza, by $X_k$ były parami
nieskorelowane, miały wspólną średnią $m$ i
$\sum_{k=1}^n\Var(X_k)=o(n^2)$.
\end{twierdzenie}
\begin{proof}
Oznaczmy $a_n=n^{-1}\sum_{k=1}^n\E X_k$.
Z liniowości wartości oczekiwanej
$\E\overline X_n=a_n$, a z jednorodności wariancji
$\Var(\overline X_n)=n^{-2}\Var(\sum_{k\leq n}X_k)$.
Dla ustalonego $\varepsilon>0$ i dostatecznie
dużego $n$ mamy $|a_n-m|<\varepsilon/2$.
Jeżeli $|\overline X_n-m|>\varepsilon$, to
$|\overline X_n-a_n|>\varepsilon/2$. Czebyszew daje
\[
\Prob(|\overline X_n-m|>\varepsilon)
\leq\frac{4}{\varepsilon^2}\Var(\overline X_n)
\longrightarrow0.
\]
Dla parami nieskorelowanych zmiennych kowariancje
między różnymi składnikami znikają, toteż wariancja
sumy jest sumą wariancji. To prowadzi do drugiej części.
\end{proof}

\begin{wniosek}[SPWL dla niezależnych zmiennych o skończonej wariancji]
Jeśli $X_1,X_2,\ldots$ są niezależne, mają wspólną
średnią $m$ i wspólną skończoną wariancję $\sigma^2$,
to $\overline X_n\xrightarrow{\Prob}m$ oraz
$\E|\overline X_n-m|^2=\sigma^2/n\to0$.
\end{wniosek}
\begin{proof}
Z niezależności wynika
$\Var(\sum_{k=1}^nX_k)=n\sigma^2$.
Po podzieleniu przez $n^2$ dostajemy
$\Var(\overline X_n)=\sigma^2/n$.
Ponieważ $\E\overline X_n=m$, ostatnia wielkość
jest równa $\E|\overline X_n-m|^2$.
Zbieżność w $L^2$ implikuje zbieżność według
prawdopodobieństwa.
\end{proof}

\section{SPWL przy samym istnieniu wartości oczekiwanej}

Nierówność Czebyszewa wymaga skończonej wariancji.
Istnieją jednak rozkłady o skończonej średniej i
nieskończonej wariancji, np. odpowiednio ciężkoogonowe
rozkłady potęgowe. Czy empiryczna średnia nadal się
stabilizuje? Dla niezależnych prób o tym samym rozkładzie
wystarcza sama całkowalność.

\begin{twierdzenie}[Słabe prawo wielkich liczb w wersji całkowalnej]
Niech $X_1,X_2,\ldots$ będą niezależnymi zmiennymi
losowymi o tym samym rozkładzie oraz
$\E|X_1|<\infty$. Wówczas, oznaczając $m=\E X_1$,
mamy $\overline X_n\xrightarrow{\Prob}m$.
\end{twierdzenie}
\begin{proof}
Dla $K>0$ obetnijmy duże wartości:
$Y_k^{(K)}=X_k\ind{\{|X_k|\leq K\}}$.
Dla ustalonego $K$ zmienne $Y_k^{(K)}$ są niezależne,
o tym samym rozkładzie i ograniczone modułem przez $K$,
a więc mają skończoną wariancję. Poprzednie twierdzenie
daje $n^{-1}\sum_{k\leq n}Y_k^{(K)}
\xrightarrow{\Prob}\E Y_1^{(K)}$.

Kontrolujemy teraz błąd obcięcia, bez zakładania
skończonej wariancji $X_1$. Liniowość wartości
oczekiwanej i nierówność trójkąta dają
\[
\E\left|\frac1n\sum_{k=1}^n(X_k-Y_k^{(K)})\right|
\leq\frac1n\sum_{k=1}^n\E|X_k-Y_k^{(K)}|
=\E[|X_1|\ind{\{|X_1|>K\}}].
\]
Ostatnie wyrażenie dąży do zera przy $K\to\infty$
z twierdzenia o zbieżności zdominowanej.
Ponadto $|\E Y_1^{(K)}-m|$ jest nie większe od tej
samej wielkości. Ustalmy $\varepsilon>0$. Wybieramy
$K$ tak duże, by $|\E Y_1^{(K)}-m|<\varepsilon/3$
oraz $\E|X_1-Y_1^{(K)}|$ było dowolnie małe.
Rozkładamy $\overline X_n-m$ na średni błąd obcięcia,
średnią zmiennych obciętych pomniejszoną o jej wartość
oczekiwaną oraz deterministyczny błąd średniej.
Dla pierwszego składnika nierówność Markowa daje
\[
\Prob\!\left(\left|\frac1n\sum_{k\leq n}
(X_k-Y_k^{(K)})\right|>\varepsilon/3\right)
\leq\frac3\varepsilon\E|X_1-Y_1^{(K)}|.
\]
Drugi składnik dąży do zera według prawdopodobieństwa
dla ustalonego $K$. Po przejściu z $n\to\infty$,
a następnie z $K\to\infty$, prawdopodobieństwo
$|\overline X_n-m|>\varepsilon$ dąży do zera.
\end{proof}

\section{Nierówność maksymalna Kołmogorowa}

Do mocnego prawa nie wystarczy kontrolować odchylenia
$\overline X_n$ dla pojedynczego indeksu $n$: trzeba
jednocześnie opanować nieskończenie wiele częściowych sum.
Nierówność maksymalna pozwala przejść od sumy wariancji
do prawdopodobieństwa, że \emph{któraś} suma częściowa
przekroczy próg.

\begin{twierdzenie}[Nierówność maksymalna Kołmogorowa]
Niech $Z_1,\ldots,Z_n$ będą niezależnymi rzeczywistymi
zmiennymi losowymi o $\E Z_k=0$ i
$\E Z_k^2<\infty$. Dla $S_j=\sum_{k=1}^j Z_k$
i $a>0$ zachodzi
\[
\Prob\!\left(\max_{1\leq j\leq n}|S_j|\geq a\right)
\leq\frac1{a^2}\sum_{k=1}^n\Var(Z_k).
\]
\end{twierdzenie}
\begin{proof}
Niech $A_k$ oznacza zdarzenie, że $|S_k|\geq a$
po raz pierwszy, tzn.
$A_k=\{|S_1|<a,\ldots,|S_{k-1}|<a,|S_k|\geq a\}$.
Zdarzenia $A_1,\ldots,A_n$ są rozłączne, a ich suma
to zdarzenie z tezy. Dla ustalonego $k$ zmienna
$S_n-S_k$ jest niezależna od $S_1,\ldots,S_k$
i ma średnią zero. Rozwijamy kwadrat na $A_k$:
\[
\E[S_n^2\ind{A_k}]
=\E[S_k^2\ind{A_k}]
+2\E[S_k(S_n-S_k)\ind{A_k}]
+\E[(S_n-S_k)^2\ind{A_k}].
\]
Środkowa całka jest równa zero dzięki niezależności
i zerowej średniej $S_n-S_k$. Pierwsza jest nie mniejsza
niż $a^2\Prob(A_k)$, a ostatnia jest nieujemna.
Sumując po $k$ i wykorzystując rozłączność zdarzeń,
otrzymujemy
$\E S_n^2\geq\sum_{k=1}^n\E[S_n^2\ind{A_k}]
\geq a^2\sum_{k=1}^n\Prob(A_k)$.
Ponieważ $\E S_n=0$ i składniki są niezależne,
$\E S_n^2=\sum_{k\leq n}\Var(Z_k)$. Dzielenie przez
$a^2$ kończy dowód.
\end{proof}

\begin{lemat}[Kryterium zbieżności szeregu niezależnych zmiennych]
Niech $(Z_n)_{n\geq1}$ będą niezależnymi rzeczywistymi
zmiennymi losowymi o zerowych średnich i skończonych
wariancjach. Jeśli $\sum_{n=1}^\infty\Var(Z_n)<\infty$,
to szereg $\sum_{n=1}^\infty Z_n$ jest zbieżny
prawie na pewno.
\end{lemat}
\begin{proof}
Wybierzmy rosnący ciąg indeksów $n_j$ tak, by
$\sum_{k>n_j}\Var(Z_k)\leq2^{-3j}$.
Zastosujmy nierówność maksymalną do skończonego
ciągu $Z_{n_j+1},\ldots,Z_m$. Dla każdego $m>n_j$
\[
\Prob\!\left(\max_{n_j<r\leq m}
\left|\sum_{k=n_j+1}^rZ_k\right|>2^{-j}\right)
\leq 2^{2j}\sum_{k>n_j}\Var(Z_k)\leq2^{-j}.
\]
Przechodząc z $m\to\infty$ przez ciągłość miary
z dołu, uzyskujemy to samo oszacowanie dla
supremum po wszystkich $r>n_j$.
Suma prawych stron jest skończona, więc z pierwszego
lematu Borela--Cantellego wynika, że prawie na pewno
dla wszystkich dostatecznie dużych $j$ mamy
$\sup_{r>n_j}|\sum_{k=n_j+1}^rZ_k|\leq2^{-j}$.
Różnica dowolnych dwóch sum częściowych po indeksie
$n_j$ ma zatem moduł co najwyżej $2^{1-j}$.
Sumy częściowe są prawie na pewno ciągiem Cauchy'ego,
a więc szereg jest zbieżny.
\end{proof}

\begin{lemat}[Lemat Kroneckera]
Niech $(a_n)_{n\geq1}$ będzie ciągiem liczb rzeczywistych.
Jeżeli szereg $\sum_{n=1}^\infty a_n/n$ jest zbieżny,
to $n^{-1}\sum_{k=1}^n a_k\to0$.
\end{lemat}
\begin{proof}
Niech $s_n=\sum_{k=1}^n a_k/k$ oraz $s_0=0$.
Z założenia $s_n\to s\in\R$.
Ponieważ $a_k=k(s_k-s_{k-1})$, sumowanie przez części
daje
$\sum_{k=1}^n a_k=n s_n-\sum_{k=1}^{n-1}s_k$.
Dzieląc przez $n$, otrzymujemy
$\frac1n\sum_{k=1}^n a_k
=s_n-\frac1n\sum_{k=1}^{n-1}s_k$.
Pierwszy składnik dąży do $s$, a drugi również do $s$
na mocy twierdzenia Cesàra o średnich ciągu zbieżnego.
Ich różnica dąży zatem do zera.
\end{proof}

\begin{twierdzenie}[MPWL przy sumowalnych wariancjach]
Niech $X_1,X_2,\ldots$ będą niezależnymi rzeczywistymi
zmiennymi losowymi o skończonych drugich momentach.
Jeśli
$\sum_{n=1}^\infty\Var(X_n)/n^2<\infty$,
to
\[
\frac1n\sum_{k=1}^n(X_k-\E X_k)\longrightarrow0
\quad\text{prawie na pewno}.
\]
Jeżeli ponadto $n^{-1}\sum_{k=1}^n\E X_k\to m$,
to $\overline X_n\to m$ prawie na pewno.
\end{twierdzenie}
\begin{proof}
Zmienna $Z_n=(X_n-\E X_n)/n$ ma średnią zero
i wariancję $\Var(X_n)/n^2$. Z założenia szereg
wariancji jest zbieżny, zatem na mocy kryterium
Kołmogorowa szereg $\sum_{n=1}^\infty Z_n$
zbiega prawie na pewno. Dla każdej trajektorii,
na której ten szereg jest zbieżny, stosujemy
lemat Kroneckera do liczb $a_n=X_n-\E X_n$.
Otrzymujemy pierwszą granicę. Dodanie granicy
średnich wartości oczekiwanych daje drugą.
\end{proof}

\section{Mocne prawo wielkich liczb dla prób niezależnych jednakowo rozłożonych}

Powyższe twierdzenie zakłada zbieżność szeregu
wariancji. Dla prób o wspólnym skończonym drugim momencie
jest to automatyczne, ponieważ $\sum n^{-2}<\infty$.
Istotne jest jednak silniejsze twierdzenie: dla
niezależnych zmiennych o identycznym rozkładzie
wystarcza \emph{skończony pierwszy moment}.

\begin{twierdzenie}[Mocne prawo wielkich liczb Kołmogorowa]
Niech $X_1,X_2,\ldots$ będą niezależnymi rzeczywistymi
zmiennymi losowymi o jednakowym rozkładzie i
$\E|X_1|<\infty$. Wtedy
\[
\frac1n\sum_{k=1}^n X_k\longrightarrow\E X_1
\qquad\text{prawie na pewno}.
\]
\end{twierdzenie}
\begin{proof}
\emph{Krok pierwszy: obcinamy rzadkie wielkie wartości.}
Dla $k\geq1$ zdefiniujmy
$Y_k=X_k\ind{\{|X_k|\leq k\}}$.
W odróżnieniu od obcięcia przy stałym progu stosowanego
w dowodzie słabego prawa tutaj próg rośnie z indeksem.
Z tożsamości warstwowej dla $|X_1|$ wynika
\[
\sum_{k=1}^\infty\Prob(|X_k|>k)
=\sum_{k=1}^\infty\Prob(|X_1|>k)
\leq\E|X_1|<\infty.
\]
Z pierwszego lematu Borela--Cantellego dostajemy,
że prawie na pewno $X_k=Y_k$ dla wszystkich
poza skończenie wieloma indeksami. Dla takiej
trajektorii suma $\sum_{k\leq n}(X_k-Y_k)$
od pewnego momentu jest stała, więc po podzieleniu
przez $n$ dąży do zera.

\emph{Krok drugi: kontrolujemy wariancje obciętych
zmiennych.} Zmienne $Y_k$ są niezależne, mają skończone
drugie momenty i $\Var(Y_k)\leq\E Y_k^2$.
Wystarczy pokazać
$\sum_{k=1}^\infty\E Y_k^2/k^2<\infty$.
Dzięki Tonellemu i jednakowemu rozkładowi $X_k$
\[
\sum_{k=1}^\infty\frac{\E Y_k^2}{k^2}
=\E\left[|X_1|^2
 \sum_{k=1}^\infty\frac{\ind{\{|X_1|\leq k\}}}{k^2}\right].
\]
Dla liczby $x\geq0$ mamy
$x^2\sum_{k\geq\max(1,\lceil x\rceil)}k^{-2}\leq Cx$
z pewną stałą uniwersalną $C$: jeśli $0\leq x\leq1$,
lewa strona jest nie większa niż
$x^2\sum_{k\geq1}k^{-2}\leq Cx$;
jeśli $x>1$, szacujemy ogon szeregu $\sum k^{-2}$
przez stałą razy $1/x$.
Wobec tego ostatnia wartość oczekiwana jest nie większa
niż $C\E|X_1|<\infty$. Możemy więc zastosować
poprzednie twierdzenie do $Y_k$ i otrzymać
$\frac1n\sum_{k=1}^n(Y_k-\E Y_k)\to0$
prawie na pewno.

\emph{Krok trzeci: wracamy do pierwotnych średnich.}
Zdominowana zbieżność daje
$\E Y_k=\E[X_1\ind{\{|X_1|\leq k\}}]\to\E X_1$.
Twierdzenie Cesàra zapewnia zatem
$n^{-1}\sum_{k=1}^n\E Y_k\to\E X_1$.
Wreszcie
\[
\frac1n\sum_{k=1}^n X_k
=\frac1n\sum_{k=1}^n(X_k-Y_k)
+\frac1n\sum_{k=1}^n(Y_k-\E Y_k)
+\frac1n\sum_{k=1}^n\E Y_k.
\]
Pierwszy i drugi składnik dążą do zera prawie na pewno,
a trzeci do $\E X_1$. To kończy dowód.
\end{proof}

\begin{przyklad}[Częstości w rzutach monetą]
Niech $(B_n)$ będą niezależnymi próbami Bernoulliego
z parametrem $p\in[0,1]$. Zmienna
$\overline B_n=n^{-1}\sum_{k=1}^n B_k$ jest częstością
sukcesów w pierwszych $n$ próbach.
Ponieważ $\E B_1=p$ i $\E|B_1|\leq1$,
MPWL daje $\overline B_n\to p$ prawie na pewno.
Twierdzenie nie stwierdza, że po dużej liczbie prób
kolejne sukcesy i porażki muszą się \emph{wyrównywać}:
możliwe są długie serie jednego wyniku, ale ich wpływ
na średnią z całej dotychczasowej próby maleje.
\end{przyklad}

\begin{przyklad}[Dlaczego dla rozkładu Cauchy'ego nie ma analogicznego MPWL]
Niech $C_1,C_2,\ldots$ będą niezależne i mają
standardowy rozkład Cauchy'ego. Z jego funkcji
charakterystycznej $\varphi_C(t)=e^{-|t|}$
wynika dla każdej liczby naturalnej $n$, że
\[
\varphi_{(C_1+\cdots+C_n)/n}(t)
=\bigl[\varphi_C(t/n)\bigr]^n=e^{-|t|}.
\]
Jednoznaczność funkcji charakterystycznej pokazuje,
że każda średnia $n^{-1}\sum_{k\leq n}C_k$ ma ten sam
rozkład co pojedyncze $C_1$. Nie może więc zbiegać
według prawdopodobieństwa do żadnej stałej, ponieważ
wtedy musiałaby również zbiegać do niej według rozkładu.
Brak całkowalności nie jest tylko techniczną luką
dowodu prawa wielkich liczb.
\end{przyklad}

\section{Alternatywny dowód MPWL: kontrola podciągu diadycznego}

Dla zmiennych o skończonym drugim momencie można udowodnić
mocne prawo także bez lematu Kroneckera. Taki dowód
wyjaśnia, dlaczego do zbieżności p.n. potrzebna jest
kontrola \emph{maksimum} sum częściowych, a nie wyłącznie
wariancji pojedynczej średniej.

\begin{twierdzenie}[MPWL dla niezależnych zmiennych o wspólnej wariancji]
Niech $X_1,X_2,\ldots$ będą niezależnymi zmiennymi
losowymi o wspólnej średniej $m\in\R$ oraz
$\sup_n\Var(X_n)\leq C<\infty$.
Wówczas $n^{-1}\sum_{k=1}^nX_k\to m$ prawie na pewno.
\end{twierdzenie}
\begin{proof}
Połóżmy $Z_k=X_k-m$ i $S_n=\sum_{k=1}^nZ_k$.
Dla $\varepsilon>0$ nierówność maksymalna
Kołmogorowa zastosowana do $2^j$ zmiennych daje
\[
\Prob\left(\max_{1\leq k\leq2^j}|S_k|>
\varepsilon2^j\right)
\leq\frac{\sum_{k\leq2^j}\Var(Z_k)}
{\varepsilon^2 2^{2j}}
\leq\frac{C}{\varepsilon^2 2^j}.
\]
Suma po $j\geq1$ jest skończona. Pierwszy lemat
Borela--Cantellego zapewnia, że dla prawie każdej
trajektorii, od pewnego $j$, wszystkie sumy
$S_1,\ldots,S_{2^j}$ mają moduł co najwyżej
$\varepsilon2^j$.
Dla $2^{j-1}<n\leq2^j$ otrzymujemy zatem
$|S_n|/n\leq2\varepsilon$.
Stosując ten argument jednocześnie dla
$\varepsilon=1,1/2,1/3,\ldots$ (przeliczalne
przecięcie zdarzeń prawdopodobieństwa jeden),
dostajemy $S_n/n\to0$ prawie na pewno.
\end{proof}

\begin{uwaga}
Założenie wspólnego rozkładu nie występuje w tym
twierdzeniu; wystarczy wspólna średnia i jednostajne
ograniczenie wariancji. Jest to mocniejsze od
wersji wariancyjnej dla prób iid w aspekcie
jednakowości rozkładu, ale słabsze od prawa Kołmogorowa
przy samym $\E|X_1|<\infty$ w aspekcie momentów.
\end{uwaga}

\begin{przyklad}[MPWL pomimo nieskończonej wariancji]
Niech $X_1,X_2,\ldots$ będą niezależne o rozkładzie
Pareto z gęstością
$f(x)=\alpha x^{-\alpha-1}\ind{\{x\geq1\}}$
dla ustalonego $1<\alpha<2$. Bezpośrednie całkowanie daje
$\E X_1=\alpha/(\alpha-1)<\infty$ oraz
$\E X_1^2=\infty$ (całka
$\int_1^\infty x^{1-\alpha}\,\dd\lambda_1(x)$
jest rozbieżna). Wariant MPWL wykorzystujący
sumowalne wariancje nie ma tu zastosowania, ale
MPWL Kołmogorowa zapewnia
$\overline X_n\to\alpha/(\alpha-1)$ p.n.
\end{przyklad}

\begin{przyklad}[Dlaczego nie można pominąć warunku niezależności]
Niech $Z$ ma rozkład Bernoulliego z parametrem $1/2$
i połóżmy $X_n=Z$ dla każdego $n$. Wszystkie $X_n$
mają ten sam rozkład i skończone momenty dowolnego
rzędu, lecz $\overline X_n=Z$ dla wszystkich $n$.
Zatem średnia nie dąży do $\E Z=1/2$ ani p.n.,
ani według prawdopodobieństwa. Kolejne obserwacje
nie wnoszą nowej informacji, jeśli są identyczną
kopią tej samej zmiennej.
\end{przyklad}

\begin{przyklad}[Różnica między zbieżnością średnich a zbieżnością pojedynczych wyników]
Dla niezależnych $B_n\sim\operatorname{Bernoulli}(1/2)$
średnie $\overline B_n\to1/2$ p.n. Jednocześnie
z drugiego lematu Borela--Cantellego prawie na pewno
występuje nieskończenie wiele jedynek i nieskończenie
wiele zer, a więc ciąg $B_n$ \emph{nie jest} zbieżny.
To stabilizacja średniej, nie stabilizacja
poszczególnych rezultatów.
\end{przyklad}

\chapter{Centralne twierdzenie graniczne}
\label{ch:clt}

Prawa wielkich liczb mówią, że średnia z wielu obserwacji
jest bliska wartości oczekiwanej, ale nie opisują
\emph{kształtu rozkładu} jej fluktuacji. Dla niezależnych
zmiennych o wspólnej skończonej wariancji wiemy już, że
$\Var(\overline X_n)=\sigma^2/n$.
Typowy rozmiar odchylenia średniej od $m$ ma więc rząd
$n^{-1/2}$. Centralne twierdzenie graniczne (CTG)
wyjaśnia, jaki rozkład otrzymamy po przemnożeniu
tych odchyleń przez $\sqrt n$.

\section{Centralne twierdzenie graniczne dla prób niezależnych jednakowo rozłożonych}

\begin{twierdzenie}[Centralne twierdzenie graniczne Lindeberga--Lévy'ego]
Niech $X_1,X_2,\ldots$ będą niezależnymi rzeczywistymi
zmiennymi losowymi o tym samym rozkładzie oraz niech
$m=\E X_1\in\R$ i
$\sigma^2=\Var(X_1)\in(0,\infty)$.
Dla $S_n=\sum_{k=1}^nX_k$ zachodzi
\[
\frac{S_n-nm}{\sigma\sqrt n}
=\frac{\sqrt n(\overline X_n-m)}{\sigma}
\xrightarrow{d}\Normal(0,1).
\]
W szczególności, dla dowolnych $a<b$,
\[
\lim_{n\to\infty}
\Prob\!\left(a<\frac{S_n-nm}{\sigma\sqrt n}\leq b\right)
=\Phi(b)-\Phi(a),
\]
gdzie $\Phi$ jest dystrybuantą standardowego rozkładu
normalnego.
\end{twierdzenie}
\begin{proof}
\emph{Krok pierwszy: sprowadzamy problem do średniej
zero i wariancji jeden.}
Zdefiniujmy $Z_k=(X_k-m)/\sigma$.
Zmienne $Z_k$ są niezależne, jednakowo rozłożone,
$\E Z_k=0$ oraz $\E Z_k^2=1$.
Badana zmienna ma postać
$T_n=(Z_1+\cdots+Z_n)/\sqrt n$.

\emph{Krok drugi: rozwijamy funkcję charakterystyczną
jednego składnika.} Niech $\varphi(t)=\E e^{itZ_1}$.
Ponieważ $\E Z_1^2<\infty$, funkcja $\varphi$ jest
dwukrotnie różniczkowalna, a
$\varphi(0)=1$, $\varphi'(0)=i\E Z_1=0$ oraz
$\varphi''(0)=-\E Z_1^2=-1$.
Ze wzoru Taylora z resztą $o(t^2)$ otrzymujemy
\[
\varphi(t)=1-\frac{t^2}{2}+o(t^2)
\qquad(t\to0).
\]
Ważne jest, że \emph{nie} założyliśmy istnienia
trzeciego momentu: drugi moment wystarcza do
ciągłości drugiej pochodnej, a więc do reszty $o(t^2)$.

\emph{Krok trzeci: wykorzystujemy niezależność.}
Funkcja charakterystyczna sumy niezależnych składników
jest iloczynem ich funkcji charakterystycznych.
Dla ustalonego $t\in\R$ mamy zatem
\[
\varphi_{T_n}(t)
=\left[\varphi\!\left(\frac{t}{\sqrt n}\right)\right]^n
=\left[1-\frac{t^2}{2n}+o\!\left(\frac1n\right)\right]^n.
\]
Jeśli zapiszemy wyrażenie w nawiasie jako
$1+u_n/n$, to $u_n\to-t^2/2$.
Elementarna granica $(1+u_n/n)^n\to e^{-t^2/2}$
pozostaje prawdziwa także dla zespolonych $u_n$:
przy dużym $n$ można użyć gałęzi logarytmu w
otoczeniu jedynki i rozwinięcia
$\log(1+z)=z+O(|z|^2)$.
W rezultacie
$\varphi_{T_n}(t)\to e^{-t^2/2}$ dla każdego $t\in\R$.

\emph{Krok czwarty: identyfikujemy rozkład graniczny.}
Wcześniej wykazaliśmy, że $e^{-t^2/2}$ jest
funkcją charakterystyczną $\Normal(0,1)$.
Jest ciągła w zerze, zatem twierdzenie Lévy'ego o
ciągłości daje $T_n\xrightarrow d\Normal(0,1)$.
Ponieważ rozkład normalny ma ciągłą dystrybuantę,
zbieżność prawdopodobieństw przedziałów $(a,b]$
wynika z różnicy dystrybuant w punktach $a,b$.
\end{proof}

\begin{uwaga}
CTG nie mówi, że $\overline X_n$ ma dla dużych $n$
dokładnie rozkład normalny. Mówi o granicy rozkładów
zmiennych \emph{wycentrowanych i wyskalowanych}.
Dla małych $n$ przybliżenie może być niedokładne,
zwłaszcza gdy rozkład $X_1$ ma znacznie skośne
lub ciężkie ogony. Sam warunek skończonej wariancji
nie dostarcza uniwersalnej szybkości zbieżności.
\end{uwaga}

\section{Przykłady zastosowania CTG}

\begin{przyklad}[Przybliżenie de Moivre'a--Laplace'a]
Niech $B_1,B_2,\ldots$ będą niezależnymi zmiennymi
Bernoulliego z parametrem $p\in(0,1)$ oraz
$S_n=\sum_{k=1}^nB_k\sim\Bin(n,p)$.
Ponieważ $\E B_1=p$ i $\Var(B_1)=p(1-p)$,
CTG daje
\[
\frac{S_n-np}{\sqrt{np(1-p)}}\xrightarrow d\Normal(0,1).
\]
Stąd dla $a<b$ prawdopodobieństwo
$\Prob(np+a\sqrt{np(1-p)}<S_n\leq
np+b\sqrt{np(1-p)})$ dąży do $\Phi(b)-\Phi(a)$.
Dla przybliżania pojedynczych dyskretnych
prawdopodobieństw stosuje się niekiedy poprawkę ciągłości,
ale samo CTG jest twierdzeniem o dystrybuantach
w punktach ciągłości granicznego rozkładu.
\end{przyklad}

\begin{przyklad}[Normalna aproksymacja Poissona]
Dla każdej liczby $\lambda>0$ niech
$N_\lambda\sim\Poi(\lambda)$. Gdy $\lambda\to\infty$,
zachodzi
$(N_\lambda-\lambda)/\sqrt\lambda\xrightarrow d\Normal(0,1)$.
Nie trzeba tu sprowadzać sprawy do sum o całkowitej
liczbie składników. Z funkcji charakterystycznej
rozkładu Poissona otrzymujemy bezpośrednio
\[
\varphi_{(N_\lambda-\lambda)/\sqrt\lambda}(t)
=\exp\!\left(\lambda\left[e^{it/\sqrt\lambda}
-1-\frac{it}{\sqrt\lambda}\right]\right).
\]
Dla ustalonego $t$ rozwinięcie wykładniczej daje
$e^{it/\sqrt\lambda}=1+it/\sqrt\lambda
-t^2/(2\lambda)+o(1/\lambda)$.
Wykładnik zbiega do $-t^2/2$, a więc
funkcja charakterystyczna do $e^{-t^2/2}$.
Z twierdzenia Lévy'ego wynika wskazana zbieżność.
\end{przyklad}

\begin{przyklad}[Od rozkładu t-Studenta do normalnego]
Niech $Z\sim\Normal(0,1)$ oraz niech
$W_1,W_2,\ldots$ będą niezależnymi standardowymi
zmiennymi normalnymi, niezależnymi również od $Z$.
Dla $n\in\N$ połóżmy
$U_n=\sum_{k=1}^nW_k^2\sim\chi_n^2$ oraz
$T_n=Z/\sqrt{U_n/n}\sim t_n$.
Ponieważ $\E W_1^2=1$, MPWL zapewnia
$U_n/n\to1$ prawie na pewno. Z ciągłości funkcji
$x\mapsto1/\sqrt x$ dla $x>0$ oraz z twierdzenia
Slutsky'ego dostajemy $T_n\xrightarrow d Z$.
W ten sposób uzasadniamy wcześniej podaną granicę
rozkładu Studenta przy rosnącej liczbie stopni swobody.
\end{przyklad}

\section{Co rozróżniają prawa graniczne?}

Dla zmiennych niezależnych jednakowo rozłożonych o
$\E|X_1|<\infty$ zarówno SPWL, jak i MPWL zapewniają
$\overline X_n\to m$ w odpowiednich sensach.
CTG wymaga ponadto $0<\Var(X_1)<\infty$ i mówi
coś innego: po pomnożeniu odchylenia $\overline X_n-m$
przez $\sqrt n$ i podzieleniu przez $\sigma$ rozkład
dąży do nieosobliwego rozkładu normalnego.
Na przykład obszar $m\pm 2\sigma/\sqrt n$ ma
przy dużym $n$ prawdopodobieństwo zbliżone do
$\Phi(2)-\Phi(-2)$, a nie po prostu „duże” w sensie
jakościowym.

Na tym kończy się pierwszy blok twierdzeń granicznych.
W kolejnych rozdziałach uzupełnimy go o warunkową wartość
oczekiwaną, koncentrację miary i granice dla
niejednakowo rozłożonych składników; powrócimy też
do problemów przewidywania i empirycznego opisu rozkładu.

\begin{twierdzenie}[Jednostajna zbieżność dystrybuant do granicy ciągłej]
Niech $F_n$ będą dystrybuantami na $\R$, a $F$
niech będzie \emph{ciągłą} dystrybuantą.
Jeżeli $F_n(x)\to F(x)$ dla każdego $x\in\R$,
to $\sup_{x\in\R}|F_n(x)-F(x)|\to0$.
\end{twierdzenie}
\begin{proof}
Ustalmy $\varepsilon>0$. Wybierzmy liczby $a<b$
tak, aby $F(a)<\varepsilon$ oraz
$1-F(b)<\varepsilon$. Ciągłość $F$ na $[a,b]$
pozwala wybrać skończony podział
$a=x_0<x_1<\cdots<x_m=b$, na którego
każdym przedziale przyrost $F$ jest mniejszy
niż $\varepsilon$. Dla dostatecznie dużego $n$
wszystkie skończenie liczne wartości
$|F_n(x_j)-F(x_j)|$ są mniejsze niż $\varepsilon$.
Jeśli $x_j\leq x\leq x_{j+1}$,
to monotoniczność $F_n$ i $F$ daje
$|F_n(x)-F(x)|\leq2\varepsilon$.
Poza $[a,b]$ stosujemy monotoniczność
i małe ogony $F(a),1-F(b)$, uzyskując
$|F_n(x)-F(x)|\leq2\varepsilon$ również tam.
Ponieważ $\varepsilon$ jest dowolne, teza wynika.
\end{proof}

\begin{uwaga}
W szczególności w CTG, ponieważ dystrybuanta
normalna jest ciągła, mamy
$\sup_x|\Prob(T_n\leq x)-\Phi(x)|\to0$.
Jest to mocniejsze sformułowanie niż zbieżność
w pojedynczych punktach, ale nadal \emph{nie podaje
szybkości} tej zbieżności.
\end{uwaga}

\section{Precyzja przybliżenia normalnego i dalsze przykłady}

\begin{przyklad}[Odchylenia losowego spaceru]
Niech $\xi_1,\xi_2,\ldots$ będą niezależne z
$\Prob(\xi_k=1)=\Prob(\xi_k=-1)=1/2$ i
$S_n=\sum_{k=1}^n\xi_k$. Mamy $\E\xi_k=0$,
$\Var(\xi_k)=1$ i $S_n/\sqrt n\xrightarrow d N(0,1)$.
Na przykład dla $a>0$
$\Prob(|S_n|\leq a\sqrt n)\to2\Phi(a)-1$.
Warto porównać ten fakt z MPWL, które daje
$S_n/n\to0$ p.n. Skale $n$ i $\sqrt n$ opisują
dwa różne pytania: zanik średniego odchylenia
i rozkład wahań po właściwym przeskalowaniu.
\end{przyklad}

\begin{przyklad}[Przybliżenie dwumianowe z poprawką ciągłości]
Niech $S_n\sim\Bin(n,p)$, $0<p<1$, i oznaczmy
$\mu=np$, $s=\sqrt{np(1-p)}$.
Dla liczb całkowitych $a\leq b$ stosuje się przybliżenie
\[
\Prob(a\leq S_n\leq b)\approx
\Phi\!\left(\frac{b+1/2-\mu}{s}\right)-
\Phi\!\left(\frac{a-1/2-\mu}{s}\right).
\]
Dodanie $1/2$ odpowiada zastąpieniu każdego punktu
lattice'owego przedziałem o jednostkowej długości.
To poprawka rachunkowa, nie dodatkowa teza samego CTG;
jej dokładność zależy od $n,p$ i położenia progów.
\end{przyklad}

\begin{przyklad}[Gdy klasyczne CTG nie ma zastosowania]
Dla niezależnych standardowych zmiennych Cauchy'ego
$C_1,C_2,\ldots$ zachodzi
$(C_1+\cdots+C_n)/n\stackrel d=C_1$, zatem nie tylko
brakuje skończonej wariancji, ale nawet same średnie
nie zbliżają się rozkładem do stałej.
Nie wolno w takim przypadku mechanicznie użyć
odchylenia standardowego ani normalnego przybliżenia.
\end{przyklad}

\begin{uwaga}[Gdzie w dowodzie CTG używa się drugiego momentu?]
Jeżeli $Z$ ma średnią zero i $\E Z^2=1$, to
$\varphi_Z'(t)=\E(iZe^{itZ})$
oraz $\varphi_Z''(t)=\E(-Z^2e^{itZ})$:
różniczkowanie pod całką uzasadnia odpowiednio
dominacja przez $|Z|$ oraz $Z^2$.
Zdominowana zbieżność implikuje ciągłość
$\varphi_Z''$ w zerze, dlatego wzór Taylora
z resztą Peana daje $\varphi_Z(t)=1-t^2/2+o(t^2)$.
W dowodzie nie ukryto założenia $\E|Z|^3<\infty$.
Jeśli chcemy oszacować szybkość aproksymacji,
zwykle potrzebne są silniejsze założenia, których
to twierdzenie nie zawiera.
\end{uwaga}

\chapter{Warunkowa wartość oczekiwana i projekcje w \texorpdfstring{$L^2$}{L2}}
\label{ch:condexp}

\section{Dlaczego warunkowanie wymaga nowej definicji?}

W elementarnych zadaniach, gdy $B\in\mathcal F$ i
$\Prob(B)>0$, używaliśmy
$\E[X\mid B]=\E[X\ind B]/\Prob(B)$.
Ten zapis opisuje jedną liczbę: średnią wartości $X$
po ograniczeniu próby do $B$. W wielu zastosowaniach
obserwujemy jednak nie jedno zdarzenie, lecz pewną
\emph{informację}: wynik części doświadczenia, wartość
pomiaru lub przynależność do jednego z wielu zbiorów.
Zanim poznamy tę informację, średnia zależna od niej
jest zmienną losową, a nie ustaloną liczbą.

Informację reprezentuje pod-$\sigma$-ciało
$\mathcal G\subseteq\mathcal F$: zdarzenia z
$\mathcal G$ są dokładnie tymi, których zajście
potrafimy rozstrzygnąć na podstawie dostępnej obserwacji.
Wartość $\E[X\mid\mathcal G]$ musi zatem
spełniać dwa warunki. Po pierwsze ma zależeć jedynie
od dostępnej informacji, czyli być
$\mathcal G$-mierzalna. Po drugie na \emph{każdym}
zdarzeniu $A\in\mathcal G$ jej całka powinna być
taka sama jak całka oryginalnej zmiennej $X$.
Ten drugi warunek wyznacza właściwą średnią bez
odwoływania się do nieformalnego „warunkowania
na wynik o zerowym prawdopodobieństwie”.

\begin{definicja}[Warunkowa wartość oczekiwana]
Niech $(\Omega,\mathcal F,\Prob)$ będzie przestrzenią
probabilistyczną, $\mathcal G\subseteq\mathcal F$
niech będzie pod-$\sigma$-ciałem, a $X\in L^1(\Omega,
\mathcal F,\Prob)$. Zmienną losową
$Y\in L^1(\Omega,\mathcal G,\Prob)$ nazywamy
\emph{warunkową wartością oczekiwaną $X$ względem
$\mathcal G$}, gdy
\[
\int_A Y\,\dd\Prob=\int_A X\,\dd\Prob
\qquad\text{dla każdego }A\in\mathcal G.
\]
Piszemy $Y=\E[X\mid\mathcal G]$.
Równość rozumiemy z dokładnością do zbioru
$\Prob$-miary zero.
\end{definicja}

\begin{twierdzenie}[Istnienie i jednoznaczność przez Radona--Nikodyma]
Dla każdej przestrzeni probabilistycznej
$(\Omega,\mathcal F,\Prob)$, pod-$\sigma$-ciała
$\mathcal G\subseteq\mathcal F$ i zmiennej
$X\in L^1(\Prob)$ istnieje warunkowa wartość
oczekiwana $\E[X\mid\mathcal G]$, która jest
jednoznaczna $\Prob$-prawie na pewno.
\end{twierdzenie}
\begin{proof}
Rozłóżmy $X=X^+-X^-$, gdzie
$X^+=\max(X,0)$ i $X^-=\max(-X,0)$.
Na $(\Omega,\mathcal G)$ definiujemy skończone
miary dodatnie $\nu_\pm(A)=\E[X^\pm\ind A]$.
Są przeliczalnie addytywne przez twierdzenie
Tonellego i spełniają $\nu_\pm\ll\Prob|_\mathcal G$:
gdy $\Prob(A)=0$, to $\nu_\pm(A)=0$.
Twierdzenie Radona--Nikodyma, zastosowane do miar
na \emph{mniejszym} $\sigma$-ciele $\mathcal G$,
daje nieujemne funkcje $Y_\pm$, mierzalne względem
$\mathcal G$, takie że
$\nu_\pm(A)=\int_A Y_\pm\,\dd\Prob$ dla
każdego $A\in\mathcal G$. Ponieważ
$\E Y_\pm=\nu_\pm(\Omega)=\E X^\pm<\infty$,
zmienna $Y=Y_+-Y_-$ należy do $L^1$ i spełnia
wymaganą tożsamość.

Jeśli $Y$ i $Z$ spełniają tę samą tożsamość,
różnica $D=Y-Z$ jest $\mathcal G$-mierzalna i
$\int_A D\,\dd\Prob=0$ dla wszystkich
$A\in\mathcal G$. W szczególności dla
$A_+=\{D>0\}$ całka $\E[D\ind{A_+}]=0$.
Nieujemna zmienna $D\ind{A_+}$ jest więc zerowa
p.n. Podobnie dla $A_-=\{D<0\}$.
Otrzymujemy $Y=Z$ p.n. Istnienie i jednoznaczność
wynikają zatem dokładnie z konstrukcji pochodnej
Radona--Nikodyma.
\end{proof}

\begin{uwaga}
Definicja nie wymaga, aby wszystkie zdarzenia
z $\mathcal G$ miały dodatnie prawdopodobieństwo.
Wartości warunkowej wartości oczekiwanej na zbiorach
zerowych nie są wyznaczone jednoznacznie. Jest to
własność definicji, a nie przeszkoda techniczna.
\end{uwaga}

\section{Pierwsze modele: od średniej liczbowej do funkcji obserwacji}

\begin{przyklad}[Brak informacji i pełna informacja]
Jeżeli $\mathcal G=\{\varnothing,\Omega\}$, to
każda $\mathcal G$-mierzalna zmienna rzeczywista jest
stała. Warunek całkowy wymusza
$\E[X\mid\mathcal G]=\E X$ p.n. Gdy
$\mathcal G=\mathcal F$, warunek spełnia sama
zmienna $X$, stąd $\E[X\mid\mathcal F]=X$ p.n.
Warunkowanie interpoluje między całkowitą niewiedzą
i pełną znajomością wyniku.
\end{przyklad}

\begin{twierdzenie}[Średnia na atomach skończonego podziału]
Niech $A_1,\ldots,A_m\in\mathcal F$ tworzą
rozłączny podział $\Omega$ i niech
$\mathcal G=\sigma(A_1,\ldots,A_m)$.
Jeśli $X\in L^1$, to wersją $\E[X\mid\mathcal G]$ jest
\[
\sum_{j:\Prob(A_j)>0}
\frac{\E[X\ind{A_j}]}{\Prob(A_j)}\ind{A_j},
\]
przy dowolnych wartościach na atomach zerowych.
\end{twierdzenie}
\begin{proof}
Zmienna z tezy jest stała na każdym atomie $A_j$,
więc jest $\mathcal G$-mierzalna. Na atomie
z $\Prob(A_j)>0$ jej całka jest równa
$\E[X\ind{A_j}]$; na atomie zerowym obie całki
wynoszą zero. Każde zdarzenie z $\mathcal G$
jest sumą pewnych atomów, zatem dodając równości
całek dla atomów otrzymujemy równość dla wszystkich
$A\in\mathcal G$.
\end{proof}

\begin{przyklad}[Jedna informacja o wyniku dwóch kostek]
Niech $D_1,D_2$ będą niezależnymi wynikami dwóch
uczciwych kostek. Obserwujemy tylko, czy $D_1$ jest
parzysta, tj. $\mathcal G=\sigma(\{D_1\in\{2,4,6\}\})$.
Dla $X=D_1+D_2$ mamy
$\E[D_1\mid D_1\text{ parzysta}]=4$,
$\E[D_1\mid D_1\text{ nieparzysta}]=3$,
a $\E D_2=7/2$. Zatem
\[
\E[X\mid\mathcal G]
=\frac{15}{2}\ind{\{D_1\text{ parzysta}\}}
+\frac{13}{2}\ind{\{D_1\text{ nieparzysta}\}}.
\]
To dwuwartościowa zmienna losowa: zanim zobaczymy
parzystość, nie wiemy, którą z dwóch średnich otrzymamy.
\end{przyklad}

\begin{przyklad}[Wiedza o wartości bezwzględnej nie określa znaku]
Niech $U\sim\Unif(-1,1)$ i
$\mathcal G=\sigma(U^2)$. Symetria rozkładu sprawia,
że $\E[U\mid\mathcal G]=0$, mimo że $U$ nie jest
zmienną stałą. Aby to uzasadnić bez odwoływania się
do intuicji, zauważmy, że dla każdego borelowskiego
$B\subseteq[0,1]$ mamy
$\E[U\ind{\{U^2\in B\}}]=0$:
całka po symetrycznym zbiorze z funkcji nieparzystej
jest zerowa. Ponieważ zdarzenia $\{U^2\in B\}$
tworzą $\sigma(U^2)$, zmienna zerowa spełnia
definicję warunkowej wartości oczekiwanej.
Jednocześnie $\E[U^2\mid\sigma(U^2)]=U^2$.
\end{przyklad}

\begin{twierdzenie}[Zmienna warunkowa jako funkcja obserwacji]
Jeżeli $V:\Omega\to\R$ jest zmienną losową,
$\mathcal G=\sigma(V)$ oraz $Y$ jest rzeczywistą
$\mathcal G$-mierzalną zmienną losową, to istnieje
funkcja borelowska $g:\R\to\R$ taka, że
$Y=g(V)$ na całej $\Omega$. W szczególności
istnieje funkcja borelowska $g$ spełniająca
$\E[X\mid\sigma(V)]=g(V)$ prawie na pewno.
\end{twierdzenie}
\begin{proof}
Każde zdarzenie $A\in\sigma(V)$ ma postać
$V^{-1}(B)$ dla pewnego $B\in\mathcal B(\R)$,
ponieważ rodzina wszystkich takich przeciwobrazów
jest dokładnie $\sigma(V)$. Dla prostej funkcji
$Y=\sum_j c_j\ind{A_j}$ otrzymujemy reprezentację
$Y=g(V)$ z borelowską funkcją
$g=\sum_j c_j\ind{B_j}$.
Dla dowolnej mierzalnej rzeczywistej $Y$
konstruujemy ciąg funkcji prostych $Y_m\to Y$
punktowo (osobno aproksymując części dodatnią
i ujemną). Każdą zapisujemy $Y_m=g_m(V)$.
Funkcja $g(t)=\lim_m g_m(t)$ w punktach, w których
granica jest skończona, oraz $g(t)=0$ gdzie indziej,
jest borelowska. Ponieważ $Y_m(\omega)\to Y(\omega)$
dla każdego $\omega$, dostajemy $g(V(\omega))=Y(\omega)$.
\end{proof}

\section{Reguły rachunkowe i twierdzenie o wieży}

\begin{twierdzenie}[Podstawowe własności warunkowej wartości oczekiwanej]
Niech $X,Y\in L^1(\Omega,\mathcal F,\Prob)$,
$\mathcal G\subseteq\mathcal F$ i $a,b\in\R$.
Wszystkie poniższe równości oraz nierówności zachodzą
prawie na pewno:
\begin{enumerate}
\item \emph{Liniowość:}
$\E[aX+bY\mid\mathcal G]=a\E[X\mid\mathcal G]+b\E[Y\mid\mathcal G]$.
\item \emph{Dodatniość i monotoniczność:} jeśli $X\geq0$,
to $\E[X\mid\mathcal G]\geq0$; jeśli $X\leq Y$,
to $\E[X\mid\mathcal G]\leq\E[Y\mid\mathcal G]$.
\item \emph{Zachowanie całki:}
$\E[\E[X\mid\mathcal G]]=\E X$.
\item \emph{Znane wielkości:} jeśli $X$ jest
$\mathcal G$-mierzalna, to
$\E[X\mid\mathcal G]=X$.
\item \emph{Wyciąganie czynnika:} jeśli $H$ jest
ograniczoną $\mathcal G$-mierzalną zmienną, to
$\E[HX\mid\mathcal G]=H\E[X\mid\mathcal G]$.
\item \emph{Niezależność:} jeśli $X$ jest niezależna
od $\mathcal G$ (tj. od każdego zdarzenia
w $\mathcal G$), to $\E[X\mid\mathcal G]=\E X$.
\end{enumerate}
\end{twierdzenie}
\begin{proof}
Liniowość wynika z liniowości całki na każdym
$A\in\mathcal G$ oraz jednoznaczności definicji.
Dla dodatniości niech $Z=\E[X\mid\mathcal G]$
i $X\geq0$. Gdyby $\Prob(Z<0)>0$, to dla
pewnego $r\in\N$ zdarzenie $A=\{Z\leq-1/r\}$
miałoby dodatnią miarę. Wówczas
$\int_A Z\,\dd\Prob\leq-\Prob(A)/r<0$,
co przeczy $\int_AZ\,\dd\Prob=\int_AX\,\dd\Prob\geq0$.
Monotoniczność otrzymujemy przez zastosowanie
dodatniości do $Y-X$.

Własność (3) to definicja dla $A=\Omega$.
Własność (4) wynika z tego, że sama zmienna $X$
spełnia oba warunki definicji. Dla (5) zaczynamy
od $H=\ind B$, $B\in\mathcal G$: dla
$A\in\mathcal G$ mamy
$\int_A\ind B\,\E[X\mid\mathcal G] \,\dd\Prob
=\int_{A\cap B}X\,\dd\Prob$.
Liniowość rozszerza to na funkcje proste, a
jednostajna aproksymacja ograniczonych funkcji
mierzalnych przez proste oraz zdominowana zbieżność
na dowolne ograniczone $H$.
Wreszcie dla niezależnego $X$ i $A\in\mathcal G$
$\E[X\ind A]=\E X\,\Prob(A)$.
Dla całkowalnych zmiennych wynika to z definicji
niezależności i aproksymacji funkcjami prostymi;
stała $\E X$ spełnia zatem wymaganą tożsamość.
\end{proof}

\begin{twierdzenie}[Własność wieży]
Jeśli $\mathcal G\subseteq\mathcal H\subseteq\mathcal F$
są $\sigma$-ciałami i $X\in L^1$, to
\[
\E[\E[X\mid\mathcal H]\mid\mathcal G]
=\E[X\mid\mathcal G],
\qquad
\E[\E[X\mid\mathcal G]\mid\mathcal H]
=\E[X\mid\mathcal G].
\]
\end{twierdzenie}
\begin{proof}
Pierwsza lewa strona jest $\mathcal G$-mierzalna.
Dla $A\in\mathcal G\subseteq\mathcal H$
na mocy dwóch kolejnych definicji mamy
\[
\int_A \E[\E[X\mid\mathcal H]\mid\mathcal G]
\,\dd\Prob
=\int_A\E[X\mid\mathcal H]\,\dd\Prob
=\int_A X\,\dd\Prob.
\]
Jednoznaczność daje pierwszą równość.
Dla drugiej zauważamy, że $\E[X\mid\mathcal G]$
jest już $\mathcal H$-mierzalna, więc jej ponowne
warunkowanie względem $\mathcal H$ nic nie zmienia.
\end{proof}

\begin{przyklad}[Trzy poziomy informacji]
Niech $D_1,D_2$ będą niezależnymi wynikami kostek,
$X=D_1+D_2$, $\mathcal H=\sigma(D_1)$,
a $\mathcal G=\sigma(\ind{\{D_1\text{ parzysta}\}})$.
Wtedy
$\E[X\mid\mathcal H]=D_1+7/2$,
a następnie $\E[D_1+7/2\mid\mathcal G]$
daje tę samą dwuwartościową zmienną co
bezpośrednie $\E[X\mid\mathcal G]$ z wcześniejszego
przykładu. Własność wieży oznacza, że można
najpierw uśrednić względem informacji dokładniejszej,
a potem grubszej; wynik nie zależy od tego
pośredniego kroku.
\end{przyklad}

\begin{przyklad}[Wieża nie działa dowolnie dla nieporównywalnych informacji]
Niech $\Omega=\{1,2,3\}$, każdy punkt z masą $1/3$,
$X=\ind{\{1\}}$, $\mathcal G=\sigma(\{1\})$
i $\mathcal H=\sigma(\{2\})$.
Zapisując zmienne jako wektory wartości na
$(1,2,3)$, dostajemy
\[
\E[\E[X\mid\mathcal G]\mid\mathcal H]
=(1/2,0,1/2),\qquad
\E[\E[X\mid\mathcal H]\mid\mathcal G]
=(1/2,1/4,1/4).
\]
Te wyniki są różne. Własność wieży zakłada
\emph{zawieranie się} odpowiednich $\sigma$-ciał;
bez niego kolejne projekcje na różne zbiory
informacji nie muszą być przemienne.
\end{przyklad}

\section{Granice pod warunkowaniem i nierówność Jensena}

\begin{twierdzenie}[Monotoniczna i zdominowana zbieżność warunkowa]
Niech $\mathcal G\subseteq\mathcal F$.
\begin{enumerate}
\item Jeśli $0\leq X_n\uparrow X$ p.n., a $X\in L^1$,
to $\E[X_n\mid\mathcal G]\uparrow\E[X\mid\mathcal G]$
prawie na pewno.
\item Jeśli $X_n\to X$ p.n., a $|X_n|\leq W$
p.n. dla wszystkich $n$, gdzie $W\in L^1$, to
$\E[X_n\mid\mathcal G]\to\E[X\mid\mathcal G]$
prawie na pewno oraz w $L^1$.
\end{enumerate}
\end{twierdzenie}
\begin{proof}
W (1) monotoniczność warunkowej wartości oczekiwanej
daje niemalejący ciąg $Y_n=\E[X_n\mid\mathcal G]$.
Jest on ograniczony z góry przez
$\E[X\mid\mathcal G]\in L^1$, więc granica
$Y=\lim_nY_n$ jest skończona p.n. Dla dowolnego
$A\in\mathcal G$ twierdzenie o zbieżności monotonicznej
daje
$\E[Y\ind A]=\lim_n\E[Y_n\ind A]
=\lim_n\E[X_n\ind A]=\E[X\ind A]$.
Jednoznaczność dowodzi (1).

Dla (2) mamy $X_n+W\geq0$ i $W-X_n\geq0$,
więc można zastosować (1) po uprzednim zastąpieniu
ogólnego ciągu przez suprema i infima ogonowe.
Dokładniej, niech
$L_n=\inf_{k\geq n}X_k$ i
$U_n=\sup_{k\geq n}X_k$. Wtedy
$L_n\uparrow X$, $U_n\downarrow X$ p.n., a
$-W\leq L_n\leq X_n\leq U_n\leq W$.
Wersja monotoniczna zastosowana do $L_n+W$ oraz
do $W-U_n$ daje
$\E[L_n\mid\mathcal G]\uparrow\E[X\mid\mathcal G]$
i $\E[U_n\mid\mathcal G]\downarrow\E[X\mid\mathcal G]$.
Monotoniczność i twierdzenie o trzech ciągach
dają zbieżność p.n. warunkowych wartości $X_n$.
Są one ograniczone modułem przez $\E[W\mid\mathcal G]$
na mocy dodatniości, a ten majorant jest całkowalny;
zdominowana zbieżność daje też zbieżność w $L^1$.
\end{proof}

\begin{twierdzenie}[Nierówność Jensena warunkowa]
Niech $X\in L^1(\Prob)$, a $\phi:\R\to\R$ będzie
skończoną funkcją wypukłą taką, że
$\phi(X)\in L^1(\Prob)$. Wówczas dla każdego
$\mathcal G\subseteq\mathcal F$
zachodzi
\[
\phi(\E[X\mid\mathcal G])
\leq\E[\phi(X)\mid\mathcal G]
\qquad\text{prawie na pewno}.
\]
\end{twierdzenie}
\begin{proof}
Dla dowolnego $r\in\mathbb Q$ wybierzmy prostą
podpierającą wykres wypukłej funkcji $\phi$ w punkcie
$r$, czyli liczbę $s_r$ taką, że
$\phi(x)\geq\phi(r)+s_r(x-r)$ dla wszystkich $x\in\R$.
Istnienie $s_r$ wynika z nierówności pomiędzy
lewą i prawą pochodną funkcji wypukłej.
Dodatniość i liniowość warunkowej wartości
oczekiwanej dają
\[
\E[\phi(X)\mid\mathcal G]
\geq\phi(r)+s_r(\E[X\mid\mathcal G]-r)
\quad\text{p.n.}
\]
Ponieważ liczb wymiernych jest przeliczalnie wiele,
możemy wybrać \emph{jeden} zbiór pełnej miary,
na którym nierówność zachodzi dla wszystkich $r$.
Dla ciągłej skończonej funkcji wypukłej zachodzi
$\phi(x)=\sup_{r\in\mathbb Q}\{\phi(r)+s_r(x-r)\}$:
z prawej strony mamy co najwyżej $\phi(x)$,
a dla wymiernych $r\to x$ wyrażenia zbliżają się
do $\phi(x)$, ponieważ lokalnie nachylenia
podpierające są ograniczone. Podstawienie
$x=\E[X\mid\mathcal G]$ daje tezę.
\end{proof}

\begin{wniosek}[Kontrakcja w $L^p$]
Niech $1\leq p<\infty$, $X\in L^p(\Prob)$ i
$\mathcal G\subseteq\mathcal F$. Wówczas
$\E[X\mid\mathcal G]\in L^p$ oraz
$\|\E[X\mid\mathcal G]\|_p\leq\|X\|_p$.
\end{wniosek}
\begin{proof}
Dla $p\geq1$ funkcja $x\mapsto|x|^p$ jest wypukła.
Stosując warunkowego Jensena, a następnie całkując,
dostajemy
$\E|\E[X\mid\mathcal G]|^p\leq
\E\E[|X|^p\mid\mathcal G]=\E|X|^p$.
\end{proof}

\begin{przyklad}[Dlaczego średnia warunkowa zmniejsza błąd średniokwadratowy]
Niech $X=Z+\varepsilon$, gdzie $Z$ jest całkowalna
i $\varepsilon$ jest niezależna od $Z$ o średniej
zero i skończonej wariancji. Jeśli obserwujemy $Z$,
to $\E[X\mid\sigma(Z)]=Z$. Zmienna $Z$ jest więc
średnią przewidywaną na podstawie obserwacji, a
$\varepsilon$ pozostaje nieprzewidywalnym błędem.
Gdy oba składniki są w $L^2$, ich zależność
z formalną projekcją zostanie wyprowadzona poniżej.
\end{przyklad}

\section{Interpretacja geometryczna: rzut ortogonalny w \texorpdfstring{$L^2$}{L2}}

Oznaczmy przez $L^2(\mathcal F)$ przestrzeń
rzeczywistych zmiennych losowych o skończonym drugim
momencie, utożsamiając zmienne równe prawie na pewno.
Iloczyn skalarny $\langle U,V\rangle=\E[UV]$
czyni tę przestrzeń przestrzenią Hilberta.
Podprzestrzeń $L^2(\mathcal G)$ składa się z klas
$\mathcal G$-mierzalnych zmiennych całkowalnych z
kwadratem. Jest domknięta w $L^2(\mathcal F)$:
jeżeli $Y_n\in L^2(\mathcal G)$ i $Y_n\to Y$
w $L^2$, to z kryterium podciągowego istnieje
podciąg zbiegający do $Y$ prawie na pewno.
Jego granica jest $\mathcal G$-mierzalna (po wyborze
odpowiedniej wersji $Y$), co dowodzi domkniętości.

\begin{twierdzenie}[Warunkowa wartość oczekiwana jako rzut]
Niech $X\in L^2(\Omega,\mathcal F,\Prob)$ i
$\mathcal G\subseteq\mathcal F$. Wtedy
$P_{\mathcal G}X:=\E[X\mid\mathcal G]$ jest
jedynym elementem $L^2(\mathcal G)$ spełniającym
\[
\E\big[(X-P_{\mathcal G}X)Z\big]=0
\qquad\text{dla każdego } Z\in L^2(\mathcal G).
\]
W szczególności dla każdego $Z\in L^2(\mathcal G)$
zachodzi tożsamość Pitagorasa
\[
\E|X-Z|^2
=\E|X-P_{\mathcal G}X|^2
+\E|P_{\mathcal G}X-Z|^2.
\]
Zmienna $P_{\mathcal G}X$ jest zatem jedynym
najlepszym przybliżeniem $X$ przez zmienne
$\mathcal G$-mierzalne w normie $L^2$.
\end{twierdzenie}
\begin{proof}
Kontrakcja w $L^2$ dowodzi, że $P_{\mathcal G}X$
należy do $L^2(\mathcal G)$. Z definicji mamy
$\E[(X-P_{\mathcal G}X)\ind A]=0$
dla wszystkich $A\in\mathcal G$.
Liniowość rozszerza równość na proste funkcje
$\mathcal G$-mierzalne. Jeżeli $Z\in L^2(\mathcal G)$,
istnieją proste $Z_j\to Z$ w $L^2$. Na mocy
nierówności Cauchy'ego--Schwarza
\[
\left|\E[(X-P_{\mathcal G}X)(Z_j-Z)]\right|
\leq\|X-P_{\mathcal G}X\|_2\,\|Z_j-Z\|_2\to0.
\]
Przechodzimy do granicy i otrzymujemy ortogonalność.

Teraz rozkładamy
$X-Z=(X-P_{\mathcal G}X)+(P_{\mathcal G}X-Z)$.
Drugi składnik należy do $L^2(\mathcal G)$,
więc wyraz mieszany w kwadracie normy jest zerowy.
To daje tożsamość Pitagorasa. Drugi składnik
po prawej jest nieujemny i znika wtedy i tylko
wtedy, gdy $Z=P_{\mathcal G}X$ w $L^2$.
Jednoznaczność najlepszego przybliżenia jest dowiedziona.
\end{proof}

\begin{wniosek}[Własności operatora projekcji]
Dla $X,Y\in L^2(\mathcal F)$ i
$\mathcal G\subseteq\mathcal F$ operator
$P_{\mathcal G}:L^2(\mathcal F)\to L^2(\mathcal F)$
jest liniowy, samosprzężony, idempotentny
($P_{\mathcal G}^2=P_{\mathcal G}$) i kontrakcyjny.
Jeśli $\mathcal G\subseteq\mathcal H$, to
$P_{\mathcal G}P_{\mathcal H}
=P_{\mathcal H}P_{\mathcal G}=P_{\mathcal G}$.
\end{wniosek}
\begin{proof}
Liniowość i kontrakcję wykazano wcześniej,
a idempotentność wynika z warunkowania zmiennej
już $\mathcal G$-mierzalnej. Z ortogonalności
$X-P_{\mathcal G}X$ do $L^2(\mathcal G)$
dostajemy, dla dowolnych $X,Y$,
$\langle P_{\mathcal G}X,Y\rangle
=\langle P_{\mathcal G}X,P_{\mathcal G}Y\rangle
=\langle X,P_{\mathcal G}Y\rangle$,
co jest samosprzężonością.
Ostatnie równości to dokładnie obie wersje
własności wieży dla zagnieżdżonych $\sigma$-ciał.
\end{proof}

\begin{przyklad}[Rzut na stałe i elementarna tożsamość wariancyjna]
Jeśli $\mathcal G=\{\varnothing,\Omega\}$ i $X\in L^2$,
to $P_{\mathcal G}X=\E X$. Równość Pitagorasa
przy $Z=c\in\R$ daje
$\E(X-c)^2=\Var(X)+(\E X-c)^2$.
Wyjaśnia to geometrycznie, dlaczego \emph{średnia}
minimalizuje oczekiwany kwadrat błędu, a nie np.
mediana (ta ostatnia minimalizuje inny rodzaj straty).
\end{przyklad}

\section{Wariancja warunkowa, dekompozycja błędu i kowariancja}

\begin{definicja}
Dla $X\in L^2(\Prob)$ i pod-$\sigma$-ciała
$\mathcal G\subseteq\mathcal F$ definiujemy
\emph{warunkową wariancję}
\[
\Var(X\mid\mathcal G)
:=\E[(X-\E[X\mid\mathcal G])^2\mid\mathcal G].
\]
Jest to nieujemna, $\mathcal G$-mierzalna
zmienna losowa.
\end{definicja}

\begin{twierdzenie}[Wzór na wariancję całkowitą]
Dla każdego $X\in L^2(\Prob)$ oraz
$\mathcal G\subseteq\mathcal F$ zachodzą
\[
\Var(X\mid\mathcal G)=
\E[X^2\mid\mathcal G]-\E[X\mid\mathcal G]^2
\quad\text{p.n.}
\]
oraz
\[
\Var(X)=\E[\Var(X\mid\mathcal G)]
+\Var(\E[X\mid\mathcal G]).
\]
\end{twierdzenie}
\begin{proof}
Oznaczmy $M=\E[X\mid\mathcal G]$.
Rozwinięcie kwadratu i reguła wyciągania
czynnika (najpierw dla ograniczonych obcięć $M$,
a potem przez zbieżność w $L^1$ wynikającą
z Cauchy'ego--Schwarza) dają
$\E[X M\mid\mathcal G]=M\E[X\mid\mathcal G]=M^2$.
Podobnie $\E[M^2\mid\mathcal G]=M^2$.
Stąd pierwszy wzór.

Zastosujmy następnie Pitagorasa do rozkładu
$X-\E X=(X-M)+(M-\E X)$.
Jego pierwszy składnik jest ortogonalny do
drugiego, bo drugi jest $\mathcal G$-mierzalny.
Otrzymujemy
$\Var(X)=\E(X-M)^2+\E(M-\E X)^2$.
Z własności zachowania wartości oczekiwanej
$\E(M)=\E X$, a z definicji wariancji
warunkowej $\E(X-M)^2=\E\Var(X\mid\mathcal G)$.
To dowodzi drugiej równości.
\end{proof}

\begin{przyklad}[Losowanie z dwóch grup]
Wybieramy grupę $G\in\{0,1\}$ z równymi
prawdopodobieństwami, a następnie obserwujemy
$X=2G+Z$, gdzie $Z$ jest niezależna od $G$,
$\E Z=0$, $\Var(Z)=1$. Wówczas
$\E[X\mid G]=2G$ i
$\Var(X\mid G)=1$. Wzór na wariancję całkowitą
daje $\Var(X)=1+\Var(2G)=1+1=2$.
Pierwszy składnik opisuje rozproszenie \emph{wewnątrz}
grup, drugi różnicę pomiędzy średnimi grupowymi.
\end{przyklad}

\begin{twierdzenie}[Kowariancja całkowita]
Niech $X,Y\in L^2(\Prob)$ i
$\mathcal G\subseteq\mathcal F$. Wówczas
\[
\Cov(X,Y)=\E[\Cov(X,Y\mid\mathcal G)]
+\Cov(\E[X\mid\mathcal G],\E[Y\mid\mathcal G]),
\]
gdzie
$\Cov(X,Y\mid\mathcal G)=
\E[XY\mid\mathcal G]-
\E[X\mid\mathcal G]\E[Y\mid\mathcal G]$.
\end{twierdzenie}
\begin{proof}
Niech $M=\E[X\mid\mathcal G]$ i
$N=\E[Y\mid\mathcal G]$. Własność wieży daje
$\E\Cov(X,Y\mid\mathcal G)=\E[XY]-\E[MN]$.
Ponieważ $\E M=\E X$ i $\E N=\E Y$,
$\Cov(M,N)=\E[MN]-\E X\E Y$.
Dodajemy obie równości i otrzymujemy tezę.
\end{proof}

\section{Warunkowanie względem zmiennej ciągłej i gęstości}

Jeżeli znamy wspólną gęstość pary $(X,V)$,
możemy jawnie obliczyć funkcję $g$ z reprezentacji
$\E[X\mid\sigma(V)]=g(V)$. Należy przy tym
oddzielić równość całkową określającą średnią
od naiwnego warunku „$V=v$”, który w modelu
ciągłym zwykle ma prawdopodobieństwo zero.

\begin{twierdzenie}[Wzór przez gęstość warunkową]
Niech $(X,V)$ będzie rzeczywistym wektorem losowym
o wspólnej gęstości $f_{X,V}$ względem miary
$\lambda_2$ oraz niech $\E|X|<\infty$.
Niech $f_V(v)=\int_\R f_{X,V}(x,v)\,\dd\lambda_1(x)$.
Dla tych $v$, dla których $0<f_V(v)<\infty$ oraz
$\int_\R |x|f_{X,V}(x,v)\,\dd\lambda_1(x)<\infty$,
połóżmy
\[
g(v)=\frac{\int_\R x f_{X,V}(x,v)\,\dd\lambda_1(x)}{f_V(v)},
\]
a w pozostałych punktach niech $g(v)=0$.
Wówczas $g(V)$ jest wersją
$\E[X\mid\sigma(V)]$.
\end{twierdzenie}
\begin{proof}
Z twierdzenia Tonellego i założenia $\E|X|<\infty$
wynika
\[
\int_\R\int_\R |x|f_{X,V}(x,v)
\,\dd\lambda_1(x)\,\dd\lambda_1(v)<\infty.
\]
Stąd licznik w definicji $g$ istnieje dla
$\lambda_1$-prawie wszystkich $v$; zbiory, na których
$f_V=0$ lub $f_V=\infty$, mają zerową masę
względem rozkładu $V$. Ponadto
$\int |g(v)|f_V(v)\,\dd\lambda_1(v)\leq\E|X|$,
więc $g(V)\in L^1$. Dla każdego zbioru borelowskiego
$B\subseteq\R$ twierdzenie Fubiniego daje
\[
\E[g(V)\ind{\{V\in B\}}]
=\int_Bg(v)f_V(v)\,\dd\lambda_1(v)
=\int_B\int_\R xf_{X,V}(x,v)
\,\dd\lambda_1(x)\,\dd\lambda_1(v)
=\E[X\ind{\{V\in B\}}].
\]
Każde zdarzenie z $\sigma(V)$ ma postać
$\{V\in B\}$, więc jest to dokładnie
tożsamość definiująca warunkową wartość oczekiwaną.
\end{proof}

\begin{przyklad}[Jednostajny punkt trójkąta]
Niech $(X,V)$ będzie jednostajny na zbiorze
$T=\{(x,v):0<v<x<1\}$, a więc jego gęstość
wynosi $2\ind T$. Dla $v\in(0,1)$ mamy
$f_V(v)=\int_v^1 2\,\dd\lambda_1(x)=2(1-v)$,
a gęstość $X$ warunkowo przy $V=v$ jest
jednostajna na $(v,1)$. Otrzymujemy
$\E[X\mid V=v]=(1+v)/2$ dla $v\in(0,1)$
i w konsekwencji
$\E[X\mid\sigma(V)]=(1+V)/2$ p.n.
Sprawdźmy to bezpośrednio: dla dowolnego
borelowskiego $B\subseteq(0,1)$
\[
\E[X\ind{\{V\in B\}}]
=2\int_B\!\int_v^1 x\,\dd\lambda_1(x)
\,\dd\lambda_1(v)
=\int_B(1-v^2)\,\dd\lambda_1(v),
\]
podczas gdy
$\E[((1+V)/2)\ind{\{V\in B\}}]
=\int_B (1+v)(1-v)\,\dd\lambda_1(v)$.
Obie całki są równe. Ten rachunek pokazuje
zarówno funkcję warunkową, jak i poprawne znaczenie
warunkowania na konkretną wartość zmiennej ciągłej.
\end{przyklad}

\begin{przyklad}[Suma niezależnych, jednakowo rozłożonych składników]
Niech $X,Y\in L^1$ będą niezależne i jednakowo
rozłożone, $S=X+Y$. Wtedy $\sigma(S)$
nie rozróżnia zamiany $X$ z $Y$. Formalnie
wspólny rozkład $(X,Y)$ jest niezmienniczy na taką
zamianę, więc dla każdego ograniczonego
borelowskiego $h$ zachodzi
$\E[Xh(S)]=\E[Yh(S)]$.
Dodając równości, dostajemy
$2\E[Xh(S)]=\E[Sh(S)]$. Definicja warunkowa
daje zatem $\E[X\mid\sigma(S)]=S/2$,
niezależnie od tego, czy rozkład składników ma gęstość.
\end{przyklad}

\begin{uwaga}[Wersje i zdarzenia zerowe]
Funkcja $g(v)$ w twierdzeniu o gęstości warunkowej
jest określona jednoznacznie tylko dla
$\Law(V)$-prawie wszystkich $v$. Nie ma żadnego
aksjomatu nakazującego konkretną wartość w punkcie,
w którym $f_V(v)=0$. Sformalizowane rozkłady
warunkowe jako rodziny miar zależnych od $v$
można konstruować na standardowych przestrzeniach
borelowskich, ale ta dodatkowa teoria nie jest
potrzebna do definicji i własności
$\E[X\mid\sigma(V)]$.
\end{uwaga}

\section{Co warto zapamiętać przed procesami stochastycznymi?}
Warunkowa wartość oczekiwana działa w dwóch językach.
W języku miary zachowuje całki na zdarzeniach
obserwowalnych. W języku geometrii jest rzutem
ortogonalnym na domkniętą przestrzeń zmiennych
zależnych od dostępnej informacji. Zmiana informacji
zmienia rzut, a własność wieży kontroluje kolejne
etapy ujawniania danych. W drugim tomie właśnie
ten mechanizm będzie podstawą definicji martyngału,
w którym zachodzi $\E[M_{n+1}\mid\mathcal F_n]=M_n$.

\section{Warunkowanie: czynnik nie musi być ograniczony}

W rozdziale~\ref{ch:condexp} reguła wyciągania czynnika została wykazana dla
ograniczonej zmiennej $H$. To założenie nie jest konieczne, ale jego
usunięcie wymaga całkowalności odpowiednich iloczynów.

\begin{twierdzenie}[Wyciąganie czynnika w pełnej wersji]
Niech $(\Omega,\mathcal F,\Prob)$ będzie przestrzenią probabilistyczną,
$\mathcal G\subseteq\mathcal F$, $X\in L^1(\Prob)$ i niech $H$
będzie rzeczywistą $\mathcal G$-mierzalną zmienną losową. Jeżeli
$HX\in L^1(\Prob)$, to $H\E[X\mid\mathcal G]\in L^1(\Prob)$ oraz
\[
\E[HX\mid\mathcal G]=H\E[X\mid\mathcal G]\quad\text{p.n.}
\]
\end{twierdzenie}
\begin{proof}
Z warunkowego Jensena mamy
$|\E[X\mid\mathcal G]|\leq \E[|X|\mid\mathcal G]$.
Dla $H_m=\min(|H|,m)$, który jest nieujemny, ograniczony
i $\mathcal G$-mierzalny, reguła wyciągania czynnika daje
$\E[H_m|X|\mid\mathcal G]=H_m\E[|X|\mid\mathcal G]$.
Po scałkowaniu i przejściu do granicy przez twierdzenie
Beppo Leviego otrzymujemy
\[
\E[|H|\E(|X|\mid\mathcal G)]=\E|HX|<\infty.
\]
W szczególności $H\E[X\mid\mathcal G]$ jest całkowalna.
Dla $H^{(m)}=(-m)\vee(H\wedge m)$ zachodzi
$H^{(m)}X\to HX$ i $H^{(m)}\E[X\mid\mathcal G]\to
H\E[X\mid\mathcal G]$ prawie na pewno, a obie rodziny
są zdominowane odpowiednio przez $|HX|$ oraz
$|H|\E[|X|\mid\mathcal G]$. Stosujemy wcześniejszą
regułę dla $H^{(m)}$, całkujemy po dowolnym $A\in\mathcal G$
i przechodzimy do granicy. Jednoznaczność warunkowej
wartości oczekiwanej kończy dowód.
\end{proof}

\begin{przyklad}[Dlaczego potrzebujemy warunku całkowalności]
Jeśli $X\equiv1$ oraz $H\geq0$ jest $\mathcal G$-mierzalna,
ale $\E H=\infty$, to formalna równość
$\E[HX\mid\mathcal G]=H\E[X\mid\mathcal G]$
nie jest równością dwóch elementów $L^1$.
Można osobno rozwijać warunkowe wartości oczekiwane
nieujemnych zmiennych o nieskończonej średniej,
lecz wymaga to definicji o wartościach w $[0,\infty]$.
\end{przyklad}

\section{Granice i jednostajna całkowalność}

\begin{twierdzenie}[Zbieżność $L^1$ zachowuje się przy warunkowaniu]
Niech $X_n,X\in L^1(\Prob)$, $\mathcal G\subseteq\mathcal F$.
Jeżeli $\E|X_n-X|\to0$, to
\[
\E\bigl|\E[X_n\mid\mathcal G]-\E[X\mid\mathcal G]\bigr|\to0.
\]
W szczególności wyniku tego nie trzeba wyprowadzać
z prawie pewnej zbieżności $X_n$: jest on bezpośrednią
konsekwencją kontrakcji w $L^1$.
\end{twierdzenie}
\begin{proof}
Z liniowości i warunkowej nierówności Jensena
\[
\E\left|\E[X_n-X\mid\mathcal G]\right|
\leq \E\E[|X_n-X|\mid\mathcal G]=\E|X_n-X|\to0.
\]
\end{proof}

\begin{przyklad}[Zbieżność prawie na pewno nie wystarcza]
Na $((0,1),\mathcal B((0,1)),\lambda_1)$ niech
$X_n=n\ind{(0,1/n)}$. Wtedy $X_n\to0$ prawie na pewno,
ale $\E X_n=1$ dla każdego $n$.
Przy trywialnym $\mathcal G$ mamy zatem
$\E[X_n\mid\mathcal G]=1\not\to0=\E[0\mid\mathcal G]$.
To wyjaśnia, dlaczego w twierdzeniu o zbieżności zdominowanej
w rozdziale~\ref{ch:condexp} rzeczywiście potrzebny jest majorant
całkowalny lub inny warunek gwarantujący $L^1$.
\end{przyklad}



\begin{przyklad}[Niezależność warunkowa nie jest niezależnością]
Niech $B_1,B_2$ będą niezależnymi próbami Bernoulliego
z parametrem $1/2$ i $S=B_1+B_2$.
Zmienne $B_1,B_2$ są niezależne, ale po poznaniu
$S=1$ wiemy, że dokładnie jedna z nich jest równa $1$.
Istotnie,
$\Prob(B_1=1,B_2=1\mid S=1)=0$, natomiast
$\Prob(B_1=1\mid S=1)=\Prob(B_2=1\mid S=1)=1/2$.
Iloczyn prawdopodobieństw brzegowych warunkowych
wynosi zatem $1/4$, nie zero. Obserwacja sumy
wprowadza zależność warunkową mimo wcześniejszej
niezależności; w szczególności nie wolno automatycznie
zastępować warunkowej wartości oczekiwanej iloczynu
iloczynem warunkowych wartości oczekiwanych.
\end{przyklad}

\section{Dwa dalsze przykłady warunkowania: rozkład dyskretny i gaussowski}

\begin{przyklad}[Suma niezależnych zmiennych Poissona]
Niech $X\sim\Poi(\lambda)$ i $Y\sim\Poi(\mu)$ będą
niezależne, gdzie $\lambda,\mu>0$, i niech $S=X+Y$.
Z rozkładu splotu wiemy, że $S\sim\Poi(\lambda+\mu)$.
Dla $n\in\N_0$ i $0\leq k\leq n$ obliczamy
\begin{align*}
\Prob(X=k\mid S=n)
&=\frac{\Prob(X=k,Y=n-k)}{\Prob(S=n)}\\
&=\binom nk\left(\frac{\lambda}{\lambda+\mu}\right)^k
\left(\frac{\mu}{\lambda+\mu}\right)^{n-k}.
\end{align*}
Zatem $X$ przy warunku $S=n$ ma rozkład
$\Bin(n,\lambda/(\lambda+\mu))$. Ponieważ
$\sigma(S)$ jest generowana przez rozłączny,
przeliczalny podział $\{S=n\}$,
\[
\E[X\mid\sigma(S)]
=\frac{\lambda}{\lambda+\mu}S,
\quad
\Var(X\mid\sigma(S))
=\frac{\lambda\mu}{(\lambda+\mu)^2}S.
\]
Intuicyjnie suma mówi, ile zdarzeń zaszło,
a parametry $\lambda$ i $\mu$ określają,
jaka część z nich pochodzi z każdego źródła.
\end{przyklad}

\begin{przyklad}[Dwie niezależne obserwacje normalne]
Niech $X,Z\sim\Normal(0,1)$ będą niezależne.
Obserwujemy tylko sumę $S=X+Z$ i chcemy
przewidzieć $X$. Z symetrii otrzymujemy
$\E[X\mid\sigma(S)]=S/2$.
Aby poznać również wariancję warunkową,
połóżmy $D=(X-Z)/2$. Wtedy
$X=S/2+D$, a odwrotną transformacją jest
$X=S/2+D$, $Z=S/2-D$, o wyznaczniku
Jacobianu o module $1$. Wspólna gęstość
$(S,D)$ wynosi
\[
f_{S,D}(s,d)=\frac1{2\pi}
\exp\left(-\frac{s^2}{4}-d^2\right),
\]
co rozkłada się na iloczyn gęstości
$\Normal(0,2)$ dla $S$ i $\Normal(0,1/2)$
dla $D$. Są więc niezależne.
Wobec $X=S/2+D$ i $\E D=0$ mamy
$\Var(X\mid\sigma(S))=\Var(D)=1/2$.
Po obserwacji sumy pozostaje połowa pierwotnej
wariancji pojedynczego składnika.
\end{przyklad}

\chapter{Dalsze twierdzenia graniczne i koncentracja prawdopodobieństwa}
\label{ch:advancedlimits}

\section{Od praw asymptotycznych do skończonych prób}

Prawo wielkich liczb odpowiada na pytanie, czy średnia
stabilizuje się w granicy $n\to\infty$, a CTG opisuje
kształt rozkładu przeskalowanych błędów. W praktyce
często chcemy wiedzieć więcej: jak szybko zanika
prawdopodobieństwo dużego odchylenia, co zrobić, jeśli
składniki mają różne rozkłady, albo czy zbieżność
średniej pozostaje prawdziwa po nieliniowym
przekształceniu. Wprowadzimy narzędzia
odpowiadające na te pytania, zachowując pełne
dowody w przedstawianych wersjach twierdzeń.

\section{Lemat Hoeffdinga i koncentracja średniej}

\begin{lemat}[Nierówność na wariancję zmiennej ograniczonej]
Niech $Z$ będzie rzeczywistą zmienną losową taką,
że $a\leq Z\leq b$ p.n., gdzie $a<b$.
Wówczas $\Var(Z)\leq(b-a)^2/4$.
\end{lemat}
\begin{proof}
Ponieważ $(Z-a)(b-Z)\geq0$, po rozwinięciu
i uśrednieniu otrzymujemy
$\E Z^2\leq(a+b)\E Z-ab$.
Niech $m=\E Z$. Wobec tego
$\Var(Z)=\E Z^2-m^2\leq(b-m)(m-a)$.
Iloczyn dwóch nieujemnych liczb o sumie $b-a$
nie przekracza $(b-a)^2/4$. To dowodzi tezy.
\end{proof}

\begin{lemat}[Lemat Hoeffdinga]
Niech $Z$ będzie rzeczywistą zmienną losową,
$a\leq Z\leq b$ p.n. dla pewnych $a<b$.
Dla każdego $t\in\R$ zachodzi
\[
\E\exp(t(Z-\E Z))
\leq\exp\left(\frac{t^2(b-a)^2}{8}\right).
\]
\end{lemat}
\begin{proof}
Zdefiniujmy $M(t)=\E e^{tZ}$ i
$h(t)=\log M(t)$. Ponieważ $Z$ jest ograniczona,
możemy różniczkować pod znakiem całki dowolną
liczbę razy, a $M(t)>0$. Obliczamy
$h'(t)=M'(t)/M(t)$ oraz
$h''(t)=M''(t)/M(t)-(M'(t)/M(t))^2$.
Dla każdego $t$ wyrażenie $h''(t)$ jest
wariancją $Z$ względem nowej miary
$\dd\Prob_t=e^{tZ}\dd\Prob/M(t)$. Nowa miara jest
probabilistyczna i nadal skupiona na $[a,b]$.
Poprzedni lemat daje więc
$0\leq h''(t)\leq(b-a)^2/4$ dla wszystkich $t$.

Funkcja $g(t)=h(t)-th'(0)$ spełnia $g(0)=g'(0)=0$,
a $g''(t)=h''(t)\leq(b-a)^2/4$.
Całkując tę nierówność dwukrotnie od zera do $t$
(dla $t<0$ całkujemy po odcinku $[t,0]$),
dostajemy $g(t)\leq t^2(b-a)^2/8$.
Ponieważ $h'(0)=\E Z$,
$g(t)=\log\E e^{t(Z-\E Z)}$.
Potęgowanie daje tezę.
\end{proof}

\begin{twierdzenie}[Nierówność Hoeffdinga dla niezależnych składników]
Niech $X_1,\ldots,X_n$ będą niezależnymi rzeczywistymi
zmiennymi losowymi oraz niech $a_k\leq X_k\leq b_k$
p.n., gdzie $a_k<b_k$. Oznaczmy
$S_n=\sum_{k=1}^n(X_k-\E X_k)$.
Dla każdego $u>0$ zachodzą nierówności
\[
\Prob(S_n\geq u)\leq
\exp\left(-\frac{2u^2}{\sum_{k=1}^n(b_k-a_k)^2}\right),
\qquad
\Prob(|S_n|\geq u)\leq
2\exp\left(-\frac{2u^2}{\sum_{k=1}^n(b_k-a_k)^2}\right).
\]
\end{twierdzenie}
\begin{proof}
Dla $t>0$ nierówność Markowa zastosowana
do $e^{tS_n}$ daje
$\Prob(S_n\geq u)\leq e^{-tu}\E e^{tS_n}$.
Z niezależności i lematu Hoeffdinga
$\E e^{tS_n}=\prod_k\E e^{t(X_k-\E X_k)}
\leq\exp(t^2V/8)$,
gdzie $V=\sum_k(b_k-a_k)^2$.
Otrzymujemy oszacowanie
$\Prob(S_n\geq u)\leq\exp(-tu+t^2V/8)$.
Minimalizując kwadratową funkcję wykładnika po $t>0$,
wybieramy $t=4u/V$, skąd $-tu+t^2V/8=-2u^2/V$.
Tę samą nierówność stosujemy do $-S_n$ i korzystamy
z sumowania prawdopodobieństw dwóch ogonów.
\end{proof}

\begin{przyklad}[Ile prób potrzeba do kontroli częstości?]
Jeżeli $B_1,\ldots,B_n$ są niezależnymi próbami
Bernoulliego z parametrem $p$, to $0\leq B_k\leq1$
i $\E\overline B_n=p$. Nierówność Hoeffdinga
przy $u=n\varepsilon$ daje
$\Prob(|\overline B_n-p|\geq\varepsilon)
\leq2e^{-2n\varepsilon^2}$.
Aby prawa strona była nie większa niż zadane
$\delta\in(0,1)$, wystarczy
$n\geq\log(2/\delta)/(2\varepsilon^2)$.
Jest to gwarancja nieasymptotyczna, ważna dla
każdego $n$, choć nie zawsze ostra.
\end{przyklad}

\begin{twierdzenie}[Prawie pewna szybkość stabilizacji średnich]
Niech $X_1,X_2,\ldots$ będą niezależnymi rzeczywistymi
zmiennymi losowymi o wspólnej średniej $m$ oraz
$a\leq X_k\leq b$ p.n. dla wszystkich $k$,
gdzie $a<b$. Dla każdego $\eta>0$ zachodzi
\[
|\overline X_n-m|\leq
(b-a)\sqrt{\frac{(1+\eta)\log n}{2n}}
\quad\text{dla wszystkich dostatecznie dużych }n
\text{ prawie na pewno}.
\]
\end{twierdzenie}
\begin{proof}
Wstawiamy
$\varepsilon_n=(b-a)\sqrt{(1+\eta)\log n/(2n)}$
do dwustronnego oszacowania Hoeffdinga. Dla $n\geq2$
\[
\Prob(|\overline X_n-m|\geq\varepsilon_n)
\leq2e^{-2n\varepsilon_n^2/(b-a)^2}
=2n^{-(1+\eta)}.
\]
Szereg tych prawdopodobieństw jest zbieżny.
Pierwszy lemat Borela--Cantellego stwierdza, że
odchylenia przekraczające próg $\varepsilon_n$
zdarzają się tylko skończenie wiele razy p.n.
\end{proof}

\begin{uwaga}
Oszacowanie rzędu $\sqrt{(\log n)/n}$ jest
wystarczającym, prostym w dowodzie tempem
prawie pewnej zbieżności. Nie należy utożsamiać
go z optymalnymi, bardziej subtelnymi wynikami
o fluktuacjach pojedynczych trajektorii.
\end{uwaga}

\section{Centralne twierdzenie graniczne dla niejednakowych składników}

Warunek jednakowego rozkładu w podstawowym CTG
można osłabić. Problemem stają się wówczas duże,
rzadkie składniki: jeden z nich może zdominować
całą sumę i zniszczyć granicę normalną.
Warunek Lindeberga wymusza, aby ich wkład do
wariancji po odpowiednim skalowaniu zanikał.

\begin{definicja}[Schemat trójkątny i warunek Lindeberga]
Dla $n\geq1$ niech $k_n\in\N$, a
$Y_{n,1},\ldots,Y_{n,k_n}$ niech będą
niezależnymi \emph{w obrębie każdego wiersza}
rzeczywistymi zmiennymi losowymi spełniającymi
$\E Y_{n,k}=0$ i $\E Y_{n,k}^2<\infty$.
Nie zakładamy niezależności między różnymi wierszami.
Połóżmy
$s_n^2=\sum_{k=1}^{k_n}\E Y_{n,k}^2>0$.
Mówimy, że zachodzi \emph{warunek Lindeberga}, jeżeli
 dla każdego $\varepsilon>0$
\[
\frac1{s_n^2}\sum_{k=1}^{k_n}
\E\left[Y_{n,k}^2\ind{\{|Y_{n,k}|>\varepsilon s_n\}}\right]
\longrightarrow0.
\]
\end{definicja}

\begin{twierdzenie}[CTG w wersji Lindeberga: kierunek wystarczający]
Jeżeli schemat trójkątny spełnia powyższe założenia
oraz warunek Lindeberga, to
\[
\frac{1}{s_n}\sum_{k=1}^{k_n}Y_{n,k}
\xrightarrow{d}\Normal(0,1).
\]
\end{twierdzenie}
\begin{proof}
\emph{Krok 1: normalizacja i znikanie największej
wariancji.} Połóżmy $Z_{n,k}=Y_{n,k}/s_n$.
Wtedy $\sum_k\E Z_{n,k}^2=1$, a warunek Lindeberga
przyjmuje postać
$\sum_k\E[Z_{n,k}^2\ind{\{|Z_{n,k}|>\varepsilon\}}]\to0$.
Dla dowolnego $\varepsilon>0$
\[
\E Z_{n,k}^2\leq\varepsilon^2+
\E[Z_{n,k}^2\ind{\{|Z_{n,k}|>\varepsilon\}}],
\]
zatem $\limsup_n\max_k\E Z_{n,k}^2\leq\varepsilon^2$.
Ponieważ $\varepsilon$ jest dowolne,
$\max_k\E Z_{n,k}^2\to0$. To jest warunek,
że żadna pojedyncza zmienna nie niesie zauważalnej
części całkowitej wariancji.

\emph{Krok 2: rozwinięcie pojedynczych funkcji
charakterystycznych.} Ustalmy $t\in\R$.
Dla każdej liczby rzeczywistej $u$ zachodzi
$|e^{iu}-1-iu|\leq u^2/2$ oraz
$|e^{iu}-1-iu+u^2/2|\leq C u^2\min(|u|,1)$
z pewną stałą bezwzględną $C$.
Dla $|u|\leq1$ druga nierówność wynika z reszty
Taylora rzędu trzeciego, a dla $|u|>1$ z oszacowania
wszystkich składników przez stałą razy $u^2$.
Oznaczmy $\psi_{n,k}(t)=\E e^{itZ_{n,k}}$.
Ponieważ $\E Z_{n,k}=0$, dostajemy
\[
\psi_{n,k}(t)=1-\frac{t^2}{2}\E Z_{n,k}^2+r_{n,k}(t),
\qquad
|\psi_{n,k}(t)-1|\leq\frac{t^2}{2}\E Z_{n,k}^2.
\]
Pozostaje sprawdzić, że $\sum_k|r_{n,k}(t)|\to0$.
Dla ustalonego $\delta>0$ rozbijmy oczekiwanie
w reszcie na $\{|Z_{n,k}|\leq\delta\}$ i dopełnienie.
Na pierwszej części mamy oszacowanie
$C_t\delta\E Z_{n,k}^2$, na drugiej
$C_t\E[Z_{n,k}^2\ind{\{|Z_{n,k}|>\delta\}}]$,
gdzie $C_t$ nie zależy od $n,k$. Po zsumowaniu
i wykorzystaniu sumy wariancji równej jeden,
\[
\limsup_n\sum_k|r_{n,k}(t)|\leq C_t\delta.
\]
Możemy puścić $\delta\downarrow0$, uzyskując
$\sum_k|r_{n,k}(t)|\to0$, a więc
$\sum_k(\psi_{n,k}(t)-1)\to-t^2/2$.

\emph{Krok 3: od sumy do iloczynu.} Niech
$z_{n,k}=\psi_{n,k}(t)-1$. Z poprzednich oszacowań
$\max_k|z_{n,k}|\to0$ oraz
$\sum_k|z_{n,k}|\leq t^2/2$.
W konsekwencji
$\sum_k|z_{n,k}|^2\leq
(\max_k|z_{n,k}|)\sum_k|z_{n,k}|\to0$.
Dla dużych $n$ wszystkie $1+z_{n,k}$ leżą
w otoczeniu jedynki, gdzie
$\log(1+z)=z+O(|z|^2)$. Stąd
\[
\log\prod_{k=1}^{k_n}\psi_{n,k}(t)
=\sum_kz_{n,k}+O\!\left(\sum_k|z_{n,k}|^2\right)
\longrightarrow-\frac{t^2}{2}.
\]
Niezależność w każdym wierszu daje
$\varphi_{\sum_kZ_{n,k}}(t)=\prod_k\psi_{n,k}(t)$.
Po potęgowaniu granicą funkcji charakterystycznych
jest $e^{-t^2/2}$. Twierdzenie Lévy'ego kończy dowód.
\end{proof}

\begin{wniosek}[Warunek Lapunowa]
Przy założeniach schematu trójkątnego załóżmy,
że dla pewnego $\delta>0$ istnieją wszystkie momenty
$\E|Y_{n,k}|^{2+\delta}$ i
\[
\frac{1}{s_n^{2+\delta}}
\sum_{k=1}^{k_n}\E|Y_{n,k}|^{2+\delta}
\longrightarrow0.
\]
Wówczas $s_n^{-1}\sum_{k=1}^{k_n}Y_{n,k}$
zbiega według rozkładu do $\Normal(0,1)$.
\end{wniosek}
\begin{proof}
Na zdarzeniu $\{|Y_{n,k}|>\varepsilon s_n\}$
mamy $|Y_{n,k}|^2\leq
|Y_{n,k}|^{2+\delta}/(\varepsilon s_n)^\delta$.
Sumując wartości oczekiwane i dzieląc przez $s_n^2$,
otrzymujemy lewą stronę warunku Lindeberga
nie większą niż $\varepsilon^{-\delta}$ razy
wyrażenie z założenia. To wyrażenie dąży do zera,
więc stosujemy poprzednie CTG.
\end{proof}

\begin{przyklad}[Zmienne niezależne, ale niejednakowo rozłożone]
Niech $X_k$ będą niezależne, $\E X_k=0$,
$|X_k|\leq C$ i
$s_n^2=\sum_{k=1}^n\Var(X_k)$.
Jeśli $s_n\to\infty$, to dla ustalonego
$\varepsilon>0$ od pewnego $n$ mamy
$|X_k|\leq C<\varepsilon s_n$ dla każdego
$k\leq n$. Warunek Lindeberga jest więc
spełniony automatycznie, a zatem
$\sum_{k=1}^nX_k/s_n\xrightarrow d N(0,1)$.
Wspólny rozkład nie jest potrzebny; ważna jest
całkowita wariancja oraz brak dominujących składników.
\end{przyklad}

\begin{przyklad}[Dlaczego warunek na duże składniki jest potrzebny]
W każdym wierszu niech $Y_{n,1}$ przyjmuje
wartości $\pm1$ z prawdopodobieństwami $1/2$,
a $Y_{n,k}=0$ dla $2\leq k\leq n$.
Wiersz jest niezależny, ma $s_n^2=1$,
lecz suma znormalizowana zawsze ma rozkład
Rademachera, a nie normalny.
Warunek Lindeberga nie zachodzi dla
$\varepsilon\in(0,1)$: jedyny niezerowy składnik
wnosi pełną jednostkę do sumy obciętych drugich
momentów.
\end{przyklad}

\section{Metoda delty: granice nieliniowych funkcji średnich}

Centralne twierdzenie graniczne mówi o przybliżeniu
normalnym dla średniej. Jeśli interesuje nas
$f(\overline X_n)$, np. logarytm, odwrotność albo
pierwiastek, nie powinniśmy od nowa analizować
całego rozkładu. Wystarczy lokalne rozwinięcie
funkcji wokół granicy średnich.

\begin{twierdzenie}[Metoda delty pierwszego rzędu]
Niech $T_n$ będą rzeczywistymi zmiennymi losowymi,
$\theta\in\R$, $a_n>0$ i $a_n\to\infty$.
Załóżmy, że $a_n(T_n-\theta)\xrightarrow d Z$
dla pewnej rzeczywistej zmiennej losowej $Z$.
Jeśli $g:\R\to\R$ jest funkcją borelowską, różniczkowalną
w punkcie $\theta$, to
\[
a_n\bigl(g(T_n)-g(\theta)\bigr)
\xrightarrow d g'(\theta)Z.
\]
\end{twierdzenie}
\begin{proof}
Najpierw zauważmy, że rodzina rozkładów
$a_n(T_n-\theta)$ jest ciasna, bo zbiega według
rozkładu. Dla dowolnych $\varepsilon,\delta>0$
\[
\Prob(|T_n-\theta|>\varepsilon)
\leq\Prob(|a_n(T_n-\theta)|>\varepsilon a_n)
\longrightarrow0,
\]
co wynika z ciasności i $\varepsilon a_n\to\infty$.
A zatem $T_n\xrightarrow{\Prob}\theta$.
Z różniczkowalności $g$ w $\theta$ wynika
reprezentacja
$g(x)-g(\theta)=(x-\theta)(g'(\theta)+r(x))$,
gdzie $r(x)\to0$ dla $x\to\theta$
i przyjmujemy $r(\theta)=0$.
Wówczas $r(T_n)\xrightarrow{\Prob}0$.
Po pomnożeniu przez $a_n$ mamy
\[
a_n(g(T_n)-g(\theta))
=[a_n(T_n-\theta)]\,[g'(\theta)+r(T_n)].
\]
Twierdzenie Slutsky'ego daje granicę
$Z\,g'(\theta)$, zgodnie z tezą.
\end{proof}

\begin{wniosek}[Metoda delty drugiego rzędu]
Przy założeniach poprzedniego twierdzenia załóżmy
dodatkowo, że $g$ jest dwukrotnie różniczkowalna
w punkcie $\theta$ oraz $g'(\theta)=0$.
Wtedy
\[
a_n^2(g(T_n)-g(\theta))
\xrightarrow d\frac{g''(\theta)}2Z^2.
\]
\end{wniosek}
\begin{proof}
Z rozwinięcia Taylora w punkcie $\theta$ mamy
$g(x)-g(\theta)=\frac12g''(\theta)(x-\theta)^2
+o((x-\theta)^2)$. Ponieważ $T_n\to\theta$
według prawdopodobieństwa, iloraz reszty przez
$(T_n-\theta)^2$ (z wartością zero w punkcie
równości) dąży do zera według prawdopodobieństwa.
Zatem
$a_n^2(g(T_n)-g(\theta))
=(\frac12g''(\theta)+o_{\Prob}(1))
[a_n(T_n-\theta)]^2$.
Z twierdzenia o ciągłym przekształceniu
$[a_n(T_n-\theta)]^2\xrightarrow d Z^2$;
Slutsky kończy dowód.
\end{proof}

\begin{przyklad}[Asymptotyka logarytmu średniej]
Niech $X_1,X_2,\ldots$ będą niezależne, jednakowo
rozłożone, $\E X_1=m>0$ i
$0<\Var(X_1)=\sigma^2<\infty$.
Przez MPWL prawdopodobieństwo zdarzenia
$\{\overline X_n>0\}$ dąży do jedności;
na jego dopełnieniu możemy umownie ustawić
$T_n=\log m$, a na nim $T_n=\log\overline X_n$.
Z metody delty dla $g(x)=\log x$ w otoczeniu $m$
(rozszerzonej dowolnie poza tym otoczeniem)
otrzymujemy
$\sqrt n(T_n-\log m)\xrightarrow d
\Normal(0,\sigma^2/m^2)$.
Wzór pokazuje, skąd bierze się czynnik $1/m$
w przybliżeniu niepewności logarytmu średniej.
\end{przyklad}

\begin{przyklad}[Gdy pochodna pierwszego rzędu znika]
Niech $Z_1,Z_2,\ldots$ będą niezależne
$\Normal(0,1)$ i $T_n=\overline Z_n$.
Mamy $\sqrt n T_n\xrightarrow d N(0,1)$.
Dla $g(x)=x^2$ zachodzi $g'(0)=0$
i $g''(0)=2$, zatem
$nT_n^2\xrightarrow d Z^2\sim\chi^2_1$.
W tym przykładzie równość rozkładów jest nawet
prawdziwa dla każdego $n$, ponieważ
$\sqrt n\,T_n\sim N(0,1)$. Zastosowanie
metody delty pierwszego rzędu dawałoby tylko
zdegenerowaną granicę zero przy skali $\sqrt n$.
\end{przyklad}

\section{Studentyzacja: nieznana wariancja w CTG}

\begin{twierdzenie}[CTG po zastąpieniu wariancji estymatorem]
Niech $X_1,X_2,\ldots$ będą niezależnymi,
jednakowo rozłożonymi rzeczywistymi zmiennymi losowymi,
$\E X_1=m$ i $0<\Var(X_1)=\sigma^2<\infty$.
Dla $n\geq2$ połóżmy
$V_n=n^{-1}\sum_{k=1}^n(X_k-\overline X_n)^2$.
Wtedy $V_n\to\sigma^2$ prawie na pewno i
\[
\frac{\sqrt n(\overline X_n-m)}{\sqrt{V_n}}
\xrightarrow d\Normal(0,1),
\]
przy dowolnym określeniu ilorazu na zdarzeniu
$\{V_n=0\}$.
\end{twierdzenie}
\begin{proof}
Rozwijając kwadrat, otrzymujemy dokładną tożsamość
$V_n=n^{-1}\sum_{k\leq n}X_k^2-\overline X_n^2$.
Ponieważ $\E X_1^2<\infty$, mocne prawo wielkich
liczb można zastosować zarówno do $X_k$, jak i $X_k^2$:
$\overline X_n\to m$ i
$n^{-1}\sum_{k\leq n}X_k^2\to\E X_1^2$
prawie na pewno. Stąd $V_n\to\sigma^2>0$ p.n.,
a więc $1/\sqrt{V_n}\to1/\sigma$ według
prawdopodobieństwa (prawdopodobieństwo $V_n=0$
równe zero nie jest wymagane dla skończonych $n$).
Z podstawowego CTG
$\sqrt n(\overline X_n-m)\xrightarrow d N(0,\sigma^2)$.
Twierdzenie Slutsky'ego daje wskazaną granicę.
\end{proof}

\begin{uwaga}
\emph{Asymptotyczna} normalność zmiennej
studentyzowanej zachodzi dla wielu rozkładów
z drugim momentem. Nie jest tym samym co
\emph{dokładny} rozkład t-Studenta, który
wynika z dodatkowych własności próby normalnej.
Rozróżnienie to jest ważne przy przejściu
do statystyki matematycznej w tomie IV.
\end{uwaga}

\section{Jednostajna zbieżność dystrybuanty empirycznej}

Dla próby $X_1,\ldots,X_n$ obserwujemy nie tylko
średnią, ale i empiryczny rozkład: każda obserwacja
otrzymuje masę $1/n$. Dystrybuanta empiryczna ma
postać
$F_n(x)=n^{-1}\sum_{k=1}^n\ind{\{X_k\leq x\}}$.
Dla ustalonego $x$ jest to średnia prób Bernoulliego
z parametrem $F(x)$, więc MPWL daje
$F_n(x)\to F(x)$ p.n. Istotne jest mocniejsze pytanie:
czy możemy jednocześnie kontrolować błąd
na \emph{całej prostej}? W tym miejscu potrzebujemy
odrobiny więcej niż punktowego MPWL.

\begin{twierdzenie}[Glivenki--Cantellego]
Niech $X_1,X_2,\ldots$ będą niezależnymi
jednakowo rozłożonymi rzeczywistymi zmiennymi losowymi
o wspólnej dystrybuancie $F$ i niech
$F_n(x)=n^{-1}\sum_{k=1}^n\ind{\{X_k\leq x\}}$.
Wówczas
\[
\sup_{x\in\R}|F_n(x)-F(x)|\longrightarrow0
\qquad\text{prawie na pewno}.
\]
\end{twierdzenie}
\begin{proof}
\emph{Krok pierwszy: rozkład jednostajny.}
Niech $U_1,U_2,\ldots$ będą niezależne
$\Unif(0,1)$ i
$H_n(t)=n^{-1}\sum_{k\leq n}\ind{\{U_k\leq t\}}$.
Dla każdego ustalonego $t\in[0,1]$ MPWL
daje $H_n(t)\to t$ p.n. Zastosujemy je
jednocześnie tylko do przeliczalnego zbioru
liczb wymiernych $t\in[0,1]$. Ich przecięcie
zdarzeń pełnej miary nadal ma miarę jeden.
Ustalmy trajektorię z tego przecięcia i
liczbę całkowitą $m\geq1$. Dla
$0\leq j\leq m$ mamy
$H_n(j/m)\to j/m$.
Jeżeli $j/m\leq t\leq(j+1)/m$, to monotoniczność
dystrybuanty empirycznej daje
\[
H_n(j/m)-\frac{j+1}{m}\leq H_n(t)-t
\leq H_n((j+1)/m)-\frac{j}{m}.
\]
Zatem
$\sup_{t\in[0,1]}|H_n(t)-t|\leq
\max_{0\leq j\leq m}|H_n(j/m)-j/m|+1/m$.
Po przejściu z $n\to\infty$ mamy limsup
nie większy niż $1/m$, dla każdego $m$.
Ponieważ $m$ jest dowolne, wynika stąd
$\sup_{t\in[0,1]}|H_n(t)-t|\to0$ p.n.

\emph{Krok drugi: dowolna dystrybuanta.}
Niech $Q(u)=\inf\{x\in\R:F(x)\geq u\}$,
$u\in(0,1)$, będzie lewostronną funkcją kwantylową.
Jest niemalejąca i borelowska, a jej podstawowa
własność, wynikająca z monotoniczności i
prawostronnej ciągłości $F$, brzmi
$Q(u)\leq x\iff u\leq F(x)$.
Wobec tego zmienne $Q(U_k)$ są niezależne i mają
wspólną dystrybuantę $F$. Dla nich odpowiednia
dystrybuanta empiryczna jest równa
$H_n(F(x))$ dla każdego $x$.
Dlatego
$\sup_x|F_n(x)-F(x)|\leq\sup_{t\in[0,1]}|H_n(t)-t|\to0$
prawie na pewno w skonstruowanym modelu.

Ostatni krok przenosimy do dowolnej próby iid o
tej samej dystrybuancie: skończone układy obserwacji
mają ten sam rozkład, a zdarzenie
$\{\sup_x|F_n(x)-F(x)|\to0\}$ jest mierzalne
względem całego ciągu obserwacji. Mierzalność
wynika z prawostronnej ciągłości $F_n,F$:
$\sup_x|F_n(x)-F(x)|=\sup_{q\in\mathbb Q}
|F_n(q)-F(q)|$. Zdarzenie ma zatem to samo
prawdopodobieństwo w obu modelach, a w modelu
kwantylowym prawdopodobieństwo to wynosi jeden.
\end{proof}

\begin{przyklad}[Dystrybuanta empiryczna próby Bernoulliego]
Niech $X_k\sim\operatorname{Bernoulli}(p)$ będą
niezależne. Dla $x<0$ obie dystrybuanty są zerowe,
a dla $x\geq1$ obie są jedynką. Na $0\leq x<1$
mamy $F_n(x)=1-\overline X_n$ oraz $F(x)=1-p$.
Stąd
$\sup_x|F_n(x)-F(x)|=|\overline X_n-p|$.
W tym szczególnym przypadku twierdzenie
Glivenki--Cantellego jest dokładnie mocnym prawem
wielkich liczb dla częstości sukcesów.
\end{przyklad}

\section{Koncentracja i funkcje Lipschitza}

\begin{przyklad}[Koncentracja odpowiedzi układu, a nie tylko średniej]
Niech $B_1,\ldots,B_n$ będą niezależnymi próbami Bernoulliego,
$S_n=\sum_{k=1}^n B_k$, a $g:\R\to\R$ będzie
funkcją $L$-Lipschitza z $L>0$, tzn.
$|g(x)-g(y)|\leq L|x-y|$ dla wszystkich $x,y$.
Z nierówności Hoeffdinga z rozdziału~\ref{ch:advancedlimits}
wynika dla każdego $\varepsilon>0$
\[
\Prob\bigl(|g(S_n/n)-g(p)|\geq L\varepsilon\bigr)
\leq 2e^{-2n\varepsilon^2}.
\]
Ograniczenie $L>0$ jest konieczne przy zapisanym
zdarzeniu z nierównością $\geq$: dla $L=0$
zdarzenie $\{|g(S_n/n)-g(p)|\geq0\}$ jest pewne,
a prawa strona może być mniejsza od jedności.
Dla $L>0$ wynik wynika z implikacji
$|g(S_n/n)-g(p)|\geq L\varepsilon\Longrightarrow
|S_n/n-p|\geq\varepsilon$.
Ten przykład przygotowuje późniejszą nierówność
Azumy--Hoeffdinga dla martyngałów: po zastąpieniu
niezależnych składników różnicami martyngałowymi
uzyskamy podobną kontrolę.
\end{przyklad}

\begin{uwaga}[Dokładny zakres twierdzenia Lindeberga]
W rozdziale~\ref{ch:advancedlimits} udowodniono
\emph{kierunek wystarczający} twierdzenia Lindeberga--Fellera.
Sama zbieżność sum do rozkładu normalnego nie implikuje
warunku Lindeberga bez dodatkowych hipotez, takich jak
asymptotyczne zanikanie największych wariancji
(w wersji twierdzenia o równoważności).
Warunek Lindeberga kontroluje nie tylko prawdopodobieństwo
rzadkiego dużego skoku, ale jego wkład do \emph{drugiego momentu}.
Zastąpienie $Y_{n,k}^2$ samym wskaźnikiem
$\ind{\{|Y_{n,k}|>\varepsilon s_n\}}$
byłoby zbyt słabe: zdarzenie rzadkie może nieść
niemal całą wariancję sumy.
\end{uwaga}

\begin{przyklad}[Rzadki skok o istotnej wariancji]
Dla każdego $n$ niech $Y_{n,1}$ przyjmuje $\sqrt n$
i $-\sqrt n$ z prawdopodobieństwami $1/(2n)$
każde, a zero z prawdopodobieństwem $1-1/n$.
Pozostałe wyrazy wiersza są zerowe. Mamy
$\E Y_{n,1}=0$, $s_n^2=\E Y_{n,1}^2=1$.
Zmienna sumy zbiega do zera według prawdopodobieństwa,
bo prawdopodobieństwo niezerowego wyniku wynosi $1/n$,
choć jej wariancja pozostaje równa jeden.
Dla dowolnego $\varepsilon>0$ i $n>\varepsilon^2$
wyrażenie z warunku Lindeberga wynosi $1$.
Nie można więc uzyskać granicy normalnej przez
argument opierający się wyłącznie na malejącej
częstości dużych skoków.
\end{przyklad}

\section{Synteza: jak wybrać właściwe twierdzenie?}

Jeśli interesuje nas wartość średniej przy dużym $n$,
stosujemy SPWL lub MPWL. Jeśli potrzebujemy rozkładu
fluktuacji o skali $n^{-1/2}$, używamy CTG, a dla
niejednakowych składników warunku Lindeberga lub
Lapunowa. Gdy pytamy o konkretną liczbę prób potrzebną
do ograniczenia ryzyka odchylenia, pomocna jest
nierówność Hoeffdinga. Metoda delty opisuje
przekształcone estymatory, studentyzacja pozwala
zastąpić nieznaną wariancję jej empiryczną wersją,
a twierdzenie Glivenki--Cantellego kontroluje
cały rozkład empiryczny, nie tylko jedną średnią.
Wszystkie te rezultaty będą miały dalsze zastosowania
w teorii procesów stochastycznych i statystyce.

\chapter{Łańcuchy Markowa w czasie dyskretnym}
\label{ch:markov}

\section{Dlaczego przeszłość można zastąpić stanem?}

Wyobraźmy sobie układ, który każdego dnia znajduje się w jednym
z ustalonych stanów. Aby przewidzieć następny dzień, zwykle
potrzebujemy pewnej informacji o przeszłości. W modelu Markowa
zakładamy, że \emph{aktualny stan zawiera całą informację o
przeszłości potrzebną do przewidywania następnego stanu}.
Nie jest to założenie o niezależności kolejnych stanów: przeciwnie,
stan jutrzejszy zazwyczaj zależy od dzisiejszego. Kluczowe jest
to, że zależność od wcześniejszych stanów znika \emph{po
uwarunkowaniu na stan obecny}.

\begin{definicja}[Jednorodny łańcuch Markowa]
Niech $(\Omega,\mathcal F,\Prob)$ będzie przestrzenią
probabilistyczną, $S$ niepustym zbiorem skończonym lub
przeliczalnym, wyposażonym w $2^S$, oraz
$X_n:\Omega\to S$ ($n\in\N_0$) ciągiem zmiennych losowych.
Mówimy, że $(X_n)_{n\geq0}$ jest \emph{jednorodnym
łańcuchem Markowa} o macierzy przejścia
$P=(p_{ij})_{i,j\in S}$, gdy
\[
 p_{ij}\geq0,\qquad \sum_{j\in S}p_{ij}=1
 \quad(i\in S),
\]
oraz dla każdego $n\geq0$ i $j\in S$ zachodzi
\[
\Prob(X_{n+1}=j\mid\sigma(X_0,\ldots,X_n))
=p_{X_nj}\quad\text{p.n.}
\]
Początkowy rozkład $\mu_i=\Prob(X_0=i)$
razem z $P$ określa prawo wszystkich skończenie
wymiarowych rozkładów procesu. Macierz $P$ nazywamy
\emph{stochastyczną wierszowo}.
\end{definicja}

\begin{uwaga}
Gdy $\Prob(X_n=i)=0$, zapis
$\Prob(X_{n+1}=j\mid X_n=i)$ jako iloraz liczb
nie jest zdefiniowany. Sformułowanie przez
warunkową wartość oczekiwaną nie ma tej wady:
$p_{X_nj}$ określamy tylko z dokładnością do
prawdopodobieństwa zero. Jeśli przejścia zależą
od $n$, otrzymujemy łańcuch niejednorodny
z macierzami $P_n$.
\end{uwaga}


\begin{uwaga}[Jak czytać prawdopodobieństwa dotyczące procesu?]
Zmienna $X_k$ opisuje stan w chwili $k$; zapis
$X_k=i_k$ jest \emph{zdarzeniem} $\{\omega:X_k(\omega)=i_k\}$.
Zapis
\[
\Prob(X_0=i_0,\ldots,X_n=i_n)
=\Prob\left(\bigcap_{k=0}^n\{X_k=i_k\}\right)
\]
oznacza prawdopodobieństwo \emph{dokładnie wskazanej historii},
a nie dowolnej historii, która kończy się stanem $i_n$.
Przecinek wewnątrz argumentu $\Prob$ oznacza więc
przecięcie zdarzeń. Natomiast
$\Prob(X_n=j)$ sumuje prawdopodobieństwa \emph{wszystkich}
historii kończących się w $j$.

Podobnie zapis
$\Prob(X_{n+1}=j\mid X_0=i_0,\ldots,X_n=i_n)$
oznacza iloraz prawdopodobieństw zdarzeń, pod warunkiem
że wskazana historia ma dodatnie prawdopodobieństwo.
Pisanie $\Prob(X_0\mid X_1=j)$ bez wskazania zdarzenia
związanego z $X_0$ nie oznacza jednej liczby.
Można natomiast mówić o \emph{rozkładzie warunkowym}
$\Law(X_0\mid X_1=j)$, podając prawdopodobieństwa
$\Prob(X_0=i\mid X_1=j)$ dla wszystkich $i\in S$,
lub o liczbie $\E[X_0\mid X_1=j]$, gdy stany są liczbami
i $X_0$ jest całkowalna. Są to różne obiekty.
\end{uwaga}

\begin{stwierdzenie}[Wersja elementarna własności Markowa]
Niech $(X_n)$ będzie jednorodnym łańcuchem Markowa
na przeliczalnym $S$ o macierzy $P$. Jeżeli
$\Prob(X_0=i_0,\ldots,X_n=i_n)>0$, to dla każdego $j\in S$
\[
\Prob(X_{n+1}=j\mid X_0=i_0,\ldots,X_n=i_n)=p_{i_nj}.
\]
Jeśli ponadto $\Prob(X_n=i)>0$, to
$\Prob(X_{n+1}=j\mid X_n=i)=p_{ij}$.
\end{stwierdzenie}
\begin{proof}
Niech $A=\{X_0=i_0,\ldots,X_n=i_n\}\in
\sigma(X_0,\ldots,X_n)$. Z definicji łańcucha
$\E[\ind{\{X_{n+1}=j\}}\mid
\sigma(X_0,\ldots,X_n)]=p_{X_nj}$.
Mnożymy tę równość przez $\ind A$ i bierzemy
wartości oczekiwane. Otrzymujemy
$\Prob(A\cap\{X_{n+1}=j\})=p_{i_nj}\Prob(A)$.
Dzielenie przez dodatnie $\Prob(A)$ daje pierwszą
równość. Druga wynika analogicznie, wybierając
$A=\{X_n=i\}$. Jeżeli $\Prob(A)=0$, dzielenie
nie jest dozwolone: wtedy pozostaje poprawne
sformułowanie przez warunkową wartość oczekiwaną.
\end{proof}

\begin{uwaga}[Co jest założeniem, a co twierdzeniem?]
Własności Markowa nie można udowodnić dla \emph{dowolnego}
procesu losowego: definiuje ona szczególną klasę procesów.
Można natomiast udowodnić, że konkretny proces ma tę własność,
a następnie wyprowadzić z niej prawo trajektorii,
własności wielokrokowe i mocną własność Markowa.
Osobne znaczenie ma \emph{zanikanie wpływu stanu początkowego}
po wielu krokach. Jest to twierdzenie asymptotyczne,
które dla skończonych łańcuchów wymaga m.in.
nieprzywiedlności i nieokresowości.
\end{uwaga}

\begin{twierdzenie}[Prawo trajektorii]
Jeżeli $(X_n)$ jest jednorodnym łańcuchem Markowa
na $S$ o rozkładzie początkowym $\mu$ i macierzy $P$,
to dla każdego $n\geq0$ i $i_0,\ldots,i_n\in S$
\[
\Prob(X_0=i_0,\ldots,X_n=i_n)
=\mu_{i_0}\prod_{k=1}^n p_{i_{k-1}i_k},
\]
przy czym iloczyn pusty dla $n=0$ jest równy $1$.
\end{twierdzenie}
\begin{proof}
Dla $n=0$ równość jest definicją $\mu$.
Niech $A=\{X_0=i_0,\ldots,X_{n-1}=i_{n-1}\}$.
Ponieważ $\ind A$ jest $\sigma(X_0,\ldots,X_{n-1})$-
mierzalna, z definicji łańcucha i własności
wyciągania czynnika otrzymujemy
\[
\Prob(A\cap\{X_n=i_n\})
=\E\left[\ind A\E[\ind{\{X_n=i_n\}}
\mid\sigma(X_0,\ldots,X_{n-1})]\right]
=p_{i_{n-1}i_n}\Prob(A).
\]
Indukcja daje wymagany iloczyn, również wtedy,
gdy $\Prob(A)=0$; nie dzielimy tu przez
prawdopodobieństwo zdarzenia warunkującego.
\end{proof}


\begin{twierdzenie}[Iloczynowe prawo trajektorii charakteryzuje łańcuch]
Niech $S$ będzie przeliczalny, $\mu$ rozkładem na $S$,
$P$ macierzą stochastyczną, a $(X_n)_{n\geq0}$ procesem
o wartościach w $S$. Załóżmy, że dla wszystkich $n\geq0$
i wszystkich $i_0,\ldots,i_n\in S$ zachodzi
\[
\Prob(X_0=i_0,\ldots,X_n=i_n)
=\mu_{i_0}\prod_{k=1}^np_{i_{k-1}i_k}.
\]
Wtedy proces jest jednorodnym łańcuchem Markowa
o rozkładzie początkowym $\mu$ i macierzy $P$.
\end{twierdzenie}
\begin{proof}
Dla historii $A=\{X_0=i_0,\ldots,X_n=i_n\}$
o dodatnim prawdopodobieństwie dzielenie wzoru
dla $n+1$ przez wzór dla $n$ daje
$\Prob(X_{n+1}=j\mid A)=p_{i_nj}$.
Każde zdarzenie z $\sigma(X_0,\ldots,X_n)$
jest sumą rozłącznych atomów postaci $A$,
ponieważ $S^{n+1}$ jest przeliczalny.
Na atomach o prawdopodobieństwie zero nie trzeba
niczego sprawdzać; na pozostałych otrzymana
równość mówi, że $p_{X_nj}$ spełnia definicyjną
tożsamość całkową warunkowej wartości oczekiwanej.
Wobec tego
$\E[\ind{\{X_{n+1}=j\}}\mid
\sigma(X_0,\ldots,X_n)]=p_{X_nj}$ p.n.
To jest własność Markowa.
\end{proof}

\begin{twierdzenie}[Warunkowanie wstecz: przeszłość i przyszłość]
Niech $(X_n)$ będzie jednorodnym łańcuchem Markowa
na przeliczalnym $S$ o rozkładzie początkowym
$\mu$ i macierzy $P$ oraz niech $n\geq1$.
Jeśli historia $\{X_1=i_1,\ldots,X_n=i_n\}$
ma dodatnie prawdopodobieństwo, to dla każdego $i\in S$
\[
\Prob(X_0=i\mid X_1=i_1,\ldots,X_n=i_n)
=\Prob(X_0=i\mid X_1=i_1)
=\frac{\mu_i p_{ii_1}}{\sum_{k\in S}\mu_kp_{ki_1}}.
\]
W szczególności, po poznaniu $X_1$ dalsze stany
$X_2,\ldots,X_n$ nie dostarczają dodatkowej informacji
o $X_0$. Nie należy mylić tego z niezależnością
$X_0$ i $X_1$.
\end{twierdzenie}
\begin{proof}
Ze wzoru na prawo trajektorii
\begin{align*}
\Prob(X_0=i, X_1=i_1,\ldots,X_n=i_n)
&=\mu_i p_{ii_1}\,c,\\
\Prob(X_1=i_1,\ldots,X_n=i_n)
&=\left(\sum_{k\in S}\mu_kp_{ki_1}\right)c,
\end{align*}
gdzie $c=\prod_{r=2}^np_{i_{r-1}i_r}$,
a dla $n=1$ przyjmujemy $c=1$. Prawdopodobieństwo
warunkujące jest dodatnie, więc $c>0$ i
$\sum_k\mu_kp_{ki_1}>0$. Po podzieleniu
obu wzorów wspólny czynnik $c$ się skraca.
Tak samo przy $n=1$ otrzymujemy drugi warunek.
To jest \emph{warunkowa niezależność} przeszłości
i dalszej przyszłości względem stanu pośredniego.
\end{proof}


\begin{uwaga}[Dlaczego można warunkować zarówno do przodu, jak i wstecz?]
Dla trzech chwil nie należy pisać
$\Prob(X_0\mid X_1,X_2)$ jako liczby.
Jeżeli $x_1,x_2$ są zaobserwowane, szukamy np.
$\Prob(X_0=i\mid X_1=x_1,X_2=x_2)$.
Po ustaleniu $X_1=x_1$ przyszły czynnik
$p_{x_1x_2}$ jest \emph{ten sam} dla każdego
kandydata $X_0=i$, więc skraca się w ilorazie.
Jest to szczególny przypadek warunkowej niezależności:
$\sigma(X_0)$ oraz $\sigma(X_2,\ldots,X_n)$
są warunkowo niezależne względem $\sigma(X_1)$.
Bez poznania $X_1$ taka niezależność nie musi zachodzić.
\end{uwaga}

\begin{uwaga}[Łańcuch odwrócony w czasie]
Jeśli $\mu_n(j)=\Prob(X_n=j)$, to dla $n\geq1$
i $\mu_n(j)>0$ zachodzi
\[
\Prob(X_{n-1}=i\mid X_n=j)
=\frac{\mu_{n-1}(i)p_{ij}}{\mu_n(j)}.
\]
Wynik jest formułą Bayesa zastosowaną do dwóch
sąsiednich stanów; w ogólności zależy od $n$,
ponieważ rozkłady $\mu_{n-1}$ i $\mu_n$ mogą się zmieniać.
Jeśli $\mu_0=\pi$ jest stacjonarny, macierz
przejścia wstecz ma stałe współrzędne
$\widetilde p_{ji}=\pi_i p_{ij}/\pi_j$
dla $\pi_j>0$. Równowaga szczegółowa oznacza
wówczas $\widetilde p_{ji}=p_{ji}$.
\end{uwaga}

\begin{przyklad}[Pogoda z pamięcią jednego dnia]
Niech $S=\{s,d\}$ oznacza odpowiednio dzień suchy
i deszczowy, a
\[
P=\begin{pmatrix} 4/5&1/5\\ 2/5&3/5\end{pmatrix}.
\]
Jeżeli dzisiaj jest sucho, szansa deszczu jutro
wynosi $1/5$. Prawdopodobieństwo sekwencji
$s,d,d,s$, jeśli zaczynamy na pewno od $s$, jest
$(1/5)(3/5)(2/5)=6/125$. Nie mnożymy
bezwarunkowych prawdopodobieństw dni: kolejne stany
nie są niezależne.
\end{przyklad}


\begin{center}
\begin{tikzpicture}[>=Stealth,
  stan/.style={draw,very thick,ellipse,minimum width=2.5cm,
  minimum height=1.2cm,align=center,font=\small}]
\node[stan,draw=blue!65!black,fill=blue!8] (s) at (0,0) {sucho\\$s$};
\node[stan,draw=orange!75!black,fill=orange!13] (d) at (6,0) {deszcz\\$d$};
\draw[->,thick,blue!65!black] (s) edge[loop above,looseness=5] node[above] {$4/5$} (s);
\draw[->,thick,orange!75!black] (d) edge[loop above,looseness=5] node[above] {$3/5$} (d);
\draw[->,thick] (s) edge[bend left=16] node[above] {$1/5$} (d);
\draw[->,thick] (d) edge[bend left=16] node[below] {$2/5$} (s);
\end{tikzpicture}

{\small\emph{Rysunek.} Graf skierowany dwustanowego łańcucha pogody.
Każda strzałka odpowiada jednej możliwości przejścia;
suma etykiet strzałek wychodzących z danego stanu wynosi $1$.
Pętla oznacza pozostanie w tym samym stanie, a nie brak kolejnego losowania.}
\end{center}

\begin{przyklad}[Jedna ścieżka, wiele możliwych ścieżek]
Dla przedstawionego łańcucha i $X_0=s$ p.n. mamy
$\Prob(X_0=s,X_1=d,X_2=d,X_3=s)=6/125$.
Natomiast $\Prob(X_2=d\mid X_0=s)$ jest sumą
prawdopodobieństw dwóch dopuszczalnych historii
$(s,s,d)$ i $(s,d,d)$:
\[
\Prob(X_2=d\mid X_0=s)
=\underbrace{\frac45\frac15}_{s\to s\to d}
+\underbrace{\frac15\frac35}_{s\to d\to d}
=\frac7{25}.
\]
To wyjaśnia różnicę między \emph{prawdopodobieństwem
trajektorii} i \emph{prawdopodobieństwem stanu w danej chwili}.
Przykładowo
$\Prob(X_2=d\mid X_0=s,X_1=d)=3/5$,
ponieważ po poznaniu stanu $X_1$ historia wcześniejsza
nie zmienia prognozy następnego kroku.
\end{przyklad}


\begin{center}
\begin{tikzpicture}[>=Stealth,scale=.94,
  every node/.style={font=\small},
  pathstate/.style={circle,draw=blue!65!black,fill=blue!8,
  minimum size=9mm,inner sep=1pt}]
\node[pathstate] (s0) at (0,0) {$s$};
\node[pathstate] (s1) at (3,1.45) {$s$};
\node[pathstate,draw=orange!80!black,fill=orange!10] (d1) at (3,-1.45) {$d$};
\node[pathstate] (ss) at (6,2.3) {$s$};
\node[pathstate,draw=orange!80!black,fill=orange!10] (sd) at (6,.6) {$d$};
\node[pathstate] (ds) at (6,-.6) {$s$};
\node[pathstate,draw=orange!80!black,fill=orange!10] (dd) at (6,-2.3) {$d$};
\draw[->,thick] (s0)--node[above]{$4/5$}(s1);
\draw[->,thick] (s0)--node[below]{$1/5$}(d1);
\draw[->] (s1)--node[above]{$4/5$}(ss);
\draw[->,very thick,orange!75!black] (s1)--node[below]{$1/5$}(sd);
\draw[->] (d1)--node[above]{$2/5$}(ds);
\draw[->,very thick,orange!75!black] (d1)--node[below]{$3/5$}(dd);
\node at (0,3.05) {$n=0$};\node at (3,3.05) {$n=1$};\node at (6.7,3.05) {$n=2$};
\end{tikzpicture}

{\small\emph{Rysunek.} Drzewo dwóch kroków z początkowego stanu suchego.
W przeciwieństwie do grafu przejść ten rysunek odróżnia
historie: dwie pomarańczowe strzałki kończą się w deszczu,
a ich prawdopodobieństwa trzeba dodać.
Kolejne krawędzie jednej historii mnożymy.}
\end{center}

\begin{przyklad}[Proces, który nie ma własności Markowa]
Najpierw losujemy ukrytą zmienną $H\in\{1,2\}$
z prawdopodobieństwami $1/2,1/2$. Jeśli $H=1$,
losowania $X_0,X_1,\ldots$ są niezależnymi próbami
Bernoulliego o parametrze $1/4$; jeśli $H=2$,
mają parametr $3/4$. Bez znajomości $H$
obserwacje nie tworzą łańcucha Markowa.
Rzeczywiście, dla $p_1=1/4$, $p_2=3/4$,
\begin{align*}
\Prob(X_2=1\mid X_1=1,X_0=1)
&=\frac{p_1^3+p_2^3}{p_1^2+p_2^2}=\frac7{10},\\
\Prob(X_2=1\mid X_1=1,X_0=0)
&=\frac{p_1^2(1-p_1)+p_2^2(1-p_2)}
{p_1(1-p_1)+p_2(1-p_2)}=\frac12.
\end{align*}
Mimo jednakowego stanu obecnego $X_1=1$ prognoza
zależy od $X_0$. Powód: przeszłe obserwacje
pozwalają wnioskować o ukrytym parametrze $H$,
którego sam $X_1$ nie opisuje. Gdyby do stanu
dołączyć $H$, uzyskalibyśmy bogatszy opis o własności Markowa.
\end{przyklad}

\begin{przyklad}[Ruch po grafie]
Niech $G=(V,E)$ będzie skończonym nieskierowanym grafem,
w którym każdy wierzchołek ma dodatni stopień $d(v)$.
Z wierzchołka $v$ przechodzimy jednostajnie do jednego
z jego sąsiadów: $p_{vw}=1/d(v)$, gdy $\{v,w\}\in E$,
i zero w przeciwnym razie. Jest to spacer losowy
po grafie. Mimo że wybór sąsiada jest jednostajny,
stany o większym stopniu będą zwykle częściej odwiedzane.
\end{przyklad}

\section{Macierze przejścia, półgrupa i własność Markowa}

\begin{twierdzenie}[Chapmana--Kołmogorowa]
Niech $(X_n)$ będzie jednorodnym łańcuchem Markowa
o macierzy $P$ na przeliczalnym $S$. Dla $m,n\geq0$
przyjmijmy $p_{ij}^{(n)}=(P^n)_{ij}$, przy czym
$P^0=I$. Wówczas
\[
\Prob(X_{m+n}=j\mid\sigma(X_0,\ldots,X_m))
=p^{(n)}_{X_mj}\quad\text{p.n.},\qquad
p_{ij}^{(m+n)}=\sum_{k\in S}p_{ik}^{(m)}p_{kj}^{(n)}.
\]
\end{twierdzenie}
\begin{proof}
Oznaczmy $\mathcal F_m^X=\sigma(X_0,\ldots,X_m)$.
Pokażemy indukcyjnie względem $n$, że
$\E[\ind{\{X_{m+n}=j\}}\mid\mathcal F_m^X]
=(P^n)_{X_mj}$. Dla $n=0$ jest to
$\ind{\{X_m=j\}}=(P^0)_{X_mj}$, a dla $n=1$
równość jest właśnie własnością Markowa.
Załóżmy teraz tezę dla $n\geq1$ i każdego
początkowego czasu. Zastosujmy ją z czasem $m+1$,
a następnie użyjmy własności wieży:
\begin{align*}
&\E[\ind{\{X_{m+n+1}=j\}}\mid\mathcal F_m^X]\\
&=\E[(P^n)_{X_{m+1}j}\mid\mathcal F_m^X]\\
&=\E\left[\sum_{k\in S}(P^n)_{kj}
 \ind{\{X_{m+1}=k\}}\mid\mathcal F_m^X\right]\\
&=\sum_{k\in S}(P^n)_{kj}\,p_{X_mk}
 =(P^{n+1})_{X_mj}.
\end{align*}
Przy przedostatniej równości wyrazy szeregu są
nieujemne, więc można przejść z sumą pod znak
warunkowej wartości oczekiwanej przez warunkową
zbieżność monotoniczną. Ostatnią równość daje
definicja mnożenia macierzy stochastycznych.

Aby odczytać z tego wynik probabilistyczny,
wystarczy wziąć $A=\{X_m=i\}$ o dodatnim
prawdopodobieństwie i pomnożyć otrzymaną
tożsamość przez $\ind A$; dostajemy
$\Prob(X_{m+n}=j\mid X_m=i)=(P^n)_{ij}$.
Nie korzystamy z takiego ilorazu, gdy $\Prob(A)=0$.

Pozostaje równość Chapmana--Kołmogorowa. Z
identyczności macierzowej $P^{m+n}=P^mP^n$
wynika dla wszystkich $i,j$:
$(P^{m+n})_{ij}=\sum_{k\in S}(P^m)_{ik}(P^n)_{kj}$.
Alternatywnie można to uzyskać, sumując po stanie
pośrednim $X_m=k$ dla łańcucha skonstruowanego
z deterministycznym $X_0=i$. Ten drugi opis
uzasadnia intuicję: każda droga długości $m+n$
przechodzi po $m$ krokach przez dokładnie jeden
stan $k$, więc rozłączne możliwości należy zsumować.
\end{proof}


\begin{twierdzenie}[Pełna własność zapominania przeszłości]
Niech $(X_n)_{n\geq0}$ będzie jednorodnym łańcuchem
Markowa na przeliczalnym zbiorze $S$, o macierzy $P$,
a $\mathcal F_n^X=\sigma(X_0,\ldots,X_n)$. Dla
$n\geq0$, $r\geq1$ i ograniczonej funkcji
$f:S^r\to\R$ mamy
\[
\E[f(X_{n+1},\ldots,X_{n+r})\mid\mathcal F_n^X]
=\sum_{j_1,\ldots,j_r\in S} f(j_1,\ldots,j_r)
 p_{X_nj_1}p_{j_1j_2}\cdots p_{j_{r-1}j_r}
\quad\text{p.n.}
\]
W szczególności, po poznaniu $X_n$ rozkład
\emph{całego skończonego fragmentu przyszłości}
nie zależy od wcześniejszych stanów $X_0,\ldots,X_{n-1}$.
\end{twierdzenie}
\begin{proof}
Wystarczy najpierw rozważyć funkcję wskaźnikową
$f=\ind{\{(j_1,\ldots,j_r)\}}$. Ustalmy historię
$A=\{X_0=i_0,\ldots,X_n=i_n\}$ o dodatnim
prawdopodobieństwie. Z prawa trajektorii
\begin{align*}
&\Prob(A, X_{n+1}=j_1,\ldots,X_{n+r}=j_r)\\
&\qquad=\Prob(A)\,p_{i_nj_1}p_{j_1j_2}\cdots
p_{j_{r-1}j_r}.
\end{align*}
Dzielenie przez $\Prob(A)$ daje odpowiednią
warunkową równość. Jeżeli $\Prob(A)=0$, obie
strony tożsamości całkowej na $A$ mają wartość zero.
Zdarzenia $A$ tworzą przeliczalny podział $\Omega$,
a każde zdarzenie z $\mathcal F_n^X$ jest sumą
niektórych elementów tego podziału. Zsumowanie
równości po tych elementach pokazuje, że prawa
strona jest wersją wskazanej warunkowej wartości
oczekiwanej. Dla dowolnej ograniczonej $f$ mnożymy
równości wskaźnikowe przez $f(j_1,\ldots,j_r)$
i sumujemy po wszystkich ciągach stanów.
Zamiana sum z wartościami oczekiwanymi jest
uprawniona, gdyż suma modułów jest ograniczona
przez $\|f\|_\infty$ po pomnożeniu przez
odpowiednie prawdopodobieństwa trajektorii.
\end{proof}

\begin{uwaga}[Dwa różne znaczenia słowa „zapominanie”]
Powyższa własność jest \emph{dokładna i natychmiastowa}:
przeszłość przestaje mieć znaczenie dla dalszej prognozy,
gdy stan obecny jest znany. Nie przesądza ona o tym,
czy $\Prob(X_n=j\mid X_0=i)$ zbliża się, gdy
$n\to\infty$, do liczby niezależnej od $i$.
Takie \emph{asymptotyczne zapominanie początku}
wynika w skończonym przypadku z nieprzywiedlności
i nieokresowości. Dla łańcucha naprzemiennego
$0\leftrightarrow1$ własność Markowa zachodzi,
a pamięć o stanie początkowym pozostaje na zawsze.
\end{uwaga}

Dla ograniczonej $f:S\to\R$ oznaczmy
$(Pf)(i)=\sum_jp_{ij}f(j)$ oraz
$(P^nf)(i)=\sum_jp_{ij}^{(n)}f(j)$.
Wtedy
$\E[f(X_{m+n})\mid\sigma(X_0,\ldots,X_m)]
=(P^nf)(X_m)$. Ta postać jest szczególnie użyteczna:
operator $P$ \emph{propaguje obserwable w przyszłość},
natomiast rozkład wierszowy $\mu$ propaguje się jako
$\mu P^n$.

\begin{przyklad}[Dwa stany: rozwiązanie jawne]
Dla $0<\alpha,\beta\leq1$ weźmy
$P=\begin{psmallmatrix}1-\alpha&\alpha\\
\beta&1-\beta\end{psmallmatrix}$.
Jeśli $q_n=\Prob(X_n=1)$, to
$q_{n+1}=\alpha(1-q_n)+(1-\beta)q_n
=\alpha+(1-\alpha-\beta)q_n$.
Po odjęciu $q_*=\alpha/(\alpha+\beta)$ otrzymujemy
$q_{n+1}-q_*=(1-\alpha-\beta)(q_n-q_*)$,
więc
\[
q_n=\frac{\alpha}{\alpha+\beta}
+(1-\alpha-\beta)^n
\left(q_0-\frac{\alpha}{\alpha+\beta}\right).
\]
Przy $\alpha=\beta=1$ łańcuch zmienia stan
w każdym kroku i nigdy nie zbiega w rozkładzie
z ustalonego stanu, mimo istnienia rozkładu
stacjonarnego $(1/2,1/2)$. W pozostałych
przypadkach $\alpha,\beta>0$ i
$\alpha+\beta<2$ mamy $|1-\alpha-\beta|<1$,
więc rozkład dąży do stacjonarnego.
\end{przyklad}


\begin{przyklad}[Ilościowe zapominanie początku a własność Markowa]
Dla łańcucha pogody
$P=\begin{psmallmatrix}4/5&1/5\\2/5&3/5\end{psmallmatrix}$
rozkład stacjonarny jest równy
$\pi=(2/3,1/3)$. Wzór dla łańcucha dwustanowego
daje
\begin{align*}
\Prob(X_n=d\mid X_0=s)&=\frac13\bigl(1-(2/5)^n\bigr),\\
\Prob(X_n=d\mid X_0=d)&=\frac13+\frac23(2/5)^n.
\end{align*}
Różnica obu prognoz wynosi $(2/5)^n$,
a zatem zależność od stanu początkowego zanika
wykładniczo. Jest to \emph{inny fakt} niż to,
że do prognozowania $X_{n+1}$ wystarcza $X_n$.
Pierwszy mówi o coraz słabszym wpływie dawnej historii;
drugi — o warunkowej niezależności w jednym kroku.
\end{przyklad}

\section{Konstrukcja i łączenie z niezależnymi losowaniami}

Zasadnicze jest pytanie, czy każda macierz $P$ i rozkład
początkowy rzeczywiście wyznaczają łańcuch. Można to
zrobić bez polegania na nieformalnych warunkowych
losowaniach.

\begin{twierdzenie}[Konstrukcja z jednostajnych zmiennych losowych]
Dla dowolnego przeliczalnego $S$, macierzy stochastycznej
$P$ na $S$ oraz rozkładu $\mu$ na $S$ istnieje
jednorodny łańcuch Markowa $(X_n)_{n\geq0}$
o rozkładzie początkowym $\mu$ i macierzy $P$.
\end{twierdzenie}
\begin{proof}
Wybierzmy enumerację $S=\{s_1,s_2,\ldots\}$
(w przypadku skończonym kończymy listę).
Dla $i\in S$ i $u\in(0,1)$ niech $F(i,u)=s_r$,
gdzie $r$ jest najmniejszym indeksem takim, że
$\sum_{\ell\leq r}p_{i,s_\ell}\geq u$.
Definicja ma sens poza ewentualnym zbiorem punktów
brzegowych miary Lebesgue'a zero, w których można
przypisać dowolny stan; przy prawdopodobieństwie
sumującym się do jedności dla każdego $u<1$
odpowiedni indeks istnieje. Funkcja $F$ jest
mierzalna względem $2^S\otimes\mathcal B((0,1))$.
Na przestrzeni iloczynowej weźmy niezależne
$U_1,U_2,\ldots\sim\Unif(0,1)$, niezależne od
$X_0$ o rozkładzie $\mu$, i rekurencyjnie połóżmy
$X_{n+1}=F(X_n,U_{n+1})$. Ponieważ
$U_{n+1}$ jest niezależna od $X_0,U_1,\ldots,U_n$,
a w szczególności od $\sigma(X_0,\ldots,X_n)$,
warunkowo przy $X_n=i$ otrzymujemy
$\Prob(F(i,U_{n+1})=j)=p_{ij}$.
Zatem $\E[\ind{\{X_{n+1}=j\}}\mid\sigma(X_0,\ldots,X_n)]
=p_{X_nj}$, jak należało wykazać.
\end{proof}

\section{Czasy wejścia, mocna własność Markowa}

\begin{definicja}[Czas wejścia i powrotu]
Dla łańcucha $(X_n)$ i zbioru $A\subseteq S$ definiujemy
$\tau_A=\inf\{n\geq0:X_n\in A\}$ oraz
$\tau_A^+=\inf\{n\geq1:X_n\in A\}$,
przyjmując $\inf\varnothing=\infty$.
Dla jednego stanu piszemy $\tau_i$ i $\tau_i^+$.
\end{definicja}

\begin{twierdzenie}[Mocna własność Markowa, wersja dyskretna]
Niech $(X_n)$ będzie jednorodnym łańcuchem Markowa
na przeliczalnym $S$, a $\tau$ czasem zatrzymania
względem jego naturalnej filtracji
$\mathcal F_n^X=\sigma(X_0,\ldots,X_n)$.
Dla każdego $r\geq0$ i ograniczonej funkcji
$f:S^{r+1}\to\R$ zachodzi
\[
\E\bigl[f(X_\tau,\ldots,X_{\tau+r})
\ind{\{\tau<\infty\}}\mid\mathcal F_\tau^X\bigr]
=\ind{\{\tau<\infty\}}\,G_f(X_\tau),
\]
gdzie $G_f(i)=\E_i f(X_0,\ldots,X_r)$, a
$\E_i$ oznacza wartość oczekiwaną dla łańcucha
rozpoczynającego w $i$. Wartości procesu na
$\{\tau=\infty\}$ w formule rozumiemy jako zero.
\end{twierdzenie}
\begin{proof}
Definicję $\mathcal F_\tau^X$ podamy dokładnie
w następnym rozdziale; tu wystarczy jej własność:
jeśli $A\in\mathcal F_\tau^X$, to
$A\cap\{\tau=n\}\in\mathcal F_n^X$.
Dla dowolnego takiego $A$ piszemy, korzystając
z rozbicia $\{\tau<\infty\}=\bigcup_{n\geq0}\{\tau=n\}$,
\begin{align*}
&\E\bigl[\ind A\ind{\{\tau<\infty\}}
 f(X_\tau,\ldots,X_{\tau+r})\bigr]\\
&=\sum_{n\geq0}\E\bigl[
\ind{A\cap\{\tau=n\}}f(X_n,\ldots,X_{n+r})\bigr]\\
&=\sum_{n\geq0}\E\bigl[
\ind{A\cap\{\tau=n\}}G_f(X_n)\bigr]
=\E\bigl[\ind A\ind{\{\tau<\infty\}}G_f(X_\tau)\bigr].
\end{align*}
Druga równość wynika ze zwykłej własności Markowa
oraz warunkowania względem $\mathcal F_n^X$.
Szeregi są absolutnie zbieżne, bo $f$ jest ograniczona
i zdarzenia $\{\tau=n\}$ są rozłączne.
Prawa strona jest $\mathcal F_\tau^X$-mierzalna,
a więc definiuje wskazaną warunkową wartość
oczekiwaną. Nie ma tu żadnego warunkowania
na zdarzenie $\{\tau=n\}$ o dodatniej miarze:
dowód działa również dla zerowych prawdopodobieństw.
\end{proof}

\begin{przyklad}[Po pierwszym deszczu]
W modelu pogody zatrzymajmy obserwację w pierwszym
dniu deszczowym $\tau_d$. Na zdarzeniu
$\{\tau_d<\infty\}$ prawdopodobieństwo, że
następnego dnia będzie sucho, wynosi $p_{ds}$,
niezależnie od liczby dni oczekiwania na deszcz.
To właśnie treść mocnej własności Markowa:
czas $\tau_d$ jest losowy, ale przyszłość po
zatrzymaniu znów zaczyna się od obecnego stanu.
\end{przyklad}

\section{Klasy komunikujące, powracanie i okresowość}

\begin{definicja}[Komunikowanie, nieprzywiedlność, okres]
Dla $i,j\in S$ piszemy $i\to j$, jeśli istnieje
$n\geq0$ z $p_{ij}^{(n)}>0$. Stany $i$ i $j$
\emph{komunikują się}, jeśli $i\to j$ i $j\to i$.
Łańcuch jest \emph{nieprzywiedlny}, gdy wszystkie
stany komunikują się. Klasa $C\subseteq S$
jest \emph{domknięta}, jeśli $p_{ij}=0$ dla
$i\in C$, $j\notin C$. Okresem stanu $i$ nazywamy
$d(i)=\gcd\{n\geq1:p_{ii}^{(n)}>0\}$;
przez $\gcd$ rozumiemy największy wspólny dzielnik
niepustego zbioru, zaś jeśli zbiór jest pusty,
stan nie ma dodatnich czasów powrotu i okres
pozostawiamy nieokreślony.
\end{definicja}

\begin{lemat}[Okres jest wspólny dla klasy komunikującej]
Jeżeli stany $i,j$ komunikują się i istnieje
$n\geq1$ z $p_{ii}^{(n)}>0$, to $d(i)=d(j)$.
\end{lemat}
\begin{proof}
Wybierzmy $r,s\geq0$ takie, że
$p_{ij}^{(r)}>0$ i $p_{ji}^{(s)}>0$.
W szczególności istnieje cykl z $i$ do $j$ i z
powrotem długości $r+s$, a więc $d(i)$ dzieli $r+s$.
Dla dowolnego $n$ z $p_{jj}^{(n)}>0$ powstaje
cykl w $i$ długości $r+n+s$, zatem $d(i)$
dzieli $r+n+s$ i $r+s$, więc dzieli $n$.
Tym samym $d(i)$ dzieli $d(j)$.
Zamieniając $i$ z $j$, otrzymujemy podzielność
odwrotną, a więc równość.
\end{proof}

\begin{definicja}[Stan powracający i przemijający]
Stan $i\in S$ jest \emph{powracający}, jeśli
$\Prob_i(\tau_i^+<\infty)=1$, a
\emph{przemijający}, jeśli
$\Prob_i(\tau_i^+<\infty)<1$.
Stan powracający jest \emph{dodatnio powracający},
jeżeli $\E_i\tau_i^+<\infty$, a \emph{zerowo
powracający}, jeśli $\E_i\tau_i^+=\infty$.
\end{definicja}

\begin{twierdzenie}[Kryterium powracania przez szereg odwiedzin]
Dla każdego stanu $i$ przeliczalnego jednorodnego
łańcucha Markowa zachodzi równoważność
\[
 i\text{ jest powracający}\quad\Longleftrightarrow
\quad\sum_{n=0}^{\infty}p_{ii}^{(n)}=\infty.
\]
Jeżeli $r_i=\Prob_i(\tau_i^+<\infty)<1$,
to $\sum_{n\geq0}p_{ii}^{(n)}=1/(1-r_i)$.
\end{twierdzenie}
\begin{proof}
Rozważmy liczbę wszystkich wizyt w $i$,
$N_i=\sum_{n\geq0}\ind{\{X_n=i\}}$, dla łańcucha
rozpoczętego w $i$. Niech $T_k$ oznacza chwilę
$k$-tego powrotu po czasie zero; $T_0=0$.
Mocna własność Markowa zastosowana w $T_k$
daje indukcyjnie
$\Prob_i(T_k<\infty)=r_i^k$.
Z monotonicznej zbieżności
\[
\sum_{n\geq0}p_{ii}^{(n)}=\E_iN_i
=\sum_{k\geq0}\Prob_i(T_k<\infty)
=\sum_{k\geq0}r_i^k.
\]
Ostatni szereg jest skończony dokładnie dla $r_i<1$;
przy $r_i=1$ jest rozbieżny. Dowód obejmuje oba
przypadki bez zakładania, że czas powrotu ma
skończoną wartość oczekiwaną.
\end{proof}

\begin{przyklad}[Spacer symetryczny na prostej całkowitej]
Niech $X_0=0$ i $X_n=\xi_1+\cdots+\xi_n$,
gdzie $\xi_k$ są niezależne i przyjmują $\pm1$
z prawdopodobieństwem $1/2$. Dla $n\geq0$
$\Prob_0(X_{2n+1}=0)=0$ i
$\Prob_0(X_{2n}=0)=\binom{2n}{n}2^{-2n}$.
Ze wzoru Stirlinga
$\binom{2n}{n}2^{-2n}\sim1/\sqrt{\pi n}$.
Szereg prawdopodobieństw powrotu jest rozbieżny,
więc spacer jest powracający. Ponadto okres wynosi $2$.
Dla spaceru z dryfem $\Prob(\xi_k=1)=p\ne1/2$
mamy
$\Prob_0(X_{2n}=0)=\binom{2n}{n}[p(1-p)]^n
\sim[4p(1-p)]^n/\sqrt{\pi n}$,
a $4p(1-p)<1$; stan zero jest zatem przemijający.
\end{przyklad}

\begin{przyklad}[Dlaczego przy powracaniu liczy się cały szereg]
Łańcuch może w jednym kroku powracać bardzo mało
prawdopodobnie, ale i tak powracać niemal na pewno,
ponieważ ma nieskończenie wiele okazji. Dlatego
samo $p_{ii}$ nie rozstrzyga o powracaniu.
Na przykład spacer symetryczny ma $p_{00}=0$,
a mimo to stan zero jest powracający.
\end{przyklad}


\begin{stwierdzenie}[Domknięcie i komunikowanie]
Dla jednorodnego łańcucha na przeliczalnym zbiorze $S$
relacja komunikowania jest relacją równoważności.
Jeśli klasa komunikująca $C$ jest domknięta i
$X_0\in C$ p.n., to $X_n\in C$ dla wszystkich
$n\geq0$ prawie na pewno.
\end{stwierdzenie}
\begin{proof}
Zwrotność wynika z $p_{ii}^{(0)}=1$,
a symetryczność jest częścią definicji.
Jeśli $i\to j$ w $m$ krokach i $j\to k$ w $n$,
to z równania Chapmana--Kołmogorowa
$p_{ik}^{(m+n)}\geq p_{ij}^{(m)}p_{jk}^{(n)}>0$.
Tak samo, łącząc drogę powrotną $k\to j\to i$,
dostajemy $k\to i$. To dowodzi przechodniości.
Jeśli $C$ jest domknięta, to dla każdego
$i\in C$ mamy $\Prob_i(X_1\in C)=1$.
Indukcja z warunkowaniem względem $X_n$
daje $\Prob(X_n\in C)=1$ przy $X_0\in C$ p.n.
Przeliczalne przecięcie zdarzeń pełnej miary
zapewnia pozostawanie w $C$ jednocześnie
we wszystkich chwilach. 
\end{proof}

\begin{przyklad}[Dwie klasy pochłaniające i stan przejściowy]
Niech $S=\{0,1,2\}$ i
$P=\begin{psmallmatrix}1&0&0\\ 1/3&0&2/3\\ 0&0&1\end{psmallmatrix}$.
Stany $0$ i $2$ są oddzielnymi domkniętymi klasami,
a stan $1$ jest przemijający. Rozkłady
$\delta_0$ i $\delta_2$ są stacjonarne;
stacjonarny jest też każdy ich wypukły
mieszany rozkład $t\delta_0+(1-t)\delta_2$.
Zatem bez nieprzywiedlności nie musi być
jednoznacznego rozkładu stacjonarnego ani
wspólnej granicy niezależnej od startu.
\end{przyklad}

\begin{uwaga}[Kiedy dowód prawa Markowa przechodzi na proces obserwowany?]
Własność Markowa zależy od wyboru \emph{stanu}.
Jeśli stan rzeczywisty ma dwie składowe $(X_n,H)$,
ale obserwujemy tylko $X_n$, to proces obserwacji
nie musi mieć własności Markowa. Rozkład warunkowy
przyszłości może zależeć od wcześniejszych obserwacji,
gdyż pomagają one odtworzyć ukrytą składową $H$.
W dowodach własności Markowa należy więc wskazać
niezależne innowacje lub zweryfikować tożsamość
warunkową, a nie wnioskować z samej znajomości
jednokrokowych rozkładów bezwarunkowych.
\end{uwaga}

\begin{przyklad}[Rozwiązane zadanie: pochłanianie a warunkowanie wstecz]
W modelu pogody z macierzą
$P=\begin{psmallmatrix}4/5&1/5\\2/5&3/5\end{psmallmatrix}$
przyjmijmy $\Prob(X_0=s)=3/4$ i
$\Prob(X_0=d)=1/4$. Wyznaczymy prawdopodobieństwo
obserwacji $X_0=s,X_1=d,X_2=d$ oraz szansę,
że pierwszy dzień był suchy, jeśli widzimy
deszcz w dniach $1$ i $2$.

Z prawa trajektorii
$\Prob(X_0=s,X_1=d,X_2=d)
=(3/4)(1/5)(3/5)=9/100$.
Prawdopodobieństwo warunkujące wynosi
\[
\Prob(X_1=d,X_2=d)
=\left(\frac34\frac15+\frac14\frac35\right)\frac35
=\frac9{50}.
\]
Stąd
$\Prob(X_0=s\mid X_1=d,X_2=d)
=(9/100)/(9/50)=1/2$.
Gdybyśmy znali jedynie $X_1=d$, wynik także
wynosiłby
$[(3/4)(1/5)]/[(3/4)(1/5)+(1/4)(3/5)]=1/2$.
To konkretne sprawdzenie warunkowej niezależności
przeszłości i dalszej przyszłości po poznaniu
stanu pośredniego.
\end{przyklad}

\section{Rozkłady stacjonarne i równowaga}

\begin{definicja}[Rozkład stacjonarny i odwracalność]
Rozkład probabilistyczny $\pi=(\pi_i)_{i\in S}$
jest \emph{stacjonarny} dla $P$, jeśli
$\pi_j=\sum_{i\in S}\pi_i p_{ij}$ dla każdego $j$,
czyli $\pi P=\pi$.
Mówimy, że $\pi$ spełnia \emph{warunek równowagi
szczegółowej}, jeśli
$\pi_i p_{ij}=\pi_jp_{ji}$ dla wszystkich $i,j\in S$.
\end{definicja}

\begin{stwierdzenie}[Równowaga szczegółowa implikuje stacjonarność]
Jeżeli $\pi$ jest rozkładem probabilistycznym
spełniającym warunek równowagi szczegółowej,
to jest stacjonarny.
\end{stwierdzenie}
\begin{proof}
Dla każdego $j$ sumujemy po $i$:
$\sum_i\pi_i p_{ij}=\sum_i\pi_jp_{ji}
=\pi_j\sum_i p_{ji}=\pi_j$.
W ostatnim wyrażeniu sumujemy \emph{wiersz}
macierzy $P$ o indeksie początkowym $j$.
\end{proof}

\begin{przyklad}[Spacer na grafie: przewaga wierzchołków o dużym stopniu]
Dla skończonego nieskierowanego grafu bez
wierzchołków izolowanych zdefiniujmy
$\pi_v=d(v)/\sum_{w\in V}d(w)$.
Gdy $v,w$ są sąsiadami,
$\pi_vp_{vw}=1/\sum_zd(z)=\pi_wp_{wv}$;
gdy nie są, obie strony znikają. Zatem $\pi$
jest stacjonarny. Wierzchołek z trzykrotnie
większą liczbą sąsiadów ma w równowadze
trzykrotnie większe prawdopodobieństwo wizyty.
\end{przyklad}

\begin{przyklad}[Stacjonarność nie jest zbieżnością]
Dla $P=\begin{psmallmatrix}0&1\\1&0\end{psmallmatrix}$
rozkład $(1/2,1/2)$ jest stacjonarny, ale startując
od pierwszego stanu, rozkłady kolejnych $X_n$
naprzemiennie wynoszą $(1,0)$ i $(0,1)$.
Sama równość $\pi P=\pi$ nie wystarcza do
$\mu P^n\to\pi$; trzeba kontrolować okresowość.
\end{przyklad}

\begin{twierdzenie}[Istnienie rozkładu stacjonarnego na skończonym zbiorze]
Każda macierz stochastyczna $P$ na niepustym
skończonym zbiorze $S$ ma co najmniej jeden
rozkład stacjonarny.
\end{twierdzenie}
\begin{proof}
Wybierzmy dowolny rozkład $\mu$ i zdefiniujmy
$\nu_N=N^{-1}\sum_{n=0}^{N-1}\mu P^n$.
Są to rozkłady probabilistyczne w skończenie
wymiarowym sympleksie
$\{x_i\geq0,\ \sum_i x_i=1\}$, który jest zwarty.
Istnieje więc podciąg $N_k$ i rozkład $\pi$ taki,
że $\nu_{N_k}\to\pi$ współrzędnie. Ponieważ
\[
\nu_NP-\nu_N=\frac{\mu P^N-\mu}{N},
\]
prawa strona zbiega do zera współrzędnie
(jej współrzędne są ograniczone modułem przez $1/N$).
Przechodząc do granicy i korzystając z ciągłości
mnożenia macierzy w skończonym wymiarze,
otrzymujemy $\pi P=\pi$.
\end{proof}

\begin{twierdzenie}[Jedyność rozkładu stacjonarnego dla skończonego łańcucha]
Jeśli $S$ jest skończony, a $P$ nieprzywiedlna,
to rozkład stacjonarny $\pi$ istnieje, jest jedyny
i spełnia $\pi_i>0$ dla każdego $i\in S$.
\end{twierdzenie}
\begin{proof}
Istnienie wynika z poprzedniego twierdzenia.
Jeżeli $\pi_i>0$ i $i\to j$, to dla pewnego $n$
$\pi_j=(\pi P^n)_j\geq\pi_i p_{ij}^{(n)}>0$.
Ponieważ pewne $\pi_i>0$ i łańcuch jest
nieprzywiedlny, wszystkie współrzędne są dodatnie.
Niech $\rho$ będzie innym rozkładem stacjonarnym;
z tego samego argumentu $\rho_i>0$.
Niech $c=\max_{i\in S}\pi_i/\rho_i$ i
$D_i=c\rho_i-\pi_i\geq0$. Co najmniej jeden $D_i=0$.
Ponieważ obie miary są stacjonarne,
$D_j=\sum_iD_i p_{ij}$ dla wszystkich $j$.
Jeśli $D_j=0$, to każdy $i$ z $p_{ij}>0$ ma $D_i=0$.
Przez iterację dostajemy $D_i=0$ dla wszystkich $i$
mających drogę do $j$, a przez nieprzywiedlność
dla wszystkich $i$. Zatem $\pi=c\rho$.
Sumowanie współrzędnych daje $c=1$, więc
$\pi=\rho$.
\end{proof}

\begin{lemat}[Wspólna dodatnia potęga w przypadku nieokresowym]
Niech $S$ będzie skończony, $P$ nieprzywiedlna
i niech okres stanów wynosi $1$. Wtedy istnieje
$N\geq1$, dla którego wszystkie współrzędne
macierzy $P^N$ są dodatnie.
\end{lemat}
\begin{proof}
Ustalmy stan $o$. Z definicji okresu liczby
$n\geq1$ z $p_{oo}^{(n)}>0$ mają największy
wspólny dzielnik równy $1$. Istnieje zatem
\emph{skończony} podzbiór długości
$n_1,\ldots,n_r$ o tym samym NWD: kolejne NWD
liczb naturalnych tworzą ciąg nierosnący,
który po skończonej liczbie zmian stabilizuje się.
Każda dostatecznie duża liczba naturalna jest
nieujemną kombinacją całkowitą $n_1,\ldots,n_r$.
Aby to uzasadnić, oznaczmy $a=n_1$.
Reszty modulo $a$ osiągane przez nieujemne sumy
$n_1,\ldots,n_r$ tworzą podmonoid skończonej
grupy $\Z/a\Z$, a więc są podgrupą;
NWD równy $1$ oznacza, że jest to cała grupa.
Wybierzmy dla każdej reszty jedną nieujemną sumę
$h_r$ o tej reszcie. Każda liczba o danej reszcie
większa lub równa $h_r$ jest postaci $h_r+ka$,
$k\geq0$, więc jest realizowana jako długość
ścieżki z $o$ do $o$.

Dla każdego $i,j\in S$ wybierzmy dodatnie
ścieżki $i\to o$ i $o\to j$ o długościach
$r_i,s_j$; można tu dopuścić długość zero.
Wobec poprzedniego ustępu, gdy $N-r_i-s_j$
jest dostatecznie duże, istnieje ścieżka
$i\to o\to o\to j$ długości dokładnie $N$.
Ponieważ par $(i,j)$ jest skończenie wiele,
wybierzmy wspólne duże $N$. Wtedy
$p_{ij}^{(N)}>0$ dla wszystkich par.
\end{proof}

\begin{twierdzenie}[Zbieżność skończonego łańcucha ergodycznego]
Niech $S$ będzie skończony, $P$ nieprzywiedlna
i nieokresowa. Wtedy istnieje jedyny rozkład
stacjonarny $\pi$ i dla każdego początkowego
rozkładu $\mu$ zachodzi
\[
\sum_{j\in S}|(\mu P^n)_j-\pi_j|\longrightarrow0.
\]
Co więcej, istnieją stałe $C<\infty$ i $r\in(0,1)$
(zależne od $P$), dla których lewa strona jest
nie większa niż $Cr^n$.
\end{twierdzenie}
\begin{proof}
Z poprzedniego lematu wybierzmy $N$ takie,
że $Q=P^N$ ma wszystkie współrzędne dodatnie.
Niech $m=|S|$ i $\varepsilon=m\min_{i,j}Q_{ij}>0$.
Gdy $\varepsilon=1$, wszystkie wiersze $Q$ są
jednostajne i po $N$ krokach następuje pełne
zapomnienie stanu początkowego.
W przeciwnym razie $0<\varepsilon<1$ i
\[
Q_{ij}=\varepsilon\,u_j+(1-\varepsilon)R_{ij},
\qquad u_j=1/m,
\]
gdzie $R$ jest macierzą stochastyczną (nieujemność
wynika z definicji $\varepsilon$).
Dla dowolnych rozkładów $\mu,\nu$
$\mu Q-\nu Q=(1-\varepsilon)(\mu R-\nu R)$.
Ponieważ przejście stochastyczne nie zwiększa
normy sumy modułów różnic,
\[
\|\mu Q-\nu Q\|_1
\leq(1-\varepsilon)\|\mu-\nu\|_1.
\]
Ustalmy stacjonarny $\pi$ i podstawmy $\nu=\pi$.
Iteracja daje
$\|\mu Q^k-\pi\|_1\leq2(1-\varepsilon)^k$.
Dla $n=kN+r$, $0\leq r<N$, kontrakcja przez
$P^r$ daje ten sam szacunek z $k=\lfloor n/N\rfloor$.
Przy $0<\varepsilon<1$ jest on geometryczny:
można wziąć $r_0=(1-\varepsilon)^{1/N}<1$
i stałą $C=2/(1-\varepsilon)$.
Przy $\varepsilon=1$ rozkład jest dokładnie
stacjonarny od pewnego kroku, więc dowolna stała
$r_0\in(0,1)$ i wystarczająco duże $C$ wystarczą.
\end{proof}

\begin{uwaga}[Odległość całkowitej wariacji]
Na zbiorze przeliczalnym definiujemy
$\|\mu-\nu\|_{\mathrm{TV}}=\sup_{A\subseteq S}
|\mu(A)-\nu(A)|=\frac12\sum_i|\mu_i-\nu_i|$.
Użyta w dowodzie norma $\ell^1$ jest dwa razy
większa od odległości całkowitej wariacji.
W przypadku przestrzeni nieskończonych rozkład
stacjonarny może nie istnieć, a okresowość nie
jest jedyną przeszkodą zbieżności.
\end{uwaga}


\begin{przyklad}[Stacjonarność bez niezależności obserwacji]
Niech $S=\{0,1\}$, a macierz przejścia wynosi
$P=\begin{psmallmatrix}3/4&1/4\\1/4&3/4\end{psmallmatrix}$.
Rozpocznijmy proces od rozkładu stacjonarnego
$\pi=(1/2,1/2)$. Wtedy dla każdego $n$
$\Prob(X_n=1)=1/2$, lecz
$\Prob(X_0=1,X_1=1)=\pi_1p_{11}=3/8$,
podczas gdy iloczyn mas brzegowych wynosi $1/4$.
Zatem ciąg jest stacjonarny, ale kolejne stany
nie są niezależne. Rozkład każdej pojedynczej obserwacji
nie wyznacza zależności czasowej procesu.
\end{przyklad}

\section{Problemy wejścia jako równania harmoniczne}

\begin{twierdzenie}[Prawdopodobieństwo trafienia]
Niech $A,B\subseteq S$ będą rozłączne, a
$\tau_A,\tau_B$ czasami wejścia.
Dla $i\in S$ połóżmy
$h(i)=\Prob_i(\tau_A<\tau_B,\ \tau_A<\infty)$.
Wtedy $0\leq h\leq1$, $h=1$ na $A$,
$h=0$ na $B$ oraz
\[
h(i)=\sum_{j\in S}p_{ij}h(j)
\quad(i\notin A\cup B).
\]
Co więcej, $h$ jest \emph{najmniejszą}
nieujemną funkcją spełniającą te warunki brzegowe
i równanie na $S\setminus(A\cup B)$.
Jeśli $\Prob_i(\tau_A\wedge\tau_B<\infty)=1$
dla wszystkich $i$, to $h$ jest jedynym
ograniczonym rozwiązaniem układu z tymi
warunkami brzegowymi.
\end{twierdzenie}
\begin{proof}
Jeśli $i\in A$, mamy $\tau_A=0<\tau_B$,
a jeśli $i\in B$, to $\tau_B=0$,
więc wartości brzegowe wynoszą odpowiednio $1$ i $0$.
Dla $i\notin A\cup B$ rozdzielamy zdarzenie trafienia
według pierwszego kroku. Z prawa całkowitego
prawdopodobieństwa i własności Markowa
\[
 h(i)=\sum_{j\in S}\Prob_i(X_1=j)
 \Prob_i(\tau_A<\tau_B,\tau_A<\infty\mid X_1=j)
 =\sum_{j\in S}p_{ij}h(j),
\]
przy czym składniki odpowiadające $p_{ij}=0$
można pominąć. Istotne jest przesunięcie czasu:
po pierwszym kroku problem zaczyna się od nowa
w stanie $j$.

Niech $g:S\to[0,\infty)$ będzie innym rozwiązaniem
z tymi samymi wartościami brzegowymi, niekoniecznie
ograniczonym. Dla $\tau=\tau_A\wedge\tau_B$
pokażemy indukcyjnie
$g(i)=\E_i[g(X_{n\wedge\tau})]$.
Dla $n=0$ jest to oczywiste. Zdarzenie
$\{\tau\leq n\}$ oznacza, że proces znajduje się
już na brzegu, więc wartość zatrzymanego procesu
nie zmienia się. Na $\{\tau>n\}$, warunkując
względem $\mathcal F_n^X$, stosujemy
$\E[g(X_{n+1})\mid\mathcal F_n^X]=(Pg)(X_n)
=g(X_n)$. Dla nieujemnego $g$ równość
warunkowych całek jest dozwolona także przed
ustaleniem całkowalności; z indukcji wynika,
że każda rozważana wartość oczekiwana jest
równa skończonemu $g(i)$.

Na $\{\tau_A\leq n,\tau_A<\tau_B\}$ mamy
$g(X_{n\wedge\tau})=1$, a poza tym z
nieujemności otrzymujemy
$g(i)\geq\Prob_i(\tau_A\leq n,\tau_A<\tau_B)$.
Zdarzenia po prawej rosną do
$\{\tau_A<\tau_B,\tau_A<\infty\}$;
ciągłość miary od dołu daje $g(i)\geq h(i)$.

Jeśli $g$ jest ograniczona i $\tau<\infty$
prawie na pewno przy każdym starcie, to
$X_{n\wedge\tau}=X_\tau$ od pewnego losowego
momentu. Zbieżność zdominowana prowadzi do
\[
 g(i)=\lim_n\E_i g(X_{n\wedge\tau})
 =\E_i g(X_\tau)
 =\Prob_i(\tau_A<\tau_B)=h(i).
\]
W ostatniej równości używamy wartości $g=1$
na $A$ i $g=0$ na $B$. To dowodzi jednoznaczności
w klasie funkcji ograniczonych. 
\end{proof}

\begin{przyklad}[Ruina gracza: model skończony]
Niech $S=\{0,\ldots,N\}$, gdzie $N\geq2$,
$0$ i $N$ są pochłaniające,
a dla $1\leq i\leq N-1$
przejścia $i\to i+1$ i $i\to i-1$ zachodzą
z prawdopodobieństwami odpowiednio $p\in(0,1)$
i $q=1-p$. Niech $h_i$ oznacza prawdopodobieństwo
dotarcia do $N$ przed zerem ze stanu $i$.
Równanie harmoniczne przyjmuje postać
$h_i=ph_{i+1}+qh_{i-1}$, z $h_0=0$, $h_N=1$.
Jeśli $p\ne q$, szukamy rozwiązań $h_i=A+B(q/p)^i$,
co po warunkach brzegowych daje
\[
h_i=\frac{1-(q/p)^i}{1-(q/p)^N}.
\]
Jeśli $p=q=1/2$, równanie oznacza stałe
pierwsze różnice, więc $h_i=i/N$.

Również oczekiwany czas pochłonięcia
$u_i=\E_i(\tau_{\{0,N\}})$ można wyprowadzić
bez zgadywania. Ponieważ z każdego stanu wewnętrznego
istnieje droga do brzegu długości co najwyżej $N-1$
o prawdopodobieństwie ograniczonym od dołu przez
dodatnią liczbę zależną od $p,q,N$, czas
pochłonięcia ma ogon geometryczny i $u_i<\infty$.
Pierwszy krok daje $u_i=1+pu_{i+1}+qu_{i-1}$,
$u_0=u_N=0$. Przy $p=q$ rozwiązanie ma postać
$u_i=i(N-i)$; podstawienie w równaniu sprawdza
je wprost. Przy $p\ne q$, funkcja $i\mapsto i$
spełnia $p(i+1)+q(i-1)-i=p-q$,
a $h$ jest harmoniczna, stąd
\[
u_i=\frac{Nh_i-i}{p-q}.
\]
Jednoznaczność rozwiązań wynika z zasady maksimum
dla różnicy dwóch rozwiązań: dodatnie maksimum
w punkcie wewnętrznym propaguje do sąsiadów,
a więc aż do brzegu, gdzie funkcja jest zerowa.
\end{przyklad}

\begin{przyklad}[Łańcuch urodzeń i śmierci na skończonym odcinku]
Na $S=\{0,\ldots,N\}$ przyjmijmy
$p_{i,i+1}=b_i>0$ dla $i<N$,
$p_{i,i-1}=d_i>0$ dla $i>0$ oraz
$p_{ii}=1-b_i-d_i\geq0$, z pominięciem
nieistniejących przejść na brzegach.
Warunek równowagi szczegółowej daje
$\pi_i b_i=\pi_{i+1}d_{i+1}$.
Stąd rekurencyjnie
\[
\pi_i=C\prod_{k=0}^{i-1}\frac{b_k}{d_{k+1}},
\qquad
C^{-1}=1+\sum_{i=1}^N\prod_{k=0}^{i-1}
\frac{b_k}{d_{k+1}}.
\]
Pusty iloczyn przy $i=0$ jest równy $1$.
Otrzymujemy dokładnie rozkład stacjonarny.
To model np. ograniczonego bufora, w którym
przyjścia zwiększają zajętość, a obsługa ją zmniejsza.
\end{przyklad}

\section{Zależność i korelacja w łańcuchu stacjonarnym}

\begin{stwierdzenie}[Kowariancja przez operator przejścia]
Niech $S$ będzie skończony, $(X_n)$ łańcuchem
o rozkładzie początkowym stacjonarnym $\pi$, a
$f,g:S\to\R$. Wtedy dla $k\geq0$
\[
\Cov(f(X_0),g(X_k))=
\sum_{i\in S}\pi_i f(i)(P^kg)(i)
-\Bigl(\sum_i\pi_i f(i)\Bigr)
 \Bigl(\sum_j\pi_jg(j)\Bigr).
\]
\end{stwierdzenie}
\begin{proof}
Warunkowanie względem $X_0$ daje
$\E[g(X_k)\mid X_0]=(P^kg)(X_0)$.
Mnożymy przez $f(X_0)$ i stosujemy wieżę.
Następnie odejmujemy iloczyn średnich,
które przy stacjonarnym rozkładzie
nie zależą od czasu.
\end{proof}

\begin{przyklad}[Autokorelacja dwustanowego łańcucha]
Niech $S=\{0,1\}$, $\alpha,\beta>0$,
$\alpha+\beta<2$ oraz rozkład początkowy
$\pi=(\beta,\alpha)/(\alpha+\beta)$.
Dla $I_n=\ind{\{X_n=1\}}$ otrzymujemy
\[
\Cov(I_0,I_k)=\pi_1(1-\pi_1)(1-\alpha-\beta)^k.
\]
Wynika to z jawnego wzoru na prawdopodobieństwo
$X_k=1$ przy starcie z $X_0=1$.
Jeśli $\alpha+\beta>1$, korelacje zmieniają
znak naprzemiennie, a jeśli $\alpha+\beta<1$,
zanikają bez zmiany znaku. Proces może zatem
mieć pamięć statystyczną mimo ustalonego rozkładu
jednowymiarowego.
\end{przyklad}

\section{Częstości odwiedzin i twierdzenie ergodyczne}

Istnienie rozkładu stacjonarnego nie oznacza samo
przez się, że \emph{rozkład $X_n$} musi do niego
zbiegać. Pojawia się jednak inne pytanie:
jaki odsetek czasu pojedyncza trajektoria
spędza w danym stanie? Dla skończonego
nieprzywiedlnego łańcucha odpowiedź jest
znacznie stabilniejsza i nie wymaga nieokresowości.

\begin{lemat}[Jednostajny ogon czasu wejścia w skończonym łańcuchu]
Niech $S$ będzie skończony, a $P$ nieprzywiedlna.
Dla dowolnego $i\in S$ istnieją liczby całkowite
$L\geq1$ oraz $\delta>0$ takie, że dla każdego
stanu początkowego $j\in S$ i $k\geq0$
\[
\Prob_j(\tau_i>kL)\leq(1-\delta)^k.
\]
Ponadto $\E_j\tau_i<\infty$ i
$\E_i\tau_i^+<\infty$.
\end{lemat}
\begin{proof}
Ustalmy $i\in S$. Dla każdego $j\neq i$
nieprzywiedlność zapewnia ścieżkę
$j=j_0,j_1,\ldots,j_{\ell_j}=i$ o dodatnim
prawdopodobieństwie
$a_j=\prod_{r=1}^{\ell_j}p_{j_{r-1}j_r}>0$.
Dla $j=i$ wybierzmy analogicznie dodatnią
ścieżkę powrotu z $i$ do $i$, tym razem o
długości co najmniej $1$; istnienie wynika z
nieprzywiedlności i skończoności $S$
(w przypadku jednego stanu bierzemy pętlę długości $1$).
Ponieważ $S$ jest skończony, możemy ustalić
$L\geq\max_j\ell_j$ oraz
$\delta=\min_j a_j\in(0,1]$.
Z każdego stanu $j\neq i$ istnieje zatem
prawdopodobieństwo co najmniej $\delta$, że
$\tau_i\leq L$; ze stanu $i$ dotyczy to
$\tau_i^+$.

Niech $j\neq i$. Na zdarzeniu $\{\tau_i>kL\}$
w chwili $kL$ znajdujemy się w pewnym stanie
$X_{kL}\neq i$. Własność Markowa i oszacowanie
jednego bloku dają
\begin{align*}
\Prob_j(\tau_i>(k+1)L)
&=\E_j\Bigl[\ind{\{\tau_i>kL\}}
  \Prob_j(\tau_i>(k+1)L\mid\mathcal F_{kL}^X)\Bigr]\\
&\leq(1-\delta)\Prob_j(\tau_i>kL).
\end{align*}
Iteracja od $k=0$ daje zadany geometryczny ogon.
Dla $j=i$ czas $\tau_i=0$, więc oszacowanie
jest trywialne.

Dla dodatniego czasu powrotu pierwszy blok
rozpoczyna się w $i$ i wykorzystujemy wybraną
ścieżkę powrotu. Jeżeli po $k\geq1$ blokach
nie było powrotu, to $X_{kL}\neq i$; od tego
momentu stosujemy drogę prowadzącą do $i$.
W obu przypadkach prawdopodobieństwo sukcesu
w następnym bloku jest co najmniej $\delta$.
Ten sam argument indukcyjny daje zatem
$\Prob_i(\tau_i^+>kL)\leq(1-\delta)^k$.
Wreszcie dla nieujemnej całkowitoliczbowej zmiennej
losowej $T$ mamy $\E T=\sum_{n\geq0}\Prob(T>n)$.
Grupując wyrazy w bloki długości $L$, otrzymujemy
\[
\E_j\tau_i\leq L\sum_{k\geq0}(1-\delta)^k
=\frac L\delta,
\qquad
\E_i\tau_i^+\leq\frac L\delta.
\]
Otrzymaliśmy jednocześnie prawie pewną skończoność
obu czasów i ich całkowalność.
\end{proof}

\begin{twierdzenie}[Częstości stanów i wzór na średni czas powrotu]
Niech $(X_n)$ będzie jednorodnym nieprzywiedlnym
łańcuchem Markowa na skończonym $S$, a $\pi$
jego jedynym rozkładem stacjonarnym. Wtedy
dla każdego $i\in S$ i dowolnego rozkładu
początkowego
\[
\frac1n\sum_{k=0}^{n-1}\ind{\{X_k=i\}}
\longrightarrow\pi_i\quad\text{prawie na pewno}.
\]
Ponadto $\E_i\tau_i^+=1/\pi_i$.
\end{twierdzenie}
\begin{proof}
Zacznijmy od prawa procesu $\Prob_i$, tzn.
$X_0=i$ prawie na pewno. Zdefiniujmy $T_0=0$ i
rekurencyjnie $T_r=\inf\{n>T_{r-1}:X_n=i\}$.
Z poprzedniego lematu i mocnej własności Markowa
wszystkie $T_r$ są skończone p.n., a odstępy
$L_r=T_r-T_{r-1}$ są niezależne i mają wspólny
rozkład $\tau_i^+$ przy $\Prob_i$. Nie jest to
stwierdzenie o niezależności stanów łańcucha:
niezależne są \emph{całe odcinki pomiędzy kolejnymi
powrotami}, po rozpoczęciu każdego z nich w stanie $i$.
Niech $m_i=\E_i\tau_i^+\in(0,\infty)$. MPWL daje
\[
 \frac{T_r}{r}=\frac1r\sum_{s=1}^rL_s
 \longrightarrow m_i\quad\text{p.n.}
\]

W celu przejścia od długości cykli do liczby
odwiedzin wprowadźmy
$r(n)=\max\{r\geq0:T_r\leq n\}$.
Ponieważ wszystkie odstępy są skończone i
$T_r\geq r$, mamy $r(n)\to\infty$ przy $n\to\infty$.
Z definicji $T_{r(n)}\leq n<T_{r(n)+1}$, skąd
\[
 \frac{T_{r(n)}}{r(n)}
 \leq\frac n{r(n)}
 <\frac{T_{r(n)+1}}{r(n)+1}
   \frac{r(n)+1}{r(n)}.
\]
Dla $r(n)\geq1$ obie skrajne strony zbiegają
do $m_i$, a więc $r(n)/n\to1/m_i$.
Liczba $N_i(n)=\sum_{k=0}^{n-1}\ind{\{X_k=i\}}$
zlicza chwile $T_r<n$, podczas gdy $r(n)$
zlicza dodatnie powroty $T_r\leq n$.
Zatem $|N_i(n)-r(n)|\leq1$, a stąd
$N_i(n)/n\to1/m_i$ prawie na pewno.

Dla dowolnego stanu początkowego $j$ mamy
$\tau_i<\infty$ p.n. (poprzedni lemat).
Po tej chwili, dzięki mocnej własności Markowa,
proces ma prawo łańcucha startującego od $i$.
Przed $\tau_i$ nie ma żadnej wizyty w $i$;
pojedyncza skończona chwila wejścia nie wpływa
na stosunek liczby wizyt do $n$. Stąd granica
$1/m_i$ jest niezależna od stanu początkowego,
a zsumowanie po rozkładzie początkowym daje ją
dla dowolnego startu. W skończonym $S$ można
jednocześnie wybrać zbiór pełnej miary dla
wszystkich stanów $i$.

Aby rozpoznać tę granicę, rozpocznijmy łańcuch
z rozkładu stacjonarnego $\pi$. Wtedy
$\E_\pi N_i(n)/n=\pi_i$ dla każdego $n$,
bo $\Prob_\pi(X_k=i)=\pi_i$. Ponieważ
$0\leq N_i(n)/n\leq1$, zbieżność zdominowana
daje $\pi_i=1/m_i$, czyli
$\E_i\tau_i^+=1/\pi_i$. Zwróćmy uwagę,
że w tym rozumowaniu \emph{nie} użyliśmy
nieokresowości: okresowy łańcuch również może
mieć stabilne częstości odwiedzin.
\end{proof}

\begin{wniosek}[Ergodyczne uśrednianie obserwabli]
Przy założeniach poprzedniego twierdzenia,
dla każdej funkcji $f:S\to\R$ i dowolnego
rozkładu początkowego zachodzi
\[
\frac1n\sum_{k=0}^{n-1}f(X_k)
\longrightarrow\sum_{i\in S}\pi_i f(i)
\quad\text{prawie na pewno}.
\]
\end{wniosek}
\begin{proof}
Ponieważ $S$ jest skończony,
$\frac1n\sum_{k<n}f(X_k)=
\sum_{i\in S}f(i)\frac1n\sum_{k<n}\ind{\{X_k=i\}}$.
Dla wszystkich stanów jednocześnie zachodzi
prawie pewna zbieżność częstości, a skończona
suma pozwala przejść do granicy.
\end{proof}

\begin{przyklad}[Naprzemienne odwiedziny mimo stabilnych częstości]
Dla macierzy
$P=\begin{psmallmatrix}0&1\\1&0\end{psmallmatrix}$
i $X_0=0$ rozkład $X_n$ nie jest zbieżny:
znajduje się na przemian w stanie $0$ lub $1$.
Jednak na każdej trajektorii
$\frac1n\sum_{k<n}\ind{\{X_k=0\}}\to1/2$
i analogicznie dla stanu $1$. Nieokresowość
potrzebna jest dla zbieżności \emph{rozkładów
w kolejnych chwilach}, ale nie dla \emph{częstości
na jednej długiej trajektorii}.
\end{przyklad}

\section{Macierz fundamentalna łańcucha pochłaniającego}

\begin{twierdzenie}[Liczba wizyt i prawdopodobieństwa pochłonięcia]
Niech $S=T\sqcup A$ będzie skończoną przestrzenią
stanów, gdzie $A$ jest niepustym zbiorem stanów
pochłaniających, a dla każdego $i\in T$ istnieje
droga o dodatnim prawdopodobieństwie do $A$.
Po uporządkowaniu stanów zapiszmy
\[
P=\begin{pmatrix}Q&R\\0&I\end{pmatrix},
\]
gdzie $Q$ opisuje przejścia wewnątrz $T$,
a $R$ z $T$ do $A$. Wtedy $Q^n\to0$, macierz
$I-Q$ jest odwracalna i
\[
N:=(I-Q)^{-1}=\sum_{n=0}^{\infty}Q^n.
\]
Współrzędna $N_{ij}$ jest oczekiwaną liczbą
wizyt w stanie $j\in T$ (łącznie z chwilą zero)
przed pochłonięciem, przy starcie w $i\in T$.
Oczekiwane czasy pochłonięcia tworzą wektor
$N\mathbf 1$, a prawdopodobieństwa pochłonięcia
w poszczególnych stanach $a\in A$ tworzą
macierz $B=NR$.
\end{twierdzenie}
\begin{proof}
Ponieważ $T$ jest skończony i z każdego stanu
istnieje dodatnia droga do $A$, znajdziemy
$L\geq1$, $\delta>0$ takie, że z dowolnego
stanu $i\in T$ prawdopodobieństwo wejścia do $A$
w najbliższych $L$ krokach jest co najmniej
$\delta$. Dzięki własności Markowa
$\Prob_i(\tau_A>kL)\leq(1-\delta)^k$.
A zatem wszystkie sumy wierszy $Q^n$ dążą do zera,
i to geometrycznie. Otrzymujemy zbieżność
szeregu macierzy $\sum_{n\geq0}Q^n$.
Dla sumy częściowej
$(I-Q)\sum_{n=0}^{m}Q^n=I-Q^{m+1}$.
Przechodząc z $m\to\infty$, dostajemy
$(I-Q)N=I$; analogicznie $N(I-Q)=I$.

Dla $i,j\in T$ prawdopodobieństwo
$\Prob_i(X_n=j,\tau_A>n)$ jest równe
$(Q^n)_{ij}$. Sumując po $n$ i stosując
monotoniczną zbieżność, otrzymujemy
$N_{ij}=\E_i\sum_{n<\tau_A}\ind{\{X_n=j\}}$.
Sumowanie po $j$ daje
$\E_i\tau_A=(N\mathbf 1)_i$.
Prawdopodobieństwo pochłonięcia w $a\in A$
dokładnie po $n+1$ krokach wynosi
$(Q^nR)_{ia}$. Sumujemy po $n$, więc
$B_{ia}=\sum_{n\geq0}(Q^nR)_{ia}=(NR)_{ia}$.
\end{proof}

\begin{przyklad}[Cztery stany, dwa cele pochłaniające]
Niech $S=\{0,1,2,3\}$, stany $0,3$ będą
pochłaniające, a przejścia z pozostałych
stanów będą następujące:
ze stanu $1$ idziemy do $0$ z prawdopodobieństwem
$1/4$ lub do $2$ z prawdopodobieństwem $3/4$;
ze stanu $2$ idziemy do $1$ z prawdopodobieństwem
$1/2$ lub do $3$ z prawdopodobieństwem $1/2$.
Dla uporządkowania $T=(1,2)$ i $A=(0,3)$ mamy
\[
Q=\begin{pmatrix}0&3/4\\1/2&0\end{pmatrix},
\qquad R=\begin{pmatrix}1/4&0\\0&1/2\end{pmatrix},
\]
oraz
\[
N=\begin{pmatrix}8/5&6/5\\4/5&8/5\end{pmatrix},
\qquad B=NR=\begin{pmatrix}2/5&3/5\\1/5&4/5\end{pmatrix}.
\]
Zaczynając w stanie $1$, spędzamy średnio
$8/5$ kroków w stanie $1$ i $6/5$ w stanie $2$,
a czas do pochłonięcia ma średnią $14/5$.
Prawdopodobieństwo zakończenia w stanie $3$
wynosi $3/5$. Zaczynając w stanie $2$, czas
ma średnią $12/5$, a szansa zakończenia w $3$
wynosi $4/5$.
\end{przyklad}


\section{Zadanie rozwiązane: trajektorie, warunkowanie i pochłanianie}

\begin{przyklad}[Urządzenie: pełne rozwiązanie jednego modelu]
Urządzenie jest w jednym z trzech stanów:
$0$ (sprawne), $1$ (awaria przejściowa),
$2$ (wyłączone na stałe). Macierz przejścia wynosi
\[
P=\begin{pmatrix}
\frac12&\frac13&\frac16\\
\frac14&\frac12&\frac14\\
0&0&1
\end{pmatrix},
\qquad
\mu=\left(\frac34,\frac14,0\right).
\]
Rozwiążemy pięć pytań, zwracając uwagę na różnicę
między warunkowaniem, sumowaniem trajektorii
i czasami zatrzymania.

\emph{(1) Jakie jest prawdopodobieństwo trajektorii
$0\to1\to0\to2$?}
Początkowy stan nie jest ustalony, dlatego
musimy uwzględnić $\mu_0=3/4$:
\[
\Prob(X_0=0,X_1=1,X_2=0,X_3=2)
=\frac34\cdot\frac13\cdot\frac14\cdot\frac16
=\frac1{96}.
\]

\emph{(2) Jak porównać przewidywanie stanu $1$
po jednej i po dwóch obserwacjach?}
Jeśli wiemy, że $X_0=0$ i $X_1=1$, własność
Markowa daje
$\Prob(X_2=1\mid X_0=0,X_1=1)=p_{11}=1/2$.
Jeżeli znamy jedynie $X_0=0$, rozważamy wszystkie
możliwe stany pośrednie:
\[
\Prob(X_2=1\mid X_0=0)
=(P^2)_{01}=\frac12\frac13+\frac13\frac12=\frac13.
\]
Te liczby są różne, ponieważ \emph{poznanie stanu
pośredniego} zmienia informację, na której opiera się prognoza.

\emph{(3) Jak odtworzyć stan początkowy z późniejszych obserwacji?}
Po zaobserwowaniu $X_1=1$ chcemy obliczyć
$\Prob(X_0=0\mid X_1=1)$. Mamy
\[
\Prob(X_1=1)=\frac34\frac13+\frac14\frac12=\frac38,
\qquad
\Prob(X_0=0\mid X_1=1)
=\frac{(3/4)(1/3)}{3/8}=\frac23.
\]
Jeśli wiemy dodatkowo, że $X_2=0$, wynik nadal
wynosi $2/3$, bo czynnik $p_{10}=1/4$ występuje
zarówno w liczniku, jak i w mianowniku. To konkretny
przykład warunkowej niezależności przeszłości
i dalszej przyszłości przy znanym stanie pośrednim.

\emph{(4) Jakie jest prawdopodobieństwo trwałego
wyłączenia w ciągu dwóch kroków, jeśli $X_0=0$?}
Stan $2$ jest pochłaniający, więc zdarzenie
$\{\tau_2\leq2\}$ jest równoważne $\{X_2=2\}$.
Można dojść do $2$ bezpośrednio w pierwszym kroku
lub w drugim kroku przez stan $0$ albo $1$:
\[
\Prob_0(\tau_2\leq2)
=\frac16+\frac12\frac16+\frac13\frac14
=\frac13.
\]
Pierwszy składnik obejmuje trajektorię, która
dotarła do stanu $2$ już w chwili $1$ i pozostała tam
do chwili $2$. Nie wolno liczyć jej ponownie jako
nowego wejścia w chwili $2$.

\emph{(5) Ile kroków średnio upłynie do trwałego wyłączenia?}
Niech $m_i=\E_i\tau_2$, gdzie $\tau_2$
jest pierwszym czasem wejścia do stanu $2$.
Mamy $m_2=0$, a analiza pierwszego kroku daje
\[
\begin{cases}
m_0=1+\frac12m_0+\frac13m_1,\\[2pt]
m_1=1+\frac14m_0+\frac12m_1.
\end{cases}
\]
Pierwsza jedynka jest kosztem wykonanego kroku,
a po nim pozostaje oczekiwany czas dalszego
oczekiwania. Przekształcając układ otrzymujemy
$3m_0-2m_1=6$ oraz $-m_0+2m_1=4$.
Dodanie równań daje $2m_0=10$, a więc
$m_0=5$ i $m_1=9/2$.
Skończoność tych średnich jest uzasadniona:
z każdego niepochłaniającego stanu można dotrzeć
do $2$ w jednym kroku z prawdopodobieństwem
co najmniej $1/6$, więc
$\Prob_i(\tau_2>n)\leq(5/6)^n$.
Wreszcie, jeśli początkowy stan ma rozkład $\mu$,
średni czas wynosi
$\E_\mu\tau_2=(3/4)\cdot5+(1/4)\cdot(9/2)=39/8$.
\end{przyklad}

\section{Dwa dalsze modele: urna i nieskończona kolejka}

\begin{przyklad}[Model Ehrenfesta: dlaczego pojawia się dwumianowy]
Rozważmy $N\geq1$ rozróżnialnych cząstek,
z których każda znajduje się w jednym z dwóch
pojemników. W każdym kroku wybieramy jednostajnie
jedną cząstkę i przenosimy do drugiego pojemnika.
Niech $X_n\in\{0,\ldots,N\}$ oznacza liczbę
cząstek w pierwszym pojemniku. Wtedy
\[
p_{k,k+1}=\frac{N-k}{N},\qquad
p_{k,k-1}=\frac{k}{N},
\]
przy czym nieistniejące przejścia na brzegach
pomijamy. Warunek równowagi szczegółowej
$\pi_k(N-k)/N=\pi_{k+1}(k+1)/N$
prowadzi do $\pi_{k+1}/\pi_k=(N-k)/(k+1)$,
więc po normalizacji
$\pi_k=2^{-N}\binom Nk$.
Jest to rozkład dwumianowy $\Bin(N,1/2)$.

Łańcuch ma jednak okres $2$, ponieważ parzystość
liczby cząstek w pierwszym pojemniku zmienia się
w każdym kroku. Rozkład $X_n$ nie zbiega z
ustalonego stanu do $\pi$, choć częstości odwiedzin
spełniają twierdzenie ergodyczne. Jeśli w każdym
kroku z prawdopodobieństwem $1/2$ pozostaniemy
bez ruchu, a z prawdopodobieństwem $1/2$ wykonamy
powyższe przejście, nowa macierz wynosi
$P'=(I+P)/2$. Jest nieokresowa, zachowuje $\pi$
i teraz rozkłady zbiegną do $\pi$.
\end{przyklad}

\begin{przyklad}[Stabilna i niestabilna kolejka w czasie dyskretnym]
Niech $S=\N_0$, $p,q\in(0,1)$, $p+q=1$.
W stanach $i\geq1$ przechodzimy do $i+1$
z prawdopodobieństwem $p$ i do $i-1$
z prawdopodobieństwem $q$.
W stanie $0$ przechodzimy do $1$ z
prawdopodobieństwem $p$ i pozostajemy w $0$
z prawdopodobieństwem $q$. Taki łańcuch
modeluje liczbę klientów w buforze z przyjściami
i odchodzeniem pojedynczych jednostek.

Równowaga szczegółowa daje
$\pi_i p=\pi_{i+1}q$, czyli
$\pi_i=\pi_0(p/q)^i$. Jeśli $\rho=p/q<1$,
możemy unormować szereg geometryczny,
otrzymując
\[
\pi_i=(1-\rho)\rho^i,
\qquad \E_\pi X=\frac{\rho}{1-\rho}.
\]
Gdy $p\geq q$, ta sama rekurencja nie może
definiować rozkładu probabilistycznego.
Nie istnieje również żaden inny rozkład
stacjonarny: równość masy wpływającej i
wypływającej ze skończonego zbioru
$\{0,\ldots,i\}$ wymusza
$\pi_i p=\pi_{i+1}q$ dla każdego $i$.
Bez sumowalności nie można spełnić
$\sum_i\pi_i=1$. Jest to przykład sytuacji,
w której dla skończonego łańcucha stacjonarność
była automatyczna, ale w przeliczalnym
tak być nie musi.
\end{przyklad}


\section{Granice zakresu przedstawionej teorii}

Zajmowaliśmy się \emph{przeliczalnymi} przestrzeniami
stanów i \emph{dyskretnym} czasem. Dla ogólnych
przestrzeni stanów macierz zastępuje jądro
przejścia $K(x,A)$, mierzalne w $x$ i będące
miarą probabilistyczną w $A$. Przy czasie ciągłym
pojawiają się półgrupy $(P_t)_{t\geq0}$,
generatory i problemy regularności trajektorii.
Systematyczną teorię tych procesów rozwiniemy
w tomie II; tutaj konstrukcje dyskretne
przygotowują formalizm filtracji i martyngałów.


\chapter{Filtracje, adaptacja i czasy zatrzymania}
\label{ch:filtracje}

\section{Informacja rosnąca w czasie}

Warunkowa wartość oczekiwana mierzy to, co można
przewidzieć na podstawie ustalonej informacji.
W procesie losowym informacja jednak rośnie:
po każdym kroku poznajemy nowy wynik. \emph{Filtracja}
porządkuje te kolejne etapy, a warunek adaptacji
gwarantuje, że proces w chwili $n$ nie odwołuje się
do przyszłości.

\begin{definicja}[Filtracja i proces adaptowany]
Niech $(\Omega,\mathcal F,\Prob)$ będzie przestrzenią
probabilistyczną. \emph{Filtracją w czasie dyskretnym}
nazywamy ciąg pod-$\sigma$-ciał
$(\mathcal F_n)_{n\geq0}$ spełniający
$\mathcal F_n\subseteq\mathcal F_{n+1}\subseteq\mathcal F$.
Proces $(X_n)_{n\geq0}$ jest \emph{adaptowany}, jeśli
$X_n$ jest $\mathcal F_n$-mierzalna dla każdego $n$.
\emph{Naturalna filtracja} procesu to
$\mathcal F_n^X=\sigma(X_0,\ldots,X_n)$.
Proces $(H_n)_{n\geq1}$ jest \emph{przewidywalny},
jeśli $H_n$ jest $\mathcal F_{n-1}$-mierzalna.
\end{definicja}

\begin{przyklad}[Historia, dodatkowa informacja i uprzedzenie przyszłości]
Dla kolejnych rzutów monetą $\xi_k\in\{-1,1\}$,
$S_n=\sum_{k=1}^n\xi_k$ i $S_0=0$, naturalna
filtracja sum $\mathcal F_n^S$ jest równa
$\sigma(\xi_1,\ldots,\xi_n)$: z trajektorii
$S_0,\ldots,S_n$ odzyskujemy wszystkie przyrosty
$\xi_k=S_k-S_{k-1}$. Możemy też rozszerzyć informację
np. o temperaturę obserwowaną w danej chwili.
Gdyby do $\mathcal F_0$ włączyć wartość
$\xi_1$, proces sum pozostałby adaptowany, ale
nie byłby już martyngałem, bo pierwszy przyrost
nie byłby nieprzewidywalny. Martyngałowość
zależy więc \emph{od filtracji}, a nie tylko od
wspólnego rozkładu procesu.
\end{przyklad}

\begin{przyklad}[Strategia przewidywalna]
Jeżeli $H_n$ oznacza liczbę jednostek aktywa,
którą wybieramy \emph{przed} zaobserwowaniem
zmiany ceny między $n-1$ i $n$, to powinna być
$\mathcal F_{n-1}$-mierzalna. Strategia
$H_n=\operatorname{sgn}(S_n-S_{n-1})$ używałaby
już wyniku kroku $n$ i dlatego nie jest
strategią przewidywalną.
\end{przyklad}


\begin{przyklad}[Filtracja rozpisana na elementarnych zdarzeniach]
Wykonujemy dwa niezależne uczciwe rzuty monetą.
Przestrzeń wyników jest
$\Omega=\{OO,OR,RO,RR\}$, gdzie pierwsza litera
opisuje pierwszy rzut. W chwili $0$ nic nie wiemy:
$\mathcal F_0=\{\varnothing,\Omega\}$. Po pierwszym
rzucie rozróżniamy atomy $\{OO,OR\}$ i $\{RO,RR\}$,
a więc
\[
\mathcal F_1=\{\varnothing,\Omega,
\{OO,OR\},\{RO,RR\}\},\qquad
\mathcal F_2=2^\Omega.
\]
Niech $Y=\ind{\{OO\}}$ oznacza zdarzenie,
że oba rzuty dały orła. Wtedy
$\E[Y\mid\mathcal F_0]=1/4$,
$\E[Y\mid\mathcal F_1]=(1/2)\ind{\{OO,OR\}}$
oraz $\E[Y\mid\mathcal F_2]=Y$.
Pierwsza prognoza jest stała, druga wykorzystuje wynik
pierwszego rzutu, a trzecia jest już dokładną
obserwacją. Pokazuje to, dlaczego filtracja
to matematyczny model \emph{narastania informacji}.
\end{przyklad}

\section{Czasy zatrzymania i informacja przy zatrzymaniu}

\begin{definicja}[Czas zatrzymania]
Zmienna losowa $\tau:\Omega\to\N_0\cup\{\infty\}$
jest \emph{czasem zatrzymania} względem filtracji
$(\mathcal F_n)$, jeśli $\{\tau\leq n\}\in\mathcal F_n$
dla każdego $n\geq0$. Wówczas także
$\{\tau=n\}\in\mathcal F_n$.
\end{definicja}

\begin{przyklad}[Czas wejścia do zbioru]
Dla procesu adaptowanego $X_n$ o wartościach w
przeliczalnym $S$ i $A\subseteq S$ czas
$\tau_A=\inf\{n\geq0:X_n\in A\}$ jest czasem
zatrzymania, bo
$\{\tau_A\leq n\}=\bigcup_{k=0}^n\{X_k\in A\}
\in\mathcal F_n$.
\end{przyklad}

\begin{przyklad}[Ostatnia wizyta na ogół nie jest czasem zatrzymania]
Dla niezależnych rzutów monetą z wartościami
$\xi_k\in\{-1,1\}$ i ustalonego $N\geq2$
rozważmy $\rho=\max\{k\in\{1,\ldots,N\}:\xi_k=1\}$,
a gdy nie ma takiego $k$, przyjmijmy $\rho=0$.
Zdarzenie $\{\rho\leq N-1\}=\{\xi_N=-1\}$
nie należy do $\sigma(\xi_1,\ldots,\xi_{N-1})$,
ponieważ ostatni rzut jest niezależny od wcześniejszych
i niedegenerowany. Dlatego $\rho$ nie jest czasem
zatrzymania. Aby rozpoznać \emph{ostatni} sukces,
trzeba znać jeszcze przyszłość.
\end{przyklad}

\begin{definicja}[Zatrzymana informacja]
Dla czasu zatrzymania $\tau$ definiujemy
\[
\mathcal F_\tau=
\{A\in\mathcal F: A\cap\{\tau\leq n\}\in\mathcal F_n
\text{ dla każdego }n\geq0\}.
\]
Jest to $\sigma$-ciało zdarzeń rozstrzygalnych
na podstawie informacji dostępnej w chwili $\tau$.
Dla procesu adaptowanego definiujemy
$X_{n\wedge\tau}$, zaś $X_\tau$ na $\{\tau<\infty\}$
ma naturalną definicję przez sumę
$\sum_n X_n\ind{\{\tau=n\}}$ w przypadku
procesu rzeczywistego (przy odpowiednim umownym
rozszerzeniu na $\{\tau=\infty\}$).
\end{definicja}

\begin{lemat}[Podstawowe własności zatrzymania]
Niech $\tau,\sigma$ będą czasami zatrzymania
względem $(\mathcal F_n)$. Wtedy:
\begin{enumerate}
\item $\tau\wedge\sigma$ oraz $\tau\vee\sigma$ są czasami zatrzymania;
\item jeśli $\sigma\leq\tau$, to $\mathcal F_\sigma\subseteq\mathcal F_\tau$;
\item jeśli $X$ jest adaptowany, to $X_\tau\ind{\{\tau<\infty\}}$
  jest $\mathcal F_\tau$-mierzalny, przy umownej wartości zero na $\{\tau=\infty\}$.
\end{enumerate}
\end{lemat}
\begin{proof}
Dla minimum mamy
$\{\tau\wedge\sigma\leq n\}=
\{\tau\leq n\}\cup\{\sigma\leq n\}$,
a dla maksimum zastępujemy sumę zbiorów
przecięciem. Oba należą do $\mathcal F_n$.
Jeśli $A\in\mathcal F_\sigma$, to przy
$\sigma\leq\tau$ zachodzi równość
\[
A\cap\{\tau\leq n\}
=(A\cap\{\sigma\leq n\})\cap\{\tau\leq n\},
\]
wynikająca z $\sigma\leq\tau$, a prawa strona
należy do $\mathcal F_n$.
Dla (3) niech $B\subseteq\R$ będzie borelowski.
Na zdarzeniu $\{\tau\leq n\}$ przeciwobraz
$\{X_\tau\in B\}$ jest sumą skończoną
$\bigcup_{k=0}^n\{\tau=k\}\cap\{X_k\in B\}$,
należącą do $\mathcal F_n$. Stąd mierzalność
względem $\mathcal F_\tau$ po wybraniu
wartości na $\{\tau=\infty\}$.
\end{proof}

\begin{uwaga}[Dwie równoważne definicje $\mathcal F_\tau$]
Dla każdego $A\in\mathcal F$ warunki
$A\cap\{\tau\leq n\}\in\mathcal F_n$ dla wszystkich $n$
oraz $A\cap\{\tau=n\}\in\mathcal F_n$ dla wszystkich $n$
są równoważne. Istotnie, $\{\tau=n\}=\{\tau\leq n\}
\setminus\{\tau\leq n-1\}$, a odwrotnie
$\{\tau\leq n\}=\bigcup_{k=0}^n\{\tau=k\}$.
Żadnego dodatkowego warunku na $\{\tau=\infty\}$
nie trzeba dopisywać do tej równoważności.
W przyjętej definicji, gdy $\tau\equiv\infty$,
mamy nawet $\mathcal F_\tau=\mathcal F$:
proces nie został zatrzymany w żadnej skończonej chwili,
a definicja nie nakłada ograniczeń na zdarzenia
poza częściami $\{\tau\leq n\}$. Jest to konwencja
wygodna dla ogólnych czasów zatrzymania;
przy zatrzymaniu skończonym p.n. nie prowadzi do
różnic w zwykłych rachunkach warunkowych.
\end{uwaga}


\begin{przyklad}[Jaka jest informacja w chwili losowego zatrzymania?]
Niech $\xi_1,\xi_2\in\{O,R\}$ będą niezależnymi,
uczciwymi rzutami monetą, a
$\tau=1$ przy $\xi_1=O$ oraz $\tau=2$ przy $\xi_1=R$.
Filtracja $\mathcal F_0,\mathcal F_1,\mathcal F_2$
pochodzi z poprzedniego przykładu czteroelementowego.
Atomami $\mathcal F_\tau$ są
$\{OO,OR\}$, $\{RO\}$ i $\{RR\}$.
Na pierwszym atomie przestajemy obserwować po
pierwszym rzucie i nie rozróżniamy wyników
$OO$ i $OR$. Na dwóch pozostałych atomach
widzieliśmy również drugi rzut.
Dla $Y=\ind{\{OO\}}$ otrzymujemy
$\E[Y\mid\mathcal F_\tau]
=\frac12\ind{\{OO,OR\}}$.
Przykład uświadamia, że $\mathcal F_\tau$ nie jest
jednym z góry ustalonym $\mathcal F_n$:
zakres obserwacji zależy od ścieżki.
\end{przyklad}

\section{Przewidywanie przy losowym czasie}

\begin{twierdzenie}[Warunkowanie po zatrzymaniu w skończonym horyzoncie]
Niech $N\in\N$, $\tau\leq N$ będzie czasem
zatrzymania, $X\in L^1(\Prob)$ oraz
$Z_n=\E[X\mid\mathcal F_n]$ dla $0\leq n\leq N$.
Wówczas
\[
\E[X\mid\mathcal F_\tau]
=\sum_{n=0}^N\ind{\{\tau=n\}}Z_n.
\]
\end{twierdzenie}
\begin{proof}
Prawa strona jest $\mathcal F_\tau$-mierzalna,
ponieważ na $\{\tau=n\}$ zależy wyłącznie
od $\mathcal F_n$. Jest całkowalna:
\[
\sum_{n=0}^N\E[|Z_n|\ind{\{\tau=n\}}]
\leq\sum_{n=0}^N\E|X|<\infty.
\]
Dla dowolnego $A\in\mathcal F_\tau$ mamy
$A\cap\{\tau=n\}\in\mathcal F_n$ i stąd
\begin{align*}
\E\Bigl[\ind A\sum_{n=0}^N
\ind{\{\tau=n\}}Z_n\Bigr]
&=\sum_{n=0}^N\E[\ind{A\cap\{\tau=n\}}Z_n]\\
&=\sum_{n=0}^N\E[\ind{A\cap\{\tau=n\}}X]
=\E[\ind A X].
\end{align*}
Jednoznaczność warunkowej wartości oczekiwanej
daje tezę.
\end{proof}

\section{Martyngał jako konsekwencja warunkowania}

\begin{definicja}[Martyngał, podmartyngał, nadmartyngał]
Niech $(\mathcal F_n)$ będzie filtracją i
$(M_n)$ rzeczywistym procesem adaptowanym
spełniającym $\E|M_n|<\infty$ dla każdego $n$.
Nazywamy go \emph{martyngałem}, jeśli
$\E[M_{n+1}\mid\mathcal F_n]=M_n$ p.n.
dla każdego $n\geq0$;
\emph{podmartyngałem}, jeśli warunkowa wartość
oczekiwana jest większa lub równa $M_n$;
\emph{nadmartyngałem}, jeśli jest mniejsza
lub równa $M_n$. Wszystkie nierówności
rozumiemy prawie na pewno.
\end{definicja}

W całej dalszej części książki słowo \emph{martyngał}
bez dopowiedzenia dotyczy czasu dyskretnego
$\N_0$. W czasie ciągłym wymagane są dodatkowe
zagadnienia regularności i modyfikacji procesu.

\begin{przyklad}[Ruch bez dryfu i z dryfem]
Niech $(\xi_k)$ będą niezależne i całkowalne,
$\mathcal F_n=\sigma(\xi_1,\ldots,\xi_n)$,
$S_n=\sum_{k=1}^n\xi_k$ i $S_0=0$.
Jeśli $\E\xi_k=0$, to $(S_n)$ jest martyngałem,
ponieważ
$\E[S_{n+1}\mid\mathcal F_n]
=S_n+\E\xi_{n+1}=S_n$.
Jeśli $\E\xi_k=m>0$, to jest podmartyngałem;
gdy $m<0$, nadmartyngałem. W każdym przypadku
$S_n-nm$ jest martyngałem. Równoważność
\emph{braku przewidywalnego przyrostu}
i martyngałowości jest tutaj dosłowna.
\end{przyklad}

\begin{twierdzenie}[Kryterium wszystkich przyszłych chwil]
Jeśli $(M_n)$ jest procesem całkowalnym i
adaptowanym, to jest martyngałem wtedy i tylko wtedy,
gdy dla wszystkich liczb $0\leq m\leq n$
$\E[M_n\mid\mathcal F_m]=M_m$ p.n.
Analogicznie podmartyngały i nadmartyngały
spełniają odpowiednią nierówność dla $m\leq n$.
\end{twierdzenie}
\begin{proof}
Przypadek $n=m+1$ jest definicją. Jeśli teza
zachodzi dla $n-1$, własność wieży daje
$\E[M_n\mid\mathcal F_m]
=\E[\E[M_n\mid\mathcal F_{n-1}]\mid\mathcal F_m]$.
Dla martyngału równa się to
$\E[M_{n-1}\mid\mathcal F_m]=M_m$.
Dla podmartyngału używamy monotoniczności
warunkowej wartości oczekiwanej, otrzymując
$\E[M_n\mid\mathcal F_m]\geq M_m$.
Dla nadmartyngału nierówności odwracamy.
\end{proof}

\section{Zatrzymywanie procesu: podstawowa konstrukcja}

\begin{twierdzenie}[Zatrzymany (nad-, pod-)martyngał]
Niech $(M_n)$ będzie martyngałem, podmartyngałem
lub nadmartyngałem względem $(\mathcal F_n)$,
a $\tau$ czasem zatrzymania. Wtedy proces
$M_n^\tau=M_{n\wedge\tau}$ ma tę samą własność.
\end{twierdzenie}
\begin{proof}
Zachodzi tożsamość
\[
M_{(n+1)\wedge\tau}-M_{n\wedge\tau}
=\ind{\{\tau>n\}}(M_{n+1}-M_n).
\]
Zdarzenie $\{\tau>n\}$ należy do $\mathcal F_n$,
więc po warunkowaniu dostajemy
\[
\E[M_{(n+1)\wedge\tau}-M_{n\wedge\tau}
\mid\mathcal F_n]
=\ind{\{\tau>n\}}\E[M_{n+1}-M_n\mid\mathcal F_n].
\]
Prawa strona jest zerowa dla martyngału,
nieujemna dla podmartyngału i niedodatnia
dla nadmartyngału. Zatrzymany proces jest
adaptowany, a dla ustalonego $n$
$|M_{n\wedge\tau}|\leq\sum_{k=0}^n|M_k|$,
więc jest całkowalny.
\end{proof}

\begin{twierdzenie}[Twierdzenie o próbkowaniu opcjonalnym]
Niech $(M_n)$ będzie martyngałem, a $\sigma,\tau$
czasami zatrzymania spełniającymi
$0\leq\sigma\leq\tau\leq N$ p.n. dla pewnego
ustalonego $N\in\N$. Wtedy
$\E[M_\tau\mid\mathcal F_\sigma]=M_\sigma$ p.n.,
a w szczególności $\E M_\tau=\E M_0$.
Dla podmartyngału (odpowiednio nadmartyngału)
równość zastępujemy nierównością $\geq$
(odpowiednio $\leq$).
\end{twierdzenie}
\begin{proof}
Niech $A\in\mathcal F_\sigma$.
Dla każdego $1\leq k\leq N$ wskaźnik
$H_k=\ind A\ind{\{\sigma<k\leq\tau\}}$
jest $\mathcal F_{k-1}$-mierzalny:
$A\cap\{\sigma<k\}=A\cap\{\sigma\leq k-1\}
\in\mathcal F_{k-1}$,
a $\{k\leq\tau\}=\{\tau>k-1\}\in\mathcal F_{k-1}$.
Ścieżkowo zachodzi teleskopowy wzór
\[
\ind A(M_\tau-M_\sigma)
=\sum_{k=1}^N H_k(M_k-M_{k-1}).
\]
Wszystkie składniki są całkowalne, więc
\[
\E[H_k(M_k-M_{k-1})]
=\E[H_k\E[M_k-M_{k-1}\mid\mathcal F_{k-1}]].
\]
Dla martyngału wyrażenie jest zerowe,
a dla pod- i nadmartyngału ma odpowiedni znak.
Wobec dowolności $A\in\mathcal F_\sigma$
i $\mathcal F_\sigma$-mierzalności $M_\sigma$
otrzymujemy przez definicję warunkowania
żądaną równość lub nierówność.
\end{proof}

\begin{uwaga}[Dlaczego ograniczenie czasu jest ważne]
Powyższe twierdzenie nie mówi, że można zatrzymywać
martyngał w \emph{dowolnym} prawie na pewno
skończonym czasie i zachować jego średnią.
Przejście $N\to\infty$ wymaga kontroli
całkowalności zatrzymanych wartości.
Bez takiej kontroli strategia „czekam aż wygram”
może dawać pozornie sprzeczny wynik.
\end{uwaga}

\begin{twierdzenie}[Zatrzymanie w czasie o skończonej średniej]
Niech $(M_n)$ będzie martyngałem, $\tau$
czasem zatrzymania o $\E\tau<\infty$,
a przyrosty spełniają
$|M_n-M_{n-1}|\leq c$ p.n. dla wszystkich $n\geq1$,
gdzie $c<\infty$. Wówczas $M_\tau\in L^1$
oraz $\E M_\tau=\E M_0$.
\end{twierdzenie}
\begin{proof}
Z nieujemności czasu i skończonej średniej
wynika $\Prob(\tau<\infty)=1$. Dla każdego $n$
\[
|M_{n\wedge\tau}|
\leq|M_0|+\sum_{k=1}^{n\wedge\tau}|M_k-M_{k-1}|
\leq |M_0|+c\tau.
\]
Prawa strona należy do $L^1$, więc
$M_{n\wedge\tau}\to M_\tau$ nie tylko p.n.,
lecz również w $L^1$ z twierdzenia o zbieżności
zdominowanej. Dla każdego $n$ proces zatrzymany
jest martyngałem i $\E M_{n\wedge\tau}=\E M_0$.
Przechodzimy do granicy w średnich.
\end{proof}


\begin{przyklad}[Zadanie rozwiązane: pierwszy orzeł, ale nie później niż po trzech rzutach]
Niech $\xi_1,\xi_2,\xi_3$ będą niezależnymi uczciwymi
znakami ($+1$ dla orła, $-1$ dla reszki),
$S_n=\sum_{k=1}^n\xi_k$ i
$\tau=\min\{\inf\{n\geq1:\xi_n=1\},3\}$.
Jest to czas zatrzymania i zawsze $\tau\leq3$.
Wszystkie możliwe \emph{zatrzymane} scenariusze
przedstawia tabela:
\begin{center}
\begin{tabular}{c|c|c|c}
zaobserwowane rzuty & $\Prob$ & $\tau$ & $S_\tau$\\\hline
$O$ & $1/2$ & $1$ & $1$\\
$RO$ & $1/4$ & $2$ & $0$\\
$RRO$ & $1/8$ & $3$ & $-1$\\
$RRR$ & $1/8$ & $3$ & $-3$
\end{tabular}
\end{center}
Zatem
\[
\E S_\tau=\frac12+0-\frac18-\frac38=0
=\E S_0,\qquad
\E\tau=\frac12+\frac12+\frac34=\frac74.
\]
Nawet jeśli wybór chwili zatrzymania zależy od
zaobserwowanych wyników, oczekiwany wynik uczciwej
gry nie zmienia się, gdy czas jest ograniczony.
Nie oznacza to, że $S_\tau=0$ na każdej trajektorii.
\end{przyklad}

\begin{przyklad}[Czekanie na jeden punkt]
Niech $S_n=\xi_1+\cdots+\xi_n$ będzie symetrycznym
spacerem losowym na $\Z$ z $S_0=0$,
a $\tau=\inf\{n\geq0:S_n=1\}$.
Z obliczenia ruiny gracza w przedziale
$[-m,1]$ otrzymujemy
$\Prob_0(\tau_1<\tau_{-m})=m/(m+1)$.
Ponieważ zdarzenia po lewej zawierają się
w $\{\tau_1<\infty\}$, po $m\to\infty$
$\Prob(\tau<\infty)=1$. Jednak $S_\tau=1$
prawie na pewno, więc $\E S_\tau=1\ne0=\E S_0$.
Przyrosty są ograniczone przez $1$;
z poprzedniego twierdzenia wynika, że
$\E\tau=\infty$. To nie jest kontrprzykład
do twierdzenia o zatrzymaniu: zawiodło
właśnie wymagane założenie.
\end{przyklad}



\begin{uwaga}[Dlaczego nie można usunąć ograniczenia czasu bez dowodu]
Niech $S_n$ będzie symetrycznym spacerem losowym i
$\tau=\inf\{n\geq0:S_n=1\}$. W poprzednim przykładzie
wykazano $\tau<\infty$ prawie na pewno, jednak
$\E S_\tau=1\ne0=\E S_0$. Dla
$\tau_m=\tau\wedge m$ mamy na mocy ograniczonego
zatrzymywania $\E S_{\tau_m}=0$ dla każdego $m$.
Jednocześnie $S_{\tau_m}\to1$ prawie na pewno.
Gdyby rodzina $(S_{\tau_m})$ była jednostajnie
całkowalna, zbieżność prawie na pewno i kryterium
Vitaliego dawałyby $\E S_{\tau_m}\to1$, co jest
sprzeczne ze stałą wartością zero.
Zatem rodzina zatrzymanych wartości nie jest
jednostajnie całkowalna. Przykład wskazuje
\emph{dokładnie}, które przejście graniczne jest
niedozwolone, zamiast poprzestawać na stwierdzeniu,
że $\tau$ jest nieograniczony.
\end{uwaga}

\chapter{Martyngały: rachunek, strategie i nierówności}
\label{ch:martyngaly}

\section{Średni przyrost równy zero nie oznacza braku zmian}

Martyngał może bardzo silnie fluktuować: nie jest
procesem stałym i kolejne wartości nie muszą być
niezależne. Warunek
$\E[M_{n+1}-M_n\mid\mathcal F_n]=0$
mówi jedynie, że \emph{po uwzględnieniu całej
aktualnie dostępnej informacji} nie ma
przewidywalnego zysku ani straty.
Przez równość zachodzi też
$\E M_n=\E M_0$ dla każdego $n$,
ale sama stałość średniej nie wystarcza do
martyngałowości.

\begin{przyklad}[Stała średnia bez martyngału]
Niech $\xi$ przyjmuje $\pm1$ z prawdopodobieństwem
$1/2$ i połóżmy $X_0=0$, $X_1=\xi$,
$X_2=-\xi$, $X_n=X_2$ dla $n\geq2$.
Względem naturalnej filtracji $\E X_n=0$
przy każdym $n$. Jednak
$\E[X_2\mid\mathcal F_1]=-\xi\ne\xi=X_1$.
Średnia bezwarunkowa gubi informację o znaku
pierwszego wyniku.
\end{przyklad}

\begin{przyklad}[Martyngał dla dowolnej zmiennej końcowej]
Niech $Y\in L^1(\Prob)$ i niech $(\mathcal F_n)$
będzie filtracją. Wówczas
$M_n=\E[Y\mid\mathcal F_n]$ jest martyngałem:
$\E[M_{n+1}\mid\mathcal F_n]=
\E[\E[Y\mid\mathcal F_{n+1}]\mid\mathcal F_n]=M_n$.
Interpretujemy $M_n$ jako aktualną prognozę
przyszłej wielkości $Y$ na podstawie informacji
do chwili $n$.
\end{przyklad}

\begin{przyklad}[Losowe mnożenie]
Niech $Z_1,Z_2,\ldots$ będą niezależne, nieujemne,
$\E Z_k=1$, i zdefiniujmy $M_0=1$,
$M_n=\prod_{k=1}^n Z_k$.
Wtedy $\E[M_{n+1}\mid\mathcal F_n]
=M_n\E Z_{n+1}=M_n$.
Niekoniecznie $M_n\to1$: możliwe jest nawet
$M_n\to0$ p.n. mimo $\E M_n=1$.
Do tej subtelności wrócimy przy zbieżności
martyngałów.
\end{przyklad}

\begin{przyklad}[Martyngał wykładniczy]
Niech $\xi_1,\xi_2,\ldots$ będą niezależne i jednakowo
rozłożone, a dla ustalonego $t\in\R$
niech $m(t)=\E e^{t\xi_1}<\infty$ i $m(t)>0$.
Dla $S_n=\sum_{k=1}^n\xi_k$ proces
\[
M_n(t)=\frac{e^{tS_n}}{m(t)^n},\qquad M_0(t)=1,
\]
jest nieujemnym martyngałem względem naturalnej
filtracji: po warunkowaniu następnego czynnika
otrzymujemy $\E e^{t\xi_{n+1}}/m(t)=1$.
Dla spaceru symetrycznego $m(t)=\cosh t$.
Takie martyngały umożliwiają szacowanie
prawdopodobieństw trafienia w odległy próg.
\end{przyklad}

\section{Różnice martyngałowe i ortogonalność}

\begin{definicja}[Różnice martyngałowe]
Niech $(M_n)$ będzie martyngałem. Dla $n\geq1$
przyjmijmy $D_n=M_n-M_{n-1}$.
Wówczas $D_n\in L^1$, jest $\mathcal F_n$-mierzalna
i $\E[D_n\mid\mathcal F_{n-1}]=0$.
Odwrotnie, jeśli proces adaptowany ma postać
$M_n=M_0+\sum_{k=1}^nD_k$,
$M_0\in L^1(\mathcal F_0)$,
$D_k\in L^1(\mathcal F_k)$ i
$\E[D_k\mid\mathcal F_{k-1}]=0$, to jest martyngałem.
\end{definicja}

\begin{twierdzenie}[Ortogonalność przyrostów w $L^2$]
Niech $(M_n)$ będzie martyngałem takim, że
$M_n\in L^2$ dla wszystkich $n$. Wtedy
$\E[D_iD_j]=0$ dla $i\ne j$ oraz
\[
\E[(M_n-M_0)^2]=\sum_{k=1}^n\E D_k^2,
\qquad
\E M_n^2=\E M_0^2+\sum_{k=1}^n\E D_k^2.
\]
\end{twierdzenie}
\begin{proof}
Załóżmy $i<j$. Ponieważ $D_i$ jest
$\mathcal F_i\subseteq\mathcal F_{j-1}$-mierzalna,
a iloczyn $D_iD_j$ jest całkowalny na mocy
nierówności Cauchy'ego--Schwarza, reguła
wyciągania czynnika z warunkowania daje
$\E[D_iD_j]=\E[D_i\E[D_j\mid\mathcal F_{j-1}]]=0$.
Podobnie
$\E[M_0D_j]=0$, gdyż $M_0$ jest
$\mathcal F_0$-mierzalna. Rozwijamy kwadrat
$M_n=M_0+\sum_{k=1}^nD_k$, a wszystkie
wyrazy mieszane znikają.
\end{proof}

\begin{uwaga}
Ortogonalność jest słabsza od niezależności:
różnice martyngałowe mogą zależeć od wszystkich
wcześniejszych obserwacji, mimo zerowych kowariancji.
\end{uwaga}

\begin{przyklad}[Zależne, lecz ortogonalne przyrosty]
Niech $\xi_k$ będą niezależnymi symetrycznymi
znakami, $D_1=\xi_1$, $D_2=\xi_1\xi_2$,
$D_3=\xi_1(1+\xi_2)\xi_3$ i
$D_n=\xi_n$ dla $n\geq4$.
Względem $\mathcal F_n=\sigma(\xi_1,\ldots,\xi_n)$
każdy przyrost ma warunkową średnią zero.
Jednak wartość bezwzględna $D_3$ zależy od
$\xi_2$, więc $D_3$ nie jest niezależna od
historii: na zdarzeniu $\{\xi_2=-1\}$ jest zerowa,
a na $\{\xi_2=1\}$ ma moduł $2$.
Mimo to wszystkie różne przyrosty są ortogonalne.
\end{przyklad}

\section{Suma przewidywalnie ważonych przyrostów}

\begin{twierdzenie}[Transformacja martyngałowa]
Niech $(M_n)$ będzie martyngałem i
$(H_n)_{n\geq1}$ przewidywalnym procesem.
Załóżmy, że każda zmienna
$H_n(M_n-M_{n-1})$ jest całkowalna
(np. wszystkie $H_n$ są ograniczone).
Wtedy
\[
(H\cdot M)_0=0,\qquad
(H\cdot M)_n=\sum_{k=1}^nH_k(M_k-M_{k-1})
\]
jest martyngałem. Jeżeli $M$ jest
podmartyngałem i $H_n\geq0$, transformacja
jest podmartyngałem; dla nadmartyngału
jest nadmartyngałem.
\end{twierdzenie}
\begin{proof}
Wartości transformacji są całkowalne i adaptowane.
Ponieważ $H_{n+1}$ jest $\mathcal F_n$-mierzalna,
reguła wyciągania czynnika, w jej wersji
z warunkiem całkowalności iloczynu, daje
\[
\E[(H\cdot M)_{n+1}-(H\cdot M)_n
\mid\mathcal F_n]
=H_{n+1}\E[M_{n+1}-M_n\mid\mathcal F_n].
\]
W martyngale prawa strona jest zerowa; w
pod- i nadmartyngale przy $H_{n+1}\geq0$
ma odpowiedni znak.
\end{proof}

\begin{przyklad}[Zmienianie stawek nie tworzy przewidywalnego zysku]
Niech $M_n$ oznacza skumulowany wynik uczciwej gry,
a $H_n$ stawkę wybraną na podstawie historii
$\mathcal F_{n-1}$. Jeżeli ryzyko jest na tyle
kontrolowane, że $H_n\Delta M_n$ jest całkowalny,
proces sumy zysków jest martyngałem.
\emph{Nie} wynika z tego, że wynik gry nie może
być dodatni na określonej trajektorii, tylko że
nie ma dodatniego \emph{warunkowego oczekiwanego}
przyrostu. Podwajanie stawki bez kontroli
całkowalności albo przy nieograniczonym czasie
nie jest objęte tym twierdzeniem.
\end{przyklad}

\section{Rozkład Dooba i energia martyngału}

\begin{twierdzenie}[Rozkład Dooba w czasie dyskretnym]
Niech $(X_n)$ będzie całkowalnym, adaptowanym
podmartyngałem. Wtedy istnieją jednoznaczne procesy
$(M_n)$ i $(A_n)$ takie, że
$X_n=M_n+A_n$, $M$ jest martyngałem,
$A_0=0$, $A_n$ jest całkowalny i adaptowany,
$\Delta A_n=A_n-A_{n-1}$ jest
$\mathcal F_{n-1}$-mierzalny i nieujemny.
Proces $A$ nazywamy przewidywalną częścią
rosnącą podmartyngału.
\end{twierdzenie}
\begin{proof}
Definiujemy
\[
A_0=0,\qquad
A_n=\sum_{k=1}^n\E[X_k-X_{k-1}\mid\mathcal F_{k-1}],
\qquad M_n=X_n-A_n.
\]
Każdy składnik sumy definiującej $A_n$ jest
$\mathcal F_{k-1}$-mierzalny, całkowalny i
nieujemny, ponieważ $X$ jest podmartyngałem.
Zatem $A$ ma wszystkie wymagane własności,
a
$\E[M_n-M_{n-1}\mid\mathcal F_{n-1}]
=\E[X_n-X_{n-1}\mid\mathcal F_{n-1}]
-(A_n-A_{n-1})=0$.
Jeśli $X=M'+A'$ byłby drugim rozkładem,
to z $\Delta A'_n=\Delta X_n-\Delta M'_n$
po warunkowaniu otrzymujemy
$\Delta A'_n=\E[\Delta X_n\mid\mathcal F_{n-1}]$.
Wraz z $A'_0=0$ wymusza to $A'=A$ i $M'=M$.
\end{proof}

\begin{przyklad}[Kwadrat martyngału]
Dla martyngału $M$ z wartościami w $L^2$ funkcja
$x\mapsto x^2$ jest wypukła, więc
$(M_n^2)$ jest podmartyngałem. Jego część rosnąca
jest dana przez
\[
\langle M\rangle_n
=\sum_{k=1}^n\E[(M_k-M_{k-1})^2\mid\mathcal F_{k-1}],
\]
przy czym $M_n^2-\langle M\rangle_n$ jest
martyngałem. Rzeczywiście
$M_n^2-M_{n-1}^2=2M_{n-1}\Delta M_n+(\Delta M_n)^2$,
a pierwszy składnik ma warunkową średnią zero.
Wielkość $\langle M\rangle_n$ jest
\emph{przewidywalną wariacją kwadratową};
różni się od obserwowanej sumy
$[M]_n=\sum_{k=1}^n(\Delta M_k)^2$.
\end{przyklad}

\begin{przyklad}[Zależna zmienność]
Niech $\xi_k\in\{-1,1\}$ będą niezależne i uczciwe,
a $\sigma_k>0$ niech będzie
$\mathcal F_{k-1}$-mierzalna oraz spełnia
$\E\sigma_k^2<\infty$. Przyjmijmy
$M_n=\sum_{k=1}^n\sigma_k\xi_k$.
Wtedy $M$ jest martyngałem, a
$\langle M\rangle_n=\sum_{k=1}^n\sigma_k^2$.
Zmienność przyrostu może zależeć od historii;
warunek martyngałowy wymaga tylko zaniku
warunkowej średniej.
\end{przyklad}

\section{Nierówności maksymalne Dooba}

\begin{twierdzenie}[Nierówność maksymalna $L^1$]
Niech $(X_k)_{0\leq k\leq n}$ będzie nieujemnym
podmartyngałem względem $(\mathcal F_k)$.
Dla każdego $\lambda>0$
\[
\lambda\Prob\left(\max_{0\leq k\leq n}X_k\geq\lambda\right)
\leq\E\left[X_n\ind{\{\max_{0\leq k\leq n}X_k\geq\lambda\}}\right]
\leq\E X_n.
\]
W szczególności dla martyngału $(M_k)$
$\Prob(\max_{k\leq n}|M_k|\geq\lambda)
\leq\E|M_n|/\lambda$.
\end{twierdzenie}
\begin{proof}
Niech $A=\{\max_{0\leq k\leq n}X_k\geq\lambda\}$
i $\tau=\inf\{k\leq n:X_k\geq\lambda\}$,
z umową $\tau=n$ na ścieżkach bez przekroczenia.
Rozbijmy $A$ na rozłączne zdarzenia
$A_k=\{X_0<\lambda,\ldots,X_{k-1}<\lambda,
X_k\geq\lambda\}$ dla $k\leq n$.
Mamy $A_k\in\mathcal F_k$ i
$X_k\geq\lambda$ na $A_k$.
Z podmartyngałowości i wieży
$\E[X_n\ind{A_k}]\geq\E[X_k\ind{A_k}]
\geq\lambda\Prob(A_k)$.
Sumujemy po $k=0,\ldots,n$, dostając pierwszą
nierówność. Druga wynika z nieujemności $X_n$.
Dla martyngału $|M_k|$ jest podmartyngałem
przez warunkową nierówność Jensena.
\end{proof}

\begin{twierdzenie}[Nierówność maksymalna $L^p$]
Niech $1<p<\infty$, $n\in\N_0$ i
$(X_k)_{k\leq n}$ będzie nieujemnym
podmartyngałem z $X_n\in L^p$. Wtedy
\[
\left\|\max_{0\leq k\leq n}X_k\right\|_p
\leq\frac{p}{p-1}\,\|X_n\|_p.
\]
Dla martyngału $M$ z $M_n\in L^p$ można
zastąpić $X_k$ przez $|M_k|$.
\end{twierdzenie}
\begin{proof}
Z nierówności podmartyngałowej $X_k\leq
\E[X_n\mid\mathcal F_k]$ oraz warunkowego
Jensena wynika najpierw, że
$\E X_k^p\leq\E X_n^p$ dla wszystkich $k\leq n$;
w szczególności każde $X_k$ ma skończony $p$-ty moment.
Oznaczmy $Y=\max_{k\leq n}X_k$.
Z poprzedniego twierdzenia, dla każdego $t>0$,
$t\Prob(Y\geq t)\leq\E[X_n\ind{\{Y\geq t\}}]$.
Najpierw zastąpmy $Y$ przez $Y\wedge L$,
aby wszystkie całki były skończone.
Wzór warstwowy i twierdzenie Tonellego dają
\begin{align*}
\E[(Y\wedge L)^p]
&=p\int_0^L t^{p-1}\Prob(Y\geq t)\,\dd\lambda_1(t)\\
&\leq p\int_0^L t^{p-2}
\E[X_n\ind{\{Y\geq t\}}] \,\dd\lambda_1(t)\\
&=\frac{p}{p-1}\E\bigl[X_n(Y\wedge L)^{p-1}\bigr]\\
&\leq\frac{p}{p-1}\|X_n\|_p
\|Y\wedge L\|_p^{p-1},
\end{align*}
gdzie w ostatnim kroku zastosowano nierówność
Höldera. Jeśli $\|Y\wedge L\|_p=0$, teza jest
oczywista; w innym przypadku dzielimy przez
$\|Y\wedge L\|_p^{p-1}$, a następnie
przechodzimy z $L\to\infty$ przez monotoniczną
zbieżność.
\end{proof}

\begin{przyklad}[Kontrola maksimum trajektorii]
Jeżeli $S_n$ jest symetrycznym spacerem losowym,
to $\E S_n^2=n$. Z nierówności Dooba dla $p=2$
otrzymujemy
$\E(\max_{k\leq n}|S_k|)^2\leq4n$.
Nierówność Markowa daje
$\Prob(\max_{k\leq n}|S_k|\geq u)\leq4n/u^2$.
Kontrolujemy w ten sposób \emph{całą drogę},
a nie tylko wartość końcową $S_n$.
\end{przyklad}

\section{Nierówność Azumy--Hoeffdinga}

\begin{twierdzenie}[Koncentracja martyngału o ograniczonych przyrostach]
Niech $(M_k)_{0\leq k\leq n}$ będzie rzeczywistym
martyngałem. Załóżmy, że istnieją stałe
$c_k\geq0$ spełniające
$|M_k-M_{k-1}|\leq c_k$ p.n. dla $1\leq k\leq n$.
Wtedy dla każdego $u>0$ i
$V=\sum_{k=1}^n c_k^2>0$ zachodzi
\[
\Prob(M_n-M_0\geq u)\leq e^{-u^2/(2V)},
\quad
\Prob(|M_n-M_0|\geq u)\leq2e^{-u^2/(2V)}.
\]
Gdy $V=0$, przyrosty są zerowe p.n. i zdarzenie
odchylenia o dodatnie $u$ ma prawdopodobieństwo zero.
\end{twierdzenie}
\begin{proof}
Dla $D_k=\Delta M_k$ mamy
$\E[D_k\mid\mathcal F_{k-1}]=0$ oraz
$-c_k\leq D_k\leq c_k$.
Warunkowa wersja lematu Hoeffdinga dowodzona
jest dokładnie jak wersja bezwarunkowa, z
miarą prawdopodobieństwa zastąpioną warunkową
średnią i tym samym oszacowaniem wariancji
zmiennej na odcinku. Daje
\[
\E[e^{tD_k}\mid\mathcal F_{k-1}]
\leq e^{t^2c_k^2/2}\qquad(t\in\R).
\]
Można to również uzasadnić bez odwoływania się
do regularnych rozkładów warunkowych:
na przedziale $[-c_k,c_k]$ funkcja $e^{tx}$
leży poniżej cięciwy łączącej końce,
więc po warunkowaniu i użyciu zerowej średniej
otrzymuje się $\cosh(tc_k)\leq e^{t^2c_k^2/2}$.
Własność wieży daje rekurencyjnie
\[
\E e^{t(M_n-M_0)}\leq
\exp\left(\frac{t^2}{2}\sum_{k=1}^n c_k^2\right).
\]
Dla $t>0$ nierówność Markowa prowadzi do
$\Prob(M_n-M_0\geq u)\leq
\exp(-tu+t^2V/2)$.
Minimalizujemy przy $t=u/V$, otrzymując
$e^{-u^2/(2V)}$. Ten sam rachunek stosujemy do
$-(M_n-M_0)$ i dodajemy dwa ogony.
\end{proof}

\begin{przyklad}[Uczciwy spacer i prawdopodobieństwo odchylenia]
Dla niezależnych symetrycznych znaków
$S_n=\sum_{k\leq n}\xi_k$ i $c_k=1$.
Zatem dla $u>0$
$\Prob(|S_n|\geq u)\leq2e^{-u^2/(2n)}$.
Nierówność Czebyszewa dałaby jedynie
$\Prob(|S_n|\geq u)\leq n/u^2$.
Oszacowanie wykładnicze jest dużo mocniejsze
przy dużych odchyleniach, ale korzysta z ograniczenia
przyrostów, nie tylko z wariancji.
\end{przyklad}

\begin{przyklad}[Odkrywanie współrzędnych funkcji wielu losowań]
Niech $\xi_1,\ldots,\xi_n$ będą niezależne, a
$f(\xi_1,\ldots,\xi_n)$ całkowalna.
Definiujemy martyngał Dooba
$M_k=\E[f(\xi_1,\ldots,\xi_n)\mid
\sigma(\xi_1,\ldots,\xi_k)]$, gdzie $M_0=\E f$.
Jeżeli zmiana tylko $k$-tej współrzędnej
może zmienić $f$ o co najwyżej $c_k$,
to dla ustalonej historii $\xi_1,\ldots,\xi_{k-1}$
możliwe wartości $M_k$ leżą w przedziale długości
co najwyżej $c_k$. W konsekwencji
$|M_k-M_{k-1}|\leq c_k$; Azuma daje
\[
\Prob(|f-\E f|\geq u)\leq
2\exp\left(-\frac{u^2}{2\sum_kc_k^2}\right).
\]
Ostrzejszą stałą $2$ w wykładniku daje
warunkowy lemat Hoeffdinga zastosowany do
przyrostu $M_k$, którego \emph{przedział wartości}
ma długość $c_k$; wtedy otrzymuje się
$\Prob(|f-\E f|\geq u)\leq
2\exp(-2u^2/\sum_kc_k^2)$.
To znana nierówność ograniczonych różnic.
\end{przyklad}


\begin{uwaga}[Dlaczego w nierówności ograniczonych różnic jest lepsza stała?]
Samo oszacowanie $|\Delta M_k|\leq c_k$ mówi,
że przyrost leży w przedziale długości $2c_k$.
W przykładzie odkrywania współrzędnych
istnieje jednak przedział $[a_k,b_k]$
o długości $b_k-a_k\leq c_k$, wyznaczony przez
informację do chwili $k-1$, w którym leży
$\Delta M_k$. Dla zmiennej o warunkowej
średniej zero warunkowy lemat Hoeffdinga daje
\[
 \E[e^{t\Delta M_k}\mid\mathcal F_{k-1}]
 \leq\exp\bigl(t^2(b_k-a_k)^2/8\bigr)
 \leq\exp(t^2c_k^2/8).
\]
Po wymnożeniu oszacowań przez wieżę i minimalizacji
$-tu+t^2\sum_kc_k^2/8$ otrzymujemy
$\exp(-2u^2/\sum_kc_k^2)$ dla jednego ogona.
\end{uwaga}

\section{Proces Galtona--Watsona i znormalizowany martyngał}

\begin{przyklad}[Rozmnażanie i zmienna populacja]
Niech $Z_0=1$, a
$Z_{n+1}=\sum_{j=1}^{Z_n}\xi_{n,j}$, gdzie
$\xi_{n,j}$ są niezależnymi, jednakowo rozłożonymi
nieujemnymi całkowitymi zmiennymi losowymi,
niezależnymi od dotychczasowej historii,
o średniej $m=\E\xi_{n,j}>0$ i wariancji
$\sigma^2<\infty$. Filtrację $\mathcal F_n$
tworzą populacje i wszystkie losowania do
chwili $n-1$. Warunkowo przy $Z_n$ mamy
\[
\E[Z_{n+1}\mid\mathcal F_n]=mZ_n,
\qquad \Var(Z_{n+1}\mid\mathcal F_n)=\sigma^2 Z_n.
\]
Zatem $W_n=Z_n/m^n$ jest nieujemnym martyngałem.
Średnia $\E Z_n=m^n$ wynika indukcyjnie z wieży,
a więc $\E W_n=1$.

Jeśli $m>1$, to przez rozkład wariancji
\[
\E[(W_{n+1}-W_n)^2]
=\frac{\sigma^2\E Z_n}{m^{2(n+1)}}
=\frac{\sigma^2}{m^{n+2}}.
\]
Ortogonalność przyrostów daje
\[
\E W_n^2=1+\frac{\sigma^2}{m(m-1)}(1-m^{-n})
\leq1+\frac{\sigma^2}{m(m-1)}.
\]
Wniosek o zbieżności w $L^2$ i p.n. uzyskamy
w następnym rozdziale. Dla $m\leq1$ oszacowanie
nie daje jednostajnej kontroli drugiego momentu;
rozstrzygnięcie granic wymaga innego argumentu.
\end{przyklad}


\chapter{Podmartyngały, nadmartyngały i twierdzenia o zbieżności}
\label{ch:nadmartyngaly}

\section{Przewidywalny kierunek zmian a trajektorie}

Podmartyngał ma nieujemne warunkowe przyrosty,
a nadmartyngał niepozytywne. Nie oznacza to,
że ich trajektorie są odpowiednio niemalejące
lub nierosnące: nawet nadmartyngał może
wiele razy rosnąć, a podmartyngał spadać.
Pojęcia te opisują \emph{warunkową tendencję},
a nie monotoniczność wyników pojedynczego losowania.

\begin{przyklad}[Proces z ujemnym dryfem może rosnąć]
Niech $\xi_k$ będą niezależne i przyjmują
$+2$ z prawdopodobieństwem $1/4$ oraz $-1$
z prawdopodobieństwem $3/4$. Wtedy
$\E\xi_k=-1/4$, więc
$X_n=\sum_{k=1}^n\xi_k$ jest nadmartyngałem.
Mimo to w każdym kroku może wzrosnąć o $2$.
Proces $M_n=X_n+n/4$ jest martyngałem.
\end{przyklad}

\begin{twierdzenie}[Wypukła funkcja martyngału]
Niech $(M_n)$ będzie martyngałem,
a $\phi:\R\to\R$ skończoną wypukłą funkcją.
Jeśli $\phi(M_n)\in L^1$ dla każdego $n$,
to $(\phi(M_n))$ jest podmartyngałem.
W szczególności procesy $|M_n|$, $|M_n|^p$
przy odpowiedniej całkowalności ($p\geq1$)
oraz $M_n^2$ dla martyngału w $L^2$
są podmartyngałami.
\end{twierdzenie}
\begin{proof}
Warunkowa nierówność Jensena daje
$\E[\phi(M_{n+1})\mid\mathcal F_n]
\geq\phi(\E[M_{n+1}\mid\mathcal F_n])
=\phi(M_n)$. Wszystkie zmienne są całkowalne
na mocy założeń.
\end{proof}

\begin{przyklad}[Konweksość nie ratuje braku całkowalności]
Dla martyngału $M$ z samym pierwszym momentem
$M_n^2$ może mieć nieskończoną wartość oczekiwaną,
a więc nie jest podmartyngałem w definicji
przyjmującej całkowalność. Dlatego przy procesach
$M_n^2$ zakładamy $M_n\in L^2$.
\end{przyklad}

\begin{przyklad}[Gdy po zatrzymaniu zostaje nieukończona próba]
Dla deterministycznej ścieżki
$(x_0,\ldots,x_6)=(0,2,1,3,0,1,2)$
i przedziału $[a,b]=[1,2]$ ukończone przekroczenia w górę to trzy odrębne pary:
$0\to2$ (kroki $0,1$), $1\to3$ (kroki $2,3$)
i $0\to1\to2$ (kroki $4,5,6$).
Aby uniknąć nieporozumień, można wypisać czasy
$\sigma_1=0$, $\tau_1=1$,
$\sigma_2=2$, $\tau_2=3$,
$\sigma_3=4$, $\tau_3=6$,
a zatem $U_6(1,2)=3$. W dowodzie nierówności
przekroczeń zliczamy właśnie takie pełne pary,
a nie wszystkie wzrosty ani liczbę przekroczeń
jednego wybranego poziomu.
\end{przyklad}

\section{Przekroczenia przedziałów i oscylacje}

Aby dowieść zbieżności prawie na pewno, nie wystarczy
kontrolować średniej procesu. Trzeba wykluczyć
nieskończone oscylacje między dwiema odległymi
wartościami. Służy temu \emph{nierówność
przekroczeń w górę} (upcrossing inequality).

\begin{definicja}[Liczba przekroczeń przedziału]
Niech $x_0,\ldots,x_N\in\R$ oraz $a<b$.
Definiujemy rekurencyjnie
$\sigma_1=\inf\{k\geq0:x_k\leq a\}$,
$\tau_1=\inf\{k>\sigma_1:x_k\geq b\}$,
a dla $r\geq2$
$\sigma_r=\inf\{k>\tau_{r-1}:x_k\leq a\}$
oraz $\tau_r=\inf\{k>\sigma_r:x_k\geq b\}$,
przyjmując $\inf\varnothing=\infty$.
Liczbą ukończonych przekroczeń w górę
od $a$ do $b$ do chwili $N$ nazywamy
$U_N(a,b)=\#\{r:\tau_r\leq N\}$.
Dla procesu $X$ stosujemy tę definicję
ścieżka po ścieżce.
\end{definicja}

\begin{twierdzenie}[Nierówność przekroczeń Dooba]
Niech $(X_n)_{n\geq0}$ będzie rzeczywistym,
całkowalnym podmartyngałem względem
$(\mathcal F_n)$. Dla liczb $a<b$ i $N\in\N_0$
\[
(b-a)\E U_N(a,b)
\leq\E[(X_N-a)^+]-\E[(X_0-a)^+].
\]
\end{twierdzenie}
\begin{proof}
Połóżmy $Y_k=(X_k-a)^+$. Ponieważ funkcja
$x\mapsto(x-a)^+$ jest wypukła i niemalejąca,
warunkowa nierówność Jensena oraz podmartyngałowość
$X$ dają
$\E[Y_k\mid\mathcal F_{k-1}]\geq Y_{k-1}$.
Zatem $Y$ jest nieujemnym podmartyngałem.

Niech $\sigma_r,\tau_r$ będą chwilami określonymi
w definicji przekroczenia. Są czasami zatrzymania:
po każdej obserwacji wiadomo, czy nastąpiło
rozpoczęcie lub zakończenie kolejnego przekroczenia,
ponieważ obie decyzje zależą wyłącznie od historii
do tej chwili. Można to sformalizować indukcyjnie
za pomocą warunków $\{\sigma_r\leq n\}$ i
$\{\tau_r\leq n\}$, będących skończonymi sumami
przecięć zdarzeń z $\mathcal F_n$.

Przed wykonaniem kroku $k$ określmy
\[
H_k=\sum_{r=1}^{N+1}
 \ind{\{\sigma_r<k\leq\tau_r\}},
 \qquad 1\leq k\leq N.
\]
Nie więcej niż jeden składnik sumy może być równy
$1$, bo przedziały $(\sigma_r,\tau_r]$ są rozłączne.
Ponadto
$\{\sigma_r<k\}=\{\sigma_r\leq k-1\}\in\mathcal F_{k-1}$
oraz $\{\tau_r\geq k\}=\{\tau_r>k-1\}\in\mathcal F_{k-1}$.
Oznacza to, że $H_k\in\{0,1\}$ jest
$\mathcal F_{k-1}$-mierzalny: decyzję o
uwzględnieniu kroku $k$ podejmujemy \emph{przed}
poznaniem przyrostu $Y_k-Y_{k-1}$.

Każde ukończone przekroczenie daje
$Y_{\tau_r}-Y_{\sigma_r}\geq b-a$.
Jeśli ostatnie rozpoczęte przekroczenie nie
zostało zakończone przed $N$, jego niedokończony
wkład $Y_N-Y_{\sigma_r}=Y_N\geq0$.
Dlatego ścieżkowo
\[
\sum_{k=1}^N H_k(Y_k-Y_{k-1})
 \geq(b-a)U_N(a,b).
\]
Z drugiej strony proces $Y$ jest podmartyngałem,
a więc dla każdego $k$
\[
\E[(1-H_k)(Y_k-Y_{k-1})]
=\E[(1-H_k)\E[Y_k-Y_{k-1}\mid\mathcal F_{k-1}]]\geq0.
\]
Rozpisujemy sumę całkowitych przyrostów $Y$ na
część wybraną przez $H$ i niewybraną:
\begin{align*}
(b-a)\E U_N(a,b)
&\leq\E\sum_{k=1}^NH_k(Y_k-Y_{k-1})\\
&=\E(Y_N-Y_0)
 -\sum_{k=1}^N\E[(1-H_k)(Y_k-Y_{k-1})]\\
&\leq\E(Y_N-Y_0)
=\E[(X_N-a)^+]-\E[(X_0-a)^+].
\end{align*}
Wszystkie sumy są skończone, a ich składniki
całkowalne; nie przechodziliśmy do nieskończonego
horyzontu bez oszacowania.
\end{proof}


\begin{uwaga}[Od liczby przekroczeń do istnienia granicy]
Jeśli ciąg rzeczywisty ma
$\liminf_nx_n<\limsup_nx_n$, to istnieją
liczby wymierne $a<b$ pomiędzy tymi granicami.
Dla dowolnie dużego czasu znajdziemy wówczas
późniejsze wartości $\leq a$ i jeszcze późniejsze
$\geq b$, czyli kolejny przekroczony przedział.
To prowadzi do nieskończenie wielu przekroczeń.
Dlatego skończoność $U_\infty(a,b)$
\emph{dla wszystkich wymiernych $a<b$ jednocześnie}
wyklucza oscylacje; pozostaje jednak osobne
zadanie wykluczenia granic $\pm\infty$.
\end{uwaga}

\begin{twierdzenie}[Twierdzenie Dooba o zbieżności podmartyngałów]
Niech $(X_n)$ będzie rzeczywistym podmartyngałem
i załóżmy, że
$\sup_{n\geq0}\E[X_n^+]<\infty$.
Wtedy istnieje skończona rzeczywista zmienna
losowa $X_\infty$ taka, że
$X_n\to X_\infty$ prawie na pewno.
W szczególności każdy nieujemny nadmartyngał
zbiega prawie na pewno do skończonej zmiennej
losowej, a każdy martyngał z
$\sup_n\E|M_n|<\infty$ zbiega p.n.
\end{twierdzenie}
\begin{proof}
Ustalmy wymierne $a<b$. Z nierówności przekroczeń
oraz tożsamości
$(X_n-a)^+\leq X_n^++|a|$
wynika
\[
\sup_N\E U_N(a,b)
\leq\frac{\sup_N\E X_N^++|a|}{b-a}<\infty.
\]
Ponieważ $U_N(a,b)$ rośnie względem $N$,
z twierdzenia o zbieżności monotonicznej
$\E U_\infty(a,b)<\infty$, a zatem
$U_\infty(a,b)<\infty$ p.n. Dla przeliczalnie
wielu wymiernych par $a<b$ wybieramy jeden
zbiór pełnej miary, na którym wszystkie te liczby
przekroczeń są skończone.

Jeśli na takiej trajektorii zachodziłoby
$\liminf_n X_n<\limsup_n X_n$, można by
wybrać wymierne $a<b$ pomiędzy obiema granicami.
Trajektoria przekraczałaby przedział $[a,b]$
w górę nieskończenie wiele razy, sprzecznie
z uzyskanym wynikiem. Zatem istnieje granica
$X_\infty$ w rozszerzonej prostej $[-\infty,\infty]$.

Pozostaje wykluczyć wartości nieskończone.
Z lematu Fatou
$\E[X_\infty^+]\leq\liminf_n\E X_n^+<\infty$,
więc $X_\infty\ne+\infty$ p.n.
Ponieważ $\E X_n$ jest niemalejące,
$\E X_n\geq\E X_0$ i
$\E X_n^-=\E X_n^+-\E X_n
\leq\sup_k\E X_k^+-\E X_0$.
Fatou zastosowany do $X_n^-$ wyklucza
$X_\infty=-\infty$ z dodatnim prawdopodobieństwem.
To daje skończoną granicę.

Dla nieujemnego nadmartyngału $Z_n$ proces
$X_n=-Z_n$ jest podmartyngałem oraz
$X_n^+=0$. Dla martyngału ograniczonego w $L^1$
mamy $\sup_n\E M_n^+<\infty$.
Oba przypadki są więc konsekwencjami już
udowodnionej wersji twierdzenia.
\end{proof}

\begin{przyklad}[Przyrosty mogą oscylować, mimo zbieżności]
Niech $Z_k$ będą niezależnymi symetrycznymi znakami,
a $M_n=\sum_{k=1}^n2^{-k}Z_k$.
Jest to martyngał, a szereg
$\sum_k2^{-k}Z_k$ zbiega nawet bezwzględnie
na każdej trajektorii. Wiele przyrostów jest
naprzemiennie dodatnich i ujemnych.
Zbieżność martyngału nie wymaga monotoniczności.
\end{przyklad}

\section{Jednostajna całkowalność i granice w \texorpdfstring{$L^1$}{L1}}

\begin{definicja}[Jednostajna całkowalność rodziny]
Rodzina $\mathcal K\subseteq L^1(\Prob)$ jest
\emph{jednostajnie całkowalna}, jeśli
\[
\lim_{K\to\infty}\sup_{X\in\mathcal K}
\E[|X|\ind{\{|X|>K\}}]=0.
\]
\end{definicja}

\begin{lemat}[Rodzina warunkowych wartości jednego $L^1$]
Dla $Y\in L^1(\Prob)$ rodzina
$\{\E[Y\mid\mathcal G]:\mathcal G\subseteq\mathcal F
\text{ jest }\sigma\text{-ciałem}\}$
jest jednostajnie całkowalna.
\end{lemat}
\begin{proof}
Ustalmy $K>0$, $\mathcal G\subseteq\mathcal F$
i połóżmy $Z=\E[Y\mid\mathcal G]$ oraz
$V=|Y|\ind{\{|Y|>K/2\}}$.
Rozdzielając $Y$ na część o module nie większym
niż $K/2$ i pozostałą, otrzymujemy z warunkowej
nierówności trójkąta
\[
|Z|\leq\frac K2+\E[V\mid\mathcal G].
\]
Na zdarzeniu $A=\{|Z|>K\}\in\mathcal G$
stąd $\E[V\mid\mathcal G]>K/2$.
Po pomnożeniu przez $\ind A$ i wzięciu średniej
wynika nierówność
$\frac K2\Prob(A)\leq\E[V\ind A]\leq\E V$.
Z drugiej strony, znów na $A$, mamy
\[
\E[|Z|\ind A]
\leq\frac K2\Prob(A)
 +\E[\E[V\mid\mathcal G]\ind A]
\leq2\E V.
\]
Zastosowaliśmy tu definiującą własność
warunkowej wartości oczekiwanej przy zdarzeniu
$A\in\mathcal G$. Po prawej znajduje się
wyrażenie niezależne od wyboru $\mathcal G$.
Ponieważ $Y\in L^1$, zbieżność zdominowana daje
$\E[|Y|\ind{\{|Y|>K/2\}}]\to0$.
Przechodząc z $K\to\infty$ otrzymujemy dokładnie
definicję jednostajnej całkowalności.
\end{proof}

\begin{twierdzenie}[Kryterium $L^1$ dla martyngałów]
Niech $(M_n)$ będzie martyngałem względem
$(\mathcal F_n)$. Następujące warunki są równoważne:
\begin{enumerate}
\item rodzina $\{M_n:n\geq0\}$ jest jednostajnie całkowalna;
\item istnieje $M_\infty\in L^1$ takie, że
$M_n\to M_\infty$ prawie na pewno oraz w $L^1$;
\item istnieje $Y\in L^1$ takie, że
$M_n=\E[Y\mid\mathcal F_n]$ dla każdego $n$.
\end{enumerate}
W warunku (3) można wybrać $Y=M_\infty$.
\end{twierdzenie}
\begin{proof}
Załóżmy (1). Jednostajna całkowalność implikuje
ograniczoność rodziny w $L^1$:
dla dowolnego ustalonego $K$ mamy
$\E|M_n|\leq K+\E[|M_n|\ind{\{|M_n|>K\}}]$.
Twierdzenie o zbieżności martyngałów daje więc
$M_n\to M_\infty$ p.n. i $M_\infty\in L^1$
na mocy Fatou. Z kryterium Vitaliego
(zbieżność według prawdopodobieństwa plus
jednostajna całkowalność) otrzymujemy zbieżność
w $L^1$. To daje (2).

Niech zachodzi (2). Dla $m\geq n$ mamy
$M_n=\E[M_m\mid\mathcal F_n]$.
Ponieważ warunkowanie jest kontrakcją w $L^1$,
przechodząc do granicy $m\to\infty$ dostajemy
$M_n=\E[M_\infty\mid\mathcal F_n]$,
czyli (3).

Jeśli zachodzi (3), to poprzedni lemat
zastosowany do tej samej zmiennej $Y$
i do $\mathcal F_n$ daje jednostajną całkowalność
rodziny $(M_n)$, czyli (1).
\end{proof}

\begin{przyklad}[Martyngał o stałej średniej, którego granica ma średnią zero]
Niech $B_1,B_2,\ldots$ będą niezależnymi
próbami Bernoulliego z parametrem $1/2$.
Połóżmy
$M_0=1$ oraz
$M_n=2^n\ind{\{B_1=\cdots=B_n=1\}}$.
Warunkowo przy historii do chwili $n$
$\E[M_{n+1}\mid\mathcal F_n]=M_n$,
a zatem $\E M_n=1$ dla każdego $n$.
Jednak prawdopodobieństwo, że \emph{wszystkie}
rzuty dały sukces, wynosi
$\lim_n2^{-n}=0$. Zatem $M_n\to0$ prawie na
pewno, choć $\E M_n=1$. Rodzina nie jest
jednostajnie całkowalna, a zbieżność nie
zachodzi w $L^1$. Pojedyncze, coraz rzadsze
i coraz większe wartości niosą całą średnią.
\end{przyklad}

\begin{twierdzenie}[Twierdzenie Lévy'ego o rosnącej informacji]
Niech $(\mathcal F_n)_{n\geq0}$ będzie filtracją
na $(\Omega,\mathcal F,\Prob)$,
$\mathcal F_\infty=\sigma(\bigcup_{n\geq0}\mathcal F_n)$
i $Y\in L^1(\Prob)$. Wtedy
\[
\E[Y\mid\mathcal F_n]
\longrightarrow\E[Y\mid\mathcal F_\infty]
\quad\text{p.n. i w }L^1.
\]
\end{twierdzenie}
\begin{proof}
Niech $M_n=\E[Y\mid\mathcal F_n]$.
Jest to martyngał jednostajnie całkowalny
na mocy poprzedniego lematu; zatem istnieje
$M_\infty\in L^1$ taki, że $M_n\to M_\infty$
p.n. i w $L^1$.
Ponieważ każda $M_n$ jest $\mathcal F_\infty$-
mierzalna, granica ma wersję
$\mathcal F_\infty$-mierzalną.
Dla dowolnego $A\in\mathcal F_k$ i $n\geq k$
\[
\E[M_n\ind A]=\E[Y\ind A].
\]
Przechodząc do granicy w $L^1$, otrzymujemy
$\E[M_\infty\ind A]=\E[Y\ind A]$
dla wszystkich $A$ z $\bigcup_k\mathcal F_k$.
Ta rodzina jest algebrą zbiorów, ponieważ
filtracja jest rosnąca. Klasa zbiorów
$A\in\mathcal F_\infty$ spełniających powyższą
równość jest klasą monotoniczną (wynika to
z zdominowanej zbieżności dla
$M_\infty,Y\in L^1$), więc twierdzenie
o klasie monotonicznej rozszerza równość
na całe $\mathcal F_\infty$.
Z definicji warunkowej wartości oczekiwanej
otrzymujemy
$M_\infty=\E[Y\mid\mathcal F_\infty]$ p.n.
\end{proof}

\begin{przyklad}[Coraz dokładniejsze odczyty obserwacji]
Niech $U\sim\Unif(0,1)$ i
$D_n=2^{-n}\lfloor2^nU\rfloor$,
$\mathcal F_n=\sigma(D_n)$.
Wiedza o $D_n$ oznacza znajomość pierwszych
$n$ cyfr binarnych $U$ (poza nieistotną
kwestią dwóch rozwinięć na zbiorze
prawdopodobieństwa zero). Mamy
$0\leq U-D_n<2^{-n}$, więc $D_n\to U$ p.n.
i $U$ jest $\mathcal F_\infty$-mierzalna.
Twierdzenie Lévy'ego daje
$\E[U\mid\mathcal F_n]\to U$ p.n. i w $L^1$.
Możemy policzyć więcej: przy $D_n=k/2^n$
zmienna $U$ jest jednostajna na przedziale
$[k/2^n,(k+1)/2^n)$, więc
$\E[U\mid\mathcal F_n]=D_n+2^{-(n+1)}$.
Prognoza z obserwacji staje się coraz dokładniejsza.
\end{przyklad}

\section{Zbieżność w \texorpdfstring{$L^2$}{L2} i ograniczona energia}

\begin{twierdzenie}[Martyngał ograniczony w $L^2$]
Niech $(M_n)$ będzie martyngałem z $M_n\in L^2$
i $\sup_n\E M_n^2<\infty$. Wtedy istnieje
$M_\infty\in L^2$ takie, że
$M_n\to M_\infty$ prawie na pewno i w $L^2$.
Ponadto
\[
\sum_{n=1}^{\infty}\E(M_n-M_{n-1})^2
=\E M_\infty^2-\E M_0^2<\infty.
\]
\end{twierdzenie}
\begin{proof}
Z ortogonalności przyrostów
$\sum_{k=1}^n\E D_k^2
=\E M_n^2-\E M_0^2$.
Ponieważ prawa strona jest niemalejąca
i ograniczona, szereg sum kwadratów jest
zbieżny. Dla $m>n$ ortogonalność daje
\[
\E|M_m-M_n|^2=\sum_{k=n+1}^m\E D_k^2\to0
\qquad(n,m\to\infty).
\]
A zatem $(M_n)$ jest ciągiem Cauchy'ego
w zupełnej przestrzeni $L^2$, więc ma granicę
$M_\infty$ w $L^2$. Osobno z nierówności
Cauchy'ego--Schwarza
$\sup_n\E|M_n|\leq\sup_n(\E M_n^2)^{1/2}<\infty$.
Twierdzenie o zbieżności martyngałów daje
zbieżność p.n. do pewnej granicy; granica ta
musi zgadzać się z granicą w $L^2$
(zbieżność w $L^2$ implikuje zbieżność według
prawdopodobieństwa, a granica jest w tym sensie
jednoznaczna). Wreszcie zbieżność w $L^2$
implikuje $\E M_n^2\to\E M_\infty^2$;
przechodzimy do granicy w skończonych
sumach energii.
\end{proof}

\begin{przyklad}[Powrót do procesu rozgałęziającego]
Dla procesu Galtona--Watsona z poprzedniego rozdziału,
przy średniej liczby potomków $m>1$ i
skończonej wariancji $\sigma^2$, martyngał
$W_n=Z_n/m^n$ spełnia
$\sup_n\E W_n^2\leq1+\sigma^2/[m(m-1)]$.
Istnieje więc $W_\infty\in L^2$ takie, że
$W_n\to W_\infty$ p.n. i w $L^2$, a
$\E W_\infty=1$. Nie można jednak na podstawie
samej tej równości stwierdzić, że
$W_\infty>0$ na każdej trajektorii,
na której populacja nigdy nie wymiera;
to wymaga osobnej analizy procesu.
\end{przyklad}

\begin{twierdzenie}[Zbieżność martyngału ograniczonego w $L^p$]
Niech $p>1$ i niech $(M_n)$ będzie martyngałem
spełniającym $\sup_n\E|M_n|^p<\infty$.
Wtedy istnieje $M_\infty\in L^p$ takie, że
$M_n\to M_\infty$ prawie na pewno i w $L^p$.
\end{twierdzenie}
\begin{proof}
Z nierówności Höldera
$\sup_n\E|M_n|\leq(\sup_n\E|M_n|^p)^{1/p}<\infty$.
Twierdzenie Dooba daje więc istnienie skończonej
granicy $M_\infty$ prawie na pewno.
Nierówność maksymalna Dooba w $L^p$ dla
podmartyngału $|M_n|$ mówi, że dla każdego $N$
\[
\E\left(\max_{0\leq k\leq N}|M_k|\right)^p
\leq\left(\frac{p}{p-1}\right)^p\E|M_N|^p.
\]
Prawa strona jest jednostajnie ograniczona.
Z monotonicznej zbieżności przy $N\to\infty$
$M^*:=\sup_{k\geq0}|M_k|$ należy do $L^p$.
Ponieważ $|M_n-M_\infty|^p\leq (2M^*)^p$
prawie na pewno i $M_n\to M_\infty$ p.n.,
zbieżność zdominowana daje
$\E|M_n-M_\infty|^p\to0$.
\end{proof}

\begin{przyklad}[Ograniczenie średniego maksimum]
Jeżeli martyngał ma $\sup_n\E|M_n|^p\leq C$
dla pewnego $p>1$, to jego całkowite
maksimum $\sup_n|M_n|$ ma skończony $p$-ty moment,
a dokładniej
$\E(\sup_n|M_n|)^p\leq(p/(p-1))^pC$.
W przeciwieństwie do samej zbieżności p.n.
otrzymujemy tu kontrolę dużych wychyleń
\emph{całej nieskończonej trajektorii}.
\end{przyklad}

\section{Opcjonalne zatrzymanie przy jednostajnej całkowalności}

\begin{twierdzenie}[Zatrzymanie martyngału domkniętego]
Niech $M_n=\E[Y\mid\mathcal F_n]$, gdzie
$Y\in L^1(\Prob)$, a $\tau$ będzie czasem
zatrzymania z $\Prob(\tau<\infty)=1$.
Wówczas $M_\tau\in L^1$,
$M_\tau=\E[Y\mid\mathcal F_\tau]$ p.n.
i $\E M_\tau=\E M_0$.
\end{twierdzenie}
\begin{proof}
Niech $\tau_n=\tau\wedge n$. Jest to ograniczony
czas zatrzymania, a wzór warunkowania w takim
czasie daje
$M_{\tau_n}=\E[Y\mid\mathcal F_{\tau_n}]$.
Rodzina wszystkich warunkowych wartości tej samej
całkowalnej zmiennej $Y$ jest jednostajnie
całkowalna, zatem $\{M_{\tau_n}:n\geq0\}$
też jest jednostajnie całkowalna.
Ponieważ $\tau<\infty$ p.n., dla niemal każdej
ścieżki od pewnego $n$ mamy $\tau_n=\tau$,
a więc $M_{\tau_n}=M_\tau$. Stąd przez kryterium
Vitaliego $M_{\tau_n}\to M_\tau$ w $L^1$.

Wskażmy dokładnie, dlaczego granica jest warunkową
wartością oczekiwaną. Dla $A\in\mathcal F_\tau$
i każdego $n$ zdarzenie
$A_n=A\cap\{\tau\leq n\}$ należy do
$\mathcal F_{\tau\wedge n}$:
na $\{\tau\wedge n\leq k\}$ przy $k<n$ mamy
$A_n\cap\{\tau\wedge n\leq k\}=A\cap\{\tau\leq k\}
\in\mathcal F_k$, a dla $k\geq n$
$A_n\in\mathcal F_n\subseteq\mathcal F_k$.
Dlatego z definicji warunkowania
\[
 \E[M_{\tau_n}\ind{A_n}]=\E[Y\ind{A_n}].
\]
Ponieważ $\tau<\infty$ p.n., mamy
$\ind{A_n}\to\ind A$ prawie na pewno.
Lewą stronę możemy przeprowadzić do granicy
dzięki zbieżności $M_{\tau_n}\to M_\tau$ w $L^1$
i oszacowaniu
$|\E[M_\tau(\ind{A_n}-\ind A)]|
\leq\E[|M_\tau|\ind{\{\tau>n\}}]\to0$.
Prawa strona zbiega przez zdominowaną zbieżność,
gdyż $Y\in L^1$. Ostatecznie
$\E[M_\tau\ind A]=\E[Y\ind A]$
dla każdego $A\in\mathcal F_\tau$.
Z osobnej własności mierzalności zatrzymanej
zmiennej $M_\tau$ wynika jej $\mathcal F_\tau$-
mierzalność. Zatem
$M_\tau=\E[Y\mid\mathcal F_\tau]$ p.n., a
w szczególności $\E M_\tau=\E Y=\E M_0$.
\end{proof}

\section{Przykłady przekrojowe i typowe pułapki}


\begin{lemat}[Maksimum nieujemnego nadmartyngału]
Niech $(Z_k)_{0\leq k\leq N}$ będzie
nieujemnym nadmartyngałem względem $(\mathcal F_k)$,
a $a>0$. Wtedy
\[
\Prob\left(\max_{0\leq k\leq N}Z_k\geq a\right)
\leq\frac{\E Z_0}{a}.
\]
\end{lemat}
\begin{proof}
Niech $\tau$ będzie pierwszą chwilą osiągnięcia
poziomu $a$, jeśli występuje przed lub w chwili $N$,
a w przeciwnym wypadku niech $\tau=N$.
Czas $\tau$ jest ograniczonym czasem zatrzymania.
Dla zdarzenia $A=\{\max_{k\leq N}Z_k\geq a\}$
mamy $Z_\tau\geq a$ na $A$, natomiast poza $A$
zmienna $Z_\tau$ pozostaje nieujemna.
Stąd $a\Prob(A)\leq\E Z_\tau$.
Z twierdzenia o opcjonalnym zatrzymaniu
nadmartyngału dla czasu ograniczonego
$\E Z_\tau\leq\E Z_0$, co kończy dowód.
\end{proof}

\begin{przyklad}[Ryzyko przekroczenia progu przed czasem $N$]
Dla martyngału o ograniczonych przyrostach
$|D_k|\leq c_k$ można uzyskać silniejsze
oszacowanie \emph{maksimum} niż przez
sumowanie nierówności dla każdej chwili.
Dla $t>0$ z dowodu Azumy wynika, że
\[
Z_k=\exp\left(t(M_k-M_0)-\frac{t^2}{2}
\sum_{j=1}^k c_j^2\right)
\]
jest nieujemnym nadmartyngałem (warunkowa średnia
czynnika jest co najwyżej jeden).
Jeżeli $\max_{k\leq N}(M_k-M_0)\geq u$, to dla
pewnego $k\leq N$ zachodzi
$Z_k\geq\exp(tu-t^2V_N/2)$, gdzie
$V_N=\sum_{j=1}^N c_j^2$.
Dla nieujemnego nadmartyngału $(Z_k)$
możemy zastosować nierówność maksymalną:
$\Prob(\max_{k\leq N}Z_k\geq a)\leq\E Z_0/a=1/a$.
Wynika ona z ograniczonego zatrzymania przy
pierwszym przekroczeniu $a$ i nieujemności
$Z_N$.
Minimalizując po $t=u/V_N$ dostajemy
\[
\Prob\left(\max_{k\leq N}(M_k-M_0)\geq u\right)
\leq e^{-u^2/(2V_N)}.
\]
Oszacowanie dotyczy \emph{całej trajektorii do $N$},
a nie tylko wartości końcowej.
\end{przyklad}

\begin{przyklad}[Dlaczego dla nadmartyngału nie trzeba oczekiwać spadku na każdej ścieżce]
Niech $S_n$ będzie symetrycznym spacerem losowym.
Proces $Z_n=S_n^2$ jest podmartyngałem i
$\E Z_n=n$ rośnie, lecz z prawdopodobieństwem
dodatnim $Z_{n+1}<Z_n$ (np. przejście z
$S_n=2$ do $S_{n+1}=1$). Analogicznie
$-S_n^2$ jest nadmartyngałem, choć
może zwiększyć swoją wartość w jednym kroku.
\end{przyklad}

\begin{przyklad}[Inne założenia, inne wnioski]
Trzy fakty trzeba wyraźnie odróżnić:
\begin{enumerate}
\item nieujemny nadmartyngał zbiega p.n.,
  lecz bez dodatkowej kontroli nie musi zbiegać w $L^1$;
\item martyngał ograniczony w $L^1$ zbiega p.n.,
  lecz może stracić wartość oczekiwaną w granicy;
\item martyngał jednostajnie całkowalny zbiega
  p.n. i w $L^1$, a ograniczony w $L^2$
  zbiega nawet w $L^2$.
\end{enumerate}
Przykład $M_n=2^n\ind{\{B_1=\cdots=B_n=1\}}$
pokazuje jednocześnie, dlaczego nie wolno
zamieniać punktów (2) i (3).
\end{przyklad}

\section{Jak te pojęcia łączą się z łańcuchami Markowa?}

\begin{twierdzenie}[Funkcja harmoniczna łańcucha daje martyngał]
Niech $(X_n)$ będzie jednorodnym łańcuchem Markowa
na przeliczalnym $S$ o macierzy $P$ i naturalnej
filtracji. Niech $h:S\to\R$ będzie ograniczona
i spełnia $Ph=h$, tzn.
$h(i)=\sum_jp_{ij}h(j)$ dla każdego $i\in S$.
Wtedy $(h(X_n))$ jest martyngałem.
Jeżeli $Ph\leq h$, jest nadmartyngałem;
gdy $Ph\geq h$, podmartyngałem.
\end{twierdzenie}
\begin{proof}
Ograniczoność $h$ zapewnia całkowalność.
Z własności Markowa
$\E[h(X_{n+1})\mid\mathcal F_n^X]
=(Ph)(X_n)$.
Podstawienie równości lub odpowiedniej nierówności
prowadzi bezpośrednio do trzech przypadków.
\end{proof}

\begin{przyklad}[Ruina gracza po raz drugi: metoda martyngałowa]
Rozważmy spacer losowy na $\{0,\ldots,N\}$
z granicami pochłaniającymi, z prawdopodobieństwem
kroku w górę $p\in(0,1)$ i w dół $q=1-p$.
Niech $\tau$ będzie chwilą pochłonięcia.
Wykazana wcześniej geometryczna kontrola ogona
$\tau$ daje $\E\tau<\infty$.
Przy $p=q=1/2$ zatrzymany proces $X_{n\wedge\tau}$
jest ograniczonym martyngałem. Zatem przez
zbieżność zdominowaną
$\E_iX_\tau=i$. Ponieważ $X_\tau\in\{0,N\}$,
$N\Prob_i(X_\tau=N)=i$, czyli $h_i=i/N$.
Aby obliczyć średni czas, zauważmy, że
$X_{n\wedge\tau}^2-(n\wedge\tau)$ jest
martyngałem, bo warunkowo przy stanie wewnętrznym
$\E[(X_{n+1}-X_n)^2\mid\mathcal F_n]=1$.
Biorąc wartości oczekiwane i przechodząc z
$n\to\infty$, dostajemy
$\E_i X_\tau^2-\E_i\tau=i^2$;
przejście uzasadnia ograniczoność $X_{n\wedge\tau}^2$
i monotoniczna zbieżność $n\wedge\tau\uparrow\tau$.
Ponieważ $\E_iX_\tau^2=N^2(i/N)=Ni$,
$\E_i\tau=i(N-i)$.

Gdy $p\ne q$, proces
$X_{n\wedge\tau}-(p-q)(n\wedge\tau)$ jest
martyngałem. Otrzymujemy
$N h_i-i=(p-q)\E_i\tau$, zgodnie z wcześniejszym
wzorem. Z kolei funkcja harmoniczna
$h_i=[1-(q/p)^i]/[1-(q/p)^N]$
prowadzi do ograniczonego martyngału $h(X_{n\wedge\tau})$
i pozwala odzyskać prawdopodobieństwo wygranej.
Te obliczenia pokazują trzy zastosowania: martyngały liniowe
służą do wyznaczania wartości oczekiwanych, kwadratowe
do obliczania średnich czasów, a harmoniczne do
wyznaczania prawdopodobieństw trafienia.
\end{przyklad}



\section{Powtórzenie: pytania kontrolne o zależność i informację}

\begin{przyklad}[Dlaczego dwie historie nie zawsze dają tę samą prognozę?]
W łańcuchu Markowa, gdy znamy $X_n=i$, każda historia
o dodatnim prawdopodobieństwie kończąca się stanem $i$
prowadzi do tej samej prognozy $X_{n+1}$.
Natomiast w mieszaninie niezależnych prób
Bernoulliego z ukrytym parametrem (przykład z
rozdziału~\ref{ch:markov}) dwa różne przebiegi
mogą kończyć się tym samym wynikiem $X_n$, ale
prowadzą do różnych prognoz, bo inaczej informują
o nieobserwowanym parametrze. To rozróżnienie
pomaga sprawdzić, czy wybrana zmienna stanu
rzeczywiście \emph{opisuje dostatecznie dużo informacji}.

Własność Markowa \emph{nie} wymaga, by $X_{n+1}$
była niezależna od $X_n$. Jedynie warunkowa
zależność od czasów wcześniejszych niż $n$
zanika po poznaniu $X_n$. Inaczej jest przy
zbieżności do stacjonarności: zanikanie zależności
od $X_0$ jest twierdzeniem o odległych czasach
oraz odpowiednich założeniach na macierz przejścia.
\end{przyklad}

\end{document}
